\documentclass[10pt,letterpaper]{article}
\usepackage{iftex}
\ifPDFTeX
  \usepackage[T1]{fontenc}
  \usepackage[utf8]{inputenc}
  \usepackage{lmodern}
\else
  \usepackage{fontspec}
\fi
\usepackage[spanish,es-tabla]{babel}
\usepackage[letterpaper,margin=2.25cm]{geometry}
\usepackage{microtype}
\usepackage{setspace}
\usepackage{enumitem}
\usepackage{longtable}
\usepackage{booktabs}
\usepackage{array}
\usepackage{tabularx}
\usepackage{ragged2e}
\usepackage{fancyhdr}
\usepackage{titlesec}
\usepackage{hyperref}
\usepackage{xurl}
\usepackage{amsmath,amssymb}
\usepackage{xcolor}
\usepackage{pdflscape}
\ifPDFTeX
  \renewcommand{\familydefault}{\sfdefault}
\else
  \setmainfont{Arimo}
  \setsansfont{Arimo}
  \setmonofont{DejaVu Sans Mono}
\fi
\hypersetup{colorlinks=true,linkcolor=black,urlcolor=blue,pdftitle={SRS TOP Racing v2.6 normativa consolidada},pdfauthor={TOP Racing}}
\setstretch{1.06}
\sloppy
\setlist{itemsep=0.18em,topsep=0.3em}
\setcounter{secnumdepth}{3}
\setcounter{tocdepth}{2}
\makeatletter
\renewcommand{\@pnumwidth}{2.2em}
\makeatother
\pagestyle{fancy}
\setlength{\headheight}{14pt}
\fancyhf{}
\lhead{SRS TOP Racing v2.6}
\rhead{\thepage}
\titleformat{\section}{\Large\bfseries}{\thesection}{0.7em}{}
\titleformat{\subsection}{\large\bfseries}{\thesubsection}{0.7em}{}
\titleformat{\subsubsection}{\normalsize\bfseries}{\thesubsubsection}{0.7em}{}
\newcolumntype{Y}{>{\RaggedRight\arraybackslash}X}
\newcolumntype{P}[1]{>{\RaggedRight\arraybackslash}p{#1}}
\begin{document}
\begin{titlepage}
\centering
\vspace*{2.2cm}
{\Huge\bfseries Especificación de Requerimientos de Software\par}
\vspace{0.5cm}
{\LARGE Sistema TOP Racing\par}
\vspace{1.1cm}
{\Large Versión 2.6 - tecnología abierta obligatoria y acceso operativo publicado\par}
\vspace{0.4cm}
{\large Especificación normativa consolidada\par}
\vfill
{\large 21 de julio de 2026\par}
\vspace{0.4cm}
{\small Fuente de trazabilidad: SRS IEEE 830 integrada heredada y registro de decisiones asociado.\par}
\end{titlepage}
\pagenumbering{roman}
\tableofcontents
\clearpage
\pagenumbering{arabic}

\section{Introducción}
\subsection{Propósito}
Esta SRS define el comportamiento requerido del sistema TOP Racing y constituye la referencia normativa para análisis, diseño, implementación, verificación y aceptación. La versión 2.6 conserva la arquitectura documental normalizada de las versiones 2.0 a 2.5 e integra las decisiones confirmadas sobre subasta sellada de vehículos, eficiencia relativa al mercado, mérito tecnodeportivo, tecnología abierta obligatoria, difusión tecnológica, competencia por peso entre eventos y límites operativos de acceso publicados por el promotor.

La versión 2.6 mantiene que la subasta de vehículos es sellada, exige el máximo nivel de verificación, admite la puja del dueño sobre su propio vehículo, permite una sola puja vigente por usuario y vehículo y solo acepta sustituciones estrictamente mayores. No publica resultados, conteos ni consecuencias derivadas antes del cierre. La línea de tendencia global modela los precios en función del rendimiento observado y la eficiencia mide la diferencia favorable o desfavorable entre el precio de tendencia y el precio oficial de subasta. Registra además que pujar el mínimo real de venta es la estrategia robusta recomendada bajo la información disponible. Finalmente, explicita como rationale institucional que la matriz de pesos produce competencia entre eventos que comparten ámbitos territoriales y temporales comparables: un promotor puede responder mejorando el P/C de su evento o diferenciando territorio y periodo, lo que puede inducir expansión descentralizada del ecosistema sin garantizar por sí mismo cobertura óptima. También corrige el estatuto de la tecnología abierta del vehículo como obligación transversal del dominio, separándola de la admisión irrestricta al evento físico, que puede quedar limitada por reglas publicadas de acceso, cupo, sede, seguridad o elegibilidad operativa dentro del alcance permitido por TOP Racing.

El objetivo de la consolidación es reducir ambigüedad y duplicación sin eliminar contenido del dominio. Los identificadores de la SRS heredada se conservan como procedencia y se relacionan con identificadores canónicos. No quedan decisiones PEND normativas activas en esta versión.

\subsection{Alcance}
TOP Racing organiza y registra eventos competitivos, actores y papeles funcionales, vehículos e inscripciones, subastas de acceso y de vehículos, resultados, clasificaciones, eficiencia relativa al mercado, premios, patrocinios, aportaciones, comisiones, pagos, devoluciones, seguridad comunitaria, disputas de entrega, publicación histórica e interacción comunitaria.

En esta especificación, el término vehículo designa el artefacto competitivo de pista definido por la categoría del evento. La plataforma no presupone que dicho artefacto sea necesariamente un vehículo motorizado de alto costo; el dominio admite categorías con tecnología casera, artesanal, reutilizada, no motorizada o de bajo costo cuando las reglas publicadas del evento y sus restricciones de seguridad lo permitan, pero exige que el vehículo satisfaga la condición transversal de tecnología abierta definida por esta SRS.

El sistema no ejecuta físicamente las carreras ni sustituye a una autoridad deportiva externa. Sí registra datos oficiales, aplica las reglas del dominio que esta SRS declara normativas y ejecuta los flujos económicos reales exigidos por esas reglas.

\subsection{Modelo normativo de la especificación}
Los elementos de la especificación se clasifican de la siguiente forma:
\begin{description}[style=nextline,leftmargin=2.8cm,labelwidth=2.45cm]
\item[DEC] decisión de origen confirmada por el interesado; conserva procedencia y contexto.
\item[BR] regla de negocio o del dominio; define invariantes y políticas conceptuales.
\item[CALC] definición matemática normativa; define variables, ecuaciones y condiciones de cálculo.
\item[RF] requisito funcional atómico; contiene una obligación observable del sistema.
\item[UI] requisito de interfaz externa; especifica información, acciones o restricciones de interacción observables sin imponer tecnología interna.
\item[RC-EXT] restricción externa o de cumplimiento del dominio.
\item[QR] requisito de calidad verificable.
\item[AC] criterio de aceptación asociado a un elemento normativo.
\end{description}

La precedencia es: una definición CALC o una regla BR específica prevalece sobre una descripción contextual de alto nivel; un RF, UI, RC-EXT o QR debe ser compatible con las reglas y cálculos que referencia. Los elementos DEC explican el origen, pero la aceptación del software se realiza contra los elementos normativos canónicos y sus AC.

\subsection{Convenciones de redacción y atomicidad}
Una obligación normativa usa ``debe'' o ``no debe''. Cada RF, UI, RC-EXT o QR expresa una obligación principal verificable. Los efectos explicativos, objetivos de negocio e hipótesis causales se registran como \emph{rationale} y no se interpretan como comportamiento obligatorio salvo que exista un elemento normativo separado.

Los requisitos de interfaz pueden fijar comportamiento observable, visibilidad, información, acciones disponibles y restricciones de interacción confirmadas. No prescriben framework, biblioteca, arquitectura interna, motor de persistencia, estructura física de datos ni disposición gráfica exacta, salvo que una restricción externa lo exija expresamente.

\subsection{Terminología esencial}

\noindent\textbf{Evento.} instancia competitiva con un ciclo de vida definido.\par
\noindent\textbf{Estado del evento.} situación oficial del evento dentro de su ciclo de vida. En esta SRS los estados oficiales son \texttt{creado}, \texttt{inscripciones}, \texttt{velocidad}, \texttt{carrera}, \texttt{subasta}, \texttt{publicado} y \texttt{cancelado}.\par
\noindent\textbf{Promotor.} usuario que convoca y administra un evento específico conforme a las reglas publicadas.\par
\noindent\textbf{Competidor.} papel funcional del usuario que accede a un evento como competidor y compite en él.\par
\noindent\textbf{Dueño de vehículo.} papel funcional del usuario al que pertenece un vehículo registrado en el sistema y que administra su información dentro del contexto del sistema y del evento.\par
\noindent\textbf{Vehículo de competición.} artefacto físico admitido por una categoría de evento para producir rendimiento medible en pista; puede ser motorizado o no motorizado, de tecnología casera, artesanal, reutilizada o de bajo costo si satisface las reglas publicadas, las restricciones de seguridad del evento y la condición transversal de tecnología abierta del vehículo.\par
\noindent\textbf{Tecnología abierta del vehículo.} condición obligatoria del vehículo de competición dentro del dominio TOP Racing por la cual la solución técnica del vehículo debe quedar documentada con información suficiente para ser consultada, estudiada, replicada o reutilizada conforme a los términos publicados aplicables. La apertura tecnológica no implica ausencia de reglas de seguridad, restricciones legales externas, derechos de terceros, control físico del evento ni admisión irrestricta de participantes por parte del promotor.\par
\noindent\textbf{Postor.} papel funcional del usuario que participa en subastas del ecosistema con arreglo a las reglas publicadas.\par
\noindent\textbf{Aportante.} usuario o tercero que realiza una o varias aportaciones económicas, ya sea directamente a la bolsa de un evento o campeonato como patrocinio, o a favor de un participante específico para compartir el premio que a ese participante corresponda.\par
\noindent\textbf{Espectador.} papel funcional del usuario que consulta información pública del ecosistema y puede participar en las funciones públicas disponibles para un evento.\par
\noindent\textbf{Admin.} único usuario administrativo del sistema.\par
\noindent\textbf{Inscripción del competidor.} registro mediante el cual el dueño de un vehículo previamente registrado lo inscribe en un evento, identifica al usuario piloto cuando exista pilotaje humano directo y deja asociados los datos económicos mínimos de esa inscripción.\par
\noindent\textbf{Inscripción cancelada.} inscripción cuyo registro se conserva, pero que queda invalidada a todos los efectos del evento por una causa permitida por las reglas aplicables.\par
\noindent\textbf{Inscripción computable.} inscripción que no está cancelada ni descalificada y cuyo participante obtuvo al menos un tiempo de vuelta válido en calificación; es elegible para la carrera y para los cálculos de eficiencia, terminación y premios dependientes de la actividad competitiva.\par
\noindent\textbf{Límite de carrera.} regla publicada que determina la terminación ordinaria de la carrera por tiempo máximo o por distancia objetivo.\par
\noindent\textbf{Rendimiento observado.} velocidad media correspondiente a la vuelta válida más rápida de calificación de un participante computable; queda fijada antes de la subasta de vehículos.\par
\noindent\textbf{Eficiencia (del dominio).} indicador económico relativo que compara el precio oficial de subasta de un vehículo con el precio de tendencia asociado a su rendimiento por la relación precio--rendimiento del evento.\par
\noindent\textbf{Eficiencia relativa al mercado.} diferencia entre el precio indicado por la línea de tendencia del evento para el rendimiento observado y el precio oficial de subasta del vehículo; es positiva cuando el mercado valoró el vehículo por debajo de la tendencia y negativa cuando lo valoró por encima.\par
\noindent\textbf{Mérito tecnodeportivo.} mérito integrado que combina desarrollo tecnológico costo--eficiente, capacidad de extraer rendimiento deportivo del vehículo y validación práctica de esa combinación mediante rendimiento observado en pista.\par
\noindent\textbf{Pozo económico oficial.} monto económico consolidado del evento a partir del pozo inicial publicado, de las aportaciones directas a la bolsa del evento y de las contribuciones válidamente aceptadas a favor de participantes dentro de la ventana permitida.\par
\noindent\textbf{Matriz de pesos.} estructura publicada por el sistema que expresa el peso relativo con el que un evento aporta puntos a los distintos niveles territoriales y temporales que le resultan aplicables.\par
\noindent\textbf{Puntaje de prioridad (P/C).} indicador publicado propio de un evento que interviene en la determinación de su matriz de pesos conforme a las reglas del dominio.\par
\noindent\textbf{Contribución de eficiencia del participante.} aportación económica asociada a un participante específico que sirve para compartir y aumentar el premio basado en eficiencia que a ese participante corresponda en un evento o en un campeonato de eficiencia.\par
\noindent\textbf{Contribución de terminación del campeonato.} aportación económica asociada a un participante específico de un campeonato de terminación que sirve para compartir y aumentar el premio de terminación que a ese participante corresponda en ese campeonato.\par
\noindent\textbf{Identidad verificada del usuario.} condición por la cual el sistema reconoce que la identidad de un usuario ha sido verificada y, por tanto, puede habilitarle funciones que impliquen participación dentro del ecosistema.\par
\noindent\textbf{Verificación reforzada del usuario.} máximo nivel de verificación operativa del dominio, gestionado por \texttt{admin}, requerido para operaciones de mayor riesgo como inscribir vehículos, pujar en subastas de vehículos y promover eventos.\par
\noindent\textbf{Usuario excluido.} usuario cuyo estado booleano de exclusión está activo y que, por ello, no puede participar operativamente en el sistema.\par
\noindent\textbf{Distribución económica posterior.} reparto posterior al resultado económico individual de un participante entre el propio participante y las personas asociadas a las contribuciones hechas a su favor, según las reglas publicadas.\par
\noindent\textbf{Campeonato de terminación.} campeonato cuya determinación económica usa resultados acumulados de terminación en lugar de resultados de eficiencia.\par
\noindent\textbf{Patrocinio.} Aportación directa a la bolsa de un evento o de un campeonato; puede vincularse a un anuncio conforme a las reglas aplicables.\par
\noindent\textbf{Premio individual.} Parte del pozo aplicable que corresponde a un participante antes de la distribución económica posterior y de los cargos que procedan.\par
\noindent\textbf{Campeonato de eficiencia.} Campeonato cuyo resultado económico computable se basa en la clasificación oficial de eficiencia y la secuencia base del dominio.\par
\subsection{Referencias y procedencia documental}
\begin{itemize}
\item IEEE 830 se conserva como referencia histórica de organización de la SRS; la numeración y el metamodelo de esta versión son propios del proyecto.
\item SRS TOP Racing IEEE 830 integrada heredada: fuente completa de requisitos, decisiones, glosario y trazabilidad previa.
\item Registro de decisiones prioritarias del cliente: consolidado como DEC-001 a DEC-180 en el Anexo B; las decisiones superadas conservan su procedencia histórica y declaran expresamente su sustitución.
\item Informe de mejoras de la SRS: insumo crítico para la refactorización, sin carácter normativo.
\end{itemize}

\section{Descripción general del producto}
\subsection{Perspectiva del producto}
TOP Racing es una plataforma operativa y comunitaria que articula competencia, mercado, liquidación económica y publicación de resultados bajo un mismo conjunto de reglas. El producto organiza y valida digitalmente información producida alrededor de eventos físicos; no ejecuta la actividad física del evento.

\subsection{Actores y papeles funcionales}
Los papeles funcionales son acumulables, excepto \texttt{admin}, que es único y exclusivo. La tabla resume responsabilidades; las autorizaciones concretas se determinan por los RF, UI y RC-EXT.
\begin{longtable}{P{2.8cm}P{11.5cm}}
\toprule
\textbf{Papel} & \textbf{Responsabilidad conceptual}\\
\midrule
Promotor & Convoca y opera eventos, captura datos oficiales, gestiona subastas y publica resultados conforme a sus atribuciones.\\
Competidor & Participa en eventos, queda asociado a resultados, premios y clasificaciones y consulta su información aplicable.\\
Dueño de vehículo & Registra vehículos de su propiedad, los inscribe cuando está habilitado y participa en las obligaciones derivadas de su subasta.\\
Postor & Presenta pujas en subastas para las que está habilitado y consulta sus pujas y resultados permitidos.\\
Aportante & Realiza aportaciones económicas a bolsas o a participantes según las reglas activas y participa en distribuciones posteriores cuando corresponda.\\
Espectador & Consulta información pública y puede participar en la subasta de acceso del evento.\\
Admin & Ejerce las funciones administrativas exclusivas, moderación y gestión global de verificación, exclusión y parámetros del modelo de negocio.\\
\bottomrule
\end{longtable}

\subsection{Mapa de capacidades}
\begin{enumerate}
\item Identidad, acceso, papeles funcionales y aceptación legal versionada.
\item Ciclo de vida, convocatoria, condiciones, reglas y operación de eventos.
\item Vehículos, inscripciones, resultados de velocidad y carrera y comprobación comunitaria de seguridad.
\item Subastas de acceso y vehículos, retenciones, adjudicación, transferencia y disputas de entrega.
\item Campeonatos, clasificaciones, P/C, matriz de pesos y puntos escalados.
\item Eficiencia relativa al mercado calculada y publicada después del cierre de la fase de subasta de vehículos.
\item Premios de evento y campeonato, terminación, contribuciones, patrocinio y distribución económica posterior.
\item Comisiones, cuotas, cobros, pagos, devoluciones y detalle económico al usuario.
\item Consulta pública, historial, foro comunitario, novedades y moderación.
\end{enumerate}

\subsection{Ciclo de vida del evento}
Los estados oficiales son \texttt{creado}, \texttt{inscripciones}, \texttt{velocidad}, \texttt{carrera}, \texttt{subasta}, \texttt{publicado} y \texttt{cancelado}. La secuencia ordinaria es:
\[
\texttt{creado}\rightarrow\texttt{inscripciones}\rightarrow\texttt{velocidad}\rightarrow\texttt{carrera}\rightarrow\texttt{subasta}\rightarrow\texttt{publicado}.
\]
El estado \texttt{cancelado} representa terminación sin publicación oficial de resultados; sus condiciones y efectos económicos se especifican en los requisitos canónicos.

\begin{longtable}{P{2.85cm}P{10.65cm}}
\toprule
\textbf{Estado} & \textbf{Propósito operativo resumido}\\
\midrule
\texttt{creado} & Preparación de datos, reglas y condiciones publicables.\\
\texttt{inscripciones} & Convocatoria publicada, registro de inscripciones y ventanas económicas permitidas.\\
\texttt{velocidad} & Captura de clasificación y datos válidos de velocidad.\\
\texttt{carrera} & Captura de lugares, vueltas y resultados de carrera.\\
\texttt{subasta} & Subasta sellada de vehículos; las pujas y sus consecuencias permanecen reservadas hasta el cierre.\\
\texttt{publicado} & Resultados oficiales, clasificaciones, liquidación y consulta histórica.\\
\texttt{cancelado} & Terminación sin resultados oficiales y ejecución de reversiones económicas aplicables.\\
\bottomrule
\end{longtable}

\subsection{Modelo económico conceptual}
El flujo económico distingue aportaciones directas a bolsa, contribuciones asociadas a participantes o inscripciones, premios individuales, distribución posterior y comisiones. Una aportación directa a la bolsa de un evento o campeonato constituye patrocinio; una aportación asociada a un participante tiene como finalidad aumentar y compartir el premio que corresponda a ese participante.

\begin{longtable}{P{3.7cm}P{9.8cm}}
\toprule
\textbf{Elemento} & \textbf{Relación conceptual}\\
\midrule
Pozo económico oficial del evento & Integra los componentes que las reglas económicas declaren computables para el premio basado en eficiencia.\\
Pozo de terminación & Fondo separado destinado al premio de terminación.\\
Pozo de campeonato & Fondo del campeonato correspondiente, integrado por aportaciones computables dentro de la ventana permitida.\\
Contribución a participante & Se vincula a un participante o inscripción y participa en la magnitud y distribución posterior del premio aplicable.\\
Patrocinio & Aportación directa a una bolsa; puede vincularse a un anuncio con imagen.\\
Distribución posterior & Reparte un premio individual entre el participante y las personas asociadas a contribuciones computables, en proporción a las cantidades aplicables.\\
Comisión o cuota & Cargo del modelo de negocio calculado y cobrado conforme a parámetros administrativos y bases positivas definidas.\\
\bottomrule
\end{longtable}

\subsection{Campeonatos, clasificaciones y matriz de pesos}
La SRS heredada define seis niveles temporales, ocho niveles territoriales y cuatro modalidades, de los que deriva una estructura de 48 campeonatos y 192 clasificaciones. Cada evento determina los campeonatos y clasificaciones aplicables mediante sus niveles territorial y temporal. El P/C propio del evento interviene en la comparación de eventos y en la matriz de pesos; la puntuación escalada de un resultado se obtiene mediante la base de puntos y el peso vigente, conforme al catálogo CALC.

\subsection{Objetivos de calidad}
La claridad operativa, auditabilidad, integridad, transparencia controlada, permanencia histórica, reproducibilidad numérica, retroalimentación de interfaz, legibilidad, contraste, consistencia y prevención de error se especifican normativamente en la sección de QR. Esta sección no introduce obligaciones adicionales.

\section{Reglas normativas del dominio}
\subsection{Reglas de negocio consolidadas}
Las siguientes BR conservan las reglas de alto nivel de la SRS heredada. Las BR siguientes consolidan invariantes y políticas confirmadas; no existen decisiones PEND normativas activas.

\subsubsection{Identidad y acceso}
\noindent\textbf{BR-IAM-001.} Los papeles funcionales son acumulables para un mismo usuario, excepto \texttt{admin}, que es único y exclusivo.\par
\noindent\textbf{BR-IAM-002.} La identidad verificada, la verificación reforzada y la exclusión son estados lógicos independientes.\par
\noindent\textbf{BR-IAM-003.} La participación operativa requiere identidad verificada, ausencia de exclusión activa y aceptación registrada de la versión vigente de términos y condiciones.\par
\noindent\textbf{BR-IAM-004.} La verificación reforzada es adicionalmente obligatoria para inscribir vehículos, pujar en subastas de vehículos y promover eventos.\par
\noindent\textbf{BR-IAM-005.} El nivel visible de acceso se deriva de los estados de verificación, exclusión y aceptación legal; el estado de sesión no modifica por sí mismo ese nivel.\par
\subsubsection{Evento y estados}
\noindent\textbf{BR-EVT-001.} Los estados oficiales del evento son \texttt{creado}, \texttt{inscripciones}, \texttt{velocidad}, \texttt{carrera}, \texttt{subasta}, \texttt{publicado} y \texttt{cancelado}.\par
\noindent\textbf{BR-EVT-002.} La secuencia ordinaria del evento avanza desde \texttt{creado} hasta \texttt{publicado}; toda transición requiere las guardas normativas aplicables.\par
\noindent\textbf{BR-EVT-003.} La convocatoria corresponde a la publicación del evento en estado \texttt{inscripciones}; las reglas, condiciones y datos públicos obligatorios deben estar publicados antes de abrir inscripciones.\par
\noindent\textbf{BR-EVT-004.} Una inscripción cancelada conserva su registro histórico, permanece visible como cancelada y no interviene en resultados, clasificaciones ni cálculos económicos posteriores.\par
\noindent\textbf{BR-EVT-005.} Un evento cancelado termina sin publicación oficial de resultados y activa las devoluciones o liberaciones económicas exigidas por las reglas aplicables.\par
\subsubsection{Vehículos e inscripciones}
\noindent\textbf{BR-VEH-001.} Un vehículo solo puede inscribirse en un evento si está previamente registrado como propiedad del usuario dueño que realiza la inscripción.\par
\noindent\textbf{BR-VEH-002.} La inscripción al evento la registra el dueño del vehículo y ese dueño es el participante de la inscripción a efectos competitivos y económicos; cuando existe pilotaje humano directo, el piloto puede ser un usuario distinto del participante dueño.\par
\noindent\textbf{BR-VEH-003.} Solo un vehículo asociado a una inscripción computable que haya participado en el evento puede ser sometido a la subasta de vehículos; antes de la subasta debe estar bajo control del evento mediante entrega al promotor y queda sujeto a transferencia obligatoria según el resultado oficial vigente.\par
\noindent\textbf{BR-VEH-004.} Una inscripción solo califica para la carrera si el participante obtiene al menos un tiempo de vuelta válido durante la calificación. La falta de tiempo válido produce condición de \texttt{no calificada}; no equivale por sí misma a cancelación ni descalificación y no genera devolución automática de los importes ya cobrados.\par
\noindent\textbf{BR-TEC-001.} TOP Racing no impone como condición de dominio que el vehículo de competición sea motorizado, costoso o perteneciente a una frontera tecnológica avanzada; cada categoría de evento debe definir las condiciones técnicas, de seguridad, medición y participación aplicables al artefacto competitivo admitido, manteniendo la obligación transversal de tecnología abierta del vehículo dentro del dominio TOP Racing.\par
\noindent\textbf{BR-TEC-002.} Las reglas de una categoría pueden admitir tecnología casera, artesanal, reutilizada, no motorizada o de bajo costo siempre que cumpla las restricciones de seguridad, medición y participación publicadas. En todos los casos, el vehículo de competición debe satisfacer la condición de tecnología abierta del vehículo definida por esta SRS.\par
\noindent\textbf{BR-TEC-003.} La finalidad económica de la eficiencia relativa al mercado es favorecer el mérito tecnodeportivo: rendimiento observado alto en relación con la valoración económica revelada por el mercado del evento, sin prometer que todo costo de desarrollo quede observado por el sistema.\par
\noindent\textbf{BR-TEC-004.} La tecnología abierta del vehículo permite la consulta y estudio de la documentación técnica publicada conforme a los términos aplicables. La transferencia de un vehículo mediante subasta puede permitir adicionalmente que el adjudicatario estudie, inspeccione o analice técnicamente el artefacto físico adquirido. TOP Racing no impone una confidencialidad tecnológica general por defecto, sin perjuicio de restricciones legales externas, seguridad, derechos de terceros o términos específicos publicados.\par
\noindent\textbf{BR-TEC-005.} La apertura tecnológica del vehículo es una obligación transversal del dominio TOP Racing y debe registrarse y publicarse conforme a las reglas aplicables del evento. La apertura tecnológica no sustituye los resultados medidos ni el cálculo de eficiencia como base obligatoria del mérito competitivo.\par
\noindent\textbf{BR-TEC-006.} La obligación de tecnología abierta del vehículo no implica acceso irrestricto a todos los eventos. El promotor puede publicar reglas de acceso, cupo, seguridad, sede, elegibilidad operativa o participación que limiten la admisión al evento cuando dichas reglas sean compatibles con las obligaciones normativas del dominio. TOP Racing debe registrar y aplicar las reglas publicadas que estén dentro de su alcance operativo, pero no garantiza por sí mismo el acceso físico o externo a un evento.\par
\subsubsection{Subastas}
\noindent\textbf{BR-AUC-001.} La subasta de acceso corresponde al acceso de espectadores; el acceso del competidor se produce por inscripción.\par
\noindent\textbf{BR-AUC-002.} La fase de subasta de vehículos es sellada: antes de su cierre no se publican pujas, conteos de pujas, orden provisional, precio provisional, adjudicatario provisional ni resultados derivados de eficiencia o premio. Cada usuario autorizado puede consultar privadamente su propia puja vigente.\par
\noindent\textbf{BR-AUC-003.} En cada subasta de vehículo existe como máximo una puja vigente por usuario; una puja vigente solo puede sustituirse por otra de importe estrictamente mayor.\par
\noindent\textbf{BR-AUC-004.} Cualquier usuario con verificación reforzada activa, aceptación vigente de términos y sin exclusión puede pujar por vehículos con independencia de sus demás papeles funcionales. El dueño puede pujar por su propio vehículo.\par
\noindent\textbf{BR-AUC-005.} La entrega, disputa y resolución de un vehículo adjudicado a un usuario distinto del dueño previo están sujetas a los plazos publicados del evento. Cuando el dueño previo resulte adjudicatario, conserva el vehículo y no se inicia un flujo de entrega o disputa consigo mismo.\par
\noindent\textbf{BR-AUC-006.} Las pujas de usuarios distintos del dueño requieren retención suficiente del principal. La puja del dueño sobre su propio vehículo está exenta de retención del principal, pero no de la comisión aplicable al comprador cuando resulte adjudicatario.\par
\subsubsection{Publicación y reserva}
\noindent\textbf{BR-PUB-001.} Toda información que el dominio declare pública debe permanecer visible con el alcance temporal aplicable; la información declarada reservada debe restringirse a los actores autorizados.\par
\noindent\textbf{BR-PUB-002.} Todo resultado oficial debe conservar trazabilidad suficiente hacia sus datos de origen y las reglas aplicadas.\par
\noindent\textbf{BR-PUB-003.} La identidad de quien presenta una protesta de seguridad permanece reservada para la consulta pública, pero puede ser visible para \texttt{admin} y el promotor del evento.\par
\subsubsection{Campeonatos}
\noindent\textbf{BR-CHP-001.} La estructura del dominio comprende seis niveles temporales, ocho niveles territoriales y cuatro modalidades, que originan 48 campeonatos y 192 clasificaciones.\par
\noindent\textbf{BR-CHP-002.} Los campeonatos y clasificaciones aplicables a un evento se determinan por sus niveles territorial y temporal y por la modalidad correspondiente.\par
\subsubsection{Matriz de pesos}
\noindent\textbf{BR-WGT-001.} La matriz de pesos del evento debe estar publicada desde \texttt{inscripciones} y permanecer visible durante el ciclo del evento.\par
\noindent\textbf{BR-WGT-002.} La matriz de pesos parte de 100; los niveles más bajos, semana y variante, conservan peso fijo 100.\par
\noindent\textbf{BR-WGT-003.} Cuando dos eventos comparables se afectan, el P/C determina cuál conserva peso y cuál desciende en la secuencia base; los empates de P/C favorecen al evento registrado antes.\par
\noindent\textbf{BR-WGT-004.} El costo por participante es el costo económico mínimo obligatorio publicado de una inscripción ordinaria: cuota obligatoria de inscripción, cargo obligatorio de renta de pista imputable a la inscripción y costo obligatorio de pesaje imputable a la inscripción. No incluye aportaciones voluntarias, patrocinios, pujas, compras de vehículos, retenciones reembolsables ni comisiones por beneficios posteriores.\par
\noindent\textbf{BR-WGT-005.} El costo por participante se fija al abrir \texttt{inscripciones} y no cambia por inscripciones posteriores, cancelaciones, descalificaciones ni variaciones del número real de participantes. Si la renta se publica únicamente como costo total del evento, su imputación analítica por participante usa el número esperado de participantes publicado para la convocatoria.\par
\noindent\textbf{BR-WGT-006.} La competencia por peso se produce únicamente entre eventos comparables que se solapan en los niveles territoriales y temporales inferiores aplicables al campeonato superior correspondiente; fuera de ese solapamiento no existe competencia directa por el mismo peso en esa dimensión.\par
\noindent\textbf{BR-WGT-007.} La matriz de pesos crea una presión de diferenciación institucional: un promotor puede mejorar el P/C de su evento para competir dentro de un ámbito congestionado o puede elegir otro territorio o periodo que reduzca el conjunto de eventos comparables, sujeto a las reglas efectivas de campeonatos y comparación.\par
\noindent\textbf{BR-WGT-008.} Los pesos fijos 100 de semana y variante preservan la relevancia de los resultados en los niveles mínimos; los descensos de peso afectan la contribución a campeonatos superiores compartidos y no invalidan el evento ni sus resultados locales.\par
\subsubsection{Eficiencia}
\noindent\textbf{BR-EFF-001.} La eficiencia relativa al mercado compara el precio oficial de subasta de cada vehículo con el precio que la línea de tendencia global del evento asocia a su rendimiento observado.\par
\noindent\textbf{BR-EFF-002.} El rendimiento observado es la velocidad media correspondiente a la vuelta válida más rápida de calificación y queda fijado antes de abrir la subasta de vehículos.\par
\noindent\textbf{BR-EFF-003.} La línea de tendencia de mercado se ajusta una sola vez, después del cierre de las subastas de vehículos, sobre todos los pares válidos de rendimiento observado y precio oficial de subasta del evento, mediante el modelo log-lineal $\ln(\widehat{C}(S))=a+bS$, equivalente a $\widehat{C}(S)=\exp(a+bS)$.\par
\noindent\textbf{BR-EFF-004.} La diferencia de eficiencia es el precio de tendencia correspondiente al rendimiento observado menos el precio oficial de subasta: $D_i=\widehat{C}(S_i)-C_i$. El orden es descendente y los empates favorecen la inscripción más antigua.\par
\noindent\textbf{BR-EFF-005.} La eficiencia, los puntos de eficiencia y el premio basado en eficiencia no se publican ni se presentan como resultados provisionales mientras la subasta de vehículos permanezca abierta; solo se determinan y publican a partir de los precios oficiales posteriores al cierre.\par
\noindent\textbf{BR-EFF-006.} La muestra de eficiencia contiene un único par válido $(S_i,C_i)$ por inscripción computable. Cada participante contribuye a la misma línea de tendencia global y es evaluado respecto de esa línea. Los precios no positivos y los registros técnicamente duplicados no forman parte de la muestra válida.\par
\subsubsection{Economía y aportaciones}
\noindent\textbf{BR-ECO-001.} Toda aportación directa a la bolsa de un evento o campeonato constituye patrocinio y puede asociarse a publicidad del patrocinador.\par
\noindent\textbf{BR-ECO-002.} Una aportación asociada a un participante o inscripción tiene como finalidad aumentar y compartir el premio que corresponda a ese beneficiario.\par
\noindent\textbf{BR-ECO-003.} Las aportaciones del evento se aceptan únicamente durante \texttt{inscripciones}; las aportaciones de campeonato se aceptan solo hasta antes del cierre de inscripciones del primer evento del campeonato.\par
\noindent\textbf{BR-ECO-004.} La función de aportaciones a participantes de evento y de campeonato puede activarse o desactivarse globalmente por \texttt{admin}; la aportación directa a bolsa permanece conceptualmente diferenciada.\par
\noindent\textbf{BR-ECO-005.} Mientras la fase de subasta de vehículos permanezca abierta, la información pública usada por un participante para valorar su premio de eficiencia se limita al pozo inicial publicado; el pozo económico oficial consolidado y los importes de aportaciones o contribuciones que lo integran no se publican hasta después del cierre de la fase.\par
\subsubsection{Premios de evento}
\noindent\textbf{BR-EPR-001.} El premio de eficiencia del evento depende de la base de puntos de eficiencia y de la contribución total de eficiencia computable del participante.\par
\noindent\textbf{BR-EPR-002.} El premio de terminación es separado del premio de eficiencia y solo incluye inscripciones computables que terminen efectivamente la carrera.\par
\noindent\textbf{BR-EPR-003.} El peso de terminación depende de la contribución total de terminación y de la relación entre vueltas completadas y vueltas esperadas.\par
\noindent\textbf{BR-EPR-004.} Cuando la suma de coeficientes de eficiencia o la suma de pesos de terminación sea cero, no existe base proporcional para adjudicar el premio correspondiente; el sistema debe restituir las cantidades computables a sus respectivos orígenes económicos y no debe cobrar comisión de beneficio sobre cantidades restituidas.\par
\subsubsection{Distribución posterior}
\noindent\textbf{BR-DST-001.} La distribución económica posterior reparte un premio individual entre el participante y las personas asociadas a contribuciones computables, en proporción a las cantidades que integran la contribución total aplicable.\par
\noindent\textbf{BR-DST-002.} Si un premio individual ya existe pero la suma de contribuciones computables usada para su distribución posterior es cero, el participante recibe el 100\% del premio individual.\par
\subsubsection{Premios de campeonato}
\noindent\textbf{BR-CPR-001.} El premio de campeonato puede corresponder a un campeonato de eficiencia o de terminación.\par
\noindent\textbf{BR-CPR-002.} En un campeonato de eficiencia, el resultado computable del premio se deriva de los lugares oficiales y la secuencia base; en uno de terminación, se deriva de la suma de relaciones vueltas completadas/vueltas esperadas de los eventos computables.\par
\subsubsection{Modelo de negocio}
\noindent\textbf{BR-BUS-001.} Las comisiones del competidor y de aportantes asociados a participantes se calculan únicamente sobre la parte positiva definida para cada caso.\par
\noindent\textbf{BR-BUS-002.} El comprador o adjudicatario de un vehículo está sujeto a la comisión de vehículo calculada sobre el precio oficial de subasta, incluso cuando el adjudicatario sea el dueño previo. El vendedor no soporta esa comisión en su condición de vendedor.\par
\noindent\textbf{BR-BUS-003.} Las aportaciones directas de patrocinio están exentas de comisión para el patrocinador.\par
\noindent\textbf{BR-BUS-004.} Las comisiones y cuotas se cobran dentro del flujo económico real en el que se produce el beneficio u operación gravada, y el usuario debe conocer previamente la regla y el parámetro aplicables.\par
\subsubsection{Seguridad}
\noindent\textbf{BR-SAF-001.} Las objeciones y protestas de seguridad se tramitan dentro de las ventanas y reglas del evento; sus resultados públicos y reservas de identidad se rigen por las reglas de publicación.\par
\subsubsection{Comunidad}
\noindent\textbf{BR-COM-001.} Los usuarios con identidad verificada pueden participar en el foro comunitario y el canal de novedades; la moderación corresponde a \texttt{admin}.\par
\subsubsection{Política numérica y monetaria}
\noindent\textbf{BR-NUM-001.} Todo importe monetario cobrado, retenido, devuelto, pagado o conservado como saldo se representa como un entero no negativo de unidades completas de la moneda declarada para el evento o campeonato; TOP Racing no admite fracciones de unidad monetaria en operaciones económicas.\par
\noindent\textbf{BR-NUM-002.} Los cálculos analíticos del dominio, incluidas regresiones, velocidades, expectativas, diferencias, razones y coeficientes, usan aritmética de punto flotante binario de 64 bits y no aplican redondeo de dominio intermedio salvo que una definición CALC lo prescriba expresamente.\par
\noindent\textbf{BR-NUM-003.} Toda distribución proporcional que termine en pagos monetarios enteros debe reconciliar exactamente el total distribuido mediante \texttt{CALC-MON-001}; ninguna operación de reparto puede crear ni destruir unidades monetarias por efecto de conversión o redondeo.\par
\subsubsection{Compatibilidad de operaciones}
\noindent\textbf{BR-OPS-001.} Dos operaciones son incompatibles cuando modifican el mismo agregado, una cambia el estado de un agregado del que depende la otra, ambas pueden reservar o debitar el mismo saldo, ambas modifican la misma puja vigente, una invalida una precondición de la otra o ambas pertenecen al mismo alcance de idempotencia y representan una posible ejecución duplicada.\par
\noindent\textbf{BR-OPS-002.} Las consultas no se bloquean por defecto durante una operación pendiente. La interfaz debe bloquear únicamente las operaciones para las que \texttt{BR-OPS-001} determine conflicto dentro del alcance afectado.\par
\subsection{Catálogo de cálculos}
Las definiciones CALC son normativas. Una implementación puede utilizar cualquier tecnología interna que reproduzca el resultado definido bajo la política numérica vigente.

\noindent\textbf{CALC-WGT-001 --- Puntaje de prioridad del evento}\par
Para un evento $e$, sea $I(e)$ la cuota obligatoria de inscripción, $R(e)$ el cargo obligatorio de renta de pista imputable a una inscripción y $W(e)$ el costo obligatorio de pesaje imputable a una inscripción. El costo por participante es $C(e)=I(e)+R(e)+W(e)$. Si la renta se publica únicamente como costo total $R_{total}(e)$, se usa analíticamente $R(e)=R_{total}(e)/N_{esperado}(e)$, donde $N_{esperado}(e)>0$ es el número esperado de participantes publicado. La contribución mínima incluida en $I(e)$ no se suma nuevamente. Cuando $C(e)>0$, $PC(e)=P(e)/C(e)$. Cuando $C(e)=0$ y $P(e)>0$, $PC(e)=+\infty$ exclusivamente para ordenamiento de prioridad; los empates favorecen al evento registrado antes.\par\vspace{0.5em}
\noindent\textbf{CALC-WGT-002 --- Escalamiento de puntos por peso}\par
Para una base de puntos $B$ y un peso vigente $W$, la puntuación escalada es $S=B\times W/100$. No se aplica redondeo intermedio no prescrito; la política numérica se define en NUM-001.\par\vspace{0.5em}
\noindent\textbf{CALC-EFF-001 --- Rendimiento observado para eficiencia}\par
Para cada participante computable $i$, sea $d_e>0$ la longitud oficial de una vuelta del evento y sea $T_i^*>0$ el menor tiempo de una vuelta válida de calificación del participante. El rendimiento observado es la velocidad media de esa vuelta: $S_i=d_e/T_i^*$. El valor $S_i$ queda fijado antes de la subasta de vehículos y no cambia por las pujas ni por el resultado de la subasta.\par\vspace{0.5em}
\noindent\textbf{CALC-EFF-002 --- Línea de tendencia log-lineal de precios en función del rendimiento}\par
Sea $n$ el número de pares válidos $(S_i,C_i)$ del evento, con $S_i>0$, $C_i>0$ y un único par por inscripción computable. Si $n=0$, no existe clasificación de eficiencia. Si $n=1$, el único participante ocupa el primer lugar y se define $D_1=0$. Si $n=2$, ambos participantes tienen diferencia de eficiencia normativa igual a cero y empatan; el orden se resuelve por la marca temporal de inscripción válida más antigua. Si $n\ge3$ y la variable $S_i$ tiene variación suficiente para estimar una pendiente única, el sistema estima por mínimos cuadrados ordinarios una única función global $\ln(C_i)=a+bS_i+\varepsilon_i$ sobre todos los participantes computables. El precio de tendencia aplicable es $\widehat{C}(S_i)=\exp(a+bS_i)$. Si la matriz del modelo no tiene rango suficiente para estimar una pendiente única, se usa el modelo de solo intercepto $\widehat{C}_i=\exp(\overline{\ln C})$, media geométrica de los precios oficiales válidos.\par\vspace{0.5em}
\noindent\textbf{CALC-EFF-003 --- Diferencia de eficiencia relativa al mercado}\par
Para $n=1$ se define $D_1=0$. Para $n=2$ se definen $D_1=D_2=0$ sin comparar residuos numéricos calculados. Para $n\ge3$, sea $\widehat{C}(S_i)$ el precio de tendencia indicado por \texttt{CALC-EFF-002} para el rendimiento observado $S_i$ y sea $C_i$ el precio oficial de subasta. La diferencia de eficiencia es $D_i=\widehat{C}(S_i)-C_i$. Un valor positivo indica valoración inferior a la tendencia para el rendimiento observado; un valor negativo indica valoración superior a la tendencia. El orden de eficiencia es descendente por $D_i$ y los empates se resuelven por la marca temporal de inscripción válida más antigua.\par\vspace{0.5em}
\noindent\textbf{CALC-PTS-001 --- Secuencia base de puntos}\par
Sean $B_0=100$, $B_1=63$, $B_2=40$, $B_3=25$ y $B_4=16$. Para una posición ordinal $r\ge1$, sean $q=\lfloor(r-1)/5\rfloor$ y $j=(r-1)\bmod5$, con $j\in\{0,1,2,3,4\}$. La base de puntos es $P(r)=B_j/10^q$. Por tanto, la secuencia comienza $100,63,40,25,16,10,6.3,4,2.5,1.6,1,0.63,0.4,0.25,0.16,\ldots$ y continúa indefinidamente por la misma recurrencia.\par\vspace{0.5em}
\noindent\textbf{CALC-EPR-001 --- Coeficiente de premio de eficiencia del evento}\par
Para cada participante computable $i$, sea $B_i$ su base de puntos de eficiencia y $E_i$ su contribución total de eficiencia computable. El coeficiente es $K_i=B_i E_i$.\par\vspace{0.5em}
\noindent\textbf{CALC-EPR-002 --- Premio individual de eficiencia del evento}\par
Sea $P$ el pozo económico oficial y $K=\sum_j K_j$ la suma de coeficientes de participantes computables. Si $K>0$, la cuota ideal de eficiencia de $i$ es $P_i^*=P(K_i/K)$ y los importes monetarios enteros se asignan mediante \texttt{CALC-MON-001}. Si $K=0$, no se adjudica premio individual de eficiencia y las cantidades computables del pozo se restituyen a sus respectivos orígenes económicos.\par\vspace{0.5em}
\noindent\textbf{CALC-EPR-003 --- Pozo computable del premio de terminación}\par
Sea $T_0$ el pozo publicado del premio de terminación y $t_i$ la contribución total de terminación computable de cada inscripción. El pozo computable es $T=T_0+\sum_i t_i$.\par\vspace{0.5em}
\noindent\textbf{CALC-EPR-004 --- Vueltas esperadas}\par
La carrera declara un único tipo de límite: \texttt{tiempo} o \texttt{distancia}. Para una carrera limitada por tiempo, sea $T_e>0$ el tiempo máximo oficial y $L_i>0$ el mejor tiempo de vuelta válido de calificación de la inscripción computable $i$, expresados en la misma unidad temporal; entonces $E_{i,e}=T_e/L_i$. Para una carrera limitada por distancia, sea $D_e>0$ la distancia objetivo y $d_e>0$ la longitud oficial de una vuelta, expresadas en la misma unidad de distancia; entonces $E_e=D_e/d_e$ para todas las inscripciones computables. Las vueltas esperadas pueden ser fraccionarias y no se redondean antes de calcular $v/E$. Una inscripción sin tiempo de vuelta válido de calificación no es computable y no entra en este cálculo.\par\vspace{0.5em}
\noindent\textbf{CALC-EPR-005 --- Peso de terminación}\par
Para una inscripción elegible $i$, con contribución total de terminación $t_i$, vueltas completadas $v_i$ y vueltas esperadas $E_i>0$, el peso de terminación es $W_i=t_i(v_i/E_i)$.\par\vspace{0.5em}
\noindent\textbf{CALC-EPR-006 --- Premio individual de terminación}\par
Sea $T$ el pozo computable de terminación y $W=\sum_j W_j$ la suma de pesos de las inscripciones elegibles. Si $W>0$, la cuota ideal de terminación de $i$ es $T_i^*=T(W_i/W)$ y los importes monetarios enteros se asignan mediante \texttt{CALC-MON-001}. Si $W=0$, no se adjudica premio de terminación y las cantidades computables se restituyen a sus respectivos orígenes económicos.\par\vspace{0.5em}
\noindent\textbf{CALC-DST-001 --- Distribución posterior proporcional}\par
Para un premio individual entero $R$ y contribuciones computables $q_1,\ldots,q_n$ vinculadas al participante y a terceros, sea $Q=\sum_j q_j$. Si $Q>0$, la cuota ideal del beneficiario $j$ es $R_j^*=R(q_j/Q)$ y los importes monetarios enteros se asignan mediante \texttt{CALC-MON-001}. Si $Q=0$, el participante recibe $R$ íntegramente. La identidad pública de terceros se rige por las reglas de reserva aplicables.\par\vspace{0.5em}
\noindent\textbf{CALC-MON-001 --- Reconciliación monetaria por mayor residuo}\par
Para distribuir un total monetario entero $M$ a partir de cuotas ideales no negativas $x_1,\ldots,x_n$ cuya suma matemática es $M$, se asigna inicialmente $m_i=\lfloor x_i\rfloor$ a cada beneficiario. Sea $L=M-\sum_i m_i$. Las $L$ unidades restantes se asignan, una por beneficiario y por ronda si fuera necesario, en orden descendente de la parte fraccionaria $x_i-\lfloor x_i\rfloor$. Los empates de residuo se resuelven por la marca temporal normativa del derecho económico aplicable y, si persiste el empate, por identificador canónico ascendente. El resultado debe satisfacer exactamente $\sum_i m_i=M$.\par\vspace{0.5em}
\noindent\textbf{CALC-CPR-001 --- Resultado computable de campeonato de terminación}\par
Para el participante $i$ en un campeonato de terminación, el resultado computable es $R_i=\sum_{e\in E_i} v_{i,e}/E_{i,e}$ sobre los eventos computables del campeonato.\par\vspace{0.5em}
\noindent\textbf{CALC-CPR-002 --- Coeficiente y premio de campeonato}\par
Para cada participante $i$, sea $R_i$ su resultado computable y $Q_i$ su contribución total computable. El coeficiente es $K_i=R_iQ_i$. Sea $K=\sum_jK_j$ y $P_C$ el pozo de campeonato; si $K>0$, la cuota ideal es $P_{C,i}^*=P_C(K_i/K)$ y los importes monetarios enteros se asignan mediante \texttt{CALC-MON-001}. Si $K=0$, no se adjudica premio de campeonato y las cantidades computables se restituyen a sus respectivos orígenes económicos.\par\vspace{0.5em}
\noindent\textbf{CALC-BUS-001 --- Base positiva de comisión}\par
Cuando una regla de comisión se aplique únicamente a una parte positiva, la base es $B^+=\max(0,R-C)$, donde $R$ es el beneficio o parte de premio definida para el caso y $C$ es el costo o aportación computable que la regla de negocio atribuye al mismo caso. El porcentaje aplicable es un parámetro administrativo vigente.\par\vspace{0.5em}
\subsection{Políticas transversales cerradas}
Las decisiones que en la versión 2.0 se identificaban como \texttt{PEND-001} a \texttt{PEND-009} se consideran cerradas por \texttt{DEC-151} a \texttt{DEC-159}. Las políticas siguientes son normativas y forman parte del oráculo de aceptación.

\subsubsection{NUM-001 --- Política numérica y monetaria}
Los importes monetarios se almacenan y transfieren como enteros no negativos de unidades completas de la moneda del evento o campeonato. Los cálculos analíticos usan punto flotante binario de 64 bits y no aplican redondeo de dominio intermedio. Las distribuciones monetarias proporcionales convierten sus cuotas ideales a importes enteros exclusivamente mediante \texttt{CALC-MON-001} y deben conservar exactamente el total monetario distribuido.

\subsubsection{OPS-001 --- Matriz de compatibilidad operativa}
Las operaciones se clasifican como consulta (Q), mutación ordinaria (M), transición de estado (T), operación financiera (F) o mutación de subasta (A). La tabla define la compatibilidad base; \emph{Cond.} exige evaluar \texttt{BR-OPS-001} y \emph{No*} significa incompatibilidad cuando las operaciones pertenecen al mismo evento o agregado dependiente.
\begin{center}
\begin{tabular}{lccccc}
\toprule
\textbf{Pendiente / nueva} & Q & M & T & F & A\\
\midrule
Q & Sí & Sí & Sí & Sí & Sí\\
M & Sí & Cond. & No* & Cond. & Cond.\\
T & Sí & No* & No* & No* & No*\\
F & Sí & Cond. & No* & Cond. & Cond.\\
A & Sí & Cond. & No* & Cond. & Cond.\\
\bottomrule
\end{tabular}
\end{center}

\subsubsection{ENV-UX-001 --- Entorno de verificación UX y accesibilidad}
La verificación de interfaz debe cubrir los viewports de referencia 320$\times$568, 390$\times$844, 768$\times$1024 y 1366$\times$768 píxeles CSS; zoom de 100\% y 200\%; teclado y puntero; y orientación vertical u horizontal cuando aplique al dispositivo. La interfaz web debe cumplir WCAG 2.2 nivel AA. Como evaluación complementaria, las pantallas críticas deben probarse en condiciones reales de pista al aire libre; esta prueba de campo no sustituye los criterios cuantitativos de accesibilidad ni presupone control sobre el hardware, brillo o iluminación ambiental del dispositivo del usuario.

\section{Requisitos funcionales y de interfaz}
\subsection{Convención de presentación}
Cada requisito canónico contiene una obligación, prioridad, procedencia y referencia a un criterio de aceptación separado. La procedencia conserva el ID heredado; un requisito heredado compuesto puede originar varios requisitos canónicos.

\subsection{Eventos, estados y convocatoria}
\noindent\textbf{RF-EVT-001 --- Crear evento}\par
\begin{description}[style=nextline,leftmargin=3.15cm,labelwidth=2.8cm,font=\normalfont\bfseries]
\item[Obligación] El sistema debe permitir registrar un evento con sus datos iniciales y sus reglas publicadas, y dejarlo en estado \texttt{creado}.
\item[Prioridad] Alta
\item[Procedencia] Secciones 1.2, 2.2 y 2.6.1.; ID heredado: \texttt{RF-001}.
\item[Verificación] \texttt{AC-RF-EVT-001}
\end{description}
\vspace{0.25em}\hrule\vspace{0.65em}
\noindent\textbf{RF-EVT-002 --- Abrir inscripciones del evento}\par
\begin{description}[style=nextline,leftmargin=3.15cm,labelwidth=2.8cm,font=\normalfont\bfseries]
\item[Obligación] El sistema debe permitir abrir las inscripciones de un evento solo cuando sus datos públicos y sus reglas aplicables ya estén publicados.
\item[Prioridad] Alta
\item[Procedencia] Secciones 2.5, 2.6 y 2.6.1.; ID heredado: \texttt{RF-002}.
\item[Verificación] \texttt{AC-RF-EVT-002}
\end{description}
\vspace{0.25em}\hrule\vspace{0.65em}
\noindent\textbf{RF-EVT-003 --- Vocabulario oficial de estados del evento}\par
\begin{description}[style=nextline,leftmargin=3.15cm,labelwidth=2.8cm,font=\normalfont\bfseries]
\item[Obligación] El sistema debe representar el estado oficial de un evento exclusivamente mediante los valores \texttt{creado}, \texttt{inscripciones}, \texttt{velocidad}, \texttt{carrera}, \texttt{subasta}, \texttt{publicado} y \texttt{cancelado}.
\item[Prioridad] Alta
\item[Procedencia] Secciones 2.2, 2.5.1 y 2.7.; ID heredado: \texttt{RF-003}.
\item[Verificación] \texttt{AC-RF-EVT-003}
\end{description}
\vspace{0.25em}\hrule\vspace{0.65em}
\noindent\textbf{RF-EVT-004 --- Bloqueo de transiciones no autorizadas}\par
\begin{description}[style=nextline,leftmargin=3.15cm,labelwidth=2.8cm,font=\normalfont\bfseries]
\item[Obligación] El sistema debe rechazar toda transición de estado que no esté autorizada por las reglas de transición vigentes.
\item[Prioridad] Alta
\item[Procedencia] Secciones 2.2, 2.5.1 y 2.7.; ID heredado: \texttt{RF-003}.
\item[Verificación] \texttt{AC-RF-EVT-004}
\end{description}
\vspace{0.25em}\hrule\vspace{0.65em}
\noindent\textbf{RF-EVT-005 --- Cerrar inscripciones}\par
\begin{description}[style=nextline,leftmargin=3.15cm,labelwidth=2.8cm,font=\normalfont\bfseries]
\item[Obligación] El sistema debe cerrar las inscripciones de un evento al ejecutar la transición \texttt{inscripciones -> velocidad}.
\item[Prioridad] Alta
\item[Procedencia] Secciones 1.2, 2.5.1 y RF-011.; ID heredado: \texttt{RF-004}.
\item[Verificación] \texttt{AC-RF-EVT-005}
\end{description}
\vspace{0.25em}\hrule\vspace{0.65em}
\noindent\textbf{RF-EVT-006 --- Precondición para la publicación oficial}\par
\begin{description}[style=nextline,leftmargin=3.15cm,labelwidth=2.8cm,font=\normalfont\bfseries]
\item[Obligación] El sistema debe exigir que se hayan registrado los datos necesarios para determinar y publicar los resultados oficiales antes de que un evento alcance el estado \texttt{publicado}.
\item[Prioridad] Alta
\item[Procedencia] Secciones 2.2, 2.6.1 y 2.7.; ID heredado: \texttt{RF-005}.
\item[Verificación] \texttt{AC-RF-EVT-006}
\end{description}
\vspace{0.25em}\hrule\vspace{0.65em}
\noindent\textbf{RF-EVT-007 --- Publicar resultados oficiales del evento}\par
\begin{description}[style=nextline,leftmargin=3.15cm,labelwidth=2.8cm,font=\normalfont\bfseries]
\item[Obligación] El sistema debe publicar los resultados oficiales de un evento con trazabilidad a los datos y reglas aplicados.
\item[Prioridad] Alta
\item[Procedencia] Secciones 1.2, 2.5, 2.6.1 y 2.7.; ID heredado: \texttt{RF-006}.
\item[Verificación] \texttt{AC-RF-EVT-007}
\end{description}
\vspace{0.25em}\hrule\vspace{0.65em}
\noindent\textbf{RF-EVT-008 --- Inmutabilidad lógica de los resultados publicados}\par
\begin{description}[style=nextline,leftmargin=3.15cm,labelwidth=2.8cm,font=\normalfont\bfseries]
\item[Obligación] El sistema debe impedir modificaciones ordinarias en los resultados oficiales de un evento que ya esté en estado \texttt{publicado}.
\item[Prioridad] Alta
\item[Procedencia] Secciones 2.5 y 2.7.; ID heredado: \texttt{RF-007}.
\item[Verificación] \texttt{AC-RF-EVT-008}
\end{description}
\vspace{0.25em}\hrule\vspace{0.65em}
\noindent\textbf{RF-EVT-009 --- Registro histórico de hitos del evento}\par
\begin{description}[style=nextline,leftmargin=3.15cm,labelwidth=2.8cm,font=\normalfont\bfseries]
\item[Obligación] El sistema debe registrar cronológicamente los hitos críticos del ciclo del evento.
\item[Prioridad] Alta
\item[Procedencia] Secciones 2.4 y 2.7.; ID heredado: \texttt{RF-008}.
\item[Verificación] \texttt{AC-RF-EVT-009}
\end{description}
\vspace{0.25em}\hrule\vspace{0.65em}
\noindent\textbf{RF-EVT-010 --- Estado inicial obligatorio}\par
\begin{description}[style=nextline,leftmargin=3.15cm,labelwidth=2.8cm,font=\normalfont\bfseries]
\item[Obligación] Todo evento nuevo debe comenzar en el estado \texttt{creado}.
\item[Prioridad] Alta
\item[Procedencia] Sección 2.5.1.; ID heredado: \texttt{RF-009}.
\item[Verificación] \texttt{AC-RF-EVT-010}
\end{description}
\vspace{0.25em}\hrule\vspace{0.65em}
\noindent\textbf{RF-EVT-011 --- Transición de apertura de inscripciones}\par
\begin{description}[style=nextline,leftmargin=3.15cm,labelwidth=2.8cm,font=\normalfont\bfseries]
\item[Obligación] El sistema debe permitir la transición \texttt{creado -> inscripciones}.
\item[Prioridad] Alta
\item[Procedencia] Secciones 2.2 y 2.5.1.; ID heredado: \texttt{RF-010}.
\item[Verificación] \texttt{AC-RF-EVT-011}
\end{description}
\vspace{0.25em}\hrule\vspace{0.65em}
\noindent\textbf{RF-EVT-012 --- Transición de inscripciones a velocidad}\par
\begin{description}[style=nextline,leftmargin=3.15cm,labelwidth=2.8cm,font=\normalfont\bfseries]
\item[Obligación] El sistema debe permitir la transición \texttt{inscripciones -> velocidad} cuando se cumplan las guardas normativas aplicables.
\item[Prioridad] Alta
\item[Procedencia] Secciones 2.2 y 2.5.1.; ID heredado: \texttt{RF-011}.
\item[Verificación] \texttt{AC-RF-EVT-012}
\end{description}
\vspace{0.25em}\hrule\vspace{0.65em}
\noindent\textbf{RF-EVT-013 --- Cierre automático de inscripciones al iniciar velocidad}\par
\begin{description}[style=nextline,leftmargin=3.15cm,labelwidth=2.8cm,font=\normalfont\bfseries]
\item[Obligación] Al ejecutar la transición \texttt{inscripciones -> velocidad}, el sistema debe cerrar automáticamente las inscripciones del evento.
\item[Prioridad] Alta
\item[Procedencia] Secciones 2.2 y 2.5.1.; ID heredado: \texttt{RF-011}.
\item[Verificación] \texttt{AC-RF-EVT-013}
\end{description}
\vspace{0.25em}\hrule\vspace{0.65em}
\noindent\textbf{RF-EVT-014 --- Verificación de identidad previa al paso a velocidad}\par
\begin{description}[style=nextline,leftmargin=3.15cm,labelwidth=2.8cm,font=\normalfont\bfseries]
\item[Obligación] El sistema no debe permitir la transición \texttt{inscripciones -> velocidad} mientras exista alguna inscripción no cancelada asociada a un usuario cuya identidad no este verificada.
\item[Prioridad] Alta
\item[Procedencia] RF-011, RF-071A y decisiones prioritarias del cliente sobre paso a \texttt{velocidad}.; ID heredado: \texttt{RF-011A}.
\item[Verificación] \texttt{AC-RF-EVT-014}
\end{description}
\vspace{0.25em}\hrule\vspace{0.65em}
\noindent\textbf{RF-EVT-015 --- Transición a la fase de carrera}\par
\begin{description}[style=nextline,leftmargin=3.15cm,labelwidth=2.8cm,font=\normalfont\bfseries]
\item[Obligación] El sistema debe permitir la transición \texttt{velocidad -> carrera}.
\item[Prioridad] Alta
\item[Procedencia] Secciones 2.2 y 2.5.1.; ID heredado: \texttt{RF-012}.
\item[Verificación] \texttt{AC-RF-EVT-015}
\end{description}
\vspace{0.25em}\hrule\vspace{0.65em}
\noindent\textbf{RF-EVT-016 --- Transición a la fase de subasta}\par
\begin{description}[style=nextline,leftmargin=3.15cm,labelwidth=2.8cm,font=\normalfont\bfseries]
\item[Obligación] El sistema debe permitir la transición \texttt{carrera -> subasta} solo cuando los vehículos sometidos a subasta ya hayan sido entregados al promotor o cuando las inscripciones correspondientes ya hayan quedado canceladas o descalificadas con los efectos del dominio aplicables.
\item[Prioridad] Alta
\item[Procedencia] Secciones 2.2, 2.4, 2.5.1, RF-046 y decisiones prioritarias del cliente sobre entrega previa del vehículo al promotor.; ID heredado: \texttt{RF-013}.
\item[Verificación] \texttt{AC-RF-EVT-016}
\end{description}
\vspace{0.25em}\hrule\vspace{0.65em}
\noindent\textbf{RF-EVT-017 --- Transición a publicado}\par
\begin{description}[style=nextline,leftmargin=3.15cm,labelwidth=2.8cm,font=\normalfont\bfseries]
\item[Obligación] El sistema debe permitir la transición \texttt{subasta -> publicado} únicamente cuando el evento cuente con la información necesaria para determinar y publicar sus resultados oficiales, sus premios, sus pagos y sus devoluciones aplicables.
\item[Prioridad] Alta
\item[Procedencia] Secciones 2.5, 2.6.4, 2.7, RF-005 y RF-006.; ID heredado: \texttt{RF-014}.
\item[Verificación] \texttt{AC-RF-EVT-017}
\end{description}
\vspace{0.25em}\hrule\vspace{0.65em}
\noindent\textbf{RF-EVT-018 --- Registro del estado cancelado}\par
\begin{description}[style=nextline,leftmargin=3.15cm,labelwidth=2.8cm,font=\normalfont\bfseries]
\item[Obligación] Cuando se cancele un evento conforme a las reglas aplicables, el sistema debe registrarlo en estado \texttt{cancelado} y no debe producir una publicación oficial de resultados para ese evento.
\item[Prioridad] Alta
\item[Procedencia] Secciones 2.5.1, 2.6.4 y decisiones prioritarias del cliente sobre devolución inmediata en cancelación del evento.; ID heredado: \texttt{RF-015}.
\item[Verificación] \texttt{AC-RF-EVT-018}
\end{description}
\vspace{0.25em}\hrule\vspace{0.65em}
\noindent\textbf{RF-EVT-019 --- Reversión económica inmediata por cancelación del evento}\par
\begin{description}[style=nextline,leftmargin=3.15cm,labelwidth=2.8cm,font=\normalfont\bfseries]
\item[Obligación] Cuando un evento pase a \texttt{cancelado}, el sistema debe ejecutar las devoluciones y liberaciones económicas exigidas por las reglas del dominio.
\item[Prioridad] Alta
\item[Procedencia] Secciones 2.5.1, 2.6.4 y decisiones prioritarias del cliente sobre devolución inmediata en cancelación del evento.; ID heredado: \texttt{RF-015}.
\item[Verificación] \texttt{AC-RF-EVT-019}
\end{description}
\vspace{0.25em}\hrule\vspace{0.65em}
\noindent\textbf{RF-EVT-020 --- Estado terminal publicado}\par
\begin{description}[style=nextline,leftmargin=3.15cm,labelwidth=2.8cm,font=\normalfont\bfseries]
\item[Obligación] Una vez que un evento alcanza \texttt{publicado}, no debe permitirse ninguna transición a otro estado operativo.
\item[Prioridad] Alta
\item[Procedencia] Secciones 2.5 y 2.7.; ID heredado: \texttt{RF-016}.
\item[Verificación] \texttt{AC-RF-EVT-020}
\end{description}
\vspace{0.25em}\hrule\vspace{0.65em}
\noindent\textbf{RF-EVT-021 --- Estado terminal cancelado}\par
\begin{description}[style=nextline,leftmargin=3.15cm,labelwidth=2.8cm,font=\normalfont\bfseries]
\item[Obligación] Una vez que un evento alcanza \texttt{cancelado}, no debe permitirse ninguna transición a otro estado operativo.
\item[Prioridad] Alta
\item[Procedencia] Secciones 2.5.1 y 2.7.; ID heredado: \texttt{RF-017}.
\item[Verificación] \texttt{AC-RF-EVT-021}
\end{description}
\vspace{0.25em}\hrule\vspace{0.65em}
\noindent\textbf{RF-EVT-022 --- Bloqueo de transiciones no definidas}\par
\begin{description}[style=nextline,leftmargin=3.15cm,labelwidth=2.8cm,font=\normalfont\bfseries]
\item[Obligación] El sistema debe rechazar cualquier transición de estado que no esté contemplada explícitamente en los requerimientos de transición válidos.
\item[Prioridad] Alta
\item[Procedencia] Secciones 2.5.1 y 2.7.; ID heredado: \texttt{RF-018}.
\item[Verificación] \texttt{AC-RF-EVT-022}
\end{description}
\vspace{0.25em}\hrule\vspace{0.65em}
\noindent\textbf{RF-EVT-023 --- Registro de transiciones de estado}\par
\begin{description}[style=nextline,leftmargin=3.15cm,labelwidth=2.8cm,font=\normalfont\bfseries]
\item[Obligación] Toda transición válida o todo intento inválido debe quedar registrado con el estado de origen, el estado de destino, el momento y el actor que ejecutó la acción.
\item[Prioridad] Alta
\item[Procedencia] Secciones 2.4, 2.7 y RF-008.; ID heredado: \texttt{RF-019}.
\item[Verificación] \texttt{AC-RF-EVT-023}
\end{description}
\vspace{0.25em}\hrule\vspace{0.65em}
\noindent\textbf{RF-EVT-024 --- Principio de operaciones habilitadas según el estado}\par
\begin{description}[style=nextline,leftmargin=3.15cm,labelwidth=2.8cm,font=\normalfont\bfseries]
\item[Obligación] El sistema debe habilitar únicamente las operaciones propias del estado actual del evento y bloquear las de los demás estados.
\item[Prioridad] Alta
\item[Procedencia] Secciones 2.5.1 y 2.7.; ID heredado: \texttt{RF-020}.
\item[Verificación] \texttt{AC-RF-EVT-024}
\end{description}
\vspace{0.25em}\hrule\vspace{0.65em}
\noindent\textbf{RF-EVT-025 --- Operaciones en estado inscripciones}\par
\begin{description}[style=nextline,leftmargin=3.15cm,labelwidth=2.8cm,font=\normalfont\bfseries]
\item[Obligación] En estado \texttt{inscripciones}, el sistema debe permitir registrar y editar inscripciones del evento, cobrar de inmediato la inscripción válida, registrar aportaciones directas a la bolsa del evento, registrar aportaciones a participantes del evento solo cuando esa función se encuentre activa y operar la subasta de acceso conforme a sus reglas económicas.
\item[Prioridad] Alta
\item[Procedencia] Reglas operativas del ciclo del evento (2.2, 2.5.1), RF-042 a RF-045B, RF-086A y decisiones prioritarias del cliente sobre cobro inmediato.; ID heredado: \texttt{RF-021}.
\item[Verificación] \texttt{AC-RF-EVT-025}
\end{description}
\vspace{0.25em}\hrule\vspace{0.65em}
\noindent\textbf{RF-EVT-026 --- Operaciones en estado velocidad}\par
\begin{description}[style=nextline,leftmargin=3.15cm,labelwidth=2.8cm,font=\normalfont\bfseries]
\item[Obligación] En estado \texttt{velocidad}, el sistema debe permitir capturar y editar resultados de velocidad del evento y ejecutar las decisiones previas a \texttt{carrera} que las reglas del dominio aún permitan.
\item[Prioridad] Alta
\item[Procedencia] Reglas operativas del ciclo del evento (2.2, 2.5.1) y decisiones prioritarias del cliente sobre vigencia hasta antes de \texttt{carrera}.; ID heredado: \texttt{RF-022}.
\item[Verificación] \texttt{AC-RF-EVT-026}
\end{description}
\vspace{0.25em}\hrule\vspace{0.65em}
\noindent\textbf{RF-EVT-027 --- Ventana de inspección comunitaria y protestas antes de carrera}\par
\begin{description}[style=nextline,leftmargin=3.15cm,labelwidth=2.8cm,font=\normalfont\bfseries]
\item[Obligación] El sistema debe permitir la inspección comunitaria de vehículos, el registro de objeciones de seguridad, las votaciones formales y las decisiones finales del promotor únicamente mientras el evento está en \texttt{inscripciones} o \texttt{velocidad}.
\item[Prioridad] Alta
\item[Procedencia] RF-057 a RF-061 y decisiones prioritarias del cliente sobre vigencia hasta antes de \texttt{carrera}.; ID heredado: \texttt{RF-022A}.
\item[Verificación] \texttt{AC-RF-EVT-027}
\end{description}
\vspace{0.25em}\hrule\vspace{0.65em}
\noindent\textbf{RF-EVT-028 --- Operaciones en estado carrera}\par
\begin{description}[style=nextline,leftmargin=3.15cm,labelwidth=2.8cm,font=\normalfont\bfseries]
\item[Obligación] En estado \texttt{carrera}, el sistema debe permitir capturar y editar resultados de carrera del evento y registrar la recepción por el promotor de los vehículos que deban pasar a subasta.
\item[Prioridad] Alta
\item[Procedencia] Reglas operativas del ciclo del evento (2.2, 2.4, 2.5.1) y decisiones prioritarias del cliente sobre entrega del vehículo al promotor al terminar la carrera.; ID heredado: \texttt{RF-023}.
\item[Verificación] \texttt{AC-RF-EVT-028}
\end{description}
\vspace{0.25em}\hrule\vspace{0.65em}
\noindent\textbf{RF-EVT-029 --- Operaciones en estado subasta}\par
\begin{description}[style=nextline,leftmargin=3.15cm,labelwidth=2.8cm,font=\normalfont\bfseries]
\item[Obligación] En estado \texttt{subasta}, el sistema debe permitir operar las subastas de vehículos, ejecutar los cobros y pagos asociados, registrar confirmaciones de entrega y tramitar disputas de entrega conforme a las reglas del evento.
\item[Prioridad] Alta
\item[Procedencia] Reglas operativas del ciclo del evento (2.2, 2.4, 2.5.1) y RF-046 a RF-050J.; ID heredado: \texttt{RF-024}.
\item[Verificación] \texttt{AC-RF-EVT-029}
\end{description}
\vspace{0.25em}\hrule\vspace{0.65em}
\noindent\textbf{RF-EVT-030 --- Operaciones en estado creado}\par
\begin{description}[style=nextline,leftmargin=3.15cm,labelwidth=2.8cm,font=\normalfont\bfseries]
\item[Obligación] En estado \texttt{creado}, el sistema debe permitir capturar y completar la información pública del evento requerida para publicar el evento en \texttt{inscripciones}.
\item[Prioridad] Media
\item[Procedencia] Secciones 2.2, 2.5.1 y 2.6.1.; ID heredado: \texttt{RF-025}.
\item[Verificación] \texttt{AC-RF-EVT-030}
\end{description}
\vspace{0.25em}\hrule\vspace{0.65em}
\noindent\textbf{RF-EVT-031 --- Consulta de resultados oficiales en estado publicado}\par
\begin{description}[style=nextline,leftmargin=3.15cm,labelwidth=2.8cm,font=\normalfont\bfseries]
\item[Obligación] En estado \texttt{publicado}, el sistema debe permitir consultar los resultados oficiales del evento.
\item[Prioridad] Alta
\item[Procedencia] Secciones 2.5, 2.7 y RF-007.; ID heredado: \texttt{RF-026}.
\item[Verificación] \texttt{AC-RF-EVT-031}
\end{description}
\vspace{0.25em}\hrule\vspace{0.65em}
\noindent\textbf{RF-EVT-032 --- Bloqueo de nuevas capturas en estado publicado}\par
\begin{description}[style=nextline,leftmargin=3.15cm,labelwidth=2.8cm,font=\normalfont\bfseries]
\item[Obligación] En estado \texttt{publicado}, el sistema debe rechazar nuevas capturas operativas del evento.
\item[Prioridad] Alta
\item[Procedencia] Secciones 2.5, 2.7 y RF-007.; ID heredado: \texttt{RF-026}.
\item[Verificación] \texttt{AC-RF-EVT-032}
\end{description}
\vspace{0.25em}\hrule\vspace{0.65em}
\noindent\textbf{RF-EVT-033 --- Operaciones en estado cancelado}\par
\begin{description}[style=nextline,leftmargin=3.15cm,labelwidth=2.8cm,font=\normalfont\bfseries]
\item[Obligación] En estado \texttt{cancelado}, el sistema no debe permitir operaciones del evento distintas de las consultas e historiales que sigan visibles conforme al dominio.
\item[Prioridad] Alta
\item[Procedencia] Secciones 2.5.1 y 2.7.; ID heredado: \texttt{RF-027}.
\item[Verificación] \texttt{AC-RF-EVT-033}
\end{description}
\vspace{0.25em}\hrule\vspace{0.65em}
\noindent\textbf{RF-EVT-034 --- Rechazo y mensaje ante operación inválida por estado}\par
\begin{description}[style=nextline,leftmargin=3.15cm,labelwidth=2.8cm,font=\normalfont\bfseries]
\item[Obligación] Si se intenta ejecutar una operación inválida para el estado actual del evento, el sistema debe rechazarla e informar explícitamente cuál es el estado vigente y cuál es la causa del rechazo.
\item[Prioridad] Alta
\item[Procedencia] Sección 2.7 y RF-020.; ID heredado: \texttt{RF-028}.
\item[Verificación] \texttt{AC-RF-EVT-034}
\end{description}
\vspace{0.25em}\hrule\vspace{0.65em}
\noindent\textbf{RF-EVT-035 --- Ventana para aportar a evento o participante del evento}\par
\begin{description}[style=nextline,leftmargin=3.15cm,labelwidth=2.8cm,font=\normalfont\bfseries]
\item[Obligación] El sistema debe permitir aportaciones directas a la bolsa del evento y aportaciones a participantes del evento únicamente mientras el evento está en estado \texttt{inscripciones}.
\item[Prioridad] Alta
\item[Procedencia] RF-021, sección 2.6.2 y decisiones prioritarias del cliente sobre cierre al pasar a \texttt{velocidad}.; ID heredado: \texttt{RF-029}.
\item[Verificación] \texttt{AC-RF-EVT-035}
\end{description}
\vspace{0.25em}\hrule\vspace{0.65em}
\noindent\textbf{RF-EVT-036 --- Cierre de aportaciones del evento al entrar en velocidad}\par
\begin{description}[style=nextline,leftmargin=3.15cm,labelwidth=2.8cm,font=\normalfont\bfseries]
\item[Obligación] Al pasar el evento a estado \texttt{velocidad}, el sistema debe cerrar de manera inmediata e irreversible la recepción de nuevas aportaciones directas a la bolsa del evento y de nuevas aportaciones a participantes del evento.
\item[Prioridad] Alta
\item[Procedencia] Integridad del proceso económico previo a la actividad competitiva (2.5.1, 2.7) y decisiones prioritarias del cliente.; ID heredado: \texttt{RF-030}.
\item[Verificación] \texttt{AC-RF-EVT-036}
\end{description}
\vspace{0.25em}\hrule\vspace{0.65em}
\noindent\textbf{RF-EVT-037 --- Integracion de aportaciones acumuladas por participante}\par
\begin{description}[style=nextline,leftmargin=3.15cm,labelwidth=2.8cm,font=\normalfont\bfseries]
\item[Obligación] Cada nueva aportación aceptada y cobrada debe sumarse a la contribución total del participante destinatario dentro del evento.
\item[Prioridad] Alta
\item[Procedencia] Reglas de pozo y reparto (1.2, 2.2, 2.6.2).; ID heredado: \texttt{RF-031}.
\item[Verificación] \texttt{AC-RF-EVT-037}
\end{description}
\vspace{0.25em}\hrule\vspace{0.65em}
\noindent\textbf{RF-EVT-038 --- Bloqueo y mensaje por aportación fuera de ventana}\par
\begin{description}[style=nextline,leftmargin=3.15cm,labelwidth=2.8cm,font=\normalfont\bfseries]
\item[Obligación] Cuando se intente anadir una aportación fuera de la ventana permitida, el sistema debe rechazarla e indicar el estado actual del evento y la causa del rechazo.
\item[Prioridad] Media
\item[Procedencia] Sección 2.7.; ID heredado: \texttt{RF-032}.
\item[Verificación] \texttt{AC-RF-EVT-038}
\end{description}
\vspace{0.25em}\hrule\vspace{0.65em}
\noindent\textbf{RF-EVT-039 --- Bloqueo preventivo de acciones no permitidas en la IU}\par
\begin{description}[style=nextline,leftmargin=3.15cm,labelwidth=2.8cm,font=\normalfont\bfseries]
\item[Obligación] La interfaz de usuario debe presentar únicamente las acciones permitidas para el estado actual del evento y ocultar o deshabilitar las acciones no permitidas.
\item[Prioridad] Alta
\item[Procedencia] Sección 2.7 y RF-020.; ID heredado: \texttt{RF-033}.
\item[Verificación] \texttt{AC-RF-EVT-039}
\end{description}
\vspace{0.25em}\hrule\vspace{0.65em}
\noindent\textbf{RF-EVT-040 --- Impedir la solicitud de acciones no permitidas}\par
\begin{description}[style=nextline,leftmargin=3.15cm,labelwidth=2.8cm,font=\normalfont\bfseries]
\item[Obligación] El sistema debe impedir que los usuarios soliciten desde la IU acciones no permitidas.
\item[Prioridad] Alta
\item[Procedencia] Sección 2.7 y RF-020.; ID heredado: \texttt{RF-034}.
\item[Verificación] \texttt{AC-RF-EVT-040}
\end{description}
\vspace{0.25em}\hrule\vspace{0.65em}
\noindent\textbf{RF-EVT-041 --- Registro de datos mínimos de convocatoria del evento}\par
\begin{description}[style=nextline,leftmargin=3.15cm,labelwidth=2.8cm,font=\normalfont\bfseries]
\item[Obligación] El sistema debe permitir registrar, como mínimo, la fecha del evento, su ubicación, su nivel territorial, su nivel temporal, el tipo de pista, el número esperado de participantes, la cuota de inscripción, el pozo inicial, el pozo del premio de terminación, el costo de renta de pista cuando corresponda, el costo fijo por pesaje y, cuando el evento contemple subasta de vehículos, el plazo máximo para iniciar disputas de entrega de vehículos adjudicados y el plazo máximo para resolverlas.
\item[Prioridad] Alta
\item[Procedencia] Secciones 2.4, 2.6.1, RIE-002, fuente documental \texttt{TOP\_Racing\_Grok\_Clarified.tex} y decisiones prioritarias del cliente sobre nivel temporal y plazos de disputa de entrega.; ID heredado: \texttt{RF-075}.
\item[Verificación] \texttt{AC-RF-EVT-041}
\end{description}
\vspace{0.25em}\hrule\vspace{0.65em}
\noindent\textbf{RF-EVT-042 --- Conservación histórica del número esperado de participantes}\par
\begin{description}[style=nextline,leftmargin=3.15cm,labelwidth=2.8cm,font=\normalfont\bfseries]
\item[Obligación] El sistema debe conservar el número esperado de participantes publicado en la convocatoria como dato histórico de esa convocatoria, sin ajustarlo por cancelaciones posteriores.
\item[Prioridad] Media
\item[Procedencia] RF-075 y fuente prioritaria confirmada por el usuario sobre conservación histórica del número esperado de participantes.; ID heredado: \texttt{RF-075A}.
\item[Verificación] \texttt{AC-RF-EVT-042}
\end{description}
\vspace{0.25em}\hrule\vspace{0.65em}
\noindent\textbf{RF-EVT-043 --- Verificación reforzada obligatoria para promover eventos}\par
\begin{description}[style=nextline,leftmargin=3.15cm,labelwidth=2.8cm,font=\normalfont\bfseries]
\item[Obligación] El sistema debe permitir registrar y publicar eventos únicamente a usuarios con verificación reforzada activa, además de las demás condiciones generales de participación operativa.
\item[Prioridad] Alta
\item[Procedencia] Secciones 2.3, 2.5, 2.6, 3.1.24 y decisiones prioritarias del cliente sobre promoción de eventos con verificación reforzada.; ID heredado: \texttt{RF-075B}.
\item[Verificación] \texttt{AC-RF-EVT-043}
\end{description}
\vspace{0.25em}\hrule\vspace{0.65em}
\noindent\textbf{RF-EVT-044 --- Publicación de la convocatoria del evento en inscripciones}\par
\begin{description}[style=nextline,leftmargin=3.15cm,labelwidth=2.8cm,font=\normalfont\bfseries]
\item[Obligación] El sistema debe publicar la convocatoria del evento al entrar en estado \texttt{inscripciones}, junto con sus condiciones y reglas públicas.
\item[Prioridad] Alta
\item[Procedencia] Secciones 2.5, 2.6.1 y decisiones prioritarias del cliente sobre convocatoria en \texttt{inscripciones}.; ID heredado: \texttt{RF-076}.
\item[Verificación] \texttt{AC-RF-EVT-044}
\end{description}
\vspace{0.25em}\hrule\vspace{0.65em}
\noindent\textbf{RF-EVT-045 --- Establecimiento obligatorio de la carga mínima del evento}\par
\begin{description}[style=nextline,leftmargin=3.15cm,labelwidth=2.8cm,font=\normalfont\bfseries]
\item[Obligación] El sistema debe exigir que el evento tenga establecida una carga mínima aplicable antes del inicio de la actividad competitiva.
\item[Prioridad] Alta
\item[Procedencia] Secciones 2.6.1 y fuente documental \texttt{TOP\_Racing\_Grok\_Clarified.tex} (línea 126).; ID heredado: \texttt{RF-077}.
\item[Verificación] \texttt{AC-RF-EVT-045}
\end{description}
\vspace{0.25em}\hrule\vspace{0.65em}
\noindent\textbf{RF-EVT-046 --- Determinación de la clasificación combinada aplicable a la salida invertida}\par
\begin{description}[style=nextline,leftmargin=3.15cm,labelwidth=2.8cm,font=\normalfont\bfseries]
\item[Obligación] El sistema debe determinar, a partir del nivel territorial y del nivel temporal definidos para el evento, la clasificación combinada aplicable a la salida invertida y conservarla como parte de la convocatoria del evento.
\item[Prioridad] Media
\item[Procedencia] Secciones 2.6.1, fuente documental \texttt{TOP\_Racing\_Grok\_Clarified.tex} y decisiones prioritarias del cliente sobre determinación de la clasificación combinada por nivel territorial y temporal.; ID heredado: \texttt{RF-078}.
\item[Verificación] \texttt{AC-RF-EVT-046}
\end{description}
\vspace{0.25em}\hrule\vspace{0.65em}
\noindent\textbf{RF-EVT-047 --- Establecimiento obligatorio de la cantidad disponible de accesos para espectadores}\par
\begin{description}[style=nextline,leftmargin=3.15cm,labelwidth=2.8cm,font=\normalfont\bfseries]
\item[Obligación] Cuando el evento contemple subasta de acceso para espectadores, el sistema debe exigir que quede establecida la cantidad disponible de accesos como parte de la convocatoria del evento.
\item[Prioridad] Alta
\item[Procedencia] Secciones 2.6.1, RF-042 y fuente documental \texttt{TOP\_Racing\_Grok\_Clarified.tex} (línea 219).; ID heredado: \texttt{RF-079}.
\item[Verificación] \texttt{AC-RF-EVT-047}
\end{description}
\vspace{0.25em}\hrule\vspace{0.65em}
\noindent\textbf{RF-EVT-048 --- Establecimiento obligatorio del plazo máximo para iniciar disputas de entrega de vehículos}\par
\begin{description}[style=nextline,leftmargin=3.15cm,labelwidth=2.8cm,font=\normalfont\bfseries]
\item[Obligación] Cuando el evento contemple subasta de vehículos, el sistema debe exigir que quede establecido y publicado el plazo máximo para iniciar disputas de entrega de vehículos adjudicados.
\item[Prioridad] Alta
\item[Procedencia] RF-075 y decisiones prioritarias del cliente sobre plazo máximo para iniciar disputas como dato del evento.; ID heredado: \texttt{RF-079A}.
\item[Verificación] \texttt{AC-RF-EVT-048}
\end{description}
\vspace{0.25em}\hrule\vspace{0.65em}
\noindent\textbf{RF-EVT-049 --- Establecimiento obligatorio del plazo máximo para resolver disputas de entrega de vehículos}\par
\begin{description}[style=nextline,leftmargin=3.15cm,labelwidth=2.8cm,font=\normalfont\bfseries]
\item[Obligación] Cuando el evento contemple subasta de vehículos, el sistema debe exigir que quede establecido y publicado el plazo máximo para resolver disputas de entrega de vehículos adjudicados.
\item[Prioridad] Alta
\item[Procedencia] RF-075 y decisiones prioritarias del cliente sobre plazo máximo para resolver disputas como dato del evento.; ID heredado: \texttt{RF-079B}.
\item[Verificación] \texttt{AC-RF-EVT-049}
\end{description}
\vspace{0.25em}\hrule\vspace{0.65em}
\noindent\textbf{RF-EVT-050 --- Publicación de la condición de entrega previa del vehículo al promotor}\par
\begin{description}[style=nextline,leftmargin=3.15cm,labelwidth=2.8cm,font=\normalfont\bfseries]
\item[Obligación] Cuando el evento contemple subasta de vehículos, el sistema debe publicar como condición del evento que la subasta solo puede iniciarse respecto de vehículos previamente entregados al promotor.
\item[Prioridad] Alta
\item[Procedencia] RF-046 y decisiones prioritarias del cliente sobre entrega previa del vehículo al promotor.; ID heredado: \texttt{RF-079C}.
\item[Verificación] \texttt{AC-RF-EVT-050}
\end{description}
\vspace{0.25em}\hrule\vspace{0.65em}
\noindent\textbf{RF-EVT-051 --- Configuración obligatoria del límite de carrera}\par
\begin{description}[style=nextline,leftmargin=3.15cm,labelwidth=2.8cm,font=\normalfont\bfseries]
\item[Obligación] Antes de abrir inscripciones, el sistema debe exigir que el evento declare la carrera como limitada por \texttt{tiempo} o por \texttt{distancia} y publique el valor positivo del límite correspondiente; para límite por distancia también debe estar publicada la longitud oficial positiva de una vuelta.
\item[Prioridad] Alta
\item[Procedencia] \texttt{DEC-155}, decisión confirmada sobre carreras limitadas por tiempo o distancia.
\item[Verificación] \texttt{AC-RF-EVT-051}
\end{description}
\vspace{0.25em}\hrule\vspace{0.65em}
\subsection{Consolidación y liquidación económica}
\noindent\textbf{RF-PAY-001 --- Consolidación del pozo económico oficial}\par
\begin{description}[style=nextline,leftmargin=3.15cm,labelwidth=2.8cm,font=\normalfont\bfseries]
\item[Obligación] El sistema debe consolidar el pozo económico oficial del evento a partir del pozo inicial publicado y de las cantidades efectivamente cobradas y computables dentro de la ventana permitida.
\item[Prioridad] Alta
\item[Procedencia] Secciones 2.6.2, 2.6.3 y 2.6.4.; ID heredado: \texttt{RF-035}.
\item[Verificación] \texttt{AC-RF-PAY-001}
\end{description}
\vspace{0.25em}\hrule\vspace{0.65em}
\noindent\textbf{RF-PAY-002 --- Conservación del registro económico oficial}\par
\begin{description}[style=nextline,leftmargin=3.15cm,labelwidth=2.8cm,font=\normalfont\bfseries]
\item[Obligación] El sistema debe conservar el detalle del pozo inicial publicado, el detalle de los cobros y retenciones aceptadas, los pagos, las devoluciones y el total oficial consolidado como parte del registro económico oficial del evento.
\item[Prioridad] Alta
\item[Procedencia] Secciones 2.6.2, 2.6.3 y 2.7.; ID heredado: \texttt{RF-036}.
\item[Verificación] \texttt{AC-RF-PAY-002}
\end{description}
\vspace{0.25em}\hrule\vspace{0.65em}
\noindent\textbf{RF-PAY-003 --- Determinación del premio individual del evento}\par
\begin{description}[style=nextline,leftmargin=3.15cm,labelwidth=2.8cm,font=\normalfont\bfseries]
\item[Obligación] El sistema debe determinar, para cada participante computable del evento, el premio individual que le corresponda a partir de su desempeño, de las reglas publicadas y de los datos oficiales del evento.
\item[Prioridad] Alta
\item[Procedencia] Sección 2.6.4.; ID heredado: \texttt{RF-037}.
\item[Verificación] \texttt{AC-RF-PAY-003}
\end{description}
\vspace{0.25em}\hrule\vspace{0.65em}
\noindent\textbf{RF-PAY-004 --- Registro trazable del premio individual del evento}\par
\begin{description}[style=nextline,leftmargin=3.15cm,labelwidth=2.8cm,font=\normalfont\bfseries]
\item[Obligación] El sistema debe registrar de manera trazable el premio individual del evento correspondiente a cada competidor computable.
\item[Prioridad] Alta
\item[Procedencia] Secciones 2.6.4 y 2.7.; ID heredado: \texttt{RF-038}.
\item[Verificación] \texttt{AC-RF-PAY-004}
\end{description}
\vspace{0.25em}\hrule\vspace{0.65em}
\noindent\textbf{RF-PAY-005 --- Distribución económica posterior del premio individual del evento}\par
\begin{description}[style=nextline,leftmargin=3.15cm,labelwidth=2.8cm,font=\normalfont\bfseries]
\item[Obligación] El sistema debe distribuir el premio individual del evento de cada competidor computable entre ese competidor y los aportantes asociados a sus contribuciones cobradas.
\item[Prioridad] Alta
\item[Procedencia] Sección 2.6.4 y fuente prioritaria confirmada por el usuario sobre contribuciones de terceros.; ID heredado: \texttt{RF-039}.
\item[Verificación] \texttt{AC-RF-PAY-005}
\end{description}
\vspace{0.25em}\hrule\vspace{0.65em}
\noindent\textbf{RF-PAY-006 --- Liquidación económica oficial del evento}\par
\begin{description}[style=nextline,leftmargin=3.15cm,labelwidth=2.8cm,font=\normalfont\bfseries]
\item[Obligación] El sistema debe producir la liquidación económica oficial del evento a partir del pozo económico oficial, de los resultados del dominio que correspondan, de la distribución económica posterior aplicable y de las devoluciones que deban ejecutarse, excluyendo toda inscripción cancelada o descalificada.
\item[Prioridad] Alta
\item[Procedencia] Secciones 2.6.3 y 2.6.4.; ID heredado: \texttt{RF-040}.
\item[Verificación] \texttt{AC-RF-PAY-006}
\end{description}
\vspace{0.25em}\hrule\vspace{0.65em}
\noindent\textbf{RF-PAY-007 --- Devolución inmediata por cancelación o descalificación previa a carrera}\par
\begin{description}[style=nextline,leftmargin=3.15cm,labelwidth=2.8cm,font=\normalfont\bfseries]
\item[Obligación] El sistema debe devolver inmediatamente las cantidades cobradas que correspondan cuando el evento se cancele o cuando una inscripción quede cancelada o descalificada antes de \texttt{carrera}, incluidas las aportaciones de terceros asociadas a esa inscripción dentro del evento.
\item[Prioridad] Alta
\item[Procedencia] Secciones 2.6.4 y decisiones prioritarias del cliente sobre devolución inmediata.; ID heredado: \texttt{RF-040A}.
\item[Verificación] \texttt{AC-RF-PAY-007}
\end{description}
\vspace{0.25em}\hrule\vspace{0.65em}
\noindent\textbf{RF-PAY-008 --- Pago automático e inmediato a todos los beneficiarios del evento}\par
\begin{description}[style=nextline,leftmargin=3.15cm,labelwidth=2.8cm,font=\normalfont\bfseries]
\item[Obligación] Una vez determinada la liquidación económica oficial del evento, el sistema debe pagar de inmediato la parte correspondiente a cada beneficiario de la distribución económica posterior del evento.
\item[Prioridad] Alta
\item[Procedencia] Sección 2.6.4 y decisiones prioritarias del cliente sobre pago automático a todos los beneficiarios.; ID heredado: \texttt{RF-040B}.
\item[Verificación] \texttt{AC-RF-PAY-008}
\end{description}
\vspace{0.25em}\hrule\vspace{0.65em}
\noindent\textbf{RF-PAY-009 --- Publicación de resultados económicos oficiales}\par
\begin{description}[style=nextline,leftmargin=3.15cm,labelwidth=2.8cm,font=\normalfont\bfseries]
\item[Obligación] El sistema debe publicar los premios individuales oficiales del evento junto con la distribución económica posterior correspondiente a cada competidor computable, mostrando sin identidad las cantidades y la parte de premio correspondientes a terceros cuando existan.
\item[Prioridad] Alta
\item[Procedencia] Secciones 2.6.3, 2.6.4 y 2.7.; ID heredado: \texttt{RF-041}.
\item[Verificación] \texttt{AC-RF-PAY-009}
\end{description}
\vspace{0.25em}\hrule\vspace{0.65em}
\subsection{Subastas y pujas}
\noindent\textbf{RF-AUC-001 --- Apertura de subasta de acceso de espectadores al evento}\par
\begin{description}[style=nextline,leftmargin=3.15cm,labelwidth=2.8cm,font=\normalfont\bfseries]
\item[Obligación] Cuando las reglas publicadas del evento contemplen subasta de acceso para espectadores, el sistema debe registrar la cantidad disponible de accesos y habilitar la recepción de pujas secretas para esa subasta.
\item[Prioridad] Alta
\item[Procedencia] Secciones 2.2, 2.4, 2.6.1 y fuente documental \texttt{TOP\_Racing\_Grok\_Clarified.tex} (líneas 146 y 219).; ID heredado: \texttt{RF-042}.
\item[Verificación] \texttt{AC-RF-AUC-001}
\end{description}
\vspace{0.25em}\hrule\vspace{0.65em}
\noindent\textbf{RF-AUC-002 --- Registro de pujas secretas en subasta de acceso con retención suficiente}\par
\begin{description}[style=nextline,leftmargin=3.15cm,labelwidth=2.8cm,font=\normalfont\bfseries]
\item[Obligación] El sistema debe registrar cada puja secreta de la subasta de acceso solo cuando quede respaldada por una retención económica suficiente para el monto pujado, sin revelar su monto individual mientras la subasta permanezca abierta.
\item[Prioridad] Alta
\item[Procedencia] Secciones 2.4, 2.6.4 y decisiones prioritarias del cliente sobre pagos reales y retención en subasta de acceso.; ID heredado: \texttt{RF-043}.
\item[Verificación] \texttt{AC-RF-AUC-002}
\end{description}
\vspace{0.25em}\hrule\vspace{0.65em}
\noindent\textbf{RF-AUC-003 --- Una sola puja vigente por usuario en subasta de acceso}\par
\begin{description}[style=nextline,leftmargin=3.15cm,labelwidth=2.8cm,font=\normalfont\bfseries]
\item[Obligación] El sistema debe mantener una sola puja vigente por usuario dentro de cada subasta de acceso de espectadores.
\item[Prioridad] Alta
\item[Procedencia] RF-043 y decisiones prioritarias del cliente sobre puja vigente única en subasta de acceso.; ID heredado: \texttt{RF-043A}.
\item[Verificación] \texttt{AC-RF-AUC-003}
\end{description}
\vspace{0.25em}\hrule\vspace{0.65em}
\noindent\textbf{RF-AUC-004 --- Sustitución solo al alza de la puja vigente de acceso}\par
\begin{description}[style=nextline,leftmargin=3.15cm,labelwidth=2.8cm,font=\normalfont\bfseries]
\item[Obligación] Mientras la subasta de acceso permanezca abierta, el sistema debe aceptar la sustitución de una puja vigente únicamente por otra de importe superior.
\item[Prioridad] Alta
\item[Procedencia] RF-043A y decisiones prioritarias del cliente sobre actualización solo al alza en subasta de acceso.; ID heredado: \texttt{RF-043B}.
\item[Verificación] \texttt{AC-RF-AUC-004}
\end{description}
\vspace{0.25em}\hrule\vspace{0.65em}
\noindent\textbf{RF-AUC-005 --- Ajuste de retención por actualización válida de la puja de acceso}\par
\begin{description}[style=nextline,leftmargin=3.15cm,labelwidth=2.8cm,font=\normalfont\bfseries]
\item[Obligación] Cuando una actualización válida reemplace la puja vigente de acceso, el sistema debe ajustar la retención económica asociada al nuevo importe.
\item[Prioridad] Alta
\item[Procedencia] RF-043A y decisiones prioritarias del cliente sobre actualización solo al alza en subasta de acceso.; ID heredado: \texttt{RF-043B}.
\item[Verificación] \texttt{AC-RF-AUC-005}
\end{description}
\vspace{0.25em}\hrule\vspace{0.65em}
\noindent\textbf{RF-AUC-006 --- Publicación del conteo de pujas vigentes en subasta de acceso}\par
\begin{description}[style=nextline,leftmargin=3.15cm,labelwidth=2.8cm,font=\normalfont\bfseries]
\item[Obligación] Mientras la subasta de acceso permanezca abierta, el sistema debe mostrar el número de pujas vigentes sin revelar sus montos individuales.
\item[Prioridad] Alta
\item[Procedencia] Secciones 2.4, 2.7 y fuente documental \texttt{TOP\_Racing\_Grok\_Clarified.tex} (línea 219).; ID heredado: \texttt{RF-044}.
\item[Verificación] \texttt{AC-RF-AUC-006}
\end{description}
\vspace{0.25em}\hrule\vspace{0.65em}
\noindent\textbf{RF-AUC-007 --- Ordenamiento de pujas válidas de acceso}\par
\begin{description}[style=nextline,leftmargin=3.15cm,labelwidth=2.8cm,font=\normalfont\bfseries]
\item[Obligación] Al cerrar la subasta de acceso, el sistema debe ordenar las pujas válidas de mayor a menor.
\item[Prioridad] Alta
\item[Procedencia] Secciones 2.4, 2.7 y fuente documental \texttt{TOP\_Racing\_Grok\_Clarified.tex} (línea 219).; ID heredado: \texttt{RF-045}.
\item[Verificación] \texttt{AC-RF-AUC-007}
\end{description}
\vspace{0.25em}\hrule\vspace{0.65em}
\noindent\textbf{RF-AUC-008 --- Adjudicación de accesos disponibles}\par
\begin{description}[style=nextline,leftmargin=3.15cm,labelwidth=2.8cm,font=\normalfont\bfseries]
\item[Obligación] Al cerrar la subasta de acceso, el sistema debe adjudicar los accesos disponibles a las mejores pujas válidas hasta agotar la cantidad publicada.
\item[Prioridad] Alta
\item[Procedencia] Secciones 2.4, 2.7 y fuente documental \texttt{TOP\_Racing\_Grok\_Clarified.tex} (línea 219).; ID heredado: \texttt{RF-045}.
\item[Verificación] \texttt{AC-RF-AUC-008}
\end{description}
\vspace{0.25em}\hrule\vspace{0.65em}
\noindent\textbf{RF-AUC-009 --- Determinación del importe de cobro del acceso adjudicado}\par
\begin{description}[style=nextline,leftmargin=3.15cm,labelwidth=2.8cm,font=\normalfont\bfseries]
\item[Obligación] Para cada acceso adjudicado, el sistema debe determinar como importe de cobro la siguiente puja válida inferior a la puja ganadora correspondiente.
\item[Prioridad] Alta
\item[Procedencia] Secciones 2.4, 2.7 y fuente documental \texttt{TOP\_Racing\_Grok\_Clarified.tex} (línea 219).; ID heredado: \texttt{RF-045}.
\item[Verificación] \texttt{AC-RF-AUC-009}
\end{description}
\vspace{0.25em}\hrule\vspace{0.65em}
\noindent\textbf{RF-AUC-010 --- Publicación del resultado de la subasta de acceso}\par
\begin{description}[style=nextline,leftmargin=3.15cm,labelwidth=2.8cm,font=\normalfont\bfseries]
\item[Obligación] Al cerrar la subasta de acceso, el sistema debe publicar la persona ganadora y el importe de cobro correspondiente a cada acceso adjudicado.
\item[Prioridad] Alta
\item[Procedencia] Secciones 2.4, 2.7 y fuente documental \texttt{TOP\_Racing\_Grok\_Clarified.tex} (línea 219).; ID heredado: \texttt{RF-045}.
\item[Verificación] \texttt{AC-RF-AUC-010}
\end{description}
\vspace{0.25em}\hrule\vspace{0.65em}
\noindent\textbf{RF-AUC-011 --- Notificación de acceso adjudicado con QR}\par
\begin{description}[style=nextline,leftmargin=3.15cm,labelwidth=2.8cm,font=\normalfont\bfseries]
\item[Obligación] Al cerrarse la subasta de acceso, el sistema debe poner a disposición de cada persona ganadora una notificación de acceso adjudicado con QR en la que consten su acceso adjudicado y el importe correspondiente.
\item[Prioridad] Alta
\item[Procedencia] RF-045, fuente documental \texttt{TOP\_Racing\_Grok\_Clarified.tex} (líneas 219 y 223) y decisiones prioritarias del cliente sobre QR en lugar de ticket.; ID heredado: \texttt{RF-045A}.
\item[Verificación] \texttt{AC-RF-AUC-011}
\end{description}
\vspace{0.25em}\hrule\vspace{0.65em}
\noindent\textbf{RF-AUC-012 --- Retencion, cobro y liberación en subasta de acceso}\par
\begin{description}[style=nextline,leftmargin=3.15cm,labelwidth=2.8cm,font=\normalfont\bfseries]
\item[Obligación] El sistema debe mantener retenido el respaldo económico suficiente de cada puja vigente de acceso, cobrar automáticamente el importe adjudicado a cada persona ganadora al cierre de la subasta y liberar de inmediato las retenciones de las pujas no ganadoras.
\item[Prioridad] Alta
\item[Procedencia] Secciones 2.6.4 y decisiones prioritarias del cliente sobre pagos reales en subasta de acceso.; ID heredado: \texttt{RF-045B}.
\item[Verificación] \texttt{AC-RF-AUC-012}
\end{description}
\vspace{0.25em}\hrule\vspace{0.65em}
\noindent\textbf{RF-AUC-013 --- Devolución inmediata de accesos adjudicados por cancelación del evento}\par
\begin{description}[style=nextline,leftmargin=3.15cm,labelwidth=2.8cm,font=\normalfont\bfseries]
\item[Obligación] Si el evento se cancela después del cobro de accesos adjudicados, el sistema debe devolver de inmediato los importes cobrados por esos accesos y liberar cualquier retención que permanezca vigente.
\item[Prioridad] Alta
\item[Procedencia] RF-045B y decisiones prioritarias del cliente sobre devolución inmediata por cancelación del evento.; ID heredado: \texttt{RF-045C}.
\item[Verificación] \texttt{AC-RF-AUC-013}
\end{description}
\vspace{0.25em}\hrule\vspace{0.65em}
\noindent\textbf{RF-AUC-014 --- Apertura de subastas de vehículos después de la carrera}\par
\begin{description}[style=nextline,leftmargin=3.15cm,labelwidth=2.8cm,font=\normalfont\bfseries]
\item[Obligación] En la fase de subasta del evento, el sistema debe habilitar la subasta solo de los vehículos asociados a inscripciones computables que hayan participado en el evento y que ya hayan sido entregados al promotor.
\item[Prioridad] Alta
\item[Procedencia] Secciones 2.2, 2.4, 2.5.1, 2.7, fuente documental \texttt{TOP\_Racing\_Grok\_Clarified.tex} (líneas 146, 221 y 223) y decisiones prioritarias del cliente sobre entrega del vehículo al promotor.; ID heredado: \texttt{RF-046}.
\item[Verificación] \texttt{AC-RF-AUC-014}
\end{description}
\vspace{0.25em}\hrule\vspace{0.65em}
\noindent\textbf{RF-AUC-015 --- Registro de pujas selladas en subastas de vehículos}\par
\begin{description}[style=nextline,leftmargin=3.15cm,labelwidth=2.8cm,font=\normalfont\bfseries]
\item[Obligación] El sistema debe registrar una puja de vehículo como información sellada y mantenerla reservada frente a todos los demás usuarios mientras la subasta permanezca abierta.
\item[Prioridad] Alta
\item[Procedencia] Secciones 2.4, 2.6.4 y decisiones sobre confidencialidad de las subastas de vehículos; ID heredado: \texttt{RF-047}.
\item[Verificación] \texttt{AC-RF-AUC-015}
\end{description}
\vspace{0.25em}\hrule\vspace{0.65em}
\noindent\textbf{RF-AUC-016 --- Verificación reforzada máxima obligatoria para pujar en vehículos}\par
\begin{description}[style=nextline,leftmargin=3.15cm,labelwidth=2.8cm,font=\normalfont\bfseries]
\item[Obligación] El sistema debe permitir pujar en subastas de vehículos a cualquier usuario con verificación reforzada activa, aceptación vigente de términos y sin exclusión, con independencia de sus demás papeles funcionales.
\item[Prioridad] Alta
\item[Procedencia] Secciones 2.3, 2.5, 2.6, 3.1.24 y decisiones sobre el máximo nivel de verificación para pujar en vehículos; ID heredado: \texttt{RF-047A}.
\item[Verificación] \texttt{AC-RF-AUC-016}
\end{description}
\vspace{0.25em}\hrule\vspace{0.65em}
\noindent\textbf{RF-AUC-017 --- Una sola puja vigente y sustitución únicamente al alza en cada vehículo}\par
\begin{description}[style=nextline,leftmargin=3.15cm,labelwidth=2.8cm,font=\normalfont\bfseries]
\item[Obligación] El sistema debe mantener como máximo una puja vigente por usuario y vehículo y aceptar su sustitución únicamente por un importe estrictamente mayor que la puja vigente anterior.
\item[Prioridad] Alta
\item[Procedencia] Secciones 2.4 y 2.7; decisión confirmada sobre pujas de vehículos únicamente crecientes; ID heredado: \texttt{RF-048}.
\item[Verificación] \texttt{AC-RF-AUC-017}
\end{description}
\vspace{0.25em}\hrule\vspace{0.65em}
\noindent\textbf{RF-AUC-018 --- Publicación del resultado oficial de la subasta solo después del cierre}\par
\begin{description}[style=nextline,leftmargin=3.15cm,labelwidth=2.8cm,font=\normalfont\bfseries]
\item[Obligación] Al cerrarse la fase de subasta de vehículos del evento, el sistema debe determinar y publicar para cada vehículo su precio oficial, el usuario adjudicatario y si el resultado implica transferencia a un tercero o conservación por el dueño previo.
\item[Prioridad] Alta
\item[Procedencia] Secciones 2.4, 2.7 y decisiones confirmadas sobre publicación exclusivamente posterior al cierre; ID heredado: \texttt{RF-049}.
\item[Verificación] \texttt{AC-RF-AUC-018}
\end{description}
\vspace{0.25em}\hrule\vspace{0.65em}
\noindent\textbf{RF-AUC-019 --- Determinación por segundo precio de la subasta sellada de vehículos}\par
\begin{description}[style=nextline,leftmargin=3.15cm,labelwidth=2.8cm,font=\normalfont\bfseries]
\item[Obligación] Al cerrar la fase de subasta de vehículos del evento, para cada vehículo con al menos dos pujas válidas el sistema debe declarar adjudicatario al usuario con la mayor puja vigente y fijar como precio oficial el segundo importe vigente más alto, incluyendo en ambos ordenamientos la puja válida del dueño sobre su propio vehículo.
\item[Prioridad] Alta
\item[Procedencia] Sección 1.3, fuente documental heredada y decisiones confirmadas sobre participación del dueño en el segundo precio; ID heredado: \texttt{RF-050}.
\item[Verificación] \texttt{AC-RF-AUC-019}
\end{description}
\vspace{0.25em}\hrule\vspace{0.65em}
\noindent\textbf{RF-AUC-020 --- Entrega obligatoria del vehículo al promotor al terminar la carrera}\par
\begin{description}[style=nextline,leftmargin=3.15cm,labelwidth=2.8cm,font=\normalfont\bfseries]
\item[Obligación] El sistema debe permitir registrar que, al terminar la carrera, el dueño entrega el vehículo al promotor para su custodia y eventual subasta.
\item[Prioridad] Alta
\item[Procedencia] Decisiones prioritarias del cliente sobre entrega del vehículo al promotor antes de subasta.; ID heredado: \texttt{RF-050A}.
\item[Verificación] \texttt{AC-RF-AUC-020}
\end{description}
\vspace{0.25em}\hrule\vspace{0.65em}
\noindent\textbf{RF-AUC-021 --- Descalificación, morosidad y exclusión del dueño por no entrega o por modificación antes de la entrega al promotor}\par
\begin{description}[style=nextline,leftmargin=3.15cm,labelwidth=2.8cm,font=\normalfont\bfseries]
\item[Obligación] Si el dueño no entrega el vehículo al promotor al terminar la carrera o si lo modifica antes de esa entrega, el sistema debe permitir al promotor descalificar a ese competidor, declararlo moroso, excluirlo del sistema y excluirlo de todos los campeonatos aplicables.
\item[Prioridad] Alta
\item[Procedencia] Decisiones prioritarias del cliente sobre no entrega o modificación del vehículo, descalificación y morosidad del dueño.; ID heredado: \texttt{RF-050B}.
\item[Verificación] \texttt{AC-RF-AUC-021}
\end{description}
\vspace{0.25em}\hrule\vspace{0.65em}
\noindent\textbf{RF-AUC-022 --- Custodia del promotor y parc ferme del vehículo subastado}\par
\begin{description}[style=nextline,leftmargin=3.15cm,labelwidth=2.8cm,font=\normalfont\bfseries]
\item[Obligación] El sistema debe conservar que el promotor mantiene bajo control del evento los vehículos recibidos para subasta, en condición de \texttt{parc ferme}, y que el dueño no puede retirarlos del control del evento en contra del resultado oficial vigente.
\item[Prioridad] Alta
\item[Procedencia] Secciones 2.4, 2.5, RF-050A y decisiones prioritarias del cliente sobre \texttt{parc ferme} y control del evento.; ID heredado: \texttt{RF-050C}.
\item[Verificación] \texttt{AC-RF-AUC-022}
\end{description}
\vspace{0.25em}\hrule\vspace{0.65em}
\noindent\textbf{RF-AUC-023 --- Retención del principal de pujas de vehículos de usuarios distintos del dueño}\par
\begin{description}[style=nextline,leftmargin=3.15cm,labelwidth=2.8cm,font=\normalfont\bfseries]
\item[Obligación] El sistema debe mantener retenido respaldo económico suficiente por el importe de cada puja vigente presentada por un usuario distinto del dueño del vehículo durante la subasta correspondiente.
\item[Prioridad] Alta
\item[Procedencia] Secciones 2.6.4 y decisiones sobre garantía real de pago con excepción de puja propia del dueño; ID heredado: \texttt{RF-050D}.
\item[Verificación] \texttt{AC-RF-AUC-023}
\end{description}
\vspace{0.25em}\hrule\vspace{0.65em}
\noindent\textbf{RF-AUC-024 --- Cobro del principal al adjudicatario externo y liberación de retenciones no ganadoras}\par
\begin{description}[style=nextline,leftmargin=3.15cm,labelwidth=2.8cm,font=\normalfont\bfseries]
\item[Obligación] Cuando el adjudicatario de un vehículo sea distinto del dueño previo, el sistema debe cobrar automáticamente el precio oficial al adjudicatario y liberar las retenciones asociadas a las pujas no ganadoras.
\item[Prioridad] Alta
\item[Procedencia] Secciones 2.6.4 y decisiones sobre cobro automático al adjudicatario externo; ID heredado: \texttt{RF-050E}.
\item[Verificación] \texttt{AC-RF-AUC-024}
\end{description}
\vspace{0.25em}\hrule\vspace{0.65em}
\noindent\textbf{RF-AUC-025 --- Pago íntegro del precio oficial al dueño cuando gana un tercero}\par
\begin{description}[style=nextline,leftmargin=3.15cm,labelwidth=2.8cm,font=\normalfont\bfseries]
\item[Obligación] Cuando el adjudicatario de un vehículo sea distinto del dueño previo y el precio oficial haya sido cobrado, el sistema debe pagar automáticamente al dueño previo el precio oficial íntegro, sin descontarle la comisión del comprador.
\item[Prioridad] Alta
\item[Procedencia] Secciones 2.6.4 y decisiones sobre pago automático íntegro al vendedor; ID heredado: \texttt{RF-050F}.
\item[Verificación] \texttt{AC-RF-AUC-025}
\end{description}
\vspace{0.25em}\hrule\vspace{0.65em}
\noindent\textbf{RF-AUC-026 --- Estado de promotor deudor únicamente en adjudicación a un tercero}\par
\begin{description}[style=nextline,leftmargin=3.15cm,labelwidth=2.8cm,font=\normalfont\bfseries]
\item[Obligación] Después del pago al dueño de un vehículo adjudicado a un usuario distinto, el sistema debe mantener al promotor en condición de deudor respecto de ese vehículo mientras el adjudicatario no confirme su entrega física.
\item[Prioridad] Alta
\item[Procedencia] Decisiones sobre condición de deudor del promotor hasta la entrega a un tercero; ID heredado: \texttt{RF-050G}.
\item[Verificación] \texttt{AC-RF-AUC-026}
\end{description}
\vspace{0.25em}\hrule\vspace{0.65em}
\noindent\textbf{RF-AUC-027 --- Confirmación de entrega por adjudicatario distinto del dueño previo}\par
\begin{description}[style=nextline,leftmargin=3.15cm,labelwidth=2.8cm,font=\normalfont\bfseries]
\item[Obligación] El sistema debe permitir que el adjudicatario distinto del dueño previo confirme en la aplicación la recepción física del vehículo adjudicado y pueda mostrar esa confirmación al promotor.
\item[Prioridad] Alta
\item[Procedencia] Decisiones sobre confirmación de entrega en la aplicación; ID heredado: \texttt{RF-050H}.
\item[Verificación] \texttt{AC-RF-AUC-027}
\end{description}
\vspace{0.25em}\hrule\vspace{0.65em}
\noindent\textbf{RF-AUC-028 --- Inicio de disputa de entrega por adjudicatario distinto del dueño previo}\par
\begin{description}[style=nextline,leftmargin=3.15cm,labelwidth=2.8cm,font=\normalfont\bfseries]
\item[Obligación] Si el promotor no entrega físicamente un vehículo adjudicado a un usuario distinto del dueño previo, el sistema debe permitir que ese adjudicatario inicie una disputa de entrega dentro del plazo máximo publicado para el evento.
\item[Prioridad] Alta
\item[Procedencia] RF-075 y decisiones sobre disputa de entrega y plazos del evento; ID heredado: \texttt{RF-050I}.
\item[Verificación] \texttt{AC-RF-AUC-028}
\end{description}
\vspace{0.25em}\hrule\vspace{0.65em}
\noindent\textbf{RF-AUC-029 --- Resolución administrativa de la disputa de entrega}\par
\begin{description}[style=nextline,leftmargin=3.15cm,labelwidth=2.8cm,font=\normalfont\bfseries]
\item[Obligación] Cuando exista una disputa de entrega abierta, el sistema debe permitir que \texttt{admin} registre si la disputa queda resuelta o si el promotor permanece moroso.
\item[Prioridad] Alta
\item[Procedencia] Decisiones prioritarias del cliente sobre resolución por \texttt{admin} y efectos de la morosidad del promotor.; ID heredado: \texttt{RF-050J}.
\item[Verificación] \texttt{AC-RF-AUC-029}
\end{description}
\vspace{0.25em}\hrule\vspace{0.65em}
\noindent\textbf{RF-AUC-030 --- Exclusión del promotor declarado moroso por disputa de entrega}\par
\begin{description}[style=nextline,leftmargin=3.15cm,labelwidth=2.8cm,font=\normalfont\bfseries]
\item[Obligación] Cuando la resolución administrativa mantenga al promotor como moroso, el sistema debe activar su exclusión.
\item[Prioridad] Alta
\item[Procedencia] Decisiones prioritarias del cliente sobre resolución por \texttt{admin} y efectos de la morosidad del promotor.; ID heredado: \texttt{RF-050J}.
\item[Verificación] \texttt{AC-RF-AUC-030}
\end{description}
\vspace{0.25em}\hrule\vspace{0.65em}
\noindent\textbf{RF-AUC-031 --- Conservación de eventos previos del promotor excluido en campeonatos}\par
\begin{description}[style=nextline,leftmargin=3.15cm,labelwidth=2.8cm,font=\normalfont\bfseries]
\item[Obligación] La exclusión de un promotor por morosidad de entrega no debe invalidar los eventos ya realizados para efectos de campeonato.
\item[Prioridad] Alta
\item[Procedencia] Decisiones prioritarias del cliente sobre resolución por \texttt{admin} y efectos de la morosidad del promotor.; ID heredado: \texttt{RF-050J}.
\item[Verificación] \texttt{AC-RF-AUC-031}
\end{description}
\vspace{0.25em}\hrule\vspace{0.65em}
\noindent\textbf{RF-AUC-032 --- Participación del dueño en la subasta de su propio vehículo}\par
\begin{description}[style=nextline,leftmargin=3.15cm,labelwidth=2.8cm,font=\normalfont\bfseries]
\item[Obligación] El sistema debe permitir que el dueño vigente de un vehículo presente una puja por ese mismo vehículo cuando cumpla las condiciones generales de elegibilidad de la subasta.
\item[Prioridad] Alta
\item[Procedencia] DEC-164; decisión prioritaria confirmada sobre participación del dueño
\item[Verificación] \texttt{AC-RF-AUC-032}
\end{description}
\vspace{0.25em}\hrule\vspace{0.65em}
\noindent\textbf{RF-AUC-033 --- Exención de retención del principal para la puja propia del dueño}\par
\begin{description}[style=nextline,leftmargin=3.15cm,labelwidth=2.8cm,font=\normalfont\bfseries]
\item[Obligación] El sistema no debe retener el importe principal de una puja presentada por el dueño vigente sobre su propio vehículo.
\item[Prioridad] Alta
\item[Procedencia] DEC-165; decisión prioritaria confirmada sobre ausencia de pago principal a sí mismo
\item[Verificación] \texttt{AC-RF-AUC-033}
\end{description}
\vspace{0.25em}\hrule\vspace{0.65em}
\noindent\textbf{RF-AUC-034 --- Autoadjudicación del vehículo al dueño previo}\par
\begin{description}[style=nextline,leftmargin=3.15cm,labelwidth=2.8cm,font=\normalfont\bfseries]
\item[Obligación] Cuando el dueño previo resulte adjudicatario de su propio vehículo, el sistema debe registrar una autoadjudicación en la que el dueño conserva el vehículo, no se ejecuta un cobro ni pago circular del precio principal y no se inicia el flujo de entrega o disputa a un comprador distinto.
\item[Prioridad] Alta
\item[Procedencia] DEC-166; decisión prioritaria confirmada sobre adjudicación al dueño actual
\item[Verificación] \texttt{AC-RF-AUC-034}
\end{description}
\vspace{0.25em}\hrule\vspace{0.65em}
\noindent\textbf{RF-AUC-035 --- Advertencia de uso, estudio e inspección técnica posterior a la adjudicación}\par
\begin{description}[style=nextline,leftmargin=3.15cm,labelwidth=2.8cm,font=\normalfont\bfseries]
\item[Obligación] Antes de aceptar la inscripción de un vehículo y antes de aceptar una puja de vehículo, el sistema debe informar que el adjudicatario legítimo puede adquirir el vehículo no solo para uso competitivo o patrimonial, sino también para inspección, estudio técnico o aprendizaje sobre su configuración física, dentro de los límites legales, de seguridad y de derechos de terceros aplicables.
\item[Prioridad] Media
\item[Procedencia] DEC-172 y DEC-173; decisión confirmada sobre difusión tecnológica y adquisición con fines de análisis técnico.
\item[Verificación] \texttt{AC-RF-AUC-035}
\end{description}
\vspace{0.25em}\hrule\vspace{0.65em}
\subsection{Resultados y clasificaciones}
\noindent\textbf{RF-CLS-001 --- Determinación de resultados del dominio del evento}\par
\begin{description}[style=nextline,leftmargin=3.15cm,labelwidth=2.8cm,font=\normalfont\bfseries]
\item[Obligación] El sistema debe determinar los resultados del dominio del evento a partir de los datos oficiales del evento y de las reglas publicadas aplicables.
\item[Prioridad] Alta
\item[Procedencia] Secciones 2.2, 2.6.4, 2.7 y fuente documental \texttt{TOP\_Racing\_Grok\_Clarified.tex} (líneas 150 y 167).; ID heredado: \texttt{RF-051}.
\item[Verificación] \texttt{AC-RF-CLS-001}
\end{description}
\vspace{0.25em}\hrule\vspace{0.65em}
\noindent\textbf{RF-CLS-002 --- Gestión de clasificaciones por modalidad}\par
\begin{description}[style=nextline,leftmargin=3.15cm,labelwidth=2.8cm,font=\normalfont\bfseries]
\item[Obligación] El sistema debe mantener clasificaciones separadas para las modalidades del dominio que correspondan al evento.
\item[Prioridad] Alta
\item[Procedencia] Secciones 1.2, 2.2 y fuente documental \texttt{TOP\_Racing\_Grok\_Clarified.tex} (líneas 150, 163 y 167).; ID heredado: \texttt{RF-052}.
\item[Verificación] \texttt{AC-RF-CLS-002}
\end{description}
\vspace{0.25em}\hrule\vspace{0.65em}
\noindent\textbf{RF-CLS-003 --- Actualización de clasificaciones por alcance territorial y temporal}\par
\begin{description}[style=nextline,leftmargin=3.15cm,labelwidth=2.8cm,font=\normalfont\bfseries]
\item[Obligación] El sistema debe reflejar el aporte de cada evento a las clasificaciones que le correspondan según su alcance territorial y temporal.
\item[Prioridad] Alta
\item[Procedencia] Secciones 1.2, 2.2, 2.4 y fuente documental \texttt{TOP\_Racing\_Grok\_Clarified.tex} (líneas 150, 163 y 171).; ID heredado: \texttt{RF-053}.
\item[Verificación] \texttt{AC-RF-CLS-003}
\end{description}
\vspace{0.25em}\hrule\vspace{0.65em}
\noindent\textbf{RF-CLS-004 --- Publicación de la matriz de pesos del evento desde inscripciones}\par
\begin{description}[style=nextline,leftmargin=3.15cm,labelwidth=2.8cm,font=\normalfont\bfseries]
\item[Obligación] El sistema debe publicar la matriz de pesos del evento desde que el evento entra en estado \texttt{inscripciones}, con sus niveles territoriales y temporales aplicables.
\item[Prioridad] Alta
\item[Procedencia] Secciones 2.2, 2.4, 2.7 y fuente documental \texttt{TOP\_Racing\_Grok\_Clarified.tex} (líneas 150 y 171).; ID heredado: \texttt{RF-054}.
\item[Verificación] \texttt{AC-RF-CLS-004}
\end{description}
\vspace{0.25em}\hrule\vspace{0.65em}
\noindent\textbf{RF-CLS-005 --- Publicación conjunta de resultados y clasificaciones del evento}\par
\begin{description}[style=nextline,leftmargin=3.15cm,labelwidth=2.8cm,font=\normalfont\bfseries]
\item[Obligación] El sistema debe publicar de manera conjunta los resultados oficiales del evento, las clasificaciones correspondientes y la matriz de pesos del evento.
\item[Prioridad] Alta
\item[Procedencia] Secciones 1.2, 2.2, 2.4, 2.7 y fuente documental \texttt{TOP\_Racing\_Grok\_Clarified.tex} (líneas 150 y 171).; ID heredado: \texttt{RF-055}.
\item[Verificación] \texttt{AC-RF-CLS-005}
\end{description}
\vspace{0.25em}\hrule\vspace{0.65em}
\noindent\textbf{RF-CLS-006 --- Exclusión de inscripciones canceladas de los resultados del evento}\par
\begin{description}[style=nextline,leftmargin=3.15cm,labelwidth=2.8cm,font=\normalfont\bfseries]
\item[Obligación] El sistema debe determinar los resultados del evento, las clasificaciones correspondientes y la publicación asociada excluyendo toda inscripción cancelada.
\item[Prioridad] Alta
\item[Procedencia] RF-051 a RF-055 y fuente prioritaria confirmada por el usuario sobre efectos de la cancelación de la inscripción.; ID heredado: \texttt{RF-055A}.
\item[Verificación] \texttt{AC-RF-CLS-006}
\end{description}
\vspace{0.25em}\hrule\vspace{0.65em}
\noindent\textbf{RF-CLS-007 --- Publicación del resultado final vigente de la subasta de vehículo}\par
\begin{description}[style=nextline,leftmargin=3.15cm,labelwidth=2.8cm,font=\normalfont\bfseries]
\item[Obligación] El sistema debe publicar en los resultados oficiales del evento únicamente el ganador y el precio finales vigentes de cada subasta de vehículo después de aplicar las decisiones y efectos públicos que correspondan sobre entrega, descalificación o morosidad.
\item[Prioridad] Alta
\item[Procedencia] RF-050B, RF-050J, RF-067E y fuentes prioritarias confirmadas por el usuario sobre resultado final vigente de la subasta.; ID heredado: \texttt{RF-055B}.
\item[Verificación] \texttt{AC-RF-CLS-007}
\end{description}
\vspace{0.25em}\hrule\vspace{0.65em}
\subsection{Seguridad y cumplimiento}
\noindent\textbf{RF-SAF-001 --- Lista estandarizada de comprobación de seguridad y cumplimiento}\par
\begin{description}[style=nextline,leftmargin=3.15cm,labelwidth=2.8cm,font=\normalfont\bfseries]
\item[Obligación] El sistema debe proporcionar una lista estandarizada de comprobación de seguridad y cumplimiento alineada con las reglas de seguridad publicadas del evento, sin sustituir ni representar el pesaje externo de entrada o de salida.
\item[Prioridad] Alta
\item[Procedencia] Secciones 2.2, 2.3 y fuente documental \texttt{TOP\_Racing\_Grok\_Clarified.tex} (líneas 131 y 152).; ID heredado: \texttt{RF-056}.
\item[Verificación] \texttt{AC-RF-SAF-001}
\end{description}
\vspace{0.25em}\hrule\vspace{0.65em}
\noindent\textbf{RF-SAF-002 --- Soporte a la inspección comunitaria previa a carrera}\par
\begin{description}[style=nextline,leftmargin=3.15cm,labelwidth=2.8cm,font=\normalfont\bfseries]
\item[Obligación] El sistema debe dar soporte a la inspección comunitaria de las condiciones visibles de seguridad de los vehículos del evento mientras el evento no haya entrado en \texttt{carrera}, sin integrar en esa inspección el pesaje externo.
\item[Prioridad] Alta
\item[Procedencia] Secciones 2.2 y 2.3, y fuente documental \texttt{TOP\_Racing\_Grok\_Clarified.tex} (líneas 95, 97 y 152), y fuente prioritaria confirmada por el usuario sobre vigencia de la inspección hasta antes de \texttt{carrera}.; ID heredado: \texttt{RF-057}.
\item[Verificación] \texttt{AC-RF-SAF-002}
\end{description}
\vspace{0.25em}\hrule\vspace{0.65em}
\noindent\textbf{RF-SAF-003 --- Registro de objeciones específicas de seguridad}\par
\begin{description}[style=nextline,leftmargin=3.15cm,labelwidth=2.8cm,font=\normalfont\bfseries]
\item[Obligación] El sistema debe permitir que un competidor registre una objeción específica de seguridad respecto de otro vehículo durante la ventana previa a \texttt{carrera} habilitada para inspección comunitaria.
\item[Prioridad] Alta
\item[Procedencia] Secciones 2.2, 2.3 y fuente documental \texttt{TOP\_Racing\_Grok\_Clarified.tex} (línea 99), y fuente prioritaria confirmada por el usuario sobre vigencia de las protestas hasta antes de \texttt{carrera}.; ID heredado: \texttt{RF-058}.
\item[Verificación] \texttt{AC-RF-SAF-003}
\end{description}
\vspace{0.25em}\hrule\vspace{0.65em}
\noindent\textbf{RF-SAF-004 --- Registro del presentante de la protesta de seguridad}\par
\begin{description}[style=nextline,leftmargin=3.15cm,labelwidth=2.8cm,font=\normalfont\bfseries]
\item[Obligación] El sistema debe conservar asociada a cada objeción de seguridad la identidad del usuario que la presentó.
\item[Prioridad] Media
\item[Procedencia] RF-058 y fuente prioritaria confirmada por el usuario sobre visibilidad restringida de quien presenta la protesta.; ID heredado: \texttt{RF-058A}.
\item[Verificación] \texttt{AC-RF-SAF-004}
\end{description}
\vspace{0.25em}\hrule\vspace{0.65em}
\noindent\textbf{RF-SAF-005 --- Registro del resultado de la comprobación de seguridad}\par
\begin{description}[style=nextline,leftmargin=3.15cm,labelwidth=2.8cm,font=\normalfont\bfseries]
\item[Obligación] El sistema debe poder registrar el resultado de la comprobación de seguridad de cada vehículo, distinguiendo al menos entre vehículo aprobado y vehículo excluido del evento por motivos de seguridad o cumplimiento.
\item[Prioridad] Media
\item[Procedencia] Secciones 2.2 y 2.3, y fuente documental \texttt{TOP\_Racing\_Grok\_Clarified.tex} (líneas 99, 101 y 131).; ID heredado: \texttt{RF-059}.
\item[Verificación] \texttt{AC-RF-SAF-005}
\end{description}
\vspace{0.25em}\hrule\vspace{0.65em}
\noindent\textbf{RF-SAF-006 --- Votación formal ante objeción no resuelta}\par
\begin{description}[style=nextline,leftmargin=3.15cm,labelwidth=2.8cm,font=\normalfont\bfseries]
\item[Obligación] Si una objeción de seguridad no se resuelve por acuerdo antes de \texttt{carrera}, el sistema debe permitir que un competidor inicie una votación formal entre los competidores registrados, registrar el total de votos emitidos y el voto de cada participante identificado por su nombre/usuario exacto, y conservar ese resultado como información para la decisión del promotor, sin que ese resultado cancele por sí mismo la inscripción afectada.
\item[Prioridad] Media
\item[Procedencia] Secciones 2.2, 2.3 y fuente documental \texttt{TOP\_Racing\_Grok\_Clarified.tex} (líneas 108 a 112), y fuentes prioritarias confirmadas por el usuario sobre carácter no vinculante de la protesta y publicidad del detalle de votos.; ID heredado: \texttt{RF-060}.
\item[Verificación] \texttt{AC-RF-SAF-006}
\end{description}
\vspace{0.25em}\hrule\vspace{0.65em}
\noindent\textbf{RF-SAF-007 --- Decisión final del promotor en disputas persistentes}\par
\begin{description}[style=nextline,leftmargin=3.15cm,labelwidth=2.8cm,font=\normalfont\bfseries]
\item[Obligación] El sistema debe permitir que el promotor registre la decisión final sobre una objeción de seguridad cuando la disputa persista sin resolución suficiente y, si esa decisión excluye a un vehículo inscrito antes de \texttt{carrera}, cancelar la inscripción correspondiente.
\item[Prioridad] Media
\item[Procedencia] Secciones 2.2, 2.3 y fuente documental \texttt{TOP\_Racing\_Grok\_Clarified.tex} (líneas 99 y 131), y fuentes prioritarias confirmadas por el usuario sobre decisión final del promotor y cancelación derivada de protesta.; ID heredado: \texttt{RF-061}.
\item[Verificación] \texttt{AC-RF-SAF-007}
\end{description}
\vspace{0.25em}\hrule\vspace{0.65em}
\noindent\textbf{RF-SAF-008 --- Registro opcional de datos de cumplimiento}\par
\begin{description}[style=nextline,leftmargin=3.15cm,labelwidth=2.8cm,font=\normalfont\bfseries]
\item[Obligación] El sistema debe permitir registrar datos de seguridad e inspección sin imponer esa captura como condición obligatoria para continuar con la gestión del evento.
\item[Prioridad] Media
\item[Procedencia] Secciones 2.2, 2.3, 2.7 y fuente documental \texttt{TOP\_Racing\_Grok\_Clarified.tex} (líneas 101 y 152).; ID heredado: \texttt{RF-062}.
\item[Verificación] \texttt{AC-RF-SAF-008}
\end{description}
\vspace{0.25em}\hrule\vspace{0.65em}
\noindent\textbf{RF-SAF-009 --- Carácter informativo de la protesta de seguridad}\par
\begin{description}[style=nextline,leftmargin=3.15cm,labelwidth=2.8cm,font=\normalfont\bfseries]
\item[Obligación] El sistema debe tratar la objeción de seguridad, su votación formal y demás elementos asociados como información para la decisión del promotor antes de \texttt{carrera}, sin efecto cancelatorio automático sobre la inscripción afectada.
\item[Prioridad] Alta
\item[Procedencia] RF-058 a RF-061 y fuente prioritaria confirmada por el usuario sobre carácter no vinculante de la protesta.; ID heredado: \texttt{RF-062A}.
\item[Verificación] \texttt{AC-RF-SAF-009}
\end{description}
\vspace{0.25em}\hrule\vspace{0.65em}
\subsection{Comunidad y novedades}
\noindent\textbf{RF-COM-001 --- Publicación en el foro comunitario}\par
\begin{description}[style=nextline,leftmargin=3.15cm,labelwidth=2.8cm,font=\normalfont\bfseries]
\item[Obligación] El sistema debe permitir que cualquier usuario con identidad verificada publique discusiones, preguntas y retroalimentación en el foro comunitario del ecosistema.
\item[Prioridad] Media
\item[Procedencia] Secciones 2.2, 2.3, 2.5, 3.4 y decisiones prioritarias del cliente sobre publicación comunitaria.; ID heredado: \texttt{RF-063}.
\item[Verificación] \texttt{AC-RF-COM-001}
\end{description}
\vspace{0.25em}\hrule\vspace{0.65em}
\noindent\textbf{RF-COM-002 --- Consulta del foro comunitario}\par
\begin{description}[style=nextline,leftmargin=3.15cm,labelwidth=2.8cm,font=\normalfont\bfseries]
\item[Obligación] El sistema debe permitir consultar las discusiones, preguntas y retroalimentación publicadas en el foro comunitario del ecosistema.
\item[Prioridad] Media
\item[Procedencia] Secciones 2.2, 2.3 y fuente documental \texttt{TOP\_Racing\_Grok\_Clarified.tex} (línea 154).; ID heredado: \texttt{RF-064}.
\item[Verificación] \texttt{AC-RF-COM-002}
\end{description}
\vspace{0.25em}\hrule\vspace{0.65em}
\noindent\textbf{RF-COM-003 --- Publicación en el canal de novedades del ecosistema}\par
\begin{description}[style=nextline,leftmargin=3.15cm,labelwidth=2.8cm,font=\normalfont\bfseries]
\item[Obligación] El sistema debe permitir que cualquier usuario con identidad verificada publique novedades del ecosistema en el canal de novedades.
\item[Prioridad] Media
\item[Procedencia] Secciones 2.1, 2.2, 2.3, 2.5, 3.4 y decisiones prioritarias del cliente sobre publicación de novedades.; ID heredado: \texttt{RF-065}.
\item[Verificación] \texttt{AC-RF-COM-003}
\end{description}
\vspace{0.25em}\hrule\vspace{0.65em}
\noindent\textbf{RF-COM-004 --- Consulta del canal de novedades del ecosistema}\par
\begin{description}[style=nextline,leftmargin=3.15cm,labelwidth=2.8cm,font=\normalfont\bfseries]
\item[Obligación] El sistema debe permitir consultar en el canal de novedades del ecosistema anuncios de eventos, resultados de carrera, resultados de subasta y destacados de la comunidad.
\item[Prioridad] Media
\item[Procedencia] Secciones 1.2, 2.2 y fuente documental \texttt{TOP\_Racing\_Grok\_Clarified.tex} (línea 154).; ID heredado: \texttt{RF-066}.
\item[Verificación] \texttt{AC-RF-COM-004}
\end{description}
\vspace{0.25em}\hrule\vspace{0.65em}
\noindent\textbf{RF-COM-005 --- Moderación del foro comunitario y del canal de novedades}\par
\begin{description}[style=nextline,leftmargin=3.15cm,labelwidth=2.8cm,font=\normalfont\bfseries]
\item[Obligación] El sistema debe permitir que \texttt{admin} modere el foro comunitario y el canal de novedades del ecosistema.
\item[Prioridad] Media
\item[Procedencia] Secciones 2.2, 2.3 y decisiones prioritarias del cliente sobre moderación.; ID heredado: \texttt{RF-066A}.
\item[Verificación] \texttt{AC-RF-COM-005}
\end{description}
\vspace{0.25em}\hrule\vspace{0.65em}
\subsection{Consulta pública e historial}
\noindent\textbf{RF-PUB-001 --- Consulta pública de la información públicada del evento}\par
\begin{description}[style=nextline,leftmargin=3.15cm,labelwidth=2.8cm,font=\normalfont\bfseries]
\item[Obligación] El sistema debe permitir la consulta pública de la información del evento que haya sido publicada conforme a las reglas aplicables.
\item[Prioridad] Alta
\item[Procedencia] Secciones 1.2, 2.2, 2.7 y fuente documental \texttt{TOP\_Racing\_Grok\_Clarified.tex} (líneas 146, 150, 154, 171 y 219).; ID heredado: \texttt{RF-067}.
\item[Verificación] \texttt{AC-RF-PUB-001}
\end{description}
\vspace{0.25em}\hrule\vspace{0.65em}
\noindent\textbf{RF-PUB-002 --- Consulta histórica de resultados oficiales del evento}\par
\begin{description}[style=nextline,leftmargin=3.15cm,labelwidth=2.8cm,font=\normalfont\bfseries]
\item[Obligación] El sistema debe conservar los resultados oficiales publicados de cada evento dentro de su historial consultable.
\item[Prioridad] Alta
\item[Procedencia] Secciones 1.2, 2.2, 2.7 y fuente documental \texttt{TOP\_Racing\_Grok\_Clarified.tex} (líneas 113, 146, 150 y 171).; ID heredado: \texttt{RF-068}.
\item[Verificación] \texttt{AC-RF-PUB-002}
\end{description}
\vspace{0.25em}\hrule\vspace{0.65em}
\noindent\textbf{RF-PUB-003 --- Trazabilidad de resultados oficiales a reglas y datos de origen}\par
\begin{description}[style=nextline,leftmargin=3.15cm,labelwidth=2.8cm,font=\normalfont\bfseries]
\item[Obligación] El sistema debe permitir consultar cada resultado oficial junto con las reglas publicadas y los datos oficiales del evento que le dan sustento.
\item[Prioridad] Alta
\item[Procedencia] Secciones 1.2, 2.7 y fuente documental \texttt{TOP\_Racing\_Grok\_Clarified.tex} (líneas 150, 171, 210, 268 y 272).; ID heredado: \texttt{RF-069}.
\item[Verificación] \texttt{AC-RF-PUB-003}
\end{description}
\vspace{0.25em}\hrule\vspace{0.65em}
\noindent\textbf{RF-PUB-004 --- Registro histórico de resultados públicos de procesos críticos del evento}\par
\begin{description}[style=nextline,leftmargin=3.15cm,labelwidth=2.8cm,font=\normalfont\bfseries]
\item[Obligación] El sistema debe incorporar al historial del evento los resultados públicos de procesos críticos del evento cuando dichos resultados hayan sido publicados por el propio sistema.
\item[Prioridad] Media
\item[Procedencia] Secciones 1.2, 2.2, 2.7 y fuente documental \texttt{TOP\_Racing\_Grok\_Clarified.tex} (líneas 113, 146, 171 y 219).; ID heredado: \texttt{RF-070}.
\item[Verificación] \texttt{AC-RF-PUB-004}
\end{description}
\vspace{0.25em}\hrule\vspace{0.65em}
\noindent\textbf{RF-PUB-005 --- Visibilidad pública de inscripciones canceladas o descalificadas}\par
\begin{description}[style=nextline,leftmargin=3.15cm,labelwidth=2.8cm,font=\normalfont\bfseries]
\item[Obligación] El sistema debe mantener visible en la consulta del evento toda inscripción cancelada o descalificada, indicando expresamente su condición y mostrando la nota correspondiente.
\item[Prioridad] Media
\item[Procedencia] RF-067, RF-088, RF-090 y fuentes prioritarias confirmadas por el usuario sobre visibilidad de inscripciones canceladas o descalificadas y de sus notas.; ID heredado: \texttt{RF-067A}.
\item[Verificación] \texttt{AC-RF-PUB-005}
\end{description}
\vspace{0.25em}\hrule\vspace{0.65em}
\noindent\textbf{RF-PUB-006 --- Visibilidad pública de protestas de seguridad y sus resultados}\par
\begin{description}[style=nextline,leftmargin=3.15cm,labelwidth=2.8cm,font=\normalfont\bfseries]
\item[Obligación] El sistema debe mantener visibles en la consulta del evento las protestas de seguridad tramitadas y sus resultados, mostrando el total de votos emitidos y el voto de cada participante identificado por su nombre o usuario exacto cuando exista votación formal, aunque el promotor decida no cancelar la inscripción afectada.
\item[Prioridad] Media
\item[Procedencia] RF-060, RF-061, RF-062A y fuente prioritaria confirmada por el usuario sobre visibilidad pública de protestas y detalle de votos.; ID heredado: \texttt{RF-067B}.
\item[Verificación] \texttt{AC-RF-PUB-006}
\end{description}
\vspace{0.25em}\hrule\vspace{0.65em}
\noindent\textbf{RF-PUB-007 --- Reserva pública de la identidad de quien presenta la protesta}\par
\begin{description}[style=nextline,leftmargin=3.15cm,labelwidth=2.8cm,font=\normalfont\bfseries]
\item[Obligación] La consulta pública del evento no debe mostrar la identidad del usuario que presentó una protesta de seguridad.
\item[Prioridad] Media
\item[Procedencia] RF-067B y fuente prioritaria confirmada por el usuario sobre reserva pública de la identidad de quien presenta la protesta.; ID heredado: \texttt{RF-067C}.
\item[Verificación] \texttt{AC-RF-PUB-007}
\end{description}
\vspace{0.25em}\hrule\vspace{0.65em}
\noindent\textbf{RF-PUB-008 --- Visibilidad pública de disputas de entrega de vehículos adjudicados}\par
\begin{description}[style=nextline,leftmargin=3.15cm,labelwidth=2.8cm,font=\normalfont\bfseries]
\item[Obligación] El sistema debe mantener visibles en la consulta del evento las disputas de entrega de vehículos adjudicados, mostrando las identidades del promotor y de la persona ganadora involucrados, los reportes de las partes y la decisión final de \texttt{admin}.
\item[Prioridad] Alta
\item[Procedencia] RF-050G a RF-050J y fuentes prioritarias confirmadas por el usuario sobre visibilidad pública de disputas de entrega.; ID heredado: \texttt{RF-067D}.
\item[Verificación] \texttt{AC-RF-PUB-008}
\end{description}
\vspace{0.25em}\hrule\vspace{0.65em}
\noindent\textbf{RF-PUB-009 --- Visibilidad pública del incumplimiento del dueño o del promotor en la entrega del vehículo}\par
\begin{description}[style=nextline,leftmargin=3.15cm,labelwidth=2.8cm,font=\normalfont\bfseries]
\item[Obligación] Cuando el dueño sea declarado moroso por no entregar el vehículo al promotor o cuando el promotor quede moroso por no entregar el vehículo al ganador, el sistema debe mantener visible esa declaración en la consulta del evento junto con los efectos públicos que correspondan.
\item[Prioridad] Alta
\item[Procedencia] RF-050B, RF-050J y fuentes prioritarias confirmadas por el usuario sobre visibilidad de la morosidad del dueño y del promotor.; ID heredado: \texttt{RF-067E}.
\item[Verificación] \texttt{AC-RF-PUB-009}
\end{description}
\vspace{0.25em}\hrule\vspace{0.65em}
\noindent\textbf{RF-PUB-010 --- Historial y trazabilidad temporal de disputas de entrega de vehículos adjudicados}\par
\begin{description}[style=nextline,leftmargin=3.15cm,labelwidth=2.8cm,font=\normalfont\bfseries]
\item[Obligación] El sistema debe incorporar al historial del evento cada disputa de entrega de vehículo adjudicado con su propia secuencia temporal consultable, incluyendo como mínimo su inicio, los reportes registrados, la decisión final de \texttt{admin} y los efectos que de ella resulten.
\item[Prioridad] Alta
\item[Procedencia] RF-050G a RF-050J, RF-067D, RF-067E y fuentes prioritarias confirmadas por el usuario sobre historial y trazabilidad temporal de la disputa de entrega.; ID heredado: \texttt{RF-070A}.
\item[Verificación] \texttt{AC-RF-PUB-010}
\end{description}
\vspace{0.25em}\hrule\vspace{0.65em}
\subsection{Identidad, acceso y exclusión}
\noindent\textbf{RF-IAM-001 --- Identificación del usuario mediante cuenta individual}\par
\begin{description}[style=nextline,leftmargin=3.15cm,labelwidth=2.8cm,font=\normalfont\bfseries]
\item[Obligación] El sistema debe identificar a cada usuario mediante una cuenta individual para el acceso a las funciones que requieran participación dentro del ecosistema.
\item[Prioridad] Alta
\item[Procedencia] Secciones 2.1, 2.3 y fuente documental \texttt{TOP\_Racing\_Grok\_Clarified.tex} (líneas 142 y 154).; ID heredado: \texttt{RF-071}.
\item[Verificación] \texttt{AC-RF-IAM-001}
\end{description}
\vspace{0.25em}\hrule\vspace{0.65em}
\noindent\textbf{RF-IAM-002 --- Verificación de identidad previa a la participación}\par
\begin{description}[style=nextline,leftmargin=3.15cm,labelwidth=2.8cm,font=\normalfont\bfseries]
\item[Obligación] El sistema debe exigir que la identidad del usuario está verificada antes de habilitarle cualquier función que implique participación operativa en el ecosistema.
\item[Prioridad] Alta
\item[Procedencia] Secciones 2.3, 2.5, REX-006 y fuente prioritaria confirmada por el usuario sobre verificación previa de identidad.; ID heredado: \texttt{RF-071A}.
\item[Verificación] \texttt{AC-RF-IAM-002}
\end{description}
\vspace{0.25em}\hrule\vspace{0.65em}
\noindent\textbf{RF-IAM-003 --- Aceptación versionada de términos y condiciones previa a la participación}\par
\begin{description}[style=nextline,leftmargin=3.15cm,labelwidth=2.8cm,font=\normalfont\bfseries]
\item[Obligación] El sistema debe exigir, antes de habilitar cualquier función de participación operativa, que el usuario tenga registrada la aceptación de la versión vigente de los términos y condiciones aplicables.
\item[Prioridad] Alta
\item[Procedencia] Secciones 2.2, 2.5, REX-007 y decisiones prioritarias del cliente sobre aceptación legal versionada.; ID heredado: \texttt{RF-071B}.
\item[Verificación] \texttt{AC-RF-IAM-003}
\end{description}
\vspace{0.25em}\hrule\vspace{0.65em}
\noindent\textbf{RF-IAM-004 --- Acceso a funciones según papeles funcionales y nivel visible de acceso del usuario}\par
\begin{description}[style=nextline,leftmargin=3.15cm,labelwidth=2.8cm,font=\normalfont\bfseries]
\item[Obligación] El sistema debe permitir a cada usuario acceder a las funciones públicas u operativas que correspondan a los papeles funcionales que tenga disponibles y al nivel visible de acceso que le resulte aplicable, distinguiendo como niveles \texttt{base}, \texttt{habilitado} y \texttt{extendido}.
\item[Prioridad] Alta
\item[Procedencia] Secciones 2.3, 2.5, 2.7, RIE-001 y fuente documental \texttt{TOP\_Racing\_Grok\_Clarified.tex} (líneas 142, 144, 146 y 154), además de decisiones prioritarias del cliente sobre verificación reforzada.; ID heredado: \texttt{RF-072}.
\item[Verificación] \texttt{AC-RF-IAM-004}
\end{description}
\vspace{0.25em}\hrule\vspace{0.65em}
\noindent\textbf{RF-IAM-005 --- Acceso restringido a la identidad de quien presenta una protesta de seguridad}\par
\begin{description}[style=nextline,leftmargin=3.15cm,labelwidth=2.8cm,font=\normalfont\bfseries]
\item[Obligación] El sistema debe permitir que solo \texttt{admin} y el promotor del evento consulten la identidad del usuario que presentó una protesta de seguridad.
\item[Prioridad] Media
\item[Procedencia] RF-058A y fuente prioritaria confirmada por el usuario sobre visibilidad restringida de quien presenta la protesta.; ID heredado: \texttt{RF-072A}.
\item[Verificación] \texttt{AC-RF-IAM-005}
\end{description}
\vspace{0.25em}\hrule\vspace{0.65em}
\noindent\textbf{RF-IAM-006 --- Restricción de acceso por declaración válida de morosidad}\par
\begin{description}[style=nextline,leftmargin=3.15cm,labelwidth=2.8cm,font=\normalfont\bfseries]
\item[Obligación] El sistema debe impedir el acceso a las funciones del sistema a todo usuario declarado moroso por decisión válida del promotor o de \texttt{admin} dentro del flujo de subasta y entrega de vehículos.
\item[Prioridad] Alta
\item[Procedencia] RF-050B, RF-050J y fuentes prioritarias confirmadas por el usuario sobre exclusión del moroso del sistema.; ID heredado: \texttt{RF-072B}.
\item[Verificación] \texttt{AC-RF-IAM-006}
\end{description}
\vspace{0.25em}\hrule\vspace{0.65em}
\noindent\textbf{RF-IAM-007 --- Habilitación reforzada de operaciones de mayor riesgo}\par
\begin{description}[style=nextline,leftmargin=3.15cm,labelwidth=2.8cm,font=\normalfont\bfseries]
\item[Obligación] El sistema debe reservar al nivel \texttt{extendido}, es decir, a usuarios con verificación reforzada activa además de las condiciones generales de participación operativa, las operaciones de inscripción de vehículos en eventos, puja en subastas de vehículos y promoción de eventos.
\item[Prioridad] Alta
\item[Procedencia] Secciones 2.3, 2.5, 2.6, 3.1.24 y decisiones prioritarias del cliente sobre verificación reforzada.; ID heredado: \texttt{RF-072C}.
\item[Verificación] \texttt{AC-RF-IAM-007}
\end{description}
\vspace{0.25em}\hrule\vspace{0.65em}
\noindent\textbf{RF-IAM-008 --- Recordatorio al intentar una operación bloqueada por nivel visible de acceso}\par
\begin{description}[style=nextline,leftmargin=3.15cm,labelwidth=2.8cm,font=\normalfont\bfseries]
\item[Obligación] Cuando una operación quede bloqueada por nivel visible de acceso, el sistema debe informar al usuario su nivel actual (\texttt{base}, \texttt{habilitado} o \texttt{extendido}), la condición requerida y la acción necesaria para habilitar la operación.
\item[Prioridad] Alta
\item[Procedencia] Secciones 2.3, 2.5, RIE-001, 3.1.25 y decisiones prioritarias del cliente sobre recordatorio al usuario.; ID heredado: \texttt{RF-072D}.
\item[Verificación] \texttt{AC-RF-IAM-008}
\end{description}
\vspace{0.25em}\hrule\vspace{0.65em}
\noindent\textbf{RF-IAM-009 --- Acumulación de papeles funcionales en un mismo usuario}\par
\begin{description}[style=nextline,leftmargin=3.15cm,labelwidth=2.8cm,font=\normalfont\bfseries]
\item[Obligación] El sistema debe permitir que un mismo usuario asuma de manera acumulativa varios papeles funcionales dentro del ecosistema, salvo la exclusión definida para \texttt{admin}.
\item[Prioridad] Alta
\item[Procedencia] Secciones 2.3 y 2.7, y fuente prioritaria confirmada por el usuario sobre papeles funcionales acumulables.; ID heredado: \texttt{RF-073}.
\item[Verificación] \texttt{AC-RF-IAM-009}
\end{description}
\vspace{0.25em}\hrule\vspace{0.65em}
\noindent\textbf{RF-IAM-010 --- Unicidad y exclusividad de admin}\par
\begin{description}[style=nextline,leftmargin=3.15cm,labelwidth=2.8cm,font=\normalfont\bfseries]
\item[Obligación] El sistema debe reconocer un único \texttt{admin} y no debe permitir que ese papel funcione como papel acumulable junto con los demás papeles funcionales.
\item[Prioridad] Alta
\item[Procedencia] Secciones 2.3 y fuente prioritaria confirmada por el usuario sobre \texttt{admin} único y exclusivo.; ID heredado: \texttt{RF-074}.
\item[Verificación] \texttt{AC-RF-IAM-010}
\end{description}
\vspace{0.25em}\hrule\vspace{0.65em}
\noindent\textbf{RF-IAM-011 --- Registro lógico del estado de identidad verificada del usuario}\par
\begin{description}[style=nextline,leftmargin=3.15cm,labelwidth=2.8cm,font=\normalfont\bfseries]
\item[Obligación] El sistema debe mantener para cada usuario un estado booleano independiente que indique si su identidad está verificada o no verificada.
\item[Prioridad] Alta
\item[Procedencia] Secciones 1.3, 2.3, 2.5 y decisiones prioritarias del cliente sobre estados booleanos del usuario.; ID heredado: \texttt{RF-147}.
\item[Verificación] \texttt{AC-RF-IAM-011}
\end{description}
\vspace{0.25em}\hrule\vspace{0.65em}
\noindent\textbf{RF-IAM-012 --- Registro lógico del estado de verificación reforzada del usuario}\par
\begin{description}[style=nextline,leftmargin=3.15cm,labelwidth=2.8cm,font=\normalfont\bfseries]
\item[Obligación] El sistema debe mantener para cada usuario un estado booleano independiente que indique si su verificación reforzada está activa o inactiva.
\item[Prioridad] Alta
\item[Procedencia] Secciones 1.3, 2.3, 2.5 y decisiones prioritarias del cliente sobre verificación reforzada.; ID heredado: \texttt{RF-147A}.
\item[Verificación] \texttt{AC-RF-IAM-012}
\end{description}
\vspace{0.25em}\hrule\vspace{0.65em}
\noindent\textbf{RF-IAM-013 --- Registro lógico del estado de exclusión del usuario}\par
\begin{description}[style=nextline,leftmargin=3.15cm,labelwidth=2.8cm,font=\normalfont\bfseries]
\item[Obligación] El sistema debe mantener para cada usuario un estado booleano independiente que indique si el usuario está excluido o no excluido.
\item[Prioridad] Alta
\item[Procedencia] Secciones 1.3, 2.3, 2.5, RF-072B y decisiones prioritarias del cliente sobre estados booleanos del usuario.; ID heredado: \texttt{RF-148}.
\item[Verificación] \texttt{AC-RF-IAM-013}
\end{description}
\vspace{0.25em}\hrule\vspace{0.65em}
\noindent\textbf{RF-IAM-014 --- Facultad global de admin sobre verificación y exclusión}\par
\begin{description}[style=nextline,leftmargin=3.15cm,labelwidth=2.8cm,font=\normalfont\bfseries]
\item[Obligación] El sistema debe permitir que \texttt{admin} verifique o desverifique la identidad de cualquier usuario y que excluya o desexcluya a cualquier usuario del sistema.
\item[Prioridad] Alta
\item[Procedencia] Secciones 2.3, 2.5, RF-074 y decisiones prioritarias del cliente sobre facultades de \texttt{admin}.; ID heredado: \texttt{RF-149}.
\item[Verificación] \texttt{AC-RF-IAM-014}
\end{description}
\vspace{0.25em}\hrule\vspace{0.65em}
\noindent\textbf{RF-IAM-015 --- Facultad global de admin sobre verificación reforzada}\par
\begin{description}[style=nextline,leftmargin=3.15cm,labelwidth=2.8cm,font=\normalfont\bfseries]
\item[Obligación] El sistema debe permitir que \texttt{admin} active o desactive la verificación reforzada de cualquier usuario.
\item[Prioridad] Alta
\item[Procedencia] Secciones 2.3, 2.5 y decisiones prioritarias del cliente sobre facultades de \texttt{admin} respecto de la verificación reforzada.; ID heredado: \texttt{RF-149A}.
\item[Verificación] \texttt{AC-RF-IAM-015}
\end{description}
\vspace{0.25em}\hrule\vspace{0.65em}
\noindent\textbf{RF-IAM-016 --- Facultad limitada del promotor sobre verificación y exclusión}\par
\begin{description}[style=nextline,leftmargin=3.15cm,labelwidth=2.8cm,font=\normalfont\bfseries]
\item[Obligación] El sistema debe permitir que el promotor verifique la identidad y excluya únicamente a usuarios vinculados operativamente a un evento que él promueve, y solo mientras ese evento permanezca dentro de su ciclo de vida.
\item[Prioridad] Alta
\item[Procedencia] Secciones 2.3, 2.5, 2.6 y decisiones prioritarias del cliente sobre alcance del promotor.; ID heredado: \texttt{RF-150}.
\item[Verificación] \texttt{AC-RF-IAM-016}
\end{description}
\vspace{0.25em}\hrule\vspace{0.65em}
\noindent\textbf{RF-IAM-017 --- Imposibilidad del promotor para modificar la verificación reforzada}\par
\begin{description}[style=nextline,leftmargin=3.15cm,labelwidth=2.8cm,font=\normalfont\bfseries]
\item[Obligación] El sistema no debe permitir que el promotor active ni desactive la verificación reforzada de ningún usuario.
\item[Prioridad] Alta
\item[Procedencia] Secciones 2.3, 2.5 y decisiones prioritarias del cliente sobre gestión de la verificación reforzada.; ID heredado: \texttt{RF-150A}.
\item[Verificación] \texttt{AC-RF-IAM-017}
\end{description}
\vspace{0.25em}\hrule\vspace{0.65em}
\noindent\textbf{RF-IAM-018 --- Efecto combinado de los estados de verificación, exclusión y aceptación vigente}\par
\begin{description}[style=nextline,leftmargin=3.15cm,labelwidth=2.8cm,font=\normalfont\bfseries]
\item[Obligación] El sistema debe habilitar la participación operativa general de un usuario únicamente cuando su identidad está verificada, su estado de exclusión no está activo y tenga registrada la aceptación vigente de términos y condiciones.
\item[Prioridad] Alta
\item[Procedencia] Secciones 2.5, 2.6, RF-071A, RF-071B, RF-072B, REX-006, REX-007 y decisiones prioritarias del cliente sobre efecto de exclusión.; ID heredado: \texttt{RF-151}.
\item[Verificación] \texttt{AC-RF-IAM-018}
\end{description}
\vspace{0.25em}\hrule\vspace{0.65em}
\noindent\textbf{RF-IAM-019 --- Efecto de la verificación reforzada sobre operaciones de mayor riesgo}\par
\begin{description}[style=nextline,leftmargin=3.15cm,labelwidth=2.8cm,font=\normalfont\bfseries]
\item[Obligación] El sistema debe exigir verificación reforzada activa, además de las condiciones generales de participación operativa, para inscribir vehículos en eventos, pujar en subastas de vehículos y promover eventos.
\item[Prioridad] Alta
\item[Procedencia] Secciones 2.3, 2.5, 2.6 y decisiones prioritarias del cliente sobre operaciones sujetas a verificación reforzada.; ID heredado: \texttt{RF-151A}.
\item[Verificación] \texttt{AC-RF-IAM-019}
\end{description}
\vspace{0.25em}\hrule\vspace{0.65em}
\noindent\textbf{RF-IAM-020 --- Cancelación de sesión en curso por exclusión}\par
\begin{description}[style=nextline,leftmargin=3.15cm,labelwidth=2.8cm,font=\normalfont\bfseries]
\item[Obligación] Cuando la exclusión de un usuario pase a estado activo, el sistema debe cancelar toda sesión vigente de ese usuario.
\item[Prioridad] Alta
\item[Procedencia] Secciones 2.5, RF-148, RF-151, 3.1.25, RIE-001 y decisión prioritaria del cliente confirmada sobre cancelación de sesión y rechazo de inicio de sesión por exclusión activa.; ID heredado: \texttt{RF-151B}.
\item[Verificación] \texttt{AC-RF-IAM-020}
\end{description}
\vspace{0.25em}\hrule\vspace{0.65em}
\noindent\textbf{RF-IAM-021 --- Rechazo de inicio de sesión durante exclusión activa}\par
\begin{description}[style=nextline,leftmargin=3.15cm,labelwidth=2.8cm,font=\normalfont\bfseries]
\item[Obligación] Mientras la exclusión de un usuario permanezca activa, el sistema debe rechazar su inicio de sesión.
\item[Prioridad] Alta
\item[Procedencia] Secciones 2.5, RF-148, RF-151, 3.1.25, RIE-001 y decisión prioritaria del cliente confirmada sobre cancelación de sesión y rechazo de inicio de sesión por exclusión activa.; ID heredado: \texttt{RF-151B}.
\item[Verificación] \texttt{AC-RF-IAM-021}
\end{description}
\vspace{0.25em}\hrule\vspace{0.65em}
\noindent\textbf{RF-IAM-022 --- Mensaje de exclusión al rechazar el inicio de sesión}\par
\begin{description}[style=nextline,leftmargin=3.15cm,labelwidth=2.8cm,font=\normalfont\bfseries]
\item[Obligación] Cuando un inicio de sesión sea rechazado por exclusión activa, la interfaz debe mostrar el mensaje de exclusión correspondiente.
\item[Prioridad] Alta
\item[Procedencia] Secciones 2.5, RF-148, RF-151, 3.1.25, RIE-001 y decisión prioritaria del cliente confirmada sobre cancelación de sesión y rechazo de inicio de sesión por exclusión activa.; ID heredado: \texttt{RF-151B}.
\item[Verificación] \texttt{AC-RF-IAM-022}
\end{description}
\vspace{0.25em}\hrule\vspace{0.65em}
\noindent\textbf{RF-IAM-023 --- Conservación trazable de cambios en verificación, verificación reforzada y exclusión}\par
\begin{description}[style=nextline,leftmargin=3.15cm,labelwidth=2.8cm,font=\normalfont\bfseries]
\item[Obligación] El sistema debe conservar trazablemente cada cambio válido en los estados de identidad verificada, verificación reforzada y exclusión del usuario, incluyendo quién lo realizó, cuándo lo realizó y sobre qué usuario recayó.
\item[Prioridad] Media
\item[Procedencia] Secciones 2.7, RNF-002 y decisiones prioritarias del cliente sobre gestión de esos estados.; ID heredado: \texttt{RF-152}.
\item[Verificación] \texttt{AC-RF-IAM-023}
\end{description}
\vspace{0.25em}\hrule\vspace{0.65em}
\noindent\textbf{RF-IAM-024 --- Exclusión de nueva identidad detectada como correspondiente a un usuario excluido}\par
\begin{description}[style=nextline,leftmargin=3.15cm,labelwidth=2.8cm,font=\normalfont\bfseries]
\item[Obligación] Si el sistema detecta que una nueva identidad corresponde a una persona ya excluida, debe excluir también esa nueva identidad.
\item[Prioridad] Alta
\item[Procedencia] Decisiones prioritarias del cliente sobre exclusión de nuevas identidades detectadas.; ID heredado: \texttt{RF-152A}.
\item[Verificación] \texttt{AC-RF-IAM-024}
\end{description}
\vspace{0.25em}\hrule\vspace{0.65em}
\noindent\textbf{RF-IAM-025 --- Consulta personal de los estados booleanos del usuario}\par
\begin{description}[style=nextline,leftmargin=3.15cm,labelwidth=2.8cm,font=\normalfont\bfseries]
\item[Obligación] El sistema debe permitir que el propio usuario consulte sus estados booleanos de identidad verificada, verificación reforzada y exclusión.
\item[Prioridad] Alta
\item[Procedencia] Secciones 1.3, 2.3, 2.5, 3.1.24 y decisiones prioritarias del cliente sobre consulta personal de estados.; ID heredado: \texttt{RF-153}.
\item[Verificación] \texttt{AC-RF-IAM-025}
\end{description}
\vspace{0.25em}\hrule\vspace{0.65em}
\noindent\textbf{RF-IAM-026 --- Consulta personal del nivel visible de acceso}\par
\begin{description}[style=nextline,leftmargin=3.15cm,labelwidth=2.8cm,font=\normalfont\bfseries]
\item[Obligación] El sistema debe mostrar al propio usuario el nivel visible de acceso derivado de sus estados vigentes usando exclusivamente \texttt{base}, \texttt{habilitado} o \texttt{extendido}.
\item[Prioridad] Alta
\item[Procedencia] Secciones 1.3, 2.3, 2.5, 3.1.24 y decisiones prioritarias del cliente sobre nivel actual visible.; ID heredado: \texttt{RF-153A}.
\item[Verificación] \texttt{AC-RF-IAM-026}
\end{description}
\vspace{0.25em}\hrule\vspace{0.65em}
\noindent\textbf{RF-IAM-027 --- Explicación de condiciones faltantes para el siguiente nivel de acceso}\par
\begin{description}[style=nextline,leftmargin=3.15cm,labelwidth=2.8cm,font=\normalfont\bfseries]
\item[Obligación] Cuando un usuario no alcance el siguiente nivel visible de acceso aplicable, el sistema debe explicar las condiciones vigentes que le faltan para alcanzarlo.
\item[Prioridad] Alta
\item[Procedencia] Secciones 1.3, 2.3, 2.5, 3.1.24 y decisiones prioritarias del cliente sobre nivel actual visible.; ID heredado: \texttt{RF-153A}.
\item[Verificación] \texttt{AC-RF-IAM-027}
\end{description}
\vspace{0.25em}\hrule\vspace{0.65em}
\noindent\textbf{RF-IAM-028 --- Recordatorio al intentar una operación bloqueada por nivel visible de acceso}\par
\begin{description}[style=nextline,leftmargin=3.15cm,labelwidth=2.8cm,font=\normalfont\bfseries]
\item[Obligación] Si un usuario intenta realizar una operación bloqueada por su nivel visible de acceso, el sistema debe informarle su nivel actual, la condición requerida y la acción que debe realizar para habilitar la operación.
\item[Prioridad] Alta
\item[Procedencia] Secciones 2.3, 2.5, 3.1.11, 3.1.24 y decisiones prioritarias del cliente sobre recordatorio ante operaciones bloqueadas.; ID heredado: \texttt{RF-153B}.
\item[Verificación] \texttt{AC-RF-IAM-028}
\end{description}
\vspace{0.25em}\hrule\vspace{0.65em}
\noindent\textbf{RF-IAM-029 --- Consulta administrativa de estados de verificación y exclusión}\par
\begin{description}[style=nextline,leftmargin=3.15cm,labelwidth=2.8cm,font=\normalfont\bfseries]
\item[Obligación] El sistema debe permitir que \texttt{admin} consulte para cualquier usuario los estados booleanos de identidad verificada, verificación reforzada y exclusión.
\item[Prioridad] Alta
\item[Procedencia] Secciones 2.3, 2.5, 3.1.24, RF-149 y RF-149A.; ID heredado: \texttt{RF-154}.
\item[Verificación] \texttt{AC-RF-IAM-029}
\end{description}
\vspace{0.25em}\hrule\vspace{0.65em}
\noindent\textbf{RF-IAM-030 --- Consulta limitada del promotor sobre estados de usuarios de su evento}\par
\begin{description}[style=nextline,leftmargin=3.15cm,labelwidth=2.8cm,font=\normalfont\bfseries]
\item[Obligación] El sistema debe permitir que el promotor consulte los estados booleanos de identidad verificada y exclusión únicamente respecto de usuarios vinculados operativamente a un evento que él promueve, y solo mientras ese evento permanezca dentro de su ciclo de vida.
\item[Prioridad] Alta
\item[Procedencia] Secciones 2.3, 2.5, RF-150 y decisiones prioritarias del cliente sobre consulta del promotor.; ID heredado: \texttt{RF-154A}.
\item[Verificación] \texttt{AC-RF-IAM-030}
\end{description}
\vspace{0.25em}\hrule\vspace{0.65em}
\noindent\textbf{RF-IAM-031 --- Reserva de estados de otros usuarios frente a usuarios no autorizados}\par
\begin{description}[style=nextline,leftmargin=3.15cm,labelwidth=2.8cm,font=\normalfont\bfseries]
\item[Obligación] El sistema no debe mostrar a un usuario los estados booleanos de identidad verificada, verificación reforzada o exclusión de otros usuarios, salvo en los casos expresamente autorizados por las reglas del dominio.
\item[Prioridad] Alta
\item[Procedencia] Secciones 2.5, 3.4, RF-143 y decisiones prioritarias del cliente sobre reserva de estados.; ID heredado: \texttt{RF-155}.
\item[Verificación] \texttt{AC-RF-IAM-031}
\end{description}
\vspace{0.25em}\hrule\vspace{0.65em}
\noindent\textbf{RF-IAM-032 --- Reserva de los estados y del nivel del usuario en la consulta pública del evento}\par
\begin{description}[style=nextline,leftmargin=3.15cm,labelwidth=2.8cm,font=\normalfont\bfseries]
\item[Obligación] La consulta pública del evento no debe mostrar los estados de identidad verificada, verificación reforzada o exclusión de los usuarios vinculados al evento, ni el nivel de verificación derivado de esos estados.
\item[Prioridad] Alta
\item[Procedencia] Secciones 1.3, 2.5, RF-067, RF-143 y decisiones prioritarias del cliente sobre consulta pública del evento.; ID heredado: \texttt{RF-156}.
\item[Verificación] \texttt{AC-RF-IAM-032}
\end{description}
\vspace{0.25em}\hrule\vspace{0.65em}
\noindent\textbf{RF-IAM-033 --- Visibilidad general de usuarios no verificados y reserva de usuarios excluidos}\par
\begin{description}[style=nextline,leftmargin=3.15cm,labelwidth=2.8cm,font=\normalfont\bfseries]
\item[Obligación] La falta de verificación de identidad no debe ocultar por sí misma a un usuario en las consultas del sistema que le correspondan; en cambio, los usuarios con exclusión activa solo deben ser visibles en las consultas generales de usuarios para \texttt{admin}, sin perjuicio de que el propio usuario pueda consultar sus estados y su nivel.
\item[Prioridad] Media
\item[Procedencia] Secciones 2.3, 2.5, RF-148, RF-154 y decisiones prioritarias del cliente sobre usuarios no verificados y usuarios excluidos.; ID heredado: \texttt{RF-157}.
\item[Verificación] \texttt{AC-RF-IAM-033}
\end{description}
\vspace{0.25em}\hrule\vspace{0.65em}
\subsection{Vehículos e inscripciones}
\noindent\textbf{RF-VEH-001 --- Registro previo obligatorio del vehículo en el sistema}\par
\begin{description}[style=nextline,leftmargin=3.15cm,labelwidth=2.8cm,font=\normalfont\bfseries]
\item[Obligación] El sistema debe exigir que un vehículo esté previamente registrado en el sistema antes de que pueda ser inscrito en un evento.
\item[Prioridad] Alta
\item[Procedencia] Fuente prioritaria confirmada por el usuario sobre registro previo del vehículo.; ID heredado: \texttt{RF-080}.
\item[Verificación] \texttt{AC-RF-VEH-001}
\end{description}
\vspace{0.25em}\hrule\vspace{0.65em}
\noindent\textbf{RF-VEH-002 --- Asociación del vehículo registrado con su usuario dueño}\par
\begin{description}[style=nextline,leftmargin=3.15cm,labelwidth=2.8cm,font=\normalfont\bfseries]
\item[Obligación] El sistema debe conservar cada vehículo registrado asociado a un único usuario dueño.
\item[Prioridad] Alta
\item[Procedencia] Fuente documental \texttt{TOP\_Racing\_Grok\_Clarified.tex} (líneas 99, 146 y 221) y fuente prioritaria confirmada por el usuario sobre propiedad del vehículo.; ID heredado: \texttt{RF-081}.
\item[Verificación] \texttt{AC-RF-VEH-002}
\end{description}
\vspace{0.25em}\hrule\vspace{0.65em}
\noindent\textbf{RF-VEH-003 --- Datos mínimos obligatorios del vehículo registrado}\par
\begin{description}[style=nextline,leftmargin=3.15cm,labelwidth=2.8cm,font=\normalfont\bfseries]
\item[Obligación] El sistema debe registrar, como mínimo, para cada vehículo su \texttt{id}, \texttt{número de chasis}, \texttt{fabricante}, \texttt{número de motor} y \texttt{usuario dueño}.
\item[Prioridad] Alta
\item[Procedencia] Fuente prioritaria confirmada por el usuario sobre datos mínimos obligatorios del vehículo registrado.; ID heredado: \texttt{RF-082}.
\item[Verificación] \texttt{AC-RF-VEH-003}
\end{description}
\vspace{0.25em}\hrule\vspace{0.65em}
\noindent\textbf{RF-VEH-004 --- Inscripción del vehículo en el evento por su usuario dueño}\par
\begin{description}[style=nextline,leftmargin=3.15cm,labelwidth=2.8cm,font=\normalfont\bfseries]
\item[Obligación] El sistema debe permitir inscribir un vehículo en un evento únicamente si la inscripción la realiza el usuario dueño asociado a ese vehículo registrado.
\item[Prioridad] Alta
\item[Procedencia] Fuente prioritaria confirmada por el usuario sobre inscripción del vehículo por su dueño.; ID heredado: \texttt{RF-083}.
\item[Verificación] \texttt{AC-RF-VEH-004}
\end{description}
\vspace{0.25em}\hrule\vspace{0.65em}
\noindent\textbf{RF-VEH-005 --- Verificación reforzada obligatoria del usuario dueño para inscribir un vehículo}\par
\begin{description}[style=nextline,leftmargin=3.15cm,labelwidth=2.8cm,font=\normalfont\bfseries]
\item[Obligación] El sistema debe permitir inscribir un vehículo en un evento únicamente si el usuario dueño que realiza la inscripción tiene verificación reforzada activa, además de las demás condiciones generales de participación.
\item[Prioridad] Alta
\item[Procedencia] Secciones 2.3, 2.5, 2.6, 3.1.24 y decisiones prioritarias del cliente sobre inscripción de vehículos con verificación reforzada.; ID heredado: \texttt{RF-083A}.
\item[Verificación] \texttt{AC-RF-VEH-005}
\end{description}
\vspace{0.25em}\hrule\vspace{0.65em}
\noindent\textbf{RF-VEH-006 --- Identificación del usuario piloto en la inscripción cuando exista pilotaje humano}\par
\begin{description}[style=nextline,leftmargin=3.15cm,labelwidth=2.8cm,font=\normalfont\bfseries]
\item[Obligación] El sistema debe exigir que la inscripción identifique al usuario piloto asociado al vehículo inscrito solo cuando la solución participante use pilotaje humano directo.
\item[Prioridad] Alta
\item[Procedencia] Fuente prioritaria confirmada por el usuario sobre datos mínimos obligatorios de la inscripción.; ID heredado: \texttt{RF-084}.
\item[Verificación] \texttt{AC-RF-VEH-006}
\end{description}
\vspace{0.25em}\hrule\vspace{0.65em}
\noindent\textbf{RF-VEH-007 --- Posibilidad de que el usuario piloto sea distinto del usuario dueño cuando exista pilotaje humano}\par
\begin{description}[style=nextline,leftmargin=3.15cm,labelwidth=2.8cm,font=\normalfont\bfseries]
\item[Obligación] Cuando exista pilotaje humano directo, el sistema debe permitir que el usuario piloto identificado en la inscripción sea distinto del usuario dueño que inscribe el vehículo.
\item[Prioridad] Alta
\item[Procedencia] Conversación fuente prioritaria consolidada y fuente prioritaria confirmada por el usuario sobre piloto distinto del dueño.; ID heredado: \texttt{RF-085}.
\item[Verificación] \texttt{AC-RF-VEH-007}
\end{description}
\vspace{0.25em}\hrule\vspace{0.65em}
\noindent\textbf{RF-VEH-008 --- Datos mínimos obligatorios de la inscripción del competidor}\par
\begin{description}[style=nextline,leftmargin=3.15cm,labelwidth=2.8cm,font=\normalfont\bfseries]
\item[Obligación] El sistema debe registrar, como mínimo, para cada inscripción el \texttt{id del vehículo}, el \texttt{id del usuario piloto} cuando exista pilotaje humano directo, la \texttt{aportación de eficiencia}, la \texttt{contribución de terminación}, el \texttt{pago de inscripción} y el \texttt{pago de renta de pista}.
\item[Prioridad] Alta
\item[Procedencia] Fuente documental \texttt{TOP\_Racing\_Grok\_Clarified.tex} (líneas 133, 144, 204, 206 y 212) y fuente prioritaria confirmada por el usuario sobre datos mínimos obligatorios de la inscripción.; ID heredado: \texttt{RF-086}.
\item[Verificación] \texttt{AC-RF-VEH-008}
\end{description}
\vspace{0.25em}\hrule\vspace{0.65em}
\noindent\textbf{RF-VEH-009 --- Aportación mínima obligatoria en la inscripción}\par
\begin{description}[style=nextline,leftmargin=3.15cm,labelwidth=2.8cm,font=\normalfont\bfseries]
\item[Obligación] El sistema debe rechazar toda inscripción cuya aportación incorporada sea inferior a 1 unidad de la moneda del evento.
\item[Prioridad] Alta
\item[Procedencia] Fuentes prioritarias confirmadas por el usuario sobre contribución mínima obligatoria en la inscripción del evento y sobre cobro real inmediato.; ID heredado: \texttt{RF-086A}.
\item[Verificación] \texttt{AC-RF-VEH-009}
\end{description}
\vspace{0.25em}\hrule\vspace{0.65em}
\noindent\textbf{RF-VEH-010 --- Cobro inmediato de importes asociados a la inscripción}\par
\begin{description}[style=nextline,leftmargin=3.15cm,labelwidth=2.8cm,font=\normalfont\bfseries]
\item[Obligación] Al aceptar una inscripción, el sistema debe cobrar de inmediato la inscripción, la aportación incorporada y la renta de pista cuando corresponda.
\item[Prioridad] Alta
\item[Procedencia] Fuentes prioritarias confirmadas por el usuario sobre contribución mínima obligatoria en la inscripción del evento y sobre cobro real inmediato.; ID heredado: \texttt{RF-086A}.
\item[Verificación] \texttt{AC-RF-VEH-010}
\end{description}
\vspace{0.25em}\hrule\vspace{0.65em}
\noindent\textbf{RF-VEH-011 --- Cobro inmediato de la inscripción del competidor}\par
\begin{description}[style=nextline,leftmargin=3.15cm,labelwidth=2.8cm,font=\normalfont\bfseries]
\item[Obligación] El sistema debe cobrar de inmediato los importes de inscripción asociados a toda inscripción registrada válidamente.
\item[Prioridad] Alta
\item[Procedencia] Fuentes prioritarias confirmadas por el usuario sobre cobro inmediato de la inscripción.; ID heredado: \texttt{RF-087}.
\item[Verificación] \texttt{AC-RF-VEH-011}
\end{description}
\vspace{0.25em}\hrule\vspace{0.65em}
\noindent\textbf{RF-VEH-012 --- Nota obligatoria en cancelación o descalificación decidida por el promotor}\par
\begin{description}[style=nextline,leftmargin=3.15cm,labelwidth=2.8cm,font=\normalfont\bfseries]
\item[Obligación] Si el promotor cancela o descalifica una inscripción por una causa permitida, el sistema debe exigir una nota escrita con la razón correspondiente.
\item[Prioridad] Alta
\item[Procedencia] Fuentes prioritarias confirmadas por el usuario sobre cancelaciones y descalificaciones decididas por el promotor.; ID heredado: \texttt{RF-088}.
\item[Verificación] \texttt{AC-RF-VEH-012}
\end{description}
\vspace{0.25em}\hrule\vspace{0.65em}
\noindent\textbf{RF-VEH-013 --- Clasificación válida para pasar a carrera}\par
\begin{description}[style=nextline,leftmargin=3.15cm,labelwidth=2.8cm,font=\normalfont\bfseries]
\item[Obligación] Al pasar el evento a \texttt{carrera}, el sistema debe marcar como \texttt{calificada} únicamente la inscripción no cancelada ni descalificada cuyo participante tenga al menos un tiempo de vuelta válido registrado durante \texttt{velocidad}; una inscripción sin tiempo válido debe quedar \texttt{no calificada} para la carrera.
\item[Prioridad] Alta
\item[Procedencia] Fuente heredada \texttt{RF-089}, precisada por \texttt{DEC-155} sobre calificación mediante tiempo de vuelta válido.; ID heredado: \texttt{RF-089}.
\item[Verificación] \texttt{AC-RF-VEH-013}
\end{description}
\vspace{0.25em}\hrule\vspace{0.65em}
\noindent\textbf{RF-VEH-014 --- Efectos de la cancelación o descalificación de la inscripción}\par
\begin{description}[style=nextline,leftmargin=3.15cm,labelwidth=2.8cm,font=\normalfont\bfseries]
\item[Obligación] Cuando una inscripción se cancele o descalifique por una causa permitida, el sistema debe conservar el registro de la inscripción afectada, invalidarla a todos los efectos del evento, excluirla de todo resultado, clasificación, cálculo económico y subasta del evento, y mantener visible su nota correspondiente en la consulta del evento.
\item[Prioridad] Alta
\item[Procedencia] Fuentes prioritarias confirmadas por el usuario sobre efectos de la cancelación de la inscripción, descalificación del competidor y visibilidad de la nota.; ID heredado: \texttt{RF-090}.
\item[Verificación] \texttt{AC-RF-VEH-014}
\end{description}
\vspace{0.25em}\hrule\vspace{0.65em}
\noindent\textbf{RF-VEH-015 --- Devolución inmediata por cancelación o descalificación previa a carrera}\par
\begin{description}[style=nextline,leftmargin=3.15cm,labelwidth=2.8cm,font=\normalfont\bfseries]
\item[Obligación] Si una inscripción queda cancelada o descalificada antes de \texttt{carrera}, el sistema debe devolver de inmediato los importes cobrados por inscripción, aportación mínima y renta de pista asociados a esa inscripción, así como las aportaciones de terceros cobradas a favor de ese participante dentro del evento.
\item[Prioridad] Alta
\item[Procedencia] Fuentes prioritarias confirmadas por el usuario sobre devoluciones inmediatas.; ID heredado: \texttt{RF-090A}.
\item[Verificación] \texttt{AC-RF-VEH-015}
\end{description}
\vspace{0.25em}\hrule\vspace{0.65em}
\noindent\textbf{RF-VEH-016 --- Descalificación y morosidad del dueño por no entrega o modificación del vehículo}\par
\begin{description}[style=nextline,leftmargin=3.15cm,labelwidth=2.8cm,font=\normalfont\bfseries]
\item[Obligación] Si el dueño no entrega el vehículo al promotor al terminar la carrera o si lo modifica antes de esa entrega, el sistema debe permitir al promotor descalificar al competidor, declararlo moroso y aplicarle la exclusión del sistema y de los campeonatos correspondientes.
\item[Prioridad] Alta
\item[Procedencia] Fuentes prioritarias confirmadas por el usuario sobre no entrega, modificación, descalificación y morosidad del dueño.; ID heredado: \texttt{RF-090B}.
\item[Verificación] \texttt{AC-RF-VEH-016}
\end{description}
\vspace{0.25em}\hrule\vspace{0.65em}
\subsection{Premio del evento}
\noindent\textbf{RF-EPR-001 --- Consulta pública del premio individual del evento}\par
\begin{description}[style=nextline,leftmargin=3.15cm,labelwidth=2.8cm,font=\normalfont\bfseries]
\item[Obligación] El sistema debe permitir consultar públicamente, para cada participante computable del evento, el premio individual que le corresponda dentro de ese evento.
\item[Prioridad] Alta
\item[Procedencia] RF-037, RF-038, RF-040, RF-041, RIE-009 y decisiones prioritarias del cliente sobre publicidad del premio individual del evento.; ID heredado: \texttt{RF-091}.
\item[Verificación] \texttt{AC-RF-EPR-001}
\end{description}
\vspace{0.25em}\hrule\vspace{0.65em}
\noindent\textbf{RF-EPR-002 --- Consulta pública de la distribución posterior asociada al premio individual del evento}\par
\begin{description}[style=nextline,leftmargin=3.15cm,labelwidth=2.8cm,font=\normalfont\bfseries]
\item[Obligación] El sistema debe permitir consultar la distribución económica posterior asociada al premio individual de cada participante computable del evento.
\item[Prioridad] Alta
\item[Procedencia] RF-039, RF-040, RF-041, RIE-009 y decisiones prioritarias del cliente sobre consulta económica del evento.; ID heredado: \texttt{RF-092}.
\item[Verificación] \texttt{AC-RF-EPR-002}
\end{description}
\vspace{0.25em}\hrule\vspace{0.65em}
\noindent\textbf{RF-EPR-003 --- Correspondencia trazable entre premio individual del evento y distribución posterior}\par
\begin{description}[style=nextline,leftmargin=3.15cm,labelwidth=2.8cm,font=\normalfont\bfseries]
\item[Obligación] El sistema debe mantener trazable la correspondencia entre el premio individual del competidor en el evento y la distribución económica posterior asociada a ese premio dentro del mismo evento.
\item[Prioridad] Media
\item[Procedencia] RF-038, RF-039, RF-040, sección 2.7 y decisiones prioritarias del cliente sobre publicación del premio individual del evento.; ID heredado: \texttt{RF-093}.
\item[Verificación] \texttt{AC-RF-EPR-003}
\end{description}
\vspace{0.25em}\hrule\vspace{0.65em}
\noindent\textbf{RF-EPR-004 --- Publicación anonimizada de las aportaciones y de la parte de premio de terceros en evento}\par
\begin{description}[style=nextline,leftmargin=3.15cm,labelwidth=2.8cm,font=\normalfont\bfseries]
\item[Obligación] Cuando existan terceros asociados a contribuciones de un competidor en un evento, el sistema debe publicar las cantidades aportadas y la parte del premio que les corresponda sin mostrar nombres ni identidad de esas personas.
\item[Prioridad] Alta
\item[Procedencia] RF-039, RF-040, sección 2.6.4 y decisiones prioritarias del cliente sobre publicación sin identidad de terceros en evento.; ID heredado: \texttt{RF-094}.
\item[Verificación] \texttt{AC-RF-EPR-004}
\end{description}
\vspace{0.25em}\hrule\vspace{0.65em}
\noindent\textbf{RF-EPR-005 --- Determinación de las aportaciones directas computables a la bolsa del evento}\par
\begin{description}[style=nextline,leftmargin=3.15cm,labelwidth=2.8cm,font=\normalfont\bfseries]
\item[Obligación] El sistema debe determinar, para cada evento, las aportaciones directas efectivamente cobradas a la bolsa del evento que integran el pozo económico oficial y que constituyen patrocinio directo del evento.
\item[Prioridad] Alta
\item[Procedencia] Secciones 2.6.2 y 2.6.4.2 y decisiones prioritarias del cliente sobre patrocinio directo de eventos.; ID heredado: \texttt{RF-115}.
\item[Verificación] \texttt{AC-RF-EPR-005}
\end{description}
\vspace{0.25em}\hrule\vspace{0.65em}
\noindent\textbf{RF-EPR-006 --- Determinación de las contribuciones de eficiencia computables por participante en el evento}\par
\begin{description}[style=nextline,leftmargin=3.15cm,labelwidth=2.8cm,font=\normalfont\bfseries]
\item[Obligación] El sistema debe determinar, para cada participante computable, la contribución total de eficiencia computable que le corresponda dentro del evento, incluida la contribución mínima obligatoria de eficiencia cobrada con la inscripción válida y las aportaciones adicionales cobradas a su favor para el premio basado en eficiencia.
\item[Prioridad] Alta
\item[Procedencia] Secciones 2.6.2 y 2.6.4.2, RF-086A y decisiones prioritarias del cliente sobre aportaciones de eficiencia a participantes.; ID heredado: \texttt{RF-116}.
\item[Verificación] \texttt{AC-RF-EPR-006}
\end{description}
\vspace{0.25em}\hrule\vspace{0.65em}
\noindent\textbf{RF-EPR-007 --- Determinación de las contribuciones de terminación computables por participante en el evento}\par
\begin{description}[style=nextline,leftmargin=3.15cm,labelwidth=2.8cm,font=\normalfont\bfseries]
\item[Obligación] El sistema debe determinar, para cada inscripción computable del evento, la contribución total de terminación que le corresponda, integrada por la contribución de terminación registrada con la inscripción válida y por las aportaciones adicionales cobradas a su favor específicamente para el premio de terminación.
\item[Prioridad] Alta
\item[Procedencia] Secciones 2.6.2, 2.6.4.2, 3.1.13 y decisiones prioritarias del cliente sobre separación entre contribuciones de eficiencia y de terminación.; ID heredado: \texttt{RF-116A}.
\item[Verificación] \texttt{AC-RF-EPR-007}
\end{description}
\vspace{0.25em}\hrule\vspace{0.65em}
\noindent\textbf{RF-EPR-008 --- Asignación base de puntos para el premio del evento}\par
\begin{description}[style=nextline,leftmargin=3.15cm,labelwidth=2.8cm,font=\normalfont\bfseries]
\item[Obligación] El sistema debe asignar la base de puntos del premio del evento conforme a la secuencia base del dominio, que comienza en 100, 63, 40, 25 y 16, y continúa con sus valores sucesivos, según el ordenamiento de eficiencia del evento.
\item[Prioridad] Alta
\item[Procedencia] Secciones 2.6.4.2 y fuente documental \texttt{TOP\_Racing\_Grok\_Clarified.tex}, además de la confirmación prioritaria del usuario.; ID heredado: \texttt{RF-117}.
\item[Verificación] \texttt{AC-RF-EPR-008}
\end{description}
\vspace{0.25em}\hrule\vspace{0.65em}
\noindent\textbf{RF-EPR-009 --- Determinación del coeficiente de premio del evento por participante}\par
\begin{description}[style=nextline,leftmargin=3.15cm,labelwidth=2.8cm,font=\normalfont\bfseries]
\item[Obligación] El sistema debe determinar, para cada participante computable, un coeficiente de premio del evento igual al producto entre la base de puntos que le corresponda por eficiencia y su contribución total de eficiencia computable dentro del evento.
\item[Prioridad] Alta
\item[Procedencia] Secciones 2.6.4.2, RF-116, RF-117 y fuente documental \texttt{TOP\_Racing\_Grok\_Clarified.tex}.; ID heredado: \texttt{RF-118}.
\item[Verificación] \texttt{AC-RF-EPR-009}
\end{description}
\vspace{0.25em}\hrule\vspace{0.65em}
\noindent\textbf{RF-EPR-010 --- Determinación de la suma de coeficientes de premio del evento}\par
\begin{description}[style=nextline,leftmargin=3.15cm,labelwidth=2.8cm,font=\normalfont\bfseries]
\item[Obligación] El sistema debe determinar la suma de los coeficientes de premio del evento correspondientes a todos los participantes computables.
\item[Prioridad] Alta
\item[Procedencia] Secciones 2.6.4.2, RF-118 y fuente documental \texttt{TOP\_Racing\_Grok\_Clarified.tex}.; ID heredado: \texttt{RF-118A}.
\item[Verificación] \texttt{AC-RF-EPR-010}
\end{description}
\vspace{0.25em}\hrule\vspace{0.65em}
\noindent\textbf{RF-EPR-011 --- Determinación operativa del premio del evento por participante}\par
\begin{description}[style=nextline,leftmargin=3.15cm,labelwidth=2.8cm,font=\normalfont\bfseries]
\item[Obligación] El sistema debe determinar, para cada participante computable, la parte del pozo económico oficial que le corresponda a partir de la fracción resultante de dividir su coeficiente de premio entre la suma de coeficientes del evento y asignar el importe monetario entero conforme a \texttt{CALC-EPR-002} y \texttt{CALC-MON-001}.
\item[Prioridad] Alta
\item[Procedencia] Secciones 2.6.4.2, RF-115, RF-118, RF-118A, fuente documental \texttt{TOP\_Racing\_Grok\_Clarified.tex} y decisiones prioritarias del cliente sobre ausencia de redondeo.; ID heredado: \texttt{RF-118B}.
\item[Verificación] \texttt{AC-RF-EPR-011}
\end{description}
\vspace{0.25em}\hrule\vspace{0.65em}
\noindent\textbf{RF-EPR-012 --- Elegibilidad para el premio de terminación}\par
\begin{description}[style=nextline,leftmargin=3.15cm,labelwidth=2.8cm,font=\normalfont\bfseries]
\item[Obligación] El sistema debe considerar para el premio de terminación únicamente a las inscripciones computables cuyos participantes hayan terminado efectivamente la carrera.
\item[Prioridad] Alta
\item[Procedencia] Secciones 2.6.2, 2.6.4.2, 3.1.12, 3.1.13, 3.1.17 y decisiones prioritarias del cliente confirmadas sobre separación entre contribuciones de eficiencia y de terminación.; ID heredado: \texttt{RF-118C}.
\item[Verificación] \texttt{AC-RF-EPR-012}
\end{description}
\vspace{0.25em}\hrule\vspace{0.65em}
\noindent\textbf{RF-EPR-013 --- Determinación del pozo computable de terminación}\par
\begin{description}[style=nextline,leftmargin=3.15cm,labelwidth=2.8cm,font=\normalfont\bfseries]
\item[Obligación] El sistema debe determinar el pozo computable del premio de terminación como la suma del pozo publicado de terminación y las contribuciones de terminación computables.
\item[Prioridad] Alta
\item[Procedencia] Secciones 2.6.2, 2.6.4.2, 3.1.12, 3.1.13, 3.1.17 y decisiones prioritarias del cliente confirmadas sobre separación entre contribuciones de eficiencia y de terminación.; ID heredado: \texttt{RF-118C}.
\item[Verificación] \texttt{AC-RF-EPR-013}
\end{description}
\vspace{0.25em}\hrule\vspace{0.65em}
\noindent\textbf{RF-EPR-014 --- Determinación de vueltas esperadas}\par
\begin{description}[style=nextline,leftmargin=3.15cm,labelwidth=2.8cm,font=\normalfont\bfseries]
\item[Obligación] El sistema debe derivar las vueltas esperadas de cada inscripción computable conforme al tipo de límite de carrera y a \texttt{CALC-EPR-004}.
\item[Prioridad] Alta
\item[Procedencia] Secciones 2.6.2, 2.6.4.2, 3.1.12, 3.1.13, 3.1.17 y decisiones prioritarias del cliente confirmadas sobre separación entre contribuciones de eficiencia y de terminación.; ID heredado: \texttt{RF-118C}.
\item[Verificación] \texttt{AC-RF-EPR-014}
\end{description}
\vspace{0.25em}\hrule\vspace{0.65em}
\noindent\textbf{RF-EPR-015 --- Determinación del peso de terminación}\par
\begin{description}[style=nextline,leftmargin=3.15cm,labelwidth=2.8cm,font=\normalfont\bfseries]
\item[Obligación] El sistema debe determinar el peso de terminación de cada inscripción como el producto entre su contribución total de terminación computable y la relación entre vueltas completadas y vueltas esperadas, conforme a \texttt{CALC-EPR-005}.
\item[Prioridad] Alta
\item[Procedencia] Secciones 2.6.2, 2.6.4.2, 3.1.12, 3.1.13, 3.1.17 y decisiones prioritarias del cliente confirmadas sobre separación entre contribuciones de eficiencia y de terminación.; ID heredado: \texttt{RF-118C}.
\item[Verificación] \texttt{AC-RF-EPR-015}
\end{description}
\vspace{0.25em}\hrule\vspace{0.65em}
\noindent\textbf{RF-EPR-016 --- Asignación proporcional del premio de terminación}\par
\begin{description}[style=nextline,leftmargin=3.15cm,labelwidth=2.8cm,font=\normalfont\bfseries]
\item[Obligación] El sistema debe asignar el pozo computable de terminación entre las inscripciones elegibles en proporción a sus pesos de terminación conforme a \texttt{CALC-EPR-006}.
\item[Prioridad] Alta
\item[Procedencia] Secciones 2.6.2, 2.6.4.2, 3.1.12, 3.1.13, 3.1.17 y decisiones prioritarias del cliente confirmadas sobre separación entre contribuciones de eficiencia y de terminación.; ID heredado: \texttt{RF-118C}.
\item[Verificación] \texttt{AC-RF-EPR-016}
\end{description}
\vspace{0.25em}\hrule\vspace{0.65em}
\noindent\textbf{RF-EPR-017 --- Objetivo de incentivo del premio de terminación}\par
\begin{description}[style=nextline,leftmargin=3.15cm,labelwidth=2.8cm,font=\normalfont\bfseries]
\item[Obligación] El mecanismo de terminación debe preservar la dependencia del premio respecto del grado relativo de terminación y de la contribución computable de la inscripción.
\item[Prioridad] Alta
\item[Procedencia] Secciones 2.6.2, 2.6.4.2, 3.1.12, 3.1.13, 3.1.17 y decisiones prioritarias del cliente confirmadas sobre separación entre contribuciones de eficiencia y de terminación.; ID heredado: \texttt{RF-118C}.
\item[Verificación] \texttt{AC-RF-EPR-017}
\end{description}
\vspace{0.25em}\hrule\vspace{0.65em}
\noindent\textbf{RF-EPR-018 --- Determinación operativa de la distribución económica posterior del premio del evento}\par
\begin{description}[style=nextline,leftmargin=3.15cm,labelwidth=2.8cm,font=\normalfont\bfseries]
\item[Obligación] El sistema debe determinar la distribución económica posterior del premio del evento de cada participante computable entre ese participante y las personas asociadas a sus contribuciones de eficiencia computables, en proporción a las cantidades que integren la contribución total de eficiencia computable de ese participante, conforme a \texttt{CALC-DST-001} y \texttt{CALC-MON-001}.
\item[Prioridad] Alta
\item[Procedencia] Secciones 2.6.4.2, RF-116, RF-118B y decisiones prioritarias del cliente sobre aportantes asociados, reparto del premio del participante y ausencia de redondeo.; ID heredado: \texttt{RF-119}.
\item[Verificación] \texttt{AC-RF-EPR-018}
\end{description}
\vspace{0.25em}\hrule\vspace{0.65em}
\noindent\textbf{RF-EPR-019 --- Determinación operativa de la distribución económica posterior del premio de terminación}\par
\begin{description}[style=nextline,leftmargin=3.15cm,labelwidth=2.8cm,font=\normalfont\bfseries]
\item[Obligación] El sistema debe determinar la distribución económica posterior del premio de terminación de cada inscripción computable entre ese participante y las personas asociadas a sus contribuciones de terminación computables, en proporción a las cantidades que integren la contribución total de terminación computable correspondiente, conforme a \texttt{CALC-DST-001} y \texttt{CALC-MON-001}.
\item[Prioridad] Alta
\item[Procedencia] Secciones 2.6.4.2, RF-116A, RF-118C y decisiones prioritarias del cliente sobre separación entre contribuciones de eficiencia y de terminación.; ID heredado: \texttt{RF-119A}.
\item[Verificación] \texttt{AC-RF-EPR-019}
\end{description}
\vspace{0.25em}\hrule\vspace{0.65em}
\noindent\textbf{RF-EPR-020 --- Publicación trazable del patrocinio y de los premios del evento}\par
\begin{description}[style=nextline,leftmargin=3.15cm,labelwidth=2.8cm,font=\normalfont\bfseries]
\item[Obligación] El sistema debe publicar de forma trazable el pozo económico oficial, las aportaciones directas cobradas a la bolsa del evento, el patrocinio publicado del evento, la contribución total de eficiencia computable de cada participante, la base de puntos por eficiencia, el coeficiente de premio del evento, el premio individual del evento por participante, el pozo del premio de terminación, la contribución total de terminación computable de cada inscripción, el premio de terminación correspondiente y las distribuciones económicas posteriores asociadas.
\item[Prioridad] Alta
\item[Procedencia] Secciones 2.6.4.2, RF-041, RF-091 a RF-094 y RF-115 a RF-119A.; ID heredado: \texttt{RF-120}.
\item[Verificación] \texttt{AC-RF-EPR-020}
\end{description}
\vspace{0.25em}\hrule\vspace{0.65em}
\noindent\textbf{RF-EPR-021 --- Pago automático e inmediato de los premios del evento a todos los beneficiarios}\par
\begin{description}[style=nextline,leftmargin=3.15cm,labelwidth=2.8cm,font=\normalfont\bfseries]
\item[Obligación] Una vez determinadas la distribución económica posterior del premio del evento y la distribución económica posterior del premio de terminación, así como las comisiones del modelo de negocio que correspondan, el sistema debe pagar de inmediato a cada beneficiario la parte neta que le corresponda y cobrar las comisiones aplicables.
\item[Prioridad] Alta
\item[Procedencia] Secciones 2.6.4.2 y decisiones prioritarias del cliente sobre pago automático a todos.; ID heredado: \texttt{RF-120A}.
\item[Verificación] \texttt{AC-RF-EPR-021}
\end{description}
\vspace{0.25em}\hrule\vspace{0.65em}
\noindent\textbf{RF-EPR-022 --- Reserva del pozo consolidado y de las aportaciones antes del cierre de la subasta de vehículos}\par
\begin{description}[style=nextline,leftmargin=3.15cm,labelwidth=2.8cm,font=\normalfont\bfseries]
\item[Obligación] Mientras la fase de subasta de vehículos permanezca abierta, el sistema no debe publicar el pozo económico oficial consolidado ni los importes agregados o individuales de las aportaciones y contribuciones que lo integran; respecto del premio de eficiencia, debe permanecer disponible públicamente únicamente el pozo inicial publicado del evento hasta el cierre de la fase.
\item[Prioridad] Alta
\item[Procedencia] \texttt{DEC-170}; decisión confirmada sobre información económica disponible al dueño al formular su puja.
\item[Verificación] \texttt{AC-RF-EPR-022}
\end{description}
\vspace{0.25em}\hrule\vspace{0.65em}
\subsection{Premio de campeonato}
\noindent\textbf{RF-CPR-001 --- Registro y cobro inmediato de aportación directa a la bolsa del campeonato}\par
\begin{description}[style=nextline,leftmargin=3.15cm,labelwidth=2.8cm,font=\normalfont\bfseries]
\item[Obligación] El sistema debe permitir registrar y cobrar de inmediato aportaciones directas destinadas a la bolsa de un campeonato específico de eficiencia o de terminación solo hasta antes de que el primer evento de ese campeonato haya cerrado inscripciones.
\item[Prioridad] Alta
\item[Procedencia] Secciones 2.2, 2.6.2 y 2.6.4.3, además de decisiones prioritarias del cliente sobre patrocinio directo y ventana de aportación al campeonato.; ID heredado: \texttt{RF-095}.
\item[Verificación] \texttt{AC-RF-CPR-001}
\end{description}
\vspace{0.25em}\hrule\vspace{0.65em}
\noindent\textbf{RF-CPR-002 --- Registro y cobro inmediato de aportación a la contribución de un participante del campeonato}\par
\begin{description}[style=nextline,leftmargin=3.15cm,labelwidth=2.8cm,font=\normalfont\bfseries]
\item[Obligación] El sistema debe permitir registrar y cobrar de inmediato aportaciones destinadas a la contribución de un participante específico dentro de un campeonato determinado, sea de eficiencia o de terminación, solo cuando la función de aportaciones a participantes de campeonato se encuentre activa y solo hasta antes de que el primer evento de ese campeonato haya cerrado inscripciones.
\item[Prioridad] Alta
\item[Procedencia] Secciones 2.6.2 y 2.6.4.3 y decisiones prioritarias del cliente sobre aportaciones a participantes en campeonatos y ventana de aportación al campeonato.; ID heredado: \texttt{RF-096}.
\item[Verificación] \texttt{AC-RF-CPR-002}
\end{description}
\vspace{0.25em}\hrule\vspace{0.65em}
\noindent\textbf{RF-CPR-003 --- Contribución mínima por defecto de participantes al premio de campeonato}\par
\begin{description}[style=nextline,leftmargin=3.15cm,labelwidth=2.8cm,font=\normalfont\bfseries]
\item[Obligación] Cuando exista premio de campeonato, el sistema debe considerar para cada participante del campeonato correspondiente una contribución mínima por defecto de 1 unidad monetaria en la familia aplicable.
\item[Prioridad] Alta
\item[Procedencia] Secciones 2.6.2 y 2.6.4.3 y fuente prioritaria confirmada por el usuario sobre contribución mínima por defecto en premio de campeonato.; ID heredado: \texttt{RF-097}.
\item[Verificación] \texttt{AC-RF-CPR-003}
\end{description}
\vspace{0.25em}\hrule\vspace{0.65em}
\noindent\textbf{RF-CPR-004 --- Determinación del pozo de campeonato}\par
\begin{description}[style=nextline,leftmargin=3.15cm,labelwidth=2.8cm,font=\normalfont\bfseries]
\item[Obligación] El sistema debe consolidar el pozo de campeonato a partir de las aportaciones directas efectivamente cobradas a la bolsa del campeonato correspondiente y de las contribuciones efectivamente cobradas a favor de participantes de ese campeonato dentro de la ventana permitida.
\item[Prioridad] Alta
\item[Procedencia] Secciones 2.6.2 y 2.6.4.3, RF-095, RF-096 y decisiones prioritarias del cliente.; ID heredado: \texttt{RF-098}.
\item[Verificación] \texttt{AC-RF-CPR-004}
\end{description}
\vspace{0.25em}\hrule\vspace{0.65em}
\noindent\textbf{RF-CPR-005 --- Registro de la finalidad económica de la aportación a participante en campeonato}\par
\begin{description}[style=nextline,leftmargin=3.15cm,labelwidth=2.8cm,font=\normalfont\bfseries]
\item[Obligación] El sistema debe conservar trazable que toda aportación realizada a la contribución de un participante dentro de un campeonato tiene como finalidad compartir el premio que le corresponda a ese participante y aumentar su magnitud.
\item[Prioridad] Alta
\item[Procedencia] Secciones 2.6.2 y 2.6.4.3 y decisiones prioritarias del cliente sobre finalidad de las aportaciones a participantes.; ID heredado: \texttt{RF-099}.
\item[Verificación] \texttt{AC-RF-CPR-005}
\end{description}
\vspace{0.25em}\hrule\vspace{0.65em}
\noindent\textbf{RF-CPR-006 --- Registro de patrocinio asociado a aportación directa de campeonato}\par
\begin{description}[style=nextline,leftmargin=3.15cm,labelwidth=2.8cm,font=\normalfont\bfseries]
\item[Obligación] El sistema debe permitir asociar a una aportación directa a la bolsa de campeonato un anuncio con imagen para su publicación como patrocinio del campeonato especificado.
\item[Prioridad] Alta
\item[Procedencia] Sección 2.6.4.3 y fuente prioritaria confirmada por el usuario sobre patrocinio público con anuncio e imagen.; ID heredado: \texttt{RF-100}.
\item[Verificación] \texttt{AC-RF-CPR-006}
\end{description}
\vspace{0.25em}\hrule\vspace{0.65em}
\noindent\textbf{RF-CPR-007 --- Publicación del patrocinio de campeonato}\par
\begin{description}[style=nextline,leftmargin=3.15cm,labelwidth=2.8cm,font=\normalfont\bfseries]
\item[Obligación] El sistema debe publicar el patrocinio asociado a cada campeonato patrocinado.
\item[Prioridad] Alta
\item[Procedencia] Secciones 2.6.2 y 2.6.4.3 y decisiones prioritarias del cliente sobre patrocinio y aportaciones a participantes.; ID heredado: \texttt{RF-101}.
\item[Verificación] \texttt{AC-RF-CPR-007}
\end{description}
\vspace{0.25em}\hrule\vspace{0.65em}
\noindent\textbf{RF-CPR-008 --- Trazabilidad económica del patrocinio de campeonato}\par
\begin{description}[style=nextline,leftmargin=3.15cm,labelwidth=2.8cm,font=\normalfont\bfseries]
\item[Obligación] El sistema debe conservar la relación trazable entre las aportaciones directas a la bolsa del campeonato, los anuncios publicados, las aportaciones a participantes y la distribución económica correspondiente.
\item[Prioridad] Alta
\item[Procedencia] Secciones 2.6.2 y 2.6.4.3 y decisiones prioritarias del cliente sobre patrocinio y aportaciones a participantes.; ID heredado: \texttt{RF-101}.
\item[Verificación] \texttt{AC-RF-CPR-008}
\end{description}
\vspace{0.25em}\hrule\vspace{0.65em}
\noindent\textbf{RF-CPR-009 --- Consulta en la página principal de campeonatos patrocinados y patrocinadores}\par
\begin{description}[style=nextline,leftmargin=3.15cm,labelwidth=2.8cm,font=\normalfont\bfseries]
\item[Obligación] El sistema debe mostrar en la página principal consultas de los campeonatos patrocinados y de los patrocinadores asociados a cada uno.
\item[Prioridad] Alta
\item[Procedencia] Secciones 2.1, 2.2, 2.6.4.3 y fuente prioritaria confirmada por el usuario sobre consulta en la página principal.; ID heredado: \texttt{RF-102}.
\item[Verificación] \texttt{AC-RF-CPR-009}
\end{description}
\vspace{0.25em}\hrule\vspace{0.65em}
\noindent\textbf{RF-CPR-010 --- Determinación de las aportaciones directas computables a la bolsa del campeonato}\par
\begin{description}[style=nextline,leftmargin=3.15cm,labelwidth=2.8cm,font=\normalfont\bfseries]
\item[Obligación] El sistema debe determinar, para cada campeonato con premio activo, las aportaciones directas efectivamente cobradas a la bolsa del campeonato que constituyen patrocinio directo de ese campeonato.
\item[Prioridad] Alta
\item[Procedencia] Secciones 2.6.2, 2.6.4.3, RF-095 y RF-098.; ID heredado: \texttt{RF-121}.
\item[Verificación] \texttt{AC-RF-CPR-010}
\end{description}
\vspace{0.25em}\hrule\vspace{0.65em}
\noindent\textbf{RF-CPR-011 --- Determinación de las contribuciones computables por participante en el campeonato}\par
\begin{description}[style=nextline,leftmargin=3.15cm,labelwidth=2.8cm,font=\normalfont\bfseries]
\item[Obligación] El sistema debe determinar, para cada participante del campeonato correspondiente, la contribución total computable que le corresponda, incluida la contribución mínima por defecto y las aportaciones adicionales efectivamente cobradas a su favor.
\item[Prioridad] Alta
\item[Procedencia] Secciones 2.6.4.3, RF-096, RF-097 y RF-099.; ID heredado: \texttt{RF-122}.
\item[Verificación] \texttt{AC-RF-CPR-011}
\end{description}
\vspace{0.25em}\hrule\vspace{0.65em}
\noindent\textbf{RF-CPR-012 --- Resultado computable de campeonato de eficiencia}\par
\begin{description}[style=nextline,leftmargin=3.15cm,labelwidth=2.8cm,font=\normalfont\bfseries]
\item[Obligación] Para un campeonato de eficiencia, el sistema debe determinar el resultado computable de cada participante a partir de los lugares oficiales del campeonato de eficiencia y la secuencia base vigente.
\item[Prioridad] Alta
\item[Procedencia] Secciones 2.6.4.3 y decisiones prioritarias del cliente sobre la lógica del premio de campeonato y su desempate.; ID heredado: \texttt{RF-123}.
\item[Verificación] \texttt{AC-RF-CPR-012}
\end{description}
\vspace{0.25em}\hrule\vspace{0.65em}
\noindent\textbf{RF-CPR-013 --- Resultado computable de campeonato de terminación}\par
\begin{description}[style=nextline,leftmargin=3.15cm,labelwidth=2.8cm,font=\normalfont\bfseries]
\item[Obligación] Para un campeonato de terminación, el sistema debe determinar el resultado computable de cada participante como la suma de las relaciones entre vueltas completadas y vueltas esperadas de sus eventos computables.
\item[Prioridad] Alta
\item[Procedencia] Secciones 2.6.4.3 y decisiones prioritarias del cliente sobre la lógica del premio de campeonato y su desempate.; ID heredado: \texttt{RF-123}.
\item[Verificación] \texttt{AC-RF-CPR-013}
\end{description}
\vspace{0.25em}\hrule\vspace{0.65em}
\noindent\textbf{RF-CPR-014 --- Desempate del resultado computable de campeonato}\par
\begin{description}[style=nextline,leftmargin=3.15cm,labelwidth=2.8cm,font=\normalfont\bfseries]
\item[Obligación] Cuando dos participantes empaten en el resultado computable para el premio de campeonato, el sistema debe favorecer a quien se haya inscrito antes en el primer evento del campeonato en el que haya participado.
\item[Prioridad] Alta
\item[Procedencia] Secciones 2.6.4.3 y decisiones prioritarias del cliente sobre la lógica del premio de campeonato y su desempate.; ID heredado: \texttt{RF-123}.
\item[Verificación] \texttt{AC-RF-CPR-014}
\end{description}
\vspace{0.25em}\hrule\vspace{0.65em}
\noindent\textbf{RF-CPR-015 --- Determinación del coeficiente de premio de campeonato por participante}\par
\begin{description}[style=nextline,leftmargin=3.15cm,labelwidth=2.8cm,font=\normalfont\bfseries]
\item[Obligación] El sistema debe determinar, para cada participante del campeonato correspondiente, un coeficiente de premio de campeonato igual al producto entre su resultado computable aplicable y su contribución total computable dentro de ese campeonato.
\item[Prioridad] Alta
\item[Procedencia] Secciones 2.6.4.3, RF-122, RF-123 y decisiones prioritarias del cliente sobre la lógica del premio de campeonato.; ID heredado: \texttt{RF-123A}.
\item[Verificación] \texttt{AC-RF-CPR-015}
\end{description}
\vspace{0.25em}\hrule\vspace{0.65em}
\noindent\textbf{RF-CPR-016 --- Determinación de la suma de coeficientes de premio de campeonato}\par
\begin{description}[style=nextline,leftmargin=3.15cm,labelwidth=2.8cm,font=\normalfont\bfseries]
\item[Obligación] El sistema debe determinar la suma de los coeficientes de premio de campeonato correspondientes a todos los participantes del campeonato con premio activo.
\item[Prioridad] Alta
\item[Procedencia] Secciones 2.6.4.3, RF-123A y decisiones prioritarias del cliente sobre la lógica del premio de campeonato.; ID heredado: \texttt{RF-123B}.
\item[Verificación] \texttt{AC-RF-CPR-016}
\end{description}
\vspace{0.25em}\hrule\vspace{0.65em}
\noindent\textbf{RF-CPR-017 --- Determinación operativa del premio de campeonato por participante}\par
\begin{description}[style=nextline,leftmargin=3.15cm,labelwidth=2.8cm,font=\normalfont\bfseries]
\item[Obligación] El sistema debe determinar, para cada participante del campeonato correspondiente, la parte del pozo de campeonato que le corresponda a partir de la fracción resultante de dividir su coeficiente de premio entre la suma de coeficientes del campeonato y asignar el importe monetario entero conforme a \texttt{CALC-CPR-002} y \texttt{CALC-MON-001}.
\item[Prioridad] Alta
\item[Procedencia] Secciones 2.6.4.3, RF-121, RF-123A, RF-123B y decisiones prioritarias del cliente sobre la lógica del premio de campeonato y ausencia de redondeo.; ID heredado: \texttt{RF-123C}.
\item[Verificación] \texttt{AC-RF-CPR-017}
\end{description}
\vspace{0.25em}\hrule\vspace{0.65em}
\noindent\textbf{RF-CPR-018 --- Determinación operativa de la distribución económica posterior del premio de campeonato}\par
\begin{description}[style=nextline,leftmargin=3.15cm,labelwidth=2.8cm,font=\normalfont\bfseries]
\item[Obligación] El sistema debe determinar la distribución económica posterior del premio de campeonato entre el participante correspondiente y las personas asociadas a sus contribuciones computables, en proporción a las cantidades que integren la contribución total computable de ese participante dentro del campeonato correspondiente, conforme a \texttt{CALC-DST-001} y \texttt{CALC-MON-001}.
\item[Prioridad] Alta
\item[Procedencia] Secciones 2.6.4.3, RF-122, RF-123C y decisiones prioritarias del cliente sobre aportaciones a participantes, reparto del premio del participante y ausencia de redondeo.; ID heredado: \texttt{RF-124}.
\item[Verificación] \texttt{AC-RF-CPR-018}
\end{description}
\vspace{0.25em}\hrule\vspace{0.65em}
\noindent\textbf{RF-CPR-019 --- Publicación y trazabilidad del patrocinio y del premio de campeonato}\par
\begin{description}[style=nextline,leftmargin=3.15cm,labelwidth=2.8cm,font=\normalfont\bfseries]
\item[Obligación] El sistema debe publicar y mantener trazable la relación entre las aportaciones directas cobradas a la bolsa del campeonato, los patrocinadores y anuncios publicados, las contribuciones cobradas a participantes, el resultado computable aplicable, los coeficientes de premio de campeonato, el premio individual del campeonato por participante y la distribución económica posterior correspondiente.
\item[Prioridad] Alta
\item[Procedencia] Secciones 2.6.4.3, RF-101, RF-102 y RF-121 a RF-124.; ID heredado: \texttt{RF-125}.
\item[Verificación] \texttt{AC-RF-CPR-019}
\end{description}
\vspace{0.25em}\hrule\vspace{0.65em}
\noindent\textbf{RF-CPR-020 --- Pago automático e inmediato del premio de campeonato a todos los beneficiarios}\par
\begin{description}[style=nextline,leftmargin=3.15cm,labelwidth=2.8cm,font=\normalfont\bfseries]
\item[Obligación] Una vez determinada la distribución económica posterior del premio de campeonato y las comisiones del modelo de negocio que correspondan, el sistema debe pagar de inmediato a cada beneficiario la parte neta que le corresponda y cobrar las comisiones aplicables.
\item[Prioridad] Alta
\item[Procedencia] Secciones 2.6.4.3 y decisiones prioritarias del cliente sobre pago automático a todos.; ID heredado: \texttt{RF-125A}.
\item[Verificación] \texttt{AC-RF-CPR-020}
\end{description}
\vspace{0.25em}\hrule\vspace{0.65em}
\noindent\textbf{RF-CPR-021 --- Consulta pública del premio individual del campeonato}\par
\begin{description}[style=nextline,leftmargin=3.15cm,labelwidth=2.8cm,font=\normalfont\bfseries]
\item[Obligación] El sistema debe permitir consultar públicamente, para cada competidor de un campeonato con premio activo, el premio individual que le corresponda dentro de ese campeonato.
\item[Prioridad] Alta
\item[Procedencia] RF-123C, RF-125, RIE-010 y decisiones prioritarias del cliente sobre publicidad del premio individual del campeonato.; ID heredado: \texttt{RF-138}.
\item[Verificación] \texttt{AC-RF-CPR-021}
\end{description}
\vspace{0.25em}\hrule\vspace{0.65em}
\noindent\textbf{RF-CPR-022 --- Consulta pública de la distribución posterior del premio de campeonato}\par
\begin{description}[style=nextline,leftmargin=3.15cm,labelwidth=2.8cm,font=\normalfont\bfseries]
\item[Obligación] El sistema debe permitir consultar la distribución económica posterior asociada al premio individual de cada competidor en un campeonato con premio activo.
\item[Prioridad] Alta
\item[Procedencia] RF-124, RF-125, RIE-010 y decisiones prioritarias del cliente sobre consulta económica de campeonato.; ID heredado: \texttt{RF-139}.
\item[Verificación] \texttt{AC-RF-CPR-022}
\end{description}
\vspace{0.25em}\hrule\vspace{0.65em}
\noindent\textbf{RF-CPR-023 --- Publicación anonimizada de las aportaciones y de la parte de premio de terceros en campeonato}\par
\begin{description}[style=nextline,leftmargin=3.15cm,labelwidth=2.8cm,font=\normalfont\bfseries]
\item[Obligación] Cuando existan terceros asociados a contribuciones de un competidor en un premio de campeonato, el sistema debe publicar las cantidades aportadas y la parte del premio que les corresponda sin mostrar nombres ni identidad de esas personas.
\item[Prioridad] Alta
\item[Procedencia] RF-124, RF-125 y decisiones prioritarias del cliente sobre publicación sin identidad de terceros en premio de campeonato.; ID heredado: \texttt{RF-140}.
\item[Verificación] \texttt{AC-RF-CPR-023}
\end{description}
\vspace{0.25em}\hrule\vspace{0.65em}
\noindent\textbf{RF-CPR-024 --- Correspondencia trazable entre premio individual del campeonato y distribución posterior}\par
\begin{description}[style=nextline,leftmargin=3.15cm,labelwidth=2.8cm,font=\normalfont\bfseries]
\item[Obligación] El sistema debe mantener trazable la correspondencia entre el premio individual del competidor en el campeonato y la distribución económica posterior asociada a ese resultado dentro de ese mismo campeonato.
\item[Prioridad] Media
\item[Procedencia] RF-123C, RF-124, RF-125 y sección 2.7.; ID heredado: \texttt{RF-141}.
\item[Verificación] \texttt{AC-RF-CPR-024}
\end{description}
\vspace{0.25em}\hrule\vspace{0.65em}
\subsection{Matriz de pesos y prioridad}
\noindent\textbf{RF-WGT-001 --- Determinación de campeonatos y clasificaciones aplicables al evento}\par
\begin{description}[style=nextline,leftmargin=3.15cm,labelwidth=2.8cm,font=\normalfont\bfseries]
\item[Obligación] El sistema debe determinar, para cada evento, los campeonatos y las clasificaciones que le resulten aplicables según su nivel territorial, su nivel temporal y las modalidades del dominio que correspondan.
\item[Prioridad] Alta
\item[Procedencia] Secciones 2.2, 2.6.3.1 y 2.6.4.; ID heredado: \texttt{RF-103}.
\item[Verificación] \texttt{AC-RF-WGT-001}
\end{description}
\vspace{0.25em}\hrule\vspace{0.65em}
\noindent\textbf{RF-WGT-002 --- Determinación del puntaje de prioridad propio del evento}\par
\begin{description}[style=nextline,leftmargin=3.15cm,labelwidth=2.8cm,font=\normalfont\bfseries]
\item[Obligación] El sistema debe determinar y conservar el puntaje de prioridad (P/C) propio del evento como el cociente entre el pozo económico oficial del evento y el costo por participante del propio evento.
\item[Prioridad] Alta
\item[Procedencia] Secciones 2.6.3, 2.6.3.1 y fuente documental TOP\_Racing\_Grok\_Clarified.tex (líneas 202 a 208).; ID heredado: \texttt{RF-104}.
\item[Verificación] \texttt{AC-RF-WGT-002}
\end{description}
\vspace{0.25em}\hrule\vspace{0.65em}
\noindent\textbf{RF-WGT-003 --- Determinación inicial de la matriz de pesos vigente del evento}\par
\begin{description}[style=nextline,leftmargin=3.15cm,labelwidth=2.8cm,font=\normalfont\bfseries]
\item[Obligación] El sistema debe determinar la matriz de pesos vigente del evento partiendo de peso 100 en toda la matriz antes de aplicar las comparaciones operativas que correspondan.
\item[Prioridad] Alta
\item[Procedencia] Secciones 2.6.3.1 y fuente documental TOP\_Racing\_Grok\_Clarified.tex (líneas 171 y 210).; ID heredado: \texttt{RF-105}.
\item[Verificación] \texttt{AC-RF-WGT-003}
\end{description}
\vspace{0.25em}\hrule\vspace{0.65em}
\noindent\textbf{RF-WGT-004 --- Aplicación del peso fijo mínimo en los niveles más bajos}\par
\begin{description}[style=nextline,leftmargin=3.15cm,labelwidth=2.8cm,font=\normalfont\bfseries]
\item[Obligación] El sistema debe asignar peso fijo de 100 a los niveles más bajos de la estructura de campeonatos, correspondientes a semana y variante, sin requerir comparación con otros eventos.
\item[Prioridad] Alta
\item[Procedencia] Secciones 2.6.3.1 y fuente documental TOP\_Racing\_Grok\_Clarified.tex (líneas 171 y 210).; ID heredado: \texttt{RF-106}.
\item[Verificación] \texttt{AC-RF-WGT-004}
\end{description}
\vspace{0.25em}\hrule\vspace{0.65em}
\noindent\textbf{RF-WGT-005 --- Comparación operativa entre eventos para ajustar pesos}\par
\begin{description}[style=nextline,leftmargin=3.15cm,labelwidth=2.8cm,font=\normalfont\bfseries]
\item[Obligación] Cuando la matriz de pesos deba ajustarse por afectación entre eventos, el sistema debe comparar entre sí los eventos que se solapen en los niveles territoriales y temporales inmediatamente inferiores aplicables al nivel superior correspondiente.
\item[Prioridad] Alta
\item[Procedencia] Secciones 2.6.3.1 y fuente documental TOP\_Racing\_Grok\_Clarified.tex (líneas 171 y 210).; ID heredado: \texttt{RF-107}.
\item[Verificación] \texttt{AC-RF-WGT-005}
\end{description}
\vspace{0.25em}\hrule\vspace{0.65em}
\noindent\textbf{RF-WGT-006 --- Conservación del peso del evento con mayor P/C}\par
\begin{description}[style=nextline,leftmargin=3.15cm,labelwidth=2.8cm,font=\normalfont\bfseries]
\item[Obligación] Cuando una comparación operativa entre eventos requiera ajustar pesos, el sistema debe conservar el peso del evento con mayor P/C.
\item[Prioridad] Alta
\item[Procedencia] Secciones 2.6.3.1 y decisiones prioritarias del cliente confirmadas sobre secuencia y desempate, además de la fuente documental TOP\_Racing\_Grok\_Clarified.tex (líneas 171 y 210).; ID heredado: \texttt{RF-108}.
\item[Verificación] \texttt{AC-RF-WGT-006}
\end{description}
\vspace{0.25em}\hrule\vspace{0.65em}
\noindent\textbf{RF-WGT-007 --- Descenso del peso del evento con menor P/C}\par
\begin{description}[style=nextline,leftmargin=3.15cm,labelwidth=2.8cm,font=\normalfont\bfseries]
\item[Obligación] Cuando una comparación operativa entre eventos requiera ajustar pesos, el sistema debe hacer descender el peso del evento con menor P/C al siguiente valor de la secuencia base vigente.
\item[Prioridad] Alta
\item[Procedencia] Secciones 2.6.3.1 y decisiones prioritarias del cliente confirmadas sobre secuencia y desempate, además de la fuente documental TOP\_Racing\_Grok\_Clarified.tex (líneas 171 y 210).; ID heredado: \texttt{RF-108}.
\item[Verificación] \texttt{AC-RF-WGT-007}
\end{description}
\vspace{0.25em}\hrule\vspace{0.65em}
\noindent\textbf{RF-WGT-008 --- Desempate de P/C por antigüedad de registro}\par
\begin{description}[style=nextline,leftmargin=3.15cm,labelwidth=2.8cm,font=\normalfont\bfseries]
\item[Obligación] Cuando dos eventos comparables tengan el mismo P/C, el sistema debe favorecer al evento registrado más temprano.
\item[Prioridad] Alta
\item[Procedencia] Secciones 2.6.3.1 y decisiones prioritarias del cliente confirmadas sobre secuencia y desempate, además de la fuente documental TOP\_Racing\_Grok\_Clarified.tex (líneas 171 y 210).; ID heredado: \texttt{RF-108}.
\item[Verificación] \texttt{AC-RF-WGT-008}
\end{description}
\vspace{0.25em}\hrule\vspace{0.65em}
\noindent\textbf{RF-WGT-009 --- Iteración del ajuste hasta estabilizar la matriz de pesos}\par
\begin{description}[style=nextline,leftmargin=3.15cm,labelwidth=2.8cm,font=\normalfont\bfseries]
\item[Obligación] El sistema debe repetir las comparaciones operativas necesarias hasta que no se requieran nuevos descensos de peso en la matriz del evento.
\item[Prioridad] Alta
\item[Procedencia] Secciones 2.6.3.1 y fuente documental TOP\_Racing\_Grok\_Clarified.tex (líneas 171 y 210).; ID heredado: \texttt{RF-108A}.
\item[Verificación] \texttt{AC-RF-WGT-009}
\end{description}
\vspace{0.25em}\hrule\vspace{0.65em}
\noindent\textbf{RF-WGT-010 --- Escalamiento de puntos por peso del evento}\par
\begin{description}[style=nextline,leftmargin=3.15cm,labelwidth=2.8cm,font=\normalfont\bfseries]
\item[Obligación] El sistema debe escalar la base de puntos obtenida en cada campeonato aplicable multiplicándola por el peso vigente del evento y dividiéndola entre 100.
\item[Prioridad] Alta
\item[Procedencia] Secciones 2.6.3.1 y fuente documental TOP\_Racing\_Grok\_Clarified.tex (línea 171).; ID heredado: \texttt{RF-108B}.
\item[Verificación] \texttt{AC-RF-WGT-010}
\end{description}
\vspace{0.25em}\hrule\vspace{0.65em}
\noindent\textbf{RF-WGT-011 --- Publicación trazable del puntaje de prioridad y de la matriz de pesos vigente}\par
\begin{description}[style=nextline,leftmargin=3.15cm,labelwidth=2.8cm,font=\normalfont\bfseries]
\item[Obligación] El sistema debe publicar de forma trazable el puntaje de prioridad propio del evento, la matriz de pesos vigente, la secuencia aplicada cuando haya ajustes y los campeonatos y clasificaciones a los que el evento aporta resultados.
\item[Prioridad] Alta
\item[Procedencia] Secciones 2.6.3, 2.6.3.1, RF-054 y RF-055.; ID heredado: \texttt{RF-109}.
\item[Verificación] \texttt{AC-RF-WGT-011}
\end{description}
\vspace{0.25em}\hrule\vspace{0.65em}
\subsection{Eficiencia relativa al mercado}
\noindent\textbf{RF-EFF-001 --- Determinación del rendimiento observado antes de la subasta}\par
\begin{description}[style=nextline,leftmargin=3.15cm,labelwidth=2.8cm,font=\normalfont\bfseries]
\item[Obligación] El sistema debe determinar el rendimiento observado de cada participante computable como la velocidad media correspondiente a su vuelta válida más rápida de calificación, conforme a \texttt{CALC-EFF-001}.
\item[Prioridad] Alta
\item[Procedencia] Secciones 2.6.4 y decisiones sobre rendimiento fijado antes de la subasta; ID heredado: \texttt{RF-110}.
\item[Verificación] \texttt{AC-RF-EFF-001}
\end{description}
\vspace{0.25em}\hrule\vspace{0.65em}
\noindent\textbf{RF-EFF-002 --- Consolidación de la muestra oficial precio--rendimiento}\par
\begin{description}[style=nextline,leftmargin=3.15cm,labelwidth=2.8cm,font=\normalfont\bfseries]
\item[Obligación] Después del cierre de las subastas de vehículos, el sistema debe consolidar un único par válido de rendimiento observado y precio oficial de subasta por inscripción computable.
\item[Prioridad] Alta
\item[Procedencia] Secciones 2.6.4.1 y decisiones sobre línea de tendencia global; ID heredado: \texttt{RF-111}.
\item[Verificación] \texttt{AC-RF-EFF-002}
\end{description}
\vspace{0.25em}\hrule\vspace{0.65em}
\noindent\textbf{RF-EFF-003 --- Ausencia de resultados provisionales de eficiencia durante la subasta}\par
\begin{description}[style=nextline,leftmargin=3.15cm,labelwidth=2.8cm,font=\normalfont\bfseries]
\item[Obligación] Mientras la subasta de vehículos permanezca abierta, el sistema no debe determinar para publicación ni presentar clasificación, puntos o premio provisionales de eficiencia.
\item[Prioridad] Alta
\item[Procedencia] DEC-162 y decisión prioritaria confirmada sobre publicación solo al cierre; ID heredado: \texttt{RF-111}.
\item[Verificación] \texttt{AC-RF-EFF-003}
\end{description}
\vspace{0.25em}\hrule\vspace{0.65em}
\noindent\textbf{RF-EFF-004 --- Reserva de información derivada de las pujas antes del cierre}\par
\begin{description}[style=nextline,leftmargin=3.15cm,labelwidth=2.8cm,font=\normalfont\bfseries]
\item[Obligación] Mientras la subasta de vehículos permanezca abierta, el sistema no debe publicar conteos de pujas, orden de pujas, precios provisionales, adjudicatarios provisionales ni datos derivados que permitan inferir la situación relativa de las pujas de terceros.
\item[Prioridad] Alta
\item[Procedencia] DEC-162 y reglas de subasta sellada; ID heredado: \texttt{RF-111}.
\item[Verificación] \texttt{AC-RF-EFF-004}
\end{description}
\vspace{0.25em}\hrule\vspace{0.65em}
\noindent\textbf{RF-EFF-005 --- Creación de la señal oficial de mercado al cierre}\par
\begin{description}[style=nextline,leftmargin=3.15cm,labelwidth=2.8cm,font=\normalfont\bfseries]
\item[Obligación] Después del cierre de las subastas de vehículos, el sistema debe utilizar los precios oficiales y los rendimientos fijados previamente para formar la muestra oficial de mercado del evento.
\item[Prioridad] Alta
\item[Procedencia] Secciones 2.6.4.1 y decisiones sobre secuencia temporal del cálculo; ID heredado: \texttt{RF-111}.
\item[Verificación] \texttt{AC-RF-EFF-005}
\end{description}
\vspace{0.25em}\hrule\vspace{0.65em}
\noindent\textbf{RF-EFF-006 --- Línea de tendencia log-lineal global de precios en función del rendimiento}\par
\begin{description}[style=nextline,leftmargin=3.15cm,labelwidth=2.8cm,font=\normalfont\bfseries]
\item[Obligación] El sistema debe ajustar una única línea de tendencia de los precios oficiales en función del rendimiento observado mediante el modelo $\ln(\widehat{C}(S))=a+bS$, equivalente a $\widehat{C}(S)=\exp(a+bS)$, y los casos de muestra definidos en \texttt{CALC-EFF-002}.
\item[Prioridad] Alta
\item[Procedencia] DEC-160, DEC-168 y corrección confirmada de la dirección de la regresión; ID heredado: \texttt{RF-112}.
\item[Verificación] \texttt{AC-RF-EFF-006}
\end{description}
\vspace{0.25em}\hrule\vspace{0.65em}
\noindent\textbf{RF-EFF-007 --- Estimación única de la línea de tendencia con todos los participantes computables}\par
\begin{description}[style=nextline,leftmargin=3.15cm,labelwidth=2.8cm,font=\normalfont\bfseries]
\item[Obligación] El sistema debe estimar una sola vez los parámetros de la línea de tendencia utilizando conjuntamente todos los pares válidos de la muestra oficial del evento.
\item[Prioridad] Alta
\item[Procedencia] DEC-160; decisión confirmada sobre regresión global; ID heredado: \texttt{RF-112}.
\item[Verificación] \texttt{AC-RF-EFF-007}
\end{description}
\vspace{0.25em}\hrule\vspace{0.65em}
\noindent\textbf{RF-EFF-008 --- Conservación trazable de la muestra y parámetros de tendencia}\par
\begin{description}[style=nextline,leftmargin=3.15cm,labelwidth=2.8cm,font=\normalfont\bfseries]
\item[Obligación] El sistema debe conservar de forma trazable la muestra oficial usada por la regresión y los parámetros ajustados de la línea de tendencia aplicable al evento.
\item[Prioridad] Alta
\item[Procedencia] Secciones 2.6.4.1 y requisitos de auditabilidad; ID heredado: \texttt{RF-112}.
\item[Verificación] \texttt{AC-RF-EFF-008}
\end{description}
\vspace{0.25em}\hrule\vspace{0.65em}
\noindent\textbf{RF-EFF-009 --- Tratamiento normativo de muestras de tamaño cero, uno y dos y de rango insuficiente}\par
\begin{description}[style=nextline,leftmargin=3.15cm,labelwidth=2.8cm,font=\normalfont\bfseries]
\item[Obligación] El sistema debe aplicar los casos $n=0$, $n=1$, $n=2$ y rango insuficiente establecidos en \texttt{CALC-EFF-002}, sin sustituirlos por comparaciones de residuos numéricos aproximados.
\item[Prioridad] Alta
\item[Procedencia] DEC-153 precisada por DEC-160; ID heredado: \texttt{RF-112}.
\item[Verificación] \texttt{AC-RF-EFF-009}
\end{description}
\vspace{0.25em}\hrule\vspace{0.65em}
\noindent\textbf{RF-EFF-010 --- Determinación de la diferencia de eficiencia precio--tendencia}\par
\begin{description}[style=nextline,leftmargin=3.15cm,labelwidth=2.8cm,font=\normalfont\bfseries]
\item[Obligación] El sistema debe determinar para cada participante computable la diferencia de eficiencia $D_i=\widehat{C}(S_i)-C_i$ conforme a \texttt{CALC-EFF-003}.
\item[Prioridad] Alta
\item[Procedencia] DEC-161 y corrección confirmada del significado de eficiencia; ID heredado: \texttt{RF-113}.
\item[Verificación] \texttt{AC-RF-EFF-010}
\end{description}
\vspace{0.25em}\hrule\vspace{0.65em}
\noindent\textbf{RF-EFF-011 --- Ordenamiento descendente por diferencia de eficiencia}\par
\begin{description}[style=nextline,leftmargin=3.15cm,labelwidth=2.8cm,font=\normalfont\bfseries]
\item[Obligación] El sistema debe ordenar la clasificación de eficiencia en sentido descendente por la diferencia $D_i$.
\item[Prioridad] Alta
\item[Procedencia] Secciones 2.6.4.1 y decisiones de clasificación; ID heredado: \texttt{RF-113A}.
\item[Verificación] \texttt{AC-RF-EFF-011}
\end{description}
\vspace{0.25em}\hrule\vspace{0.65em}
\noindent\textbf{RF-EFF-012 --- Asignación de base de puntos de eficiencia}\par
\begin{description}[style=nextline,leftmargin=3.15cm,labelwidth=2.8cm,font=\normalfont\bfseries]
\item[Obligación] El sistema debe asignar al orden de eficiencia la secuencia base definida por \texttt{CALC-PTS-001}.
\item[Prioridad] Alta
\item[Procedencia] Secciones 2.6.4.1 y DEC-152; ID heredado: \texttt{RF-113A}.
\item[Verificación] \texttt{AC-RF-EFF-012}
\end{description}
\vspace{0.25em}\hrule\vspace{0.65em}
\noindent\textbf{RF-EFF-013 --- Desempate de eficiencia por antigüedad de inscripción}\par
\begin{description}[style=nextline,leftmargin=3.15cm,labelwidth=2.8cm,font=\normalfont\bfseries]
\item[Obligación] Cuando dos participantes tengan la misma diferencia de eficiencia normativa, el sistema debe colocar primero al participante cuya inscripción válida tenga la marca temporal más antigua.
\item[Prioridad] Alta
\item[Procedencia] Secciones 2.6.4.1 y decisión de desempate; ID heredado: \texttt{RF-113A}.
\item[Verificación] \texttt{AC-RF-EFF-013}
\end{description}
\vspace{0.25em}\hrule\vspace{0.65em}
\noindent\textbf{RF-EFF-014 --- Publicación de la cadena de trazabilidad de eficiencia después del cierre}\par
\begin{description}[style=nextline,leftmargin=3.15cm,labelwidth=2.8cm,font=\normalfont\bfseries]
\item[Obligación] Después del cierre, el sistema debe permitir consultar la relación entre rendimiento observado, precio oficial de subasta, línea de tendencia, precio de tendencia, diferencia de eficiencia, base de puntos y posición resultante.
\item[Prioridad] Alta
\item[Procedencia] Secciones 2.6.4.1 y requisitos de transparencia controlada; ID heredado: \texttt{RF-114}.
\item[Verificación] \texttt{AC-RF-EFF-014}
\end{description}
\vspace{0.25em}\hrule\vspace{0.65em}
\noindent\textbf{RF-EFF-015 --- Publicación exclusiva de resultados oficiales de eficiencia}\par
\begin{description}[style=nextline,leftmargin=3.15cm,labelwidth=2.8cm,font=\normalfont\bfseries]
\item[Obligación] El sistema debe publicar la clasificación, los puntos y el premio basado en eficiencia únicamente cuando se hayan determinado a partir de la muestra oficial posterior al cierre de la subasta.
\item[Prioridad] Alta
\item[Procedencia] DEC-162 y decisión confirmada de no publicar resultados provisionales; ID heredado: \texttt{RF-114}.
\item[Verificación] \texttt{AC-RF-EFF-015}
\end{description}
\vspace{0.25em}\hrule\vspace{0.65em}
\noindent\textbf{RF-EFF-016 --- Interpretación del signo de la diferencia de eficiencia}\par
\begin{description}[style=nextline,leftmargin=3.15cm,labelwidth=2.8cm,font=\normalfont\bfseries]
\item[Obligación] El sistema debe presentar un valor positivo de $D_i$ como valoración inferior a la tendencia para el rendimiento observado y un valor negativo como valoración superior a esa tendencia.
\item[Prioridad] Alta
\item[Procedencia] DEC-161 y semántica confirmada de subvaloración y sobrevaloración; ID heredado: \texttt{RF-114}.
\item[Verificación] \texttt{AC-RF-EFF-016}
\end{description}
\vspace{0.25em}\hrule\vspace{0.65em}
\noindent\textbf{RF-EFF-017 --- Uso del precio oficial de autoadjudicación en la línea de tendencia}\par
\begin{description}[style=nextline,leftmargin=3.15cm,labelwidth=2.8cm,font=\normalfont\bfseries]
\item[Obligación] Cuando el dueño resulte adjudicatario de su propio vehículo, el sistema debe usar el precio oficial de segundo precio de esa autoadjudicación como $C_i$ en la muestra de eficiencia bajo las mismas reglas aplicables a los demás vehículos.
\item[Prioridad] Alta
\item[Procedencia] DEC-167 y decisión confirmada sobre señal económica de la puja del dueño; ID heredado: \texttt{RF-114}.
\item[Verificación] \texttt{AC-RF-EFF-017}
\end{description}
\vspace{0.25em}\hrule\vspace{0.65em}
\noindent\textbf{RF-EFF-018 --- Explicación posterior del incentivo de eficiencia}\par
\begin{description}[style=nextline,leftmargin=3.15cm,labelwidth=2.8cm,font=\normalfont\bfseries]
\item[Obligación] Después del cierre, la información explicativa de eficiencia debe permitir comprender que un vehículo competitivo valorado por debajo de la tendencia obtiene una diferencia favorable y que un precio superior a la tendencia reduce esa diferencia y puede reducir el premio de eficiencia.
\item[Prioridad] Alta
\item[Procedencia] DEC-161, DEC-167 y rationale del mecanismo de eficiencia; ID heredado: \texttt{RF-114}.
\item[Verificación] \texttt{AC-RF-EFF-018}
\end{description}
\vspace{0.25em}\hrule\vspace{0.65em}
\subsection{Patrocinio y aportaciones de evento}
\noindent\textbf{RF-SPN-001 --- Registro y cobro inmediato de aportación directa a la bolsa del evento}\par
\begin{description}[style=nextline,leftmargin=3.15cm,labelwidth=2.8cm,font=\normalfont\bfseries]
\item[Obligación] El sistema debe permitir registrar y cobrar de inmediato aportaciones directas destinadas a la bolsa de un evento específico solo mientras el evento permanezca en \texttt{inscripciones}.
\item[Prioridad] Alta
\item[Procedencia] Secciones 2.2, 2.6.2, 2.6.4.2 y decisiones prioritarias del cliente sobre patrocinio directo de eventos.; ID heredado: \texttt{RF-126}.
\item[Verificación] \texttt{AC-RF-SPN-001}
\end{description}
\vspace{0.25em}\hrule\vspace{0.65em}
\noindent\textbf{RF-SPN-002 --- Registro y cobro inmediato de aportación a la contribución de eficiencia de un participante del evento}\par
\begin{description}[style=nextline,leftmargin=3.15cm,labelwidth=2.8cm,font=\normalfont\bfseries]
\item[Obligación] El sistema debe permitir registrar y cobrar de inmediato aportaciones destinadas a la contribución de eficiencia de un participante específico dentro de un evento solo mientras el evento permanezca en \texttt{inscripciones} y solo cuando la función de aportaciones a participantes del evento se encuentre activa.
\item[Prioridad] Alta
\item[Procedencia] Secciones 2.6.2, 2.6.4.2 y decisiones prioritarias del cliente sobre aportaciones de eficiencia a participantes del evento.; ID heredado: \texttt{RF-127}.
\item[Verificación] \texttt{AC-RF-SPN-002}
\end{description}
\vspace{0.25em}\hrule\vspace{0.65em}
\noindent\textbf{RF-SPN-003 --- Registro y cobro inmediato de aportación a la contribución de terminación de un participante del evento}\par
\begin{description}[style=nextline,leftmargin=3.15cm,labelwidth=2.8cm,font=\normalfont\bfseries]
\item[Obligación] El sistema debe permitir registrar y cobrar de inmediato aportaciones destinadas a la contribución de terminación de una inscripción específica dentro de un evento solo mientras el evento permanezca en \texttt{inscripciones} y solo cuando la función de aportaciones a participantes del evento se encuentre activa.
\item[Prioridad] Alta
\item[Procedencia] Secciones 2.6.2, 2.6.4.2 y decisiones prioritarias del cliente sobre aportaciones de terminación a participantes del evento.; ID heredado: \texttt{RF-127A}.
\item[Verificación] \texttt{AC-RF-SPN-003}
\end{description}
\vspace{0.25em}\hrule\vspace{0.65em}
\noindent\textbf{RF-SPN-004 --- Registro de patrocinio asociado a aportación directa de evento}\par
\begin{description}[style=nextline,leftmargin=3.15cm,labelwidth=2.8cm,font=\normalfont\bfseries]
\item[Obligación] El sistema debe permitir asociar a una aportación directa a la bolsa del evento un anuncio con imagen para su publicación como patrocinio del evento.
\item[Prioridad] Alta
\item[Procedencia] Sección 2.6.4.2 y decisiones prioritarias del cliente sobre publicidad asociada al patrocinio directo.; ID heredado: \texttt{RF-128}.
\item[Verificación] \texttt{AC-RF-SPN-004}
\end{description}
\vspace{0.25em}\hrule\vspace{0.65em}
\noindent\textbf{RF-SPN-005 --- Publicación del patrocinio directo del evento}\par
\begin{description}[style=nextline,leftmargin=3.15cm,labelwidth=2.8cm,font=\normalfont\bfseries]
\item[Obligación] El sistema debe publicar el patrocinio directo asociado a la bolsa del evento.
\item[Prioridad] Alta
\item[Procedencia] Secciones 2.6.2, 2.6.4.2, RF-126, RF-127, RF-127A y RF-128.; ID heredado: \texttt{RF-129}.
\item[Verificación] \texttt{AC-RF-SPN-005}
\end{description}
\vspace{0.25em}\hrule\vspace{0.65em}
\noindent\textbf{RF-SPN-006 --- Publicación distinguida del patrocinio de campeonato aplicable}\par
\begin{description}[style=nextline,leftmargin=3.15cm,labelwidth=2.8cm,font=\normalfont\bfseries]
\item[Obligación] Cuando un evento pertenezca a campeonatos patrocinados, el sistema debe publicar el patrocinio aplicable por campeonato de forma distinguida del patrocinio directo del evento.
\item[Prioridad] Alta
\item[Procedencia] Secciones 2.6.2, 2.6.4.2, RF-126, RF-127, RF-127A y RF-128.; ID heredado: \texttt{RF-129}.
\item[Verificación] \texttt{AC-RF-SPN-006}
\end{description}
\vspace{0.25em}\hrule\vspace{0.65em}
\noindent\textbf{RF-SPN-007 --- Trazabilidad económica del patrocinio y aportaciones del evento}\par
\begin{description}[style=nextline,leftmargin=3.15cm,labelwidth=2.8cm,font=\normalfont\bfseries]
\item[Obligación] El sistema debe conservar la relación trazable entre aportaciones directas cobradas a la bolsa, anuncios publicados, aportaciones cobradas de eficiencia o terminación y distribuciones económicas correspondientes.
\item[Prioridad] Alta
\item[Procedencia] Secciones 2.6.2, 2.6.4.2, RF-126, RF-127, RF-127A y RF-128.; ID heredado: \texttt{RF-129}.
\item[Verificación] \texttt{AC-RF-SPN-007}
\end{description}
\vspace{0.25em}\hrule\vspace{0.65em}
\noindent\textbf{RF-SPN-008 --- Devolución inmediata de aportaciones del evento por cancelación del evento}\par
\begin{description}[style=nextline,leftmargin=3.15cm,labelwidth=2.8cm,font=\normalfont\bfseries]
\item[Obligación] Si el evento se cancela, el sistema debe devolver inmediatamente las aportaciones directas cobradas a la bolsa del evento, las aportaciones de eficiencia cobradas a participantes del evento y las aportaciones de terminación cobradas a inscripciones del evento.
\item[Prioridad] Alta
\item[Procedencia] Secciones 2.6.4 y decisiones prioritarias del cliente sobre devolución inmediata en cancelación del evento.; ID heredado: \texttt{RF-129A}.
\item[Verificación] \texttt{AC-RF-SPN-008}
\end{description}
\vspace{0.25em}\hrule\vspace{0.65em}
\noindent\textbf{RF-SPN-009 --- Devolución inmediata de aportaciones a participante por cancelación o descalificación previa a carrera}\par
\begin{description}[style=nextline,leftmargin=3.15cm,labelwidth=2.8cm,font=\normalfont\bfseries]
\item[Obligación] Si un participante del evento queda cancelado o descalificado antes de \texttt{carrera}, el sistema debe devolver inmediatamente las aportaciones de terceros cobradas a favor de ese participante tanto para eficiencia como para terminación dentro del evento.
\item[Prioridad] Alta
\item[Procedencia] Secciones 2.6.4 y decisiones prioritarias del cliente sobre devolución inmediata al cancelarse o descalificarse un participante antes de \texttt{carrera}.; ID heredado: \texttt{RF-129B}.
\item[Verificación] \texttt{AC-RF-SPN-009}
\end{description}
\vspace{0.25em}\hrule\vspace{0.65em}
\subsection{Puntos y actualización de clasificaciones}
\noindent\textbf{RF-PTS-001 --- Base de puntos de velocidad}\par
\begin{description}[style=nextline,leftmargin=3.15cm,labelwidth=2.8cm,font=\normalfont\bfseries]
\item[Obligación] Para cada campeonato aplicable en modalidad de velocidad, el sistema debe asignar la base de puntos conforme al ordenamiento oficial de velocidad y a la secuencia base vigente.
\item[Prioridad] Alta
\item[Procedencia] Secciones 2.6.3.1, 3.1.7, 3.1.16, fuente documental TOP\_Racing\_Grok\_Clarified.tex (líneas 167 a 171) y decisiones prioritarias del cliente sobre desempate.; ID heredado: \texttt{RF-130}.
\item[Verificación] \texttt{AC-RF-PTS-001}
\end{description}
\vspace{0.25em}\hrule\vspace{0.65em}
\noindent\textbf{RF-PTS-002 --- Desempate del orden de velocidad para puntos}\par
\begin{description}[style=nextline,leftmargin=3.15cm,labelwidth=2.8cm,font=\normalfont\bfseries]
\item[Obligación] Cuando dos participantes empaten en el ordenamiento oficial de velocidad, el sistema debe favorecer a quien se haya inscrito antes en el evento.
\item[Prioridad] Alta
\item[Procedencia] Secciones 2.6.3.1, 3.1.7, 3.1.16, fuente documental TOP\_Racing\_Grok\_Clarified.tex (líneas 167 a 171) y decisiones prioritarias del cliente sobre desempate.; ID heredado: \texttt{RF-130}.
\item[Verificación] \texttt{AC-RF-PTS-002}
\end{description}
\vspace{0.25em}\hrule\vspace{0.65em}
\noindent\textbf{RF-PTS-003 --- Base de puntos de carrera}\par
\begin{description}[style=nextline,leftmargin=3.15cm,labelwidth=2.8cm,font=\normalfont\bfseries]
\item[Obligación] Para cada campeonato aplicable en modalidad de carrera, el sistema debe asignar la base de puntos conforme al ordenamiento oficial de carrera y a la secuencia base vigente.
\item[Prioridad] Alta
\item[Procedencia] Secciones 2.6.3.1, 3.1.7, 3.1.16, fuente documental TOP\_Racing\_Grok\_Clarified.tex (líneas 167 a 171) y decisiones prioritarias del cliente sobre desempate.; ID heredado: \texttt{RF-131}.
\item[Verificación] \texttt{AC-RF-PTS-003}
\end{description}
\vspace{0.25em}\hrule\vspace{0.65em}
\noindent\textbf{RF-PTS-004 --- Desempate del orden de carrera para puntos}\par
\begin{description}[style=nextline,leftmargin=3.15cm,labelwidth=2.8cm,font=\normalfont\bfseries]
\item[Obligación] Cuando dos participantes empaten en el ordenamiento oficial de carrera, el sistema debe favorecer a quien se haya inscrito antes en el evento.
\item[Prioridad] Alta
\item[Procedencia] Secciones 2.6.3.1, 3.1.7, 3.1.16, fuente documental TOP\_Racing\_Grok\_Clarified.tex (líneas 167 a 171) y decisiones prioritarias del cliente sobre desempate.; ID heredado: \texttt{RF-131}.
\item[Verificación] \texttt{AC-RF-PTS-004}
\end{description}
\vspace{0.25em}\hrule\vspace{0.65em}
\noindent\textbf{RF-PTS-005 --- Base de puntos de eficiencia para campeonato}\par
\begin{description}[style=nextline,leftmargin=3.15cm,labelwidth=2.8cm,font=\normalfont\bfseries]
\item[Obligación] Para cada campeonato aplicable en modalidad de eficiencia, el sistema debe asignar la base de puntos conforme al ordenamiento oficial de eficiencia y a la secuencia base vigente.
\item[Prioridad] Alta
\item[Procedencia] Secciones 2.6.3.1, 3.1.7, 3.1.17, fuente documental TOP\_Racing\_Grok\_Clarified.tex (líneas 167 a 171 y 242) y decisiones prioritarias del cliente sobre desempate.; ID heredado: \texttt{RF-132}.
\item[Verificación] \texttt{AC-RF-PTS-005}
\end{description}
\vspace{0.25em}\hrule\vspace{0.65em}
\noindent\textbf{RF-PTS-006 --- Desempate del orden de eficiencia para campeonato}\par
\begin{description}[style=nextline,leftmargin=3.15cm,labelwidth=2.8cm,font=\normalfont\bfseries]
\item[Obligación] Cuando dos participantes empaten en el ordenamiento oficial de eficiencia, el sistema debe favorecer a quien se haya inscrito antes en el evento.
\item[Prioridad] Alta
\item[Procedencia] Secciones 2.6.3.1, 3.1.7, 3.1.17, fuente documental TOP\_Racing\_Grok\_Clarified.tex (líneas 167 a 171 y 242) y decisiones prioritarias del cliente sobre desempate.; ID heredado: \texttt{RF-132}.
\item[Verificación] \texttt{AC-RF-PTS-006}
\end{description}
\vspace{0.25em}\hrule\vspace{0.65em}
\noindent\textbf{RF-PTS-007 --- Escalamiento de puntos por modalidad en cada campeonato aplicable}\par
\begin{description}[style=nextline,leftmargin=3.15cm,labelwidth=2.8cm,font=\normalfont\bfseries]
\item[Obligación] El sistema debe determinar, para cada campeonato aplicable y para cada modalidad de velocidad, carrera y eficiencia, la puntuación escalada del evento aplicando a la base de puntos correspondiente el peso vigente del evento en ese campeonato, sin aplicar redondeo.
\item[Prioridad] Alta
\item[Procedencia] Secciones 2.6.3.1, RF-108B, RF-130, RF-131, RF-132 y decisiones prioritarias del cliente sobre ausencia de redondeo.; ID heredado: \texttt{RF-133}.
\item[Verificación] \texttt{AC-RF-PTS-007}
\end{description}
\vspace{0.25em}\hrule\vspace{0.65em}
\noindent\textbf{RF-PTS-008 --- Actualización de clasificaciones de velocidad, carrera y eficiencia}\par
\begin{description}[style=nextline,leftmargin=3.15cm,labelwidth=2.8cm,font=\normalfont\bfseries]
\item[Obligación] El sistema debe actualizar, para cada campeonato aplicable, las clasificaciones de velocidad, carrera y eficiencia agregando la puntuación escalada del evento a la clasificación correspondiente.
\item[Prioridad] Alta
\item[Procedencia] Secciones 2.6.3.1, 3.1.7, RF-133 y fuente documental TOP\_Racing\_Grok\_Clarified.tex (líneas 150, 167 y 171).; ID heredado: \texttt{RF-134}.
\item[Verificación] \texttt{AC-RF-PTS-008}
\end{description}
\vspace{0.25em}\hrule\vspace{0.65em}
\noindent\textbf{RF-PTS-009 --- Actualización de la clasificación combinada del campeonato}\par
\begin{description}[style=nextline,leftmargin=3.15cm,labelwidth=2.8cm,font=\normalfont\bfseries]
\item[Obligación] El sistema debe actualizar, para cada campeonato aplicable en modalidad combinada, la clasificación combinada agregando separadamente los puntos escalados del evento obtenidos en velocidad y en carrera.
\item[Prioridad] Alta
\item[Procedencia] Secciones 2.6.3.1, RF-133, fuente documental TOP\_Racing\_Grok\_Clarified.tex (líneas 167 a 171) y decisiones prioritarias del cliente sobre clasificación combinada.; ID heredado: \texttt{RF-135}.
\item[Verificación] \texttt{AC-RF-PTS-009}
\end{description}
\vspace{0.25em}\hrule\vspace{0.65em}
\noindent\textbf{RF-PTS-010 --- Exclusión de inscripciones canceladas de la actualización de puntos y clasificaciones}\par
\begin{description}[style=nextline,leftmargin=3.15cm,labelwidth=2.8cm,font=\normalfont\bfseries]
\item[Obligación] El sistema debe actualizar las clasificaciones y los puntos por modalidad excluyendo toda inscripción cancelada del evento.
\item[Prioridad] Alta
\item[Procedencia] RF-055A, RF-090, RF-130 a RF-135 y decisiones prioritarias del cliente sobre efectos de la cancelación.; ID heredado: \texttt{RF-136}.
\item[Verificación] \texttt{AC-RF-PTS-010}
\end{description}
\vspace{0.25em}\hrule\vspace{0.65em}
\noindent\textbf{RF-PTS-011 --- Publicación trazable de la actualización de puntos y clasificaciones por modalidad}\par
\begin{description}[style=nextline,leftmargin=3.15cm,labelwidth=2.8cm,font=\normalfont\bfseries]
\item[Obligación] El sistema debe publicar de forma trazable, para cada campeonato aplicable, la relación entre los resultados oficiales del evento, la base de puntos por modalidad, el peso vigente del evento, la puntuación escalada resultante y la actualización de las clasificaciones correspondientes.
\item[Prioridad] Alta
\item[Procedencia] Secciones 2.6.3.1, 3.1.7, RF-108B y RF-130 a RF-136.; ID heredado: \texttt{RF-137}.
\item[Verificación] \texttt{AC-RF-PTS-011}
\end{description}
\vspace{0.25em}\hrule\vspace{0.65em}
\subsection{Publicación, moderación y reserva}
\noindent\textbf{RF-GOV-001 --- Publicación transversal de la información declarada pública}\par
\begin{description}[style=nextline,leftmargin=3.15cm,labelwidth=2.8cm,font=\normalfont\bfseries]
\item[Obligación] El sistema debe publicar, en las consultas que correspondan, toda información que las reglas del dominio declaren visible para el público, con el alcance específico que corresponda a cada caso.
\item[Prioridad] Alta
\item[Procedencia] Secciones 2.2, 2.5, 2.7, 3.4 y decisiones prioritarias del cliente sobre publicidad del dominio.; ID heredado: \texttt{RF-142}.
\item[Verificación] \texttt{AC-RF-GOV-001}
\end{description}
\vspace{0.25em}\hrule\vspace{0.65em}
\noindent\textbf{RF-GOV-002 --- Reserva selectiva de identidades sujetas a restricción pública}\par
\begin{description}[style=nextline,leftmargin=3.15cm,labelwidth=2.8cm,font=\normalfont\bfseries]
\item[Obligación] El sistema debe mantener reservada en la consulta pública toda identidad que el dominio declare restringida y debe mostrarla solo a los papeles funcionales expresamente autorizados.
\item[Prioridad] Alta
\item[Procedencia] Secciones 2.5, 2.7, 3.4, RF-058A, RF-067C, RF-153 a RF-157 y decisiones prioritarias del cliente sobre reserva de identidad.; ID heredado: \texttt{RF-143}.
\item[Verificación] \texttt{AC-RF-GOV-002}
\end{description}
\vspace{0.25em}\hrule\vspace{0.65em}
\noindent\textbf{RF-GOV-003 --- Publicación anonimizada de terceros en distribuciones económicas}\par
\begin{description}[style=nextline,leftmargin=3.15cm,labelwidth=2.8cm,font=\normalfont\bfseries]
\item[Obligación] El sistema debe publicar sin nombres ni identidad las cantidades aportadas por terceros y la parte del premio que les corresponda cuando las reglas del dominio exijan anonimización de terceros en premios de evento o de campeonato.
\item[Prioridad] Alta
\item[Procedencia] Secciones 2.6.4.2, 2.6.4.3, 2.7, RIE-009, RIE-010 y decisiones prioritarias del cliente sobre publicación de premios.; ID heredado: \texttt{RF-144}.
\item[Verificación] \texttt{AC-RF-GOV-003}
\end{description}
\vspace{0.25em}\hrule\vspace{0.65em}
\noindent\textbf{RF-GOV-004 --- Moderación administrativa de espacios comunitarios}\par
\begin{description}[style=nextline,leftmargin=3.15cm,labelwidth=2.8cm,font=\normalfont\bfseries]
\item[Obligación] El sistema debe permitir que \texttt{admin} ejerza moderación sobre el foro comunitario y sobre el canal de novedades del ecosistema sin alterar por ello las demás reglas del dominio sobre publicidad, reserva o trazabilidad.
\item[Prioridad] Media
\item[Procedencia] Secciones 2.2, 2.3, 2.7, RF-063 a RF-066A y decisiones prioritarias del cliente sobre moderación.; ID heredado: \texttt{RF-145}.
\item[Verificación] \texttt{AC-RF-GOV-004}
\end{description}
\vspace{0.25em}\hrule\vspace{0.65em}
\noindent\textbf{RF-GOV-005 --- Permanencia histórica de publicaciones y procesos críticos visibles}\par
\begin{description}[style=nextline,leftmargin=3.15cm,labelwidth=2.8cm,font=\normalfont\bfseries]
\item[Obligación] El sistema debe conservar dentro del historial correspondiente toda publicación y todo proceso crítico visible cuya permanencia histórica exijan las reglas del dominio.
\item[Prioridad] Alta
\item[Procedencia] Secciones 2.2, 2.7, RF-067 a RF-070A, RF-101, RF-102, RF-129 y decisiones prioritarias del cliente sobre permanencia histórica.; ID heredado: \texttt{RF-146}.
\item[Verificación] \texttt{AC-RF-GOV-005}
\end{description}
\vspace{0.25em}\hrule\vspace{0.65em}
\subsection{Reputación del promotor}
\noindent\textbf{RF-REP-001 --- Registro de calificación del promotor por participantes}\par
\begin{description}[style=nextline,leftmargin=3.15cm,labelwidth=2.8cm,font=\normalfont\bfseries]
\item[Obligación] El sistema debe permitir que los participantes de un evento registren una calificación del promotor después del evento.
\item[Prioridad] Media
\item[Procedencia] Secciones 1.2, 2.2 y fuente documental \texttt{TOP\_Racing\_Grok\_Clarified.tex} (línea 133).; ID heredado: \texttt{RF-158}.
\item[Verificación] \texttt{AC-RF-REP-001}
\end{description}
\vspace{0.25em}\hrule\vspace{0.65em}
\noindent\textbf{RF-REP-002 --- Registro de calificación del promotor por espectadores}\par
\begin{description}[style=nextline,leftmargin=3.15cm,labelwidth=2.8cm,font=\normalfont\bfseries]
\item[Obligación] El sistema debe permitir que los espectadores de un evento registren una calificación del promotor después del evento.
\item[Prioridad] Media
\item[Procedencia] Secciones 1.2, 2.2 y fuente documental \texttt{TOP\_Racing\_Grok\_Clarified.tex} (línea 133).; ID heredado: \texttt{RF-159}.
\item[Verificación] \texttt{AC-RF-REP-002}
\end{description}
\vspace{0.25em}\hrule\vspace{0.65em}
\noindent\textbf{RF-REP-003 --- Determinación de la reputación del promotor}\par
\begin{description}[style=nextline,leftmargin=3.15cm,labelwidth=2.8cm,font=\normalfont\bfseries]
\item[Obligación] El sistema debe determinar y conservar la reputación del promotor a partir de las calificaciones del promotor registradas en eventos anteriores.
\item[Prioridad] Media
\item[Procedencia] Secciones 2.2 y fuente documental \texttt{TOP\_Racing\_Grok\_Clarified.tex} (línea 133).; ID heredado: \texttt{RF-160}.
\item[Verificación] \texttt{AC-RF-REP-003}
\end{description}
\vspace{0.25em}\hrule\vspace{0.65em}
\noindent\textbf{RF-REP-004 --- Aplicación de la reputación del promotor a eventos futuros}\par
\begin{description}[style=nextline,leftmargin=3.15cm,labelwidth=2.8cm,font=\normalfont\bfseries]
\item[Obligación] El sistema debe vincular la reputación vigente del promotor con los eventos futuros organizados por ese promotor.
\item[Prioridad] Media
\item[Procedencia] Secciones 1.2, 2.2 y fuente documental \texttt{TOP\_Racing\_Grok\_Clarified.tex} (línea 133).; ID heredado: \texttt{RF-161}.
\item[Verificación] \texttt{AC-RF-REP-004}
\end{description}
\vspace{0.25em}\hrule\vspace{0.65em}
\subsection{Modelo de negocio}
\noindent\textbf{RF-BUS-001 --- Configuración administrativa de parámetros de comisión y cuota fija}\par
\begin{description}[style=nextline,leftmargin=3.15cm,labelwidth=2.8cm,font=\normalfont\bfseries]
\item[Obligación] El sistema debe permitir que \texttt{admin} configure los parámetros de comisión del competidor por evento, del competidor por campeonato, del aportante por evento, del aportante por campeonato, del comprador de vehículo adjudicado, la cuota fija del promotor por evento, el umbral anual de gratuidad del promotor y los estados globales de activación o inactivación de aportaciones a participantes para eventos y campeonatos.
\item[Prioridad] Alta
\item[Procedencia] Secciones 2.2, 2.6.2.2 y decisiones sobre configuración administrativa del modelo de negocio; ID heredado: \texttt{RF-162}.
\item[Verificación] \texttt{AC-RF-BUS-001}
\end{description}
\vspace{0.25em}\hrule\vspace{0.65em}
\noindent\textbf{RF-BUS-002 --- Base positiva de comisión del competidor por evento}\par
\begin{description}[style=nextline,leftmargin=3.15cm,labelwidth=2.8cm,font=\normalfont\bfseries]
\item[Obligación] El sistema debe calcular la comisión del competidor por evento exclusivamente sobre la base positiva definida por CALC-BUS-001, usando como $R$ el premio individual del evento y como $C$ los pagos del dominio atribuibles al mismo evento, excluidas las compras de vehículos.
\item[Prioridad] Alta
\item[Procedencia] Secciones 2.6.2.2, 2.6.4.2 y decisiones prioritarias del cliente sobre comisión del competidor por evento.; ID heredado: \texttt{RF-163}.
\item[Verificación] \texttt{AC-RF-BUS-002}
\end{description}
\vspace{0.25em}\hrule\vspace{0.65em}
\noindent\textbf{RF-BUS-003 --- Base positiva de comisión del competidor por campeonato}\par
\begin{description}[style=nextline,leftmargin=3.15cm,labelwidth=2.8cm,font=\normalfont\bfseries]
\item[Obligación] El sistema debe calcular la comisión del competidor por campeonato exclusivamente sobre la base positiva definida por CALC-BUS-001, usando como $R$ el premio individual del campeonato y como $C$ las aportaciones computables del propio competidor en ese campeonato.
\item[Prioridad] Alta
\item[Procedencia] Secciones 2.6.2.2, 2.6.4.3 y decisiones prioritarias del cliente sobre comisión del competidor por campeonato.; ID heredado: \texttt{RF-164}.
\item[Verificación] \texttt{AC-RF-BUS-003}
\end{description}
\vspace{0.25em}\hrule\vspace{0.65em}
\noindent\textbf{RF-BUS-004 --- Base positiva de comisión del aportante asociado a un participante}\par
\begin{description}[style=nextline,leftmargin=3.15cm,labelwidth=2.8cm,font=\normalfont\bfseries]
\item[Obligación] El sistema debe calcular por separado en eventos y campeonatos la comisión del aportante asociado a un participante exclusivamente sobre la base positiva definida por CALC-BUS-001, usando como $R$ la parte de premio del aportante y como $C$ su aportación en el mismo caso.
\item[Prioridad] Alta
\item[Procedencia] Secciones 2.6.2.2, 2.6.4.2, 2.6.4.3 y decisiones prioritarias del cliente sobre comisión del aportante.; ID heredado: \texttt{RF-165}.
\item[Verificación] \texttt{AC-RF-BUS-004}
\end{description}
\vspace{0.25em}\hrule\vspace{0.65em}
\noindent\textbf{RF-BUS-005 --- Determinación de la comisión del comprador de vehículo adjudicado}\par
\begin{description}[style=nextline,leftmargin=3.15cm,labelwidth=2.8cm,font=\normalfont\bfseries]
\item[Obligación] El sistema debe calcular la comisión aplicable al comprador sobre toda adjudicación oficial de vehículo en subasta, usando como base el precio oficial, incluso cuando el comprador adjudicatario sea el dueño previo, sin cargar esa comisión al vendedor en su condición de vendedor.
\item[Prioridad] Alta
\item[Procedencia] Secciones 2.6.2.2, 3.1.6 y decisiones confirmadas sobre comisión a cargo del comprador; ID heredado: \texttt{RF-166}.
\item[Verificación] \texttt{AC-RF-BUS-005}
\end{description}
\vspace{0.25em}\hrule\vspace{0.65em}
\noindent\textbf{RF-BUS-006 --- Exención de comisión para patrocinadores directos}\par
\begin{description}[style=nextline,leftmargin=3.15cm,labelwidth=2.8cm,font=\normalfont\bfseries]
\item[Obligación] El sistema no debe cobrar comisión al patrocinador por las aportaciones directas a la bolsa de un evento o de un campeonato.
\item[Prioridad] Alta
\item[Procedencia] Secciones 2.6.2.2, 3.1.15, 3.1.20 y decisiones prioritarias del cliente sobre patrocinadores sin comisión.; ID heredado: \texttt{RF-167}.
\item[Verificación] \texttt{AC-RF-BUS-006}
\end{description}
\vspace{0.25em}\hrule\vspace{0.65em}
\noindent\textbf{RF-BUS-007 --- Cobro de cuota fija al promotor a partir del umbral anual configurado}\par
\begin{description}[style=nextline,leftmargin=3.15cm,labelwidth=2.8cm,font=\normalfont\bfseries]
\item[Obligación] El sistema debe cobrar al promotor la cuota fija vigente al crear cada evento nuevo una vez superado o alcanzado, según la configuración vigente, el umbral anual de gratuidad aplicable a ese promotor.
\item[Prioridad] Alta
\item[Procedencia] Secciones 2.6.2.2 y decisiones prioritarias del cliente sobre cuota fija del promotor.; ID heredado: \texttt{RF-168}.
\item[Verificación] \texttt{AC-RF-BUS-007}
\end{description}
\vspace{0.25em}\hrule\vspace{0.65em}
\noindent\textbf{RF-BUS-008 --- Cobro automático de comisiones y cuotas del modelo de negocio}\par
\begin{description}[style=nextline,leftmargin=3.15cm,labelwidth=2.8cm,font=\normalfont\bfseries]
\item[Obligación] El sistema debe cobrar automáticamente las comisiones y cuotas que correspondan según la configuración vigente, incluida la comisión del comprador en una autoadjudicación, dentro del flujo económico real en el que se produzca el beneficio, adjudicación u operación gravada.
\item[Prioridad] Alta
\item[Procedencia] Secciones 2.6.2.2, 2.6.4 y decisiones sobre pagos reales y comisiones; ID heredado: \texttt{RF-169}.
\item[Verificación] \texttt{AC-RF-BUS-008}
\end{description}
\vspace{0.25em}\hrule\vspace{0.65em}
\noindent\textbf{RF-BUS-009 --- Información previa de comisión o cuota aplicable}\par
\begin{description}[style=nextline,leftmargin=3.15cm,labelwidth=2.8cm,font=\normalfont\bfseries]
\item[Obligación] Antes de que un usuario participe en una operación que pueda causar comisión o cuota, el sistema debe informar la regla aplicable y el parámetro vigente; cuando el importe dependa de un resultado posterior de subasta, el sistema debe informar previamente la fórmula y publicar el importe efectivo después del cierre.
\item[Prioridad] Alta
\item[Procedencia] Secciones 2.6.2.2, 3.2 y decisiones sobre publicación de comisiones; ID heredado: \texttt{RF-170}.
\item[Verificación] \texttt{AC-RF-BUS-009}
\end{description}
\vspace{0.25em}\hrule\vspace{0.65em}
\noindent\textbf{RF-BUS-010 --- Publicación posterior del detalle de comisión o cuota cobrada}\par
\begin{description}[style=nextline,leftmargin=3.15cm,labelwidth=2.8cm,font=\normalfont\bfseries]
\item[Obligación] Una vez ejecutado el cobro de una comisión o cuota del modelo de negocio, el sistema debe poner a disposición del usuario afectado el detalle correspondiente, incluyendo la base gravable o el hecho generador, el parámetro aplicado, el importe cobrado y el importe neto resultante.
\item[Prioridad] Alta
\item[Procedencia] Secciones 2.6.2.2, 2.6.4 y decisiones prioritarias del cliente sobre detalle posterior de comisiones.; ID heredado: \texttt{RF-171}.
\item[Verificación] \texttt{AC-RF-BUS-010}
\end{description}
\vspace{0.25em}\hrule\vspace{0.65em}
\noindent\textbf{RF-BUS-011 --- Información expresa de exención de comisión cuando corresponda}\par
\begin{description}[style=nextline,leftmargin=3.15cm,labelwidth=2.8cm,font=\normalfont\bfseries]
\item[Obligación] Cuando una operación relevante para el modelo de negocio esté expresamente exenta de comisión, el sistema debe informar esa exención al usuario afectado al momento de la operación.
\item[Prioridad] Media
\item[Procedencia] Secciones 2.6.2.2 y decisiones prioritarias del cliente sobre vendedor y patrocinador sin comisión.; ID heredado: \texttt{RF-172}.
\item[Verificación] \texttt{AC-RF-BUS-011}
\end{description}
\vspace{0.25em}\hrule\vspace{0.65em}
\subsection{Requisitos de interfaz externa}
Los requisitos UI especifican comportamiento observable de interacción. No prescriben componentes gráficos, framework ni disposición exacta.

\noindent\textbf{UI-001 --- Acceso a funciones según habilitación del usuario}\par
\begin{description}[style=nextline,leftmargin=3.15cm,labelwidth=2.8cm,font=\normalfont\bfseries]
\item[Obligación] La interfaz principal debe exponer a cada usuario únicamente las funciones operativas para las que cumpla las condiciones de participación y de papel funcional aplicables, sin impedir la consulta de funciones públicas.
\item[Prioridad] Alta
\item[Procedencia] Secciones 1.2, 2.1, 2.2, 2.3, 2.4, 3.1.11, 3.1.24, 3.1.25 y decisiones prioritarias del cliente sobre página principal, niveles de verificación y aceptación legal previa y versionada.; ID heredado: \texttt{RIE-001}.
\item[Verificación] \texttt{AC-UI-001}
\end{description}
\vspace{0.25em}\hrule\vspace{0.65em}
\noindent\textbf{UI-002 --- Consulta personal de estados y nivel visible de acceso}\par
\begin{description}[style=nextline,leftmargin=3.15cm,labelwidth=2.8cm,font=\normalfont\bfseries]
\item[Obligación] La interfaz principal debe permitir al usuario consultar sus estados booleanos de identidad, verificación reforzada y exclusión, así como su nivel visible de acceso \texttt{base}, \texttt{habilitado} o \texttt{extendido}.
\item[Prioridad] Alta
\item[Procedencia] Secciones 1.2, 2.1, 2.2, 2.3, 2.4, 3.1.11, 3.1.24, 3.1.25 y decisiones prioritarias del cliente sobre página principal, niveles de verificación y aceptación legal previa y versionada.; ID heredado: \texttt{RIE-001}.
\item[Verificación] \texttt{AC-UI-002}
\end{description}
\vspace{0.25em}\hrule\vspace{0.65em}
\noindent\textbf{UI-003 --- Explicación de condiciones faltantes de acceso}\par
\begin{description}[style=nextline,leftmargin=3.15cm,labelwidth=2.8cm,font=\normalfont\bfseries]
\item[Obligación] Cuando un usuario no alcance el siguiente nivel visible de acceso, la interfaz principal debe indicar las condiciones faltantes y debe mostrar el recordatorio aplicable al intentar una operación bloqueada por nivel.
\item[Prioridad] Alta
\item[Procedencia] Secciones 1.2, 2.1, 2.2, 2.3, 2.4, 3.1.11, 3.1.24, 3.1.25 y decisiones prioritarias del cliente sobre página principal, niveles de verificación y aceptación legal previa y versionada.; ID heredado: \texttt{RIE-001}.
\item[Verificación] \texttt{AC-UI-003}
\end{description}
\vspace{0.25em}\hrule\vspace{0.65em}
\noindent\textbf{UI-004 --- Consulta de campeonatos patrocinados desde la interfaz principal}\par
\begin{description}[style=nextline,leftmargin=3.15cm,labelwidth=2.8cm,font=\normalfont\bfseries]
\item[Obligación] La interfaz principal debe permitir consultar los campeonatos patrocinados y sus patrocinadores asociados.
\item[Prioridad] Alta
\item[Procedencia] Secciones 1.2, 2.1, 2.2, 2.3, 2.4, 3.1.11, 3.1.24, 3.1.25 y decisiones prioritarias del cliente sobre página principal, niveles de verificación y aceptación legal previa y versionada.; ID heredado: \texttt{RIE-001}.
\item[Verificación] \texttt{AC-UI-004}
\end{description}
\vspace{0.25em}\hrule\vspace{0.65em}
\noindent\textbf{UI-005 --- Cancelación inmediata de sesión por exclusión}\par
\begin{description}[style=nextline,leftmargin=3.15cm,labelwidth=2.8cm,font=\normalfont\bfseries]
\item[Obligación] Cuando se active la exclusión de un usuario con una sesión en curso, el sistema debe cancelar esa sesión.
\item[Prioridad] Alta
\item[Procedencia] Secciones 1.2, 2.1, 2.2, 2.3, 2.4, 3.1.11, 3.1.24, 3.1.25 y decisiones prioritarias del cliente sobre página principal, niveles de verificación y aceptación legal previa y versionada.; ID heredado: \texttt{RIE-001}.
\item[Verificación] \texttt{AC-UI-005}
\end{description}
\vspace{0.25em}\hrule\vspace{0.65em}
\noindent\textbf{UI-006 --- Bloqueo de inicio de sesión de usuario excluido}\par
\begin{description}[style=nextline,leftmargin=3.15cm,labelwidth=2.8cm,font=\normalfont\bfseries]
\item[Obligación] El sistema debe rechazar el inicio de sesión de un usuario con exclusión activa y debe mostrar el mensaje de exclusión correspondiente.
\item[Prioridad] Alta
\item[Procedencia] Secciones 1.2, 2.1, 2.2, 2.3, 2.4, 3.1.11, 3.1.24, 3.1.25 y decisiones prioritarias del cliente sobre página principal, niveles de verificación y aceptación legal previa y versionada.; ID heredado: \texttt{RIE-001}.
\item[Verificación] \texttt{AC-UI-006}
\end{description}
\vspace{0.25em}\hrule\vspace{0.65em}
\noindent\textbf{UI-007 --- Interfaz para convocar y registrar eventos}\par
\begin{description}[style=nextline,leftmargin=3.15cm,labelwidth=2.8cm,font=\normalfont\bfseries]
\item[Obligación] La interfaz de convocatoria del evento debe permitir capturar y publicar, como mínimo, la fecha, la ubicación, el nivel territorial del evento, el nivel temporal del evento, el tipo de pista, el número esperado de participantes, la cuota de inscripción, el pozo inicial, el pozo del premio de terminación, el costo de renta de pista cuando corresponda, el costo fijo por pesaje y, cuando el evento contemple subasta de vehículos, el plazo máximo para iniciar disputas de entrega de vehículos adjudicados y el plazo máximo para resolverlas.
\item[Prioridad] Alta
\item[Procedencia] Secciones 1.2, 2.2, 2.3, 2.6.1, fuente documental \texttt{TOP\_Racing\_Grok\_Clarified.tex} (línea 144) y decisiones prioritarias del cliente sobre nivel temporal, convocatoria en \texttt{inscripciones} y plazos de disputa de entrega.; ID heredado: \texttt{RIE-002}.
\item[Verificación] \texttt{AC-UI-007}
\end{description}
\vspace{0.25em}\hrule\vspace{0.65em}
\noindent\textbf{UI-008 --- Captura de datos mínimos de inscripción}\par
\begin{description}[style=nextline,leftmargin=3.15cm,labelwidth=2.8cm,font=\normalfont\bfseries]
\item[Obligación] La interfaz de inscripción debe permitir capturar, como mínimo, el vehículo inscrito, el usuario piloto asociado cuando exista pilotaje humano directo, la aportación de eficiencia, la contribución de terminación y los importes de inscripción y renta de pista aplicables.
\item[Prioridad] Alta
\item[Procedencia] Secciones 1.2, 2.2, 2.3, 2.6.2, 3.1.13 y fuente documental \texttt{TOP\_Racing\_Grok\_Clarified.tex} (línea 144), además de decisiones prioritarias del cliente sobre cobro inmediato de la inscripción.; ID heredado: \texttt{RIE-003}.
\item[Verificación] \texttt{AC-UI-008}
\end{description}
\vspace{0.25em}\hrule\vspace{0.65em}
\noindent\textbf{UI-009 --- Visualización del cobro inmediato de la inscripción}\par
\begin{description}[style=nextline,leftmargin=3.15cm,labelwidth=2.8cm,font=\normalfont\bfseries]
\item[Obligación] Al aceptar una inscripción, la interfaz debe reflejar el cobro inmediato de los importes aplicables a esa inscripción.
\item[Prioridad] Alta
\item[Procedencia] Secciones 1.2, 2.2, 2.3, 2.6.2, 3.1.13 y fuente documental \texttt{TOP\_Racing\_Grok\_Clarified.tex} (línea 144), además de decisiones prioritarias del cliente sobre cobro inmediato de la inscripción.; ID heredado: \texttt{RIE-003}.
\item[Verificación] \texttt{AC-UI-009}
\end{description}
\vspace{0.25em}\hrule\vspace{0.65em}
\noindent\textbf{UI-010 --- Publicación de la cantidad de accesos disponibles}\par
\begin{description}[style=nextline,leftmargin=3.15cm,labelwidth=2.8cm,font=\normalfont\bfseries]
\item[Obligación] La interfaz de subasta de acceso debe publicar la cantidad de accesos de espectadores disponibles para el evento.
\item[Prioridad] Alta
\item[Procedencia] Secciones 1.2, 2.2, 2.3, 3.1.6 y fuente documental \texttt{TOP\_Racing\_Grok\_Clarified.tex} (líneas 146 y 219), además de decisiones prioritarias del cliente sobre QR, actualización al alza y pagos reales con retención.; ID heredado: \texttt{RIE-004}.
\item[Verificación] \texttt{AC-UI-010}
\end{description}
\vspace{0.25em}\hrule\vspace{0.65em}
\noindent\textbf{UI-011 --- Captura de pujas secretas de acceso con respaldo suficiente}\par
\begin{description}[style=nextline,leftmargin=3.15cm,labelwidth=2.8cm,font=\normalfont\bfseries]
\item[Obligación] La interfaz de subasta de acceso debe aceptar una puja únicamente cuando la retención económica suficiente exigida por las reglas vigentes pueda constituirse.
\item[Prioridad] Alta
\item[Procedencia] Secciones 1.2, 2.2, 2.3, 3.1.6 y fuente documental \texttt{TOP\_Racing\_Grok\_Clarified.tex} (líneas 146 y 219), además de decisiones prioritarias del cliente sobre QR, actualización al alza y pagos reales con retención.; ID heredado: \texttt{RIE-004}.
\item[Verificación] \texttt{AC-UI-011}
\end{description}
\vspace{0.25em}\hrule\vspace{0.65em}
\noindent\textbf{UI-012 --- Unicidad de la puja vigente de acceso por usuario}\par
\begin{description}[style=nextline,leftmargin=3.15cm,labelwidth=2.8cm,font=\normalfont\bfseries]
\item[Obligación] Mientras la subasta de acceso permanezca abierta, la interfaz debe mantener como máximo una puja vigente de acceso por usuario.
\item[Prioridad] Alta
\item[Procedencia] Secciones 1.2, 2.2, 2.3, 3.1.6 y fuente documental \texttt{TOP\_Racing\_Grok\_Clarified.tex} (líneas 146 y 219), además de decisiones prioritarias del cliente sobre QR, actualización al alza y pagos reales con retención.; ID heredado: \texttt{RIE-004}.
\item[Verificación] \texttt{AC-UI-012}
\end{description}
\vspace{0.25em}\hrule\vspace{0.65em}
\noindent\textbf{UI-013 --- Actualización solo al alza de la puja de acceso}\par
\begin{description}[style=nextline,leftmargin=3.15cm,labelwidth=2.8cm,font=\normalfont\bfseries]
\item[Obligación] Mientras la subasta de acceso permanezca abierta, la interfaz debe permitir actualizar una puja vigente únicamente a un importe superior.
\item[Prioridad] Alta
\item[Procedencia] Secciones 1.2, 2.2, 2.3, 3.1.6 y fuente documental \texttt{TOP\_Racing\_Grok\_Clarified.tex} (líneas 146 y 219), además de decisiones prioritarias del cliente sobre QR, actualización al alza y pagos reales con retención.; ID heredado: \texttt{RIE-004}.
\item[Verificación] \texttt{AC-UI-013}
\end{description}
\vspace{0.25em}\hrule\vspace{0.65em}
\noindent\textbf{UI-014 --- Conteo público de pujas de acceso sin importes individuales}\par
\begin{description}[style=nextline,leftmargin=3.15cm,labelwidth=2.8cm,font=\normalfont\bfseries]
\item[Obligación] Mientras la subasta de acceso permanezca abierta, la interfaz debe mostrar el número de pujas vigentes sin revelar los importes individuales de terceros.
\item[Prioridad] Alta
\item[Procedencia] Secciones 1.2, 2.2, 2.3, 3.1.6 y fuente documental \texttt{TOP\_Racing\_Grok\_Clarified.tex} (líneas 146 y 219), además de decisiones prioritarias del cliente sobre QR, actualización al alza y pagos reales con retención.; ID heredado: \texttt{RIE-004}.
\item[Verificación] \texttt{AC-UI-014}
\end{description}
\vspace{0.25em}\hrule\vspace{0.65em}
\noindent\textbf{UI-015 --- Publicación del resultado de acceso al cierre}\par
\begin{description}[style=nextline,leftmargin=3.15cm,labelwidth=2.8cm,font=\normalfont\bfseries]
\item[Obligación] Al cerrar la subasta de acceso, la interfaz debe publicar los accesos adjudicados, las personas ganadoras y el importe adjudicado de cada acceso.
\item[Prioridad] Alta
\item[Procedencia] Secciones 1.2, 2.2, 2.3, 3.1.6 y fuente documental \texttt{TOP\_Racing\_Grok\_Clarified.tex} (líneas 146 y 219), además de decisiones prioritarias del cliente sobre QR, actualización al alza y pagos reales con retención.; ID heredado: \texttt{RIE-004}.
\item[Verificación] \texttt{AC-UI-015}
\end{description}
\vspace{0.25em}\hrule\vspace{0.65em}
\noindent\textbf{UI-016 --- Cobro automático de accesos adjudicados}\par
\begin{description}[style=nextline,leftmargin=3.15cm,labelwidth=2.8cm,font=\normalfont\bfseries]
\item[Obligación] Al cerrar la subasta de acceso, el sistema debe cobrar automáticamente a cada persona ganadora el importe adjudicado aplicable.
\item[Prioridad] Alta
\item[Procedencia] Secciones 1.2, 2.2, 2.3, 3.1.6 y fuente documental \texttt{TOP\_Racing\_Grok\_Clarified.tex} (líneas 146 y 219), además de decisiones prioritarias del cliente sobre QR, actualización al alza y pagos reales con retención.; ID heredado: \texttt{RIE-004}.
\item[Verificación] \texttt{AC-UI-016}
\end{description}
\vspace{0.25em}\hrule\vspace{0.65em}
\noindent\textbf{UI-017 --- Liberación de retenciones de pujas de acceso no ganadoras}\par
\begin{description}[style=nextline,leftmargin=3.15cm,labelwidth=2.8cm,font=\normalfont\bfseries]
\item[Obligación] Al cerrar la subasta de acceso, el sistema debe liberar las retenciones asociadas a las pujas no ganadoras.
\item[Prioridad] Alta
\item[Procedencia] Secciones 1.2, 2.2, 2.3, 3.1.6 y fuente documental \texttt{TOP\_Racing\_Grok\_Clarified.tex} (líneas 146 y 219), además de decisiones prioritarias del cliente sobre QR, actualización al alza y pagos reales con retención.; ID heredado: \texttt{RIE-004}.
\item[Verificación] \texttt{AC-UI-017}
\end{description}
\vspace{0.25em}\hrule\vspace{0.65em}
\noindent\textbf{UI-018 --- Disponibilidad de la notificación QR del acceso ganado}\par
\begin{description}[style=nextline,leftmargin=3.15cm,labelwidth=2.8cm,font=\normalfont\bfseries]
\item[Obligación] Después de adjudicar un acceso, la interfaz debe poner a disposición de la persona ganadora la notificación con el QR correspondiente.
\item[Prioridad] Alta
\item[Procedencia] Secciones 1.2, 2.2, 2.3, 3.1.6 y fuente documental \texttt{TOP\_Racing\_Grok\_Clarified.tex} (líneas 146 y 219), además de decisiones prioritarias del cliente sobre QR, actualización al alza y pagos reales con retención.; ID heredado: \texttt{RIE-004}.
\item[Verificación] \texttt{AC-UI-018}
\end{description}
\vspace{0.25em}\hrule\vspace{0.65em}
\noindent\textbf{UI-019 --- Listado de vehículos computables en la subasta}\par
\begin{description}[style=nextline,leftmargin=3.15cm,labelwidth=2.8cm,font=\normalfont\bfseries]
\item[Obligación] La interfaz de subasta de vehículos debe listar los vehículos computables del evento que hayan sido entregados al promotor, junto con los datos públicos asociados a su inscripción.
\item[Prioridad] Alta
\item[Procedencia] Secciones 1.2, 2.2, 2.3, 2.4, 3.1.6 y fuente documental \texttt{TOP\_Racing\_Grok\_Clarified.tex} (líneas 146, 221 y 223), además de decisiones prioritarias del cliente sobre garantía real de pago, custodia del promotor y transferencia obligatoria del vehículo subastado.; ID heredado: \texttt{RIE-005}.
\item[Verificación] \texttt{AC-UI-019}
\end{description}
\vspace{0.25em}\hrule\vspace{0.65em}
\noindent\textbf{UI-020 --- Captura de pujas selladas con la regla de retención aplicable}\par
\begin{description}[style=nextline,leftmargin=3.15cm,labelwidth=2.8cm,font=\normalfont\bfseries]
\item[Obligación] La interfaz de subasta debe permitir registrar una puja únicamente cuando se cumpla la regla económica de respaldo aplicable al postor: retención suficiente del principal para usuarios distintos del dueño del vehículo y ausencia de retención del principal para la puja del dueño sobre su propio vehículo.
\item[Prioridad] Alta
\item[Procedencia] RIE-005, BR-AUC-006, DEC-165 y reglas económicas de retención; ID heredado: \texttt{RIE-005}.
\item[Verificación] \texttt{AC-UI-020}
\end{description}
\vspace{0.25em}\hrule\vspace{0.65em}
\noindent\textbf{UI-021 --- Reserva total de información agregada de pujas antes del cierre}\par
\begin{description}[style=nextline,leftmargin=3.15cm,labelwidth=2.8cm,font=\normalfont\bfseries]
\item[Obligación] Mientras la subasta de un vehículo permanezca abierta, la interfaz no debe mostrar conteos de pujas, importes de terceros, orden relativo, precio provisional, adjudicatario provisional ni indicadores derivados de las pujas.
\item[Prioridad] Alta
\item[Procedencia] RIE-005 y DEC-162; decisión confirmada de no publicar información antes del cierre; ID heredado: \texttt{RIE-005}.
\item[Verificación] \texttt{AC-UI-021}
\end{description}
\vspace{0.25em}\hrule\vspace{0.65em}
\noindent\textbf{UI-022 --- Consulta y aumento de la puja propia vigente}\par
\begin{description}[style=nextline,leftmargin=3.15cm,labelwidth=2.8cm,font=\normalfont\bfseries]
\item[Obligación] La interfaz debe permitir que cada usuario consulte su propia puja vigente y la sustituya únicamente por un importe estrictamente mayor, sin acceso a pujas individuales o información relativa de otros usuarios.
\item[Prioridad] Alta
\item[Procedencia] RIE-005 y DEC-163; regla de puja sellada creciente; ID heredado: \texttt{RIE-005}.
\item[Verificación] \texttt{AC-UI-022}
\end{description}
\vspace{0.25em}\hrule\vspace{0.65em}
\noindent\textbf{UI-023 --- Ausencia de eficiencia, puntos y premios provisionales durante la subasta}\par
\begin{description}[style=nextline,leftmargin=3.15cm,labelwidth=2.8cm,font=\normalfont\bfseries]
\item[Obligación] Mientras la subasta de vehículos permanezca abierta, la interfaz no debe mostrar clasificación de eficiencia, puntos de eficiencia ni premio de eficiencia calculados a partir de pujas vigentes.
\item[Prioridad] Alta
\item[Procedencia] RIE-005 y DEC-162; decisión confirmada de publicación solo al cierre; ID heredado: \texttt{RIE-005}.
\item[Verificación] \texttt{AC-UI-023}
\end{description}
\vspace{0.25em}\hrule\vspace{0.65em}
\noindent\textbf{UI-024 --- Publicación posterior del resultado oficial y su naturaleza}\par
\begin{description}[style=nextline,leftmargin=3.15cm,labelwidth=2.8cm,font=\normalfont\bfseries]
\item[Obligación] Después del cierre de la fase de subasta de vehículos, la interfaz debe mostrar para cada vehículo el precio oficial, el adjudicatario y si será transferido a un tercero o conservado por autoadjudicación del dueño previo.
\item[Prioridad] Alta
\item[Procedencia] RIE-005 y DEC-166; resultado oficial de la subasta de vehículos; ID heredado: \texttt{RIE-005}.
\item[Verificación] \texttt{AC-UI-024}
\end{description}
\vspace{0.25em}\hrule\vspace{0.65em}
\noindent\textbf{UI-025 --- Publicación del precio final y ganador al cierre}\par
\begin{description}[style=nextline,leftmargin=3.15cm,labelwidth=2.8cm,font=\normalfont\bfseries]
\item[Obligación] Al cerrar la subasta, la interfaz debe publicar el precio final y el ganador vigente de cada vehículo adjudicado.
\item[Prioridad] Alta
\item[Procedencia] Secciones 1.2, 2.2, 2.3, 2.4, 3.1.6 y fuente documental \texttt{TOP\_Racing\_Grok\_Clarified.tex} (líneas 146, 221 y 223), además de decisiones prioritarias del cliente sobre garantía real de pago, custodia del promotor y transferencia obligatoria del vehículo subastado.; ID heredado: \texttt{RIE-005}.
\item[Verificación] \texttt{AC-UI-025}
\end{description}
\vspace{0.25em}\hrule\vspace{0.65em}
\noindent\textbf{UI-026 --- Información sobre transferencia obligatoria}\par
\begin{description}[style=nextline,leftmargin=3.15cm,labelwidth=2.8cm,font=\normalfont\bfseries]
\item[Obligación] La interfaz debe mostrar que el vehículo subastado queda sujeto a transferencia obligatoria conforme al resultado oficial vigente.
\item[Prioridad] Alta
\item[Procedencia] Secciones 1.2, 2.2, 2.3, 2.4, 3.1.6 y fuente documental \texttt{TOP\_Racing\_Grok\_Clarified.tex} (líneas 146, 221 y 223), además de decisiones prioritarias del cliente sobre garantía real de pago, custodia del promotor y transferencia obligatoria del vehículo subastado.; ID heredado: \texttt{RIE-005}.
\item[Verificación] \texttt{AC-UI-026}
\end{description}
\vspace{0.25em}\hrule\vspace{0.65em}
\noindent\textbf{UI-027 --- Gestión de entrega y disputa del vehículo adjudicado}\par
\begin{description}[style=nextline,leftmargin=3.15cm,labelwidth=2.8cm,font=\normalfont\bfseries]
\item[Obligación] La interfaz debe permitir registrar y consultar la confirmación de entrega, la disputa de entrega y su resolución dentro de los plazos aplicables.
\item[Prioridad] Alta
\item[Procedencia] Secciones 1.2, 2.2, 2.3, 2.4, 3.1.6 y fuente documental \texttt{TOP\_Racing\_Grok\_Clarified.tex} (líneas 146, 221 y 223), además de decisiones prioritarias del cliente sobre garantía real de pago, custodia del promotor y transferencia obligatoria del vehículo subastado.; ID heredado: \texttt{RIE-005}.
\item[Verificación] \texttt{AC-UI-027}
\end{description}
\vspace{0.25em}\hrule\vspace{0.65em}
\noindent\textbf{UI-028 --- Consulta pública de resultados y clasificaciones del evento}\par
\begin{description}[style=nextline,leftmargin=3.15cm,labelwidth=2.8cm,font=\normalfont\bfseries]
\item[Obligación] La interfaz de consulta pública debe mostrar los resultados oficiales y las clasificaciones asociadas al evento publicado.
\item[Prioridad] Alta
\item[Procedencia] Secciones 1.2, 2.2, 2.3, 2.4, 3.1.16, 3.1.21, fuente documental TOP\_Racing\_Grok\_Clarified.tex (líneas 150 y 171) y análisis funcional confirmado de la app legacy sobre la consulta de standings por selección de nivel temporal y territorial.; ID heredado: \texttt{RIE-006}.
\item[Verificación] \texttt{AC-UI-028}
\end{description}
\vspace{0.25em}\hrule\vspace{0.65em}
\noindent\textbf{UI-029 --- Consulta pública de la matriz de pesos del evento}\par
\begin{description}[style=nextline,leftmargin=3.15cm,labelwidth=2.8cm,font=\normalfont\bfseries]
\item[Obligación] La interfaz de consulta pública debe mostrar la matriz de pesos aplicable al evento con sus niveles territoriales y temporales.
\item[Prioridad] Alta
\item[Procedencia] Secciones 1.2, 2.2, 2.3, 2.4, 3.1.16, 3.1.21, fuente documental TOP\_Racing\_Grok\_Clarified.tex (líneas 150 y 171) y análisis funcional confirmado de la app legacy sobre la consulta de standings por selección de nivel temporal y territorial.; ID heredado: \texttt{RIE-006}.
\item[Verificación] \texttt{AC-UI-029}
\end{description}
\vspace{0.25em}\hrule\vspace{0.65em}
\noindent\textbf{UI-030 --- Selección del corte territorial y temporal de standings}\par
\begin{description}[style=nextline,leftmargin=3.15cm,labelwidth=2.8cm,font=\normalfont\bfseries]
\item[Obligación] La consulta pública de standings debe permitir seleccionar un nivel territorial y un nivel temporal.
\item[Prioridad] Alta
\item[Procedencia] Secciones 1.2, 2.2, 2.3, 2.4, 3.1.16, 3.1.21, fuente documental TOP\_Racing\_Grok\_Clarified.tex (líneas 150 y 171) y análisis funcional confirmado de la app legacy sobre la consulta de standings por selección de nivel temporal y territorial.; ID heredado: \texttt{RIE-006}.
\item[Verificación] \texttt{AC-UI-030}
\end{description}
\vspace{0.25em}\hrule\vspace{0.65em}
\noindent\textbf{UI-031 --- Publicación de campeonatos y puntajes del corte seleccionado}\par
\begin{description}[style=nextline,leftmargin=3.15cm,labelwidth=2.8cm,font=\normalfont\bfseries]
\item[Obligación] Para el corte territorial y temporal seleccionado, la interfaz debe mostrar los campeonatos concretos comprendidos y los puntajes publicados aplicables por modalidad y tipo de participante cuando corresponda.
\item[Prioridad] Alta
\item[Procedencia] Secciones 1.2, 2.2, 2.3, 2.4, 3.1.16, 3.1.21, fuente documental TOP\_Racing\_Grok\_Clarified.tex (líneas 150 y 171) y análisis funcional confirmado de la app legacy sobre la consulta de standings por selección de nivel temporal y territorial.; ID heredado: \texttt{RIE-006}.
\item[Verificación] \texttt{AC-UI-031}
\end{description}
\vspace{0.25em}\hrule\vspace{0.65em}
\noindent\textbf{UI-032 --- Lista estandarizada de comprobación de seguridad}\par
\begin{description}[style=nextline,leftmargin=3.15cm,labelwidth=2.8cm,font=\normalfont\bfseries]
\item[Obligación] La interfaz de seguridad y cumplimiento debe proporcionar la lista estandarizada de comprobación aplicable al evento.
\item[Prioridad] Media
\item[Procedencia] Secciones 2.2, 2.3 y fuente documental \texttt{TOP\_Racing\_Grok\_Clarified.tex} (línea 152).; ID heredado: \texttt{RIE-007}.
\item[Verificación] \texttt{AC-UI-032}
\end{description}
\vspace{0.25em}\hrule\vspace{0.65em}
\noindent\textbf{UI-033 --- Registro del proceso formal de objeciones de seguridad}\par
\begin{description}[style=nextline,leftmargin=3.15cm,labelwidth=2.8cm,font=\normalfont\bfseries]
\item[Obligación] La interfaz de seguridad y cumplimiento debe permitir registrar objeciones de seguridad, votaciones formales y decisiones del promotor asociadas al evento.
\item[Prioridad] Media
\item[Procedencia] Secciones 2.2, 2.3 y fuente documental \texttt{TOP\_Racing\_Grok\_Clarified.tex} (línea 152).; ID heredado: \texttt{RIE-007}.
\item[Verificación] \texttt{AC-UI-033}
\end{description}
\vspace{0.25em}\hrule\vspace{0.65em}
\noindent\textbf{UI-034 --- Captura opcional de datos de seguridad e inspección}\par
\begin{description}[style=nextline,leftmargin=3.15cm,labelwidth=2.8cm,font=\normalfont\bfseries]
\item[Obligación] La interfaz debe permitir registrar datos de seguridad e inspección sin convertir esa captura, por sí sola, en condición obligatoria para continuar el ciclo del evento.
\item[Prioridad] Media
\item[Procedencia] Secciones 2.2, 2.3 y fuente documental \texttt{TOP\_Racing\_Grok\_Clarified.tex} (línea 152).; ID heredado: \texttt{RIE-007}.
\item[Verificación] \texttt{AC-UI-034}
\end{description}
\vspace{0.25em}\hrule\vspace{0.65em}
\noindent\textbf{UI-035 --- Exclusión del pesaje externo de la interfaz de seguridad}\par
\begin{description}[style=nextline,leftmargin=3.15cm,labelwidth=2.8cm,font=\normalfont\bfseries]
\item[Obligación] La interfaz de seguridad y cumplimiento no debe solicitar ni registrar el pesaje externo de entrada o de salida.
\item[Prioridad] Media
\item[Procedencia] Secciones 2.2, 2.3 y fuente documental \texttt{TOP\_Racing\_Grok\_Clarified.tex} (línea 152).; ID heredado: \texttt{RIE-007}.
\item[Verificación] \texttt{AC-UI-035}
\end{description}
\vspace{0.25em}\hrule\vspace{0.65em}
\noindent\textbf{UI-036 --- Interfaz comunitaria de discusión y novedades}\par
\begin{description}[style=nextline,leftmargin=3.15cm,labelwidth=2.8cm,font=\normalfont\bfseries]
\item[Obligación] La interfaz comunitaria debe poner a disposición un foro comunitario y un canal de novedades del ecosistema, permitir que cualquier usuario con identidad verificada publique en ambos y permitir que \texttt{admin} los modere.
\item[Prioridad] Media
\item[Procedencia] Secciones 1.2, 2.1, 2.2, 2.3, fuente documental \texttt{TOP\_Racing\_Grok\_Clarified.tex} (línea 154) y decisiones prioritarias del cliente sobre publicación por cualquier usuario con identidad verificada y moderación por \texttt{admin}.; ID heredado: \texttt{RIE-008}.
\item[Verificación] \texttt{AC-UI-036}
\end{description}
\vspace{0.25em}\hrule\vspace{0.65em}
\noindent\textbf{UI-037 --- Interfaz de premio individual del evento}\par
\begin{description}[style=nextline,leftmargin=3.15cm,labelwidth=2.8cm,font=\normalfont\bfseries]
\item[Obligación] La interfaz debe mostrar públicamente, para cada participante computable, su premio individual del evento y la distribución económica posterior que le corresponda, publicando sin identidad las cantidades y la parte de premio correspondientes a terceros cuando existan.
\item[Prioridad] Alta
\item[Procedencia] Secciones 1.2, 2.2, 2.3, 2.6.3, 2.6.4 y fuente documental \texttt{TOP\_Racing\_Grok\_Clarified.tex} (línea 148).; ID heredado: \texttt{RIE-009}.
\item[Verificación] \texttt{AC-UI-037}
\end{description}
\vspace{0.25em}\hrule\vspace{0.65em}
\noindent\textbf{UI-038 --- Selección del campeonato para operaciones económicas}\par
\begin{description}[style=nextline,leftmargin=3.15cm,labelwidth=2.8cm,font=\normalfont\bfseries]
\item[Obligación] La interfaz de premio de campeonato debe permitir seleccionar un campeonato específico de eficiencia o de terminación.
\item[Prioridad] Alta
\item[Procedencia] Secciones 2.2, 2.6.4.3 y decisiones prioritarias del cliente sobre premio de campeonato, patrocinio, aportaciones a participantes y cobro inmediato.; ID heredado: \texttt{RIE-010}.
\item[Verificación] \texttt{AC-UI-038}
\end{description}
\vspace{0.25em}\hrule\vspace{0.65em}
\noindent\textbf{UI-039 --- Registro de aportaciones directas a la bolsa del campeonato}\par
\begin{description}[style=nextline,leftmargin=3.15cm,labelwidth=2.8cm,font=\normalfont\bfseries]
\item[Obligación] La interfaz debe permitir registrar aportaciones directas a la bolsa del campeonato seleccionado dentro de la ventana temporal aplicable.
\item[Prioridad] Alta
\item[Procedencia] Secciones 2.2, 2.6.4.3 y decisiones prioritarias del cliente sobre premio de campeonato, patrocinio, aportaciones a participantes y cobro inmediato.; ID heredado: \texttt{RIE-010}.
\item[Verificación] \texttt{AC-UI-039}
\end{description}
\vspace{0.25em}\hrule\vspace{0.65em}
\noindent\textbf{UI-040 --- Registro condicionado de aportaciones a participantes de campeonato}\par
\begin{description}[style=nextline,leftmargin=3.15cm,labelwidth=2.8cm,font=\normalfont\bfseries]
\item[Obligación] La interfaz debe permitir registrar una aportación a la contribución de un participante del campeonato únicamente cuando la función global correspondiente esté activa.
\item[Prioridad] Alta
\item[Procedencia] Secciones 2.2, 2.6.4.3 y decisiones prioritarias del cliente sobre premio de campeonato, patrocinio, aportaciones a participantes y cobro inmediato.; ID heredado: \texttt{RIE-010}.
\item[Verificación] \texttt{AC-UI-040}
\end{description}
\vspace{0.25em}\hrule\vspace{0.65em}
\noindent\textbf{UI-041 --- Cobro inmediato de aportaciones de campeonato aceptadas}\par
\begin{description}[style=nextline,leftmargin=3.15cm,labelwidth=2.8cm,font=\normalfont\bfseries]
\item[Obligación] Al aceptar una aportación de campeonato, el sistema debe cobrar de inmediato el importe correspondiente.
\item[Prioridad] Alta
\item[Procedencia] Secciones 2.2, 2.6.4.3 y decisiones prioritarias del cliente sobre premio de campeonato, patrocinio, aportaciones a participantes y cobro inmediato.; ID heredado: \texttt{RIE-010}.
\item[Verificación] \texttt{AC-UI-041}
\end{description}
\vspace{0.25em}\hrule\vspace{0.65em}
\noindent\textbf{UI-042 --- Publicación del anuncio de patrocinio directo del campeonato}\par
\begin{description}[style=nextline,leftmargin=3.15cm,labelwidth=2.8cm,font=\normalfont\bfseries]
\item[Obligación] Cuando exista una aportación directa a la bolsa del campeonato con anuncio asociado, la interfaz debe publicar la imagen del patrocinio correspondiente.
\item[Prioridad] Alta
\item[Procedencia] Secciones 2.2, 2.6.4.3 y decisiones prioritarias del cliente sobre premio de campeonato, patrocinio, aportaciones a participantes y cobro inmediato.; ID heredado: \texttt{RIE-010}.
\item[Verificación] \texttt{AC-UI-042}
\end{description}
\vspace{0.25em}\hrule\vspace{0.65em}
\noindent\textbf{UI-043 --- Consulta pública del premio individual de campeonato}\par
\begin{description}[style=nextline,leftmargin=3.15cm,labelwidth=2.8cm,font=\normalfont\bfseries]
\item[Obligación] La interfaz debe permitir consultar públicamente el premio individual de campeonato de cada competidor cuando haya sido determinado oficialmente.
\item[Prioridad] Alta
\item[Procedencia] Secciones 2.2, 2.6.4.3 y decisiones prioritarias del cliente sobre premio de campeonato, patrocinio, aportaciones a participantes y cobro inmediato.; ID heredado: \texttt{RIE-010}.
\item[Verificación] \texttt{AC-UI-043}
\end{description}
\vspace{0.25em}\hrule\vspace{0.65em}
\noindent\textbf{UI-044 --- Consulta pública de la distribución posterior del premio de campeonato}\par
\begin{description}[style=nextline,leftmargin=3.15cm,labelwidth=2.8cm,font=\normalfont\bfseries]
\item[Obligación] La interfaz debe permitir consultar la distribución posterior del premio individual de campeonato conforme a las reglas de publicación y reserva aplicables.
\item[Prioridad] Alta
\item[Procedencia] Secciones 2.2, 2.6.4.3 y decisiones prioritarias del cliente sobre premio de campeonato, patrocinio, aportaciones a participantes y cobro inmediato.; ID heredado: \texttt{RIE-010}.
\item[Verificación] \texttt{AC-UI-044}
\end{description}
\vspace{0.25em}\hrule\vspace{0.65em}
\noindent\textbf{UI-045 --- Registro de aportaciones directas a la bolsa del evento}\par
\begin{description}[style=nextline,leftmargin=3.15cm,labelwidth=2.8cm,font=\normalfont\bfseries]
\item[Obligación] Mientras el evento permanezca en \texttt{inscripciones}, la interfaz debe permitir registrar aportaciones directas a su bolsa.
\item[Prioridad] Alta
\item[Procedencia] Secciones 1.2, 2.2, 2.6.2, 2.6.4.2 y decisiones prioritarias del cliente sobre patrocinio directo de eventos, aportaciones a participantes y cobro inmediato.; ID heredado: \texttt{RIE-011}.
\item[Verificación] \texttt{AC-UI-045}
\end{description}
\vspace{0.25em}\hrule\vspace{0.65em}
\noindent\textbf{UI-046 --- Registro condicionado de aportaciones de eficiencia a participantes del evento}\par
\begin{description}[style=nextline,leftmargin=3.15cm,labelwidth=2.8cm,font=\normalfont\bfseries]
\item[Obligación] Mientras el evento permanezca en \texttt{inscripciones}, la interfaz debe permitir registrar aportaciones a la contribución de eficiencia de un participante únicamente cuando la función global correspondiente esté activa.
\item[Prioridad] Alta
\item[Procedencia] Secciones 1.2, 2.2, 2.6.2, 2.6.4.2 y decisiones prioritarias del cliente sobre patrocinio directo de eventos, aportaciones a participantes y cobro inmediato.; ID heredado: \texttt{RIE-011}.
\item[Verificación] \texttt{AC-UI-046}
\end{description}
\vspace{0.25em}\hrule\vspace{0.65em}
\noindent\textbf{UI-047 --- Registro condicionado de aportaciones de terminación a inscripciones}\par
\begin{description}[style=nextline,leftmargin=3.15cm,labelwidth=2.8cm,font=\normalfont\bfseries]
\item[Obligación] Mientras el evento permanezca en \texttt{inscripciones}, la interfaz debe permitir registrar aportaciones a la contribución de terminación de una inscripción únicamente cuando la función global correspondiente esté activa.
\item[Prioridad] Alta
\item[Procedencia] Secciones 1.2, 2.2, 2.6.2, 2.6.4.2 y decisiones prioritarias del cliente sobre patrocinio directo de eventos, aportaciones a participantes y cobro inmediato.; ID heredado: \texttt{RIE-011}.
\item[Verificación] \texttt{AC-UI-047}
\end{description}
\vspace{0.25em}\hrule\vspace{0.65em}
\noindent\textbf{UI-048 --- Cobro inmediato de aportaciones del evento aceptadas}\par
\begin{description}[style=nextline,leftmargin=3.15cm,labelwidth=2.8cm,font=\normalfont\bfseries]
\item[Obligación] Al aceptar una aportación al evento o a una inscripción del evento, el sistema debe cobrar de inmediato el importe correspondiente.
\item[Prioridad] Alta
\item[Procedencia] Secciones 1.2, 2.2, 2.6.2, 2.6.4.2 y decisiones prioritarias del cliente sobre patrocinio directo de eventos, aportaciones a participantes y cobro inmediato.; ID heredado: \texttt{RIE-011}.
\item[Verificación] \texttt{AC-UI-048}
\end{description}
\vspace{0.25em}\hrule\vspace{0.65em}
\noindent\textbf{UI-049 --- Publicación del anuncio de patrocinio directo del evento}\par
\begin{description}[style=nextline,leftmargin=3.15cm,labelwidth=2.8cm,font=\normalfont\bfseries]
\item[Obligación] Cuando exista una aportación directa a la bolsa del evento con anuncio asociado, la interfaz debe publicar la imagen del patrocinio correspondiente.
\item[Prioridad] Alta
\item[Procedencia] Secciones 1.2, 2.2, 2.6.2, 2.6.4.2 y decisiones prioritarias del cliente sobre patrocinio directo de eventos, aportaciones a participantes y cobro inmediato.; ID heredado: \texttt{RIE-011}.
\item[Verificación] \texttt{AC-UI-049}
\end{description}
\vspace{0.25em}\hrule\vspace{0.65em}
\noindent\textbf{UI-050 --- Distinción del patrocinio de campeonatos aplicable al evento}\par
\begin{description}[style=nextline,leftmargin=3.15cm,labelwidth=2.8cm,font=\normalfont\bfseries]
\item[Obligación] Cuando el evento pertenezca a campeonatos patrocinados, la interfaz del evento debe mostrar el patrocinio aplicable por campeonato de forma diferenciada del patrocinio directo del evento.
\item[Prioridad] Alta
\item[Procedencia] Secciones 1.2, 2.2, 2.6.2, 2.6.4.2 y decisiones prioritarias del cliente sobre patrocinio directo de eventos, aportaciones a participantes y cobro inmediato.; ID heredado: \texttt{RIE-011}.
\item[Verificación] \texttt{AC-UI-050}
\end{description}
\vspace{0.25em}\hrule\vspace{0.65em}
\noindent\textbf{UI-051 --- Interfaz de calificación del promotor}\par
\begin{description}[style=nextline,leftmargin=3.15cm,labelwidth=2.8cm,font=\normalfont\bfseries]
\item[Obligación] La interfaz del sistema debe permitir que participantes y espectadores registren una calificación del promotor después del evento correspondiente.
\item[Prioridad] Media
\item[Procedencia] Secciones 1.2, 2.2, 3.1.26 y fuente documental \texttt{TOP\_Racing\_Grok\_Clarified.tex} (línea 133).; ID heredado: \texttt{RIE-012}.
\item[Verificación] \texttt{AC-UI-051}
\end{description}
\vspace{0.25em}\hrule\vspace{0.65em}
\noindent\textbf{UI-052 --- Configuración administrativa de parámetros del modelo de negocio}\par
\begin{description}[style=nextline,leftmargin=3.15cm,labelwidth=2.8cm,font=\normalfont\bfseries]
\item[Obligación] La interfaz administrativa debe permitir a \texttt{admin} consultar y actualizar los porcentajes de comisión, la cuota fija por evento del promotor, el umbral anual de gratuidad y los estados globales de activación de aportaciones a participantes definidos por el modelo de negocio.
\item[Prioridad] Alta
\item[Procedencia] Secciones 1.2, 2.2, 2.6.2.2, 3.1.27 y decisiones prioritarias del cliente sobre configuración y publicación de comisiones.; ID heredado: \texttt{RIE-013}.
\item[Verificación] \texttt{AC-UI-052}
\end{description}
\vspace{0.25em}\hrule\vspace{0.65em}
\noindent\textbf{UI-053 --- Información previa sobre comisión o cuota aplicable}\par
\begin{description}[style=nextline,leftmargin=3.15cm,labelwidth=2.8cm,font=\normalfont\bfseries]
\item[Obligación] Antes de confirmar una operación afectada por el modelo de negocio, la interfaz operativa debe informar al usuario si la operación causa comisión o cuota y debe identificar la regla y el parámetro vigentes aplicables.
\item[Prioridad] Alta
\item[Procedencia] Secciones 1.2, 2.2, 2.6.2.2, 3.1.27 y decisiones prioritarias del cliente sobre configuración y publicación de comisiones.; ID heredado: \texttt{RIE-013}.
\item[Verificación] \texttt{AC-UI-053}
\end{description}
\vspace{0.25em}\hrule\vspace{0.65em}
\noindent\textbf{UI-054 --- Consulta posterior del cobro o exención del modelo de negocio}\par
\begin{description}[style=nextline,leftmargin=3.15cm,labelwidth=2.8cm,font=\normalfont\bfseries]
\item[Obligación] Después de ejecutar una operación afectada por el modelo de negocio, la interfaz debe permitir al usuario consultar el detalle del cobro o de la exención correspondiente.
\item[Prioridad] Alta
\item[Procedencia] Secciones 1.2, 2.2, 2.6.2.2, 3.1.27 y decisiones prioritarias del cliente sobre configuración y publicación de comisiones.; ID heredado: \texttt{RIE-013}.
\item[Verificación] \texttt{AC-UI-054}
\end{description}
\vspace{0.25em}\hrule\vspace{0.65em}
\noindent\textbf{UI-055 --- Conjunto de pantallas críticas para evaluación UX}\par
\begin{description}[style=nextline,leftmargin=3.15cm,labelwidth=2.8cm,font=\normalfont\bfseries]
\item[Obligación] El sistema debe considerar como pantallas críticas del dominio, para efectos de evaluación UX de percepción inicial positiva y de jerarquía visual de información crítica, al menos la interfaz principal del sistema, la interfaz de inscripción de competidores, la interfaz de acceso de espectadores a la subasta del evento, la interfaz de subasta de vehículos, la interfaz de premio individual del evento, la interfaz de premio de campeonato y patrocinio, y la interfaz del evento para patrocinio y aportaciones a participantes.
\item[Prioridad] Alta
\item[Procedencia] Secciones 2.2, 2.7, 3.5, RIE-001, RIE-003, RIE-004, RIE-005, RIE-009, RIE-010, RIE-011 y decisiones prioritarias del cliente sobre importancia de la estética y la UX para la adopción.; ID heredado: \texttt{RIE-014}.
\item[Verificación] \texttt{AC-UI-055}
\end{description}
\vspace{0.25em}\hrule\vspace{0.65em}
\noindent\textbf{UI-056 --- Contexto común de inscripciones en el detalle del evento}\par
\begin{description}[style=nextline,leftmargin=3.15cm,labelwidth=2.8cm,font=\normalfont\bfseries]
\item[Obligación] La interfaz de detalle operativo del evento debe presentar los datos públicos y operativos del evento junto con un contexto común de sus inscripciones y los datos asociados a cada una.
\item[Prioridad] Alta
\item[Procedencia] Secciones 1.2, 2.2, 2.3, 3.1.4, 3.1.6, 3.1.13, 3.1.17, 3.1.18, 3.1.21, 3.2 RIE-003, RIE-005 y decisiones prioritarias del cliente confirmadas a partir del análisis funcional de la app legacy sobre la lista común de inscripciones con captura condicionada de resultados, premio de terminación y pujas.; ID heredado: \texttt{RIE-016}.
\item[Verificación] \texttt{AC-UI-056}
\end{description}
\vspace{0.25em}\hrule\vspace{0.65em}
\noindent\textbf{UI-057 --- Captura de mejor vuelta válida durante velocidad}\par
\begin{description}[style=nextline,leftmargin=3.15cm,labelwidth=2.8cm,font=\normalfont\bfseries]
\item[Obligación] Durante \texttt{velocidad}, la interfaz debe permitir al promotor registrar el mejor tiempo de vuelta válido de cada inscripción no cancelada ni descalificada.
\item[Prioridad] Alta
\item[Procedencia] Secciones 1.2, 2.2, 2.3, 3.1.4, 3.1.6, 3.1.13, 3.1.17, 3.1.18, 3.1.21, 3.2 RIE-003, RIE-005 y decisiones prioritarias del cliente confirmadas a partir del análisis funcional de la app legacy sobre la lista común de inscripciones con captura condicionada de resultados, premio de terminación y pujas.; ID heredado: \texttt{RIE-016}.
\item[Verificación] \texttt{AC-UI-057}
\end{description}
\vspace{0.25em}\hrule\vspace{0.65em}
\noindent\textbf{UI-058 --- Ausencia de velocidad válida por incumplimiento de pesaje decidido fuera del sistema}\par
\begin{description}[style=nextline,leftmargin=3.15cm,labelwidth=2.8cm,font=\normalfont\bfseries]
\item[Obligación] Durante \texttt{velocidad}, la interfaz debe permitir que el promotor deje una inscripción sin velocidad válida capturada cuando aplique la consecuencia de un incumplimiento de pesaje decidido fuera del sistema.
\item[Prioridad] Alta
\item[Procedencia] Secciones 1.2, 2.2, 2.3, 3.1.4, 3.1.6, 3.1.13, 3.1.17, 3.1.18, 3.1.21, 3.2 RIE-003, RIE-005 y decisiones prioritarias del cliente confirmadas a partir del análisis funcional de la app legacy sobre la lista común de inscripciones con captura condicionada de resultados, premio de terminación y pujas.; ID heredado: \texttt{RIE-016}.
\item[Verificación] \texttt{AC-UI-058}
\end{description}
\vspace{0.25em}\hrule\vspace{0.65em}
\noindent\textbf{UI-059 --- Captura de lugar oficial durante carrera}\par
\begin{description}[style=nextline,leftmargin=3.15cm,labelwidth=2.8cm,font=\normalfont\bfseries]
\item[Obligación] Durante \texttt{carrera}, la interfaz debe permitir al promotor registrar el lugar oficial de cada competidor computable.
\item[Prioridad] Alta
\item[Procedencia] Secciones 1.2, 2.2, 2.3, 3.1.4, 3.1.6, 3.1.13, 3.1.17, 3.1.18, 3.1.21, 3.2 RIE-003, RIE-005 y decisiones prioritarias del cliente confirmadas a partir del análisis funcional de la app legacy sobre la lista común de inscripciones con captura condicionada de resultados, premio de terminación y pujas.; ID heredado: \texttt{RIE-016}.
\item[Verificación] \texttt{AC-UI-059}
\end{description}
\vspace{0.25em}\hrule\vspace{0.65em}
\noindent\textbf{UI-060 --- Captura de vueltas durante carrera}\par
\begin{description}[style=nextline,leftmargin=3.15cm,labelwidth=2.8cm,font=\normalfont\bfseries]
\item[Obligación] Durante \texttt{carrera}, la interfaz debe permitir al promotor registrar el número oficial de vueltas de cada competidor como insumo del premio de terminación.
\item[Prioridad] Alta
\item[Procedencia] Secciones 1.2, 2.2, 2.3, 3.1.4, 3.1.6, 3.1.13, 3.1.17, 3.1.18, 3.1.21, 3.2 RIE-003, RIE-005 y decisiones prioritarias del cliente confirmadas a partir del análisis funcional de la app legacy sobre la lista común de inscripciones con captura condicionada de resultados, premio de terminación y pujas.; ID heredado: \texttt{RIE-016}.
\item[Verificación] \texttt{AC-UI-060}
\end{description}
\vspace{0.25em}\hrule\vspace{0.65em}
\noindent\textbf{UI-061 --- Descalificación con causa textual}\par
\begin{description}[style=nextline,leftmargin=3.15cm,labelwidth=2.8cm,font=\normalfont\bfseries]
\item[Obligación] Durante \texttt{velocidad} o \texttt{carrera}, la interfaz no debe permitir que el promotor confirme una descalificación sin una causa textual.
\item[Prioridad] Alta
\item[Procedencia] Secciones 1.2, 2.2, 2.3, 3.1.4, 3.1.6, 3.1.13, 3.1.17, 3.1.18, 3.1.21, 3.2 RIE-003, RIE-005 y decisiones prioritarias del cliente confirmadas a partir del análisis funcional de la app legacy sobre la lista común de inscripciones con captura condicionada de resultados, premio de terminación y pujas.; ID heredado: \texttt{RIE-016}.
\item[Verificación] \texttt{AC-UI-061}
\end{description}
\vspace{0.25em}\hrule\vspace{0.65em}
\noindent\textbf{UI-062 --- Consulta económica de terminación en el detalle del evento}\par
\begin{description}[style=nextline,leftmargin=3.15cm,labelwidth=2.8cm,font=\normalfont\bfseries]
\item[Obligación] Cuando el premio de terminación haya sido determinado, la interfaz debe permitir consultar el pozo de terminación, la contribución total computable y el premio de terminación resultante por inscripción.
\item[Prioridad] Alta
\item[Procedencia] Secciones 1.2, 2.2, 2.3, 3.1.4, 3.1.6, 3.1.13, 3.1.17, 3.1.18, 3.1.21, 3.2 RIE-003, RIE-005 y decisiones prioritarias del cliente confirmadas a partir del análisis funcional de la app legacy sobre la lista común de inscripciones con captura condicionada de resultados, premio de terminación y pujas.; ID heredado: \texttt{RIE-016}.
\item[Verificación] \texttt{AC-UI-062}
\end{description}
\vspace{0.25em}\hrule\vspace{0.65em}
\noindent\textbf{UI-063 --- Reserva de datos privados en el detalle operativo}\par
\begin{description}[style=nextline,leftmargin=3.15cm,labelwidth=2.8cm,font=\normalfont\bfseries]
\item[Obligación] La interfaz de detalle operativo debe permitir que cada usuario capture o consulte únicamente sus propios datos privados cuando la operación esté habilitada y no debe exponer datos privados de otros usuarios.
\item[Prioridad] Alta
\item[Procedencia] Secciones 1.2, 2.2, 2.3, 3.1.4, 3.1.6, 3.1.13, 3.1.17, 3.1.18, 3.1.21, 3.2 RIE-003, RIE-005 y decisiones prioritarias del cliente confirmadas a partir del análisis funcional de la app legacy sobre la lista común de inscripciones con captura condicionada de resultados, premio de terminación y pujas.; ID heredado: \texttt{RIE-016}.
\item[Verificación] \texttt{AC-UI-063}
\end{description}
\vspace{0.25em}\hrule\vspace{0.65em}
\noindent\textbf{UI-064 --- Bloqueo de edición fuera de condiciones habilitantes}\par
\begin{description}[style=nextline,leftmargin=3.15cm,labelwidth=2.8cm,font=\normalfont\bfseries]
\item[Obligación] Fuera del estado o de las condiciones de autorización aplicables, la interfaz debe conservar la consulta permitida y bloquear la edición no autorizada.
\item[Prioridad] Alta
\item[Procedencia] Secciones 1.2, 2.2, 2.3, 3.1.4, 3.1.6, 3.1.13, 3.1.17, 3.1.18, 3.1.21, 3.2 RIE-003, RIE-005 y decisiones prioritarias del cliente confirmadas a partir del análisis funcional de la app legacy sobre la lista común de inscripciones con captura condicionada de resultados, premio de terminación y pujas.; ID heredado: \texttt{RIE-016}.
\item[Verificación] \texttt{AC-UI-064}
\end{description}
\vspace{0.25em}\hrule\vspace{0.65em}
\noindent\textbf{UI-065 --- Selección del corte territorial y temporal de prioridades}\par
\begin{description}[style=nextline,leftmargin=3.15cm,labelwidth=2.8cm,font=\normalfont\bfseries]
\item[Obligación] La interfaz pública de prioridades de eventos debe permitir seleccionar un nivel territorial y un nivel temporal para consultar el corte correspondiente de eventos comparables.
\item[Prioridad] Alta
\item[Procedencia] Secciones 1.2, 2.2, 2.4, 3.1.16, RF-104, RF-109 y análisis funcional confirmado de la app legacy sobre la consulta de prioridades de eventos por selección de nivel temporal y territorial.; ID heredado: \texttt{RIE-017}.
\item[Verificación] \texttt{AC-UI-065}
\end{description}
\vspace{0.25em}\hrule\vspace{0.65em}
\noindent\textbf{UI-066 --- Publicación de discriminadores y P/C por evento}\par
\begin{description}[style=nextline,leftmargin=3.15cm,labelwidth=2.8cm,font=\normalfont\bfseries]
\item[Obligación] Para cada evento del corte consultado, la interfaz debe publicar sus discriminadores territoriales y temporales, estado, organizador, fecha, pista o variante aplicable y P/C propio.
\item[Prioridad] Alta
\item[Procedencia] Secciones 1.2, 2.2, 2.4, 3.1.16, RF-104, RF-109 y análisis funcional confirmado de la app legacy sobre la consulta de prioridades de eventos por selección de nivel temporal y territorial.; ID heredado: \texttt{RIE-017}.
\item[Verificación] \texttt{AC-UI-066}
\end{description}
\vspace{0.25em}\hrule\vspace{0.65em}
\noindent\textbf{UI-067 --- Ordenamiento consultable por P/C}\par
\begin{description}[style=nextline,leftmargin=3.15cm,labelwidth=2.8cm,font=\normalfont\bfseries]
\item[Obligación] La interfaz pública de prioridades debe permitir ordenar los eventos consultados por su P/C publicado.
\item[Prioridad] Alta
\item[Procedencia] Secciones 1.2, 2.2, 2.4, 3.1.16, RF-104, RF-109 y análisis funcional confirmado de la app legacy sobre la consulta de prioridades de eventos por selección de nivel temporal y territorial.; ID heredado: \texttt{RIE-017}.
\item[Verificación] \texttt{AC-UI-067}
\end{description}
\vspace{0.25em}\hrule\vspace{0.65em}
\noindent\textbf{UI-068 --- Navegación al detalle público del evento}\par
\begin{description}[style=nextline,leftmargin=3.15cm,labelwidth=2.8cm,font=\normalfont\bfseries]
\item[Obligación] La interfaz pública de prioridades debe permitir abrir el detalle público del evento seleccionado.
\item[Prioridad] Alta
\item[Procedencia] Secciones 1.2, 2.2, 2.4, 3.1.16, RF-104, RF-109 y análisis funcional confirmado de la app legacy sobre la consulta de prioridades de eventos por selección de nivel temporal y territorial.; ID heredado: \texttt{RIE-017}.
\item[Verificación] \texttt{AC-UI-068}
\end{description}
\vspace{0.25em}\hrule\vspace{0.65em}
\noindent\textbf{UI-069 --- Información pública del pozo durante la subasta de vehículos}\par
\begin{description}[style=nextline,leftmargin=3.15cm,labelwidth=2.8cm,font=\normalfont\bfseries]
\item[Obligación] Mientras la fase de subasta de vehículos permanezca abierta, la interfaz pública del evento debe mostrar el pozo inicial publicado como referencia económica del premio de eficiencia y no debe mostrar el pozo económico oficial consolidado ni importes de aportaciones o contribuciones que permitan reconstruirlo.
\item[Prioridad] Alta
\item[Procedencia] \texttt{DEC-170} y \texttt{RF-EPR-022}; decisión confirmada sobre información económica disponible antes del cierre.
\item[Verificación] \texttt{AC-UI-069}
\end{description}
\vspace{0.25em}\hrule\vspace{0.65em}
\noindent\textbf{UI-070 --- Declaración de categoría tecnológica y documentación abierta obligatoria}\par
\begin{description}[style=nextline,leftmargin=3.15cm,labelwidth=2.8cm,font=\normalfont\bfseries]
\item[Obligación] La interfaz de publicación de evento debe permitir declarar la categoría tecnológica admitida, incluyendo si acepta vehículos motorizados, no motorizados, caseros, artesanales, reutilizados o de bajo costo, y debe exigir el registro de la documentación abierta obligatoria del vehículo conforme a la condición de tecnología abierta del vehículo definida por esta SRS.
\item[Prioridad] Alta
\item[Procedencia] DEC-173, DEC-174, DEC-175 y DEC-180; decisión confirmada sobre bajo costo, tecnología abierta obligatoria, acceso operativo publicado y generalización a vehículos de pista con cualquier tecnología.
\item[Verificación] \texttt{AC-UI-070}
\end{description}
\vspace{0.25em}\hrule\vspace{0.65em}
\section{Restricciones externas}

\noindent\textbf{RC-EXT-001 --- Procesamiento obligatorio de retenciones, cobros, pagos y devoluciones reales}\par
\begin{description}[style=nextline,leftmargin=3.15cm,labelwidth=2.8cm,font=\normalfont\bfseries]
\item[Obligación] El sistema debe ejecutar retenciones, cobros, liberaciones, pagos y devoluciones reales cuando las reglas del dominio lo exijan.
\item[Prioridad] Alta
\item[Procedencia] Sección 2.5 y decisiones prioritarias del cliente sobre pagos reales.; ID heredado: \texttt{REX-001}.
\item[Verificación] \texttt{AC-RC-EXT-001}
\end{description}
\vspace{0.25em}\hrule\vspace{0.65em}
\noindent\textbf{RC-EXT-002 --- Sujeción a plazos publicados del evento en disputas de entrega de vehículos}\par
\begin{description}[style=nextline,leftmargin=3.15cm,labelwidth=2.8cm,font=\normalfont\bfseries]
\item[Obligación] Cuando un evento contemple subasta de vehículos, el sistema debe operar las disputas de entrega del vehículo adjudicado únicamente con arreglo al plazo máximo para iniciarlas y al plazo máximo para resolverlas publicados para ese evento.
\item[Prioridad] Alta
\item[Procedencia] Sección 2.5, RF-075, RF-079A, RF-079B y decisiones prioritarias del cliente.; ID heredado: \texttt{REX-002}.
\item[Verificación] \texttt{AC-RC-EXT-002}
\end{description}
\vspace{0.25em}\hrule\vspace{0.65em}
\noindent\textbf{RC-EXT-003 --- Custodia previa obligatoria del vehículo antes de su subasta}\par
\begin{description}[style=nextline,leftmargin=3.15cm,labelwidth=2.8cm,font=\normalfont\bfseries]
\item[Obligación] Ningún vehículo debe entrar válidamente a la subasta del evento sin haber sido previamente entregado al promotor y quedar bajo control del evento.
\item[Prioridad] Alta
\item[Procedencia] Sección 2.5, RF-046 y decisiones prioritarias del cliente sobre entrega previa del vehículo al promotor.; ID heredado: \texttt{REX-003}.
\item[Verificación] \texttt{AC-RC-EXT-003}
\end{description}
\vspace{0.25em}\hrule\vspace{0.65em}
\noindent\textbf{RC-EXT-004 --- Régimen obligatorio de publicidad del dominio}\par
\begin{description}[style=nextline,leftmargin=3.15cm,labelwidth=2.8cm,font=\normalfont\bfseries]
\item[Obligación] El sistema debe respetar la publicidad obligatoria del dominio para patrocinios, inscripciones canceladas o descalificadas, protestas de seguridad, disputas de entrega, decisiones del promotor y demás publicaciones críticas que las reglas del dominio declaren visibles.
\item[Prioridad] Alta
\item[Procedencia] Sección 2.5, RF-067A a RF-067E, RF-070A, RF-101, RF-102, RF-129 y decisiones prioritarias del cliente.; ID heredado: \texttt{REX-004}.
\item[Verificación] \texttt{AC-RC-EXT-004}
\end{description}
\vspace{0.25em}\hrule\vspace{0.65em}
\noindent\textbf{RC-EXT-005 --- Régimen obligatorio de reserva selectiva de identidad}\par
\begin{description}[style=nextline,leftmargin=3.15cm,labelwidth=2.8cm,font=\normalfont\bfseries]
\item[Obligación] El sistema debe mantener reservada para la consulta pública la identidad de quien presenta una protesta de seguridad, sin impedir que esa identidad sea visible para \texttt{admin} y para el promotor del evento.
\item[Prioridad] Alta
\item[Procedencia] Sección 2.5, RF-058A, RF-067C y DEC-027.; ID heredado: \texttt{REX-005}.
\item[Verificación] \texttt{AC-RC-EXT-005}
\end{description}
\vspace{0.25em}\hrule\vspace{0.65em}
\noindent\textbf{RC-EXT-006 --- Verificación de identidad y no exclusión antes de la participación}\par
\begin{description}[style=nextline,leftmargin=3.15cm,labelwidth=2.8cm,font=\normalfont\bfseries]
\item[Obligación] El sistema debe exigir que la identidad del usuario esté verificada y que su estado de exclusión no esté activo antes de permitirle cualquier participación operativa en el ecosistema TOP Racing.
\item[Prioridad] Alta
\item[Procedencia] Sección 2.5, sección 2.3 y DEC-069.; ID heredado: \texttt{REX-006}.
\item[Verificación] \texttt{AC-RC-EXT-006}
\end{description}
\vspace{0.25em}\hrule\vspace{0.65em}
\noindent\textbf{RC-EXT-007 --- Verificación reforzada obligatoria para operaciones de mayor riesgo}\par
\begin{description}[style=nextline,leftmargin=3.15cm,labelwidth=2.8cm,font=\normalfont\bfseries]
\item[Obligación] El sistema debe exigir verificación reforzada activa para inscribir vehículos en eventos, pujar en subastas de vehículos y promover eventos.
\item[Prioridad] Alta
\item[Procedencia] Sección 2.5, sección 2.6, RF-151A y decisiones prioritarias del cliente sobre verificación reforzada.; ID heredado: \texttt{REX-006A}.
\item[Verificación] \texttt{AC-RC-EXT-007}
\end{description}
\vspace{0.25em}\hrule\vspace{0.65em}
\noindent\textbf{RC-EXT-008 --- Aceptación de la versión vigente de términos y condiciones antes de la participación}\par
\begin{description}[style=nextline,leftmargin=3.15cm,labelwidth=2.8cm,font=\normalfont\bfseries]
\item[Obligación] El sistema debe exigir la aceptación registrada de la versión vigente de términos y condiciones antes de permitir cualquier participación operativa en el ecosistema TOP Racing.
\item[Prioridad] Alta
\item[Procedencia] Sección 2.5, RF-071B y decisiones prioritarias del cliente sobre aceptación versionada.; ID heredado: \texttt{REX-007}.
\item[Verificación] \texttt{AC-RC-EXT-008}
\end{description}
\vspace{0.25em}\hrule\vspace{0.65em}
\noindent\textbf{RC-EXT-009 --- Sujeción del vehículo subastado a transferencia obligatoria y prohibición de retiro no autorizado}\par
\begin{description}[style=nextline,leftmargin=3.15cm,labelwidth=2.8cm,font=\normalfont\bfseries]
\item[Obligación] Todo vehículo sometido a la subasta de un evento debe quedar sujeto a transferencia obligatoria según el resultado oficial vigente de esa subasta y no puede retirarse del control del evento en contra de ese resultado.
\item[Prioridad] Alta
\item[Procedencia] Sección 2.5, RF-050C y DEC-088 a DEC-090.; ID heredado: \texttt{REX-008}.
\item[Verificación] \texttt{AC-RC-EXT-009}
\end{description}
\vspace{0.25em}\hrule\vspace{0.65em}
\section{Requisitos de calidad}
Esta sección sustituye la duplicación heredada entre necesidades de calidad, RNF, RAC y la antigua RIE-015. Cada QR es la única formulación normativa de la propiedad correspondiente en esta versión.

\noindent\textbf{QR-001 --- Claridad operativa de estados y acciones}\par
\begin{description}[style=nextline,leftmargin=3.15cm,labelwidth=2.8cm,font=\normalfont\bfseries]
\item[Obligación] La interfaz debe identificar el estado oficial vigente y diferenciar las operaciones habilitadas de las bloqueadas; toda operación bloqueada debe presentar una causa comprensible.
\item[Prioridad] Alta
\item[Procedencia] Consolidación de RNF-001; ID heredado: \texttt{RNF-001}.
\item[Verificación] \texttt{AC-QR-001}
\end{description}
\vspace{0.25em}\hrule\vspace{0.65em}
\noindent\textbf{QR-002 --- Auditabilidad y trazabilidad de operaciones críticas}\par
\begin{description}[style=nextline,leftmargin=3.15cm,labelwidth=2.8cm,font=\normalfont\bfseries]
\item[Obligación] El sistema debe conservar, para cada operación crítica y resultado oficial publicado, actor cuando corresponda, marca temporal, datos de origen y reglas aplicadas suficientes para reconstruir la derivación o secuencia relevante.
\item[Prioridad] Alta
\item[Procedencia] Consolidación de RNF-002, RNF-004, RAC-002; ID heredado: \texttt{RNF-002, RNF-004, RAC-002}.
\item[Verificación] \texttt{AC-QR-002}
\end{description}
\vspace{0.25em}\hrule\vspace{0.65em}
\noindent\textbf{QR-003 --- Integridad de derivaciones del dominio}\par
\begin{description}[style=nextline,leftmargin=3.15cm,labelwidth=2.8cm,font=\normalfont\bfseries]
\item[Obligación] El sistema debe impedir que datos invalidados por cancelación, descalificación, exclusión o recálculo válido continúen interviniendo en resultados, clasificaciones, subastas o determinaciones económicas posteriores.
\item[Prioridad] Alta
\item[Procedencia] Consolidación de RNF-003, RAC-003; ID heredado: \texttt{RNF-003, RAC-003}.
\item[Verificación] \texttt{AC-QR-003}
\end{description}
\vspace{0.25em}\hrule\vspace{0.65em}
\noindent\textbf{QR-004 --- Transparencia pública controlada}\par
\begin{description}[style=nextline,leftmargin=3.15cm,labelwidth=2.8cm,font=\normalfont\bfseries]
\item[Obligación] El sistema debe publicar la información declarada pública por las reglas del dominio y debe reservar la información declarada restringida.
\item[Prioridad] Alta
\item[Procedencia] Consolidación de RAC-005; ID heredado: \texttt{RAC-005}.
\item[Verificación] \texttt{AC-QR-004}
\end{description}
\vspace{0.25em}\hrule\vspace{0.65em}
\noindent\textbf{QR-005 --- Permanencia y consultabilidad histórica}\par
\begin{description}[style=nextline,leftmargin=3.15cm,labelwidth=2.8cm,font=\normalfont\bfseries]
\item[Obligación] La información obligatoria publicada del evento no debe retirarse durante el periodo de vigencia definido por las reglas del dominio y debe permanecer consultable en el historial cuando corresponda.
\item[Prioridad] Alta
\item[Procedencia] Consolidación de RNF-005, RAC-004; ID heredado: \texttt{RNF-005, RAC-004}.
\item[Verificación] \texttt{AC-QR-005}
\end{description}
\vspace{0.25em}\hrule\vspace{0.65em}
\noindent\textbf{QR-006 --- Reproducibilidad numérica sin redondeo intermedio no prescrito}\par
\begin{description}[style=nextline,leftmargin=3.15cm,labelwidth=2.8cm,font=\normalfont\bfseries]
\item[Obligación] Los cálculos del dominio deben obedecer \texttt{NUM-001}: punto flotante binario de 64 bits para magnitudes analíticas, importes monetarios enteros en unidades completas de la moneda aplicable y reconciliación exacta de toda distribución monetaria mediante \texttt{CALC-MON-001}; no deben aplicar redondeo de dominio intermedio salvo prescripción CALC expresa.
\item[Prioridad] Alta
\item[Procedencia] Consolidación de RNF-006, RAC-001; ID heredado: \texttt{RNF-006, RAC-001}.
\item[Verificación] \texttt{AC-QR-006}
\end{description}
\vspace{0.25em}\hrule\vspace{0.65em}
\noindent\textbf{QR-007 --- Retroalimentación perceptible al inicio de acciones interactivas}\par
\begin{description}[style=nextline,leftmargin=3.15cm,labelwidth=2.8cm,font=\normalfont\bfseries]
\item[Obligación] En una muestra representativa de acciones interactivas aceptadas, el tiempo desde la activación hasta la primera señal perceptible de recepción debe ser como máximo 100 ms en el percentil 95 y ninguna medición individual debe exceder 250 ms en \texttt{ENV-UX-001}.
\item[Prioridad] Alta
\item[Procedencia] Consolidación de RNF-007; ID heredado: \texttt{RNF-007}.
\item[Verificación] \texttt{AC-QR-007}
\end{description}
\vspace{0.25em}\hrule\vspace{0.65em}
\noindent\textbf{QR-008 --- Visibilidad de espera y bloqueo por incompatibilidad operativa}\par
\begin{description}[style=nextline,leftmargin=3.15cm,labelwidth=2.8cm,font=\normalfont\bfseries]
\item[Obligación] Cuando una operación aceptada no termine dentro de 500 ms, la interfaz debe mantener un estado de espera explícito hasta recibir resultado y debe bloquear únicamente las operaciones incompatibles conforme a \texttt{OPS-001}; las consultas permanecen habilitadas por defecto.
\item[Prioridad] Alta
\item[Procedencia] Consolidación de RNF-008; ID heredado: \texttt{RNF-008}.
\item[Verificación] \texttt{AC-QR-008}
\end{description}
\vspace{0.25em}\hrule\vspace{0.65em}
\noindent\textbf{QR-009 --- Percepción inicial positiva en pantallas críticas}\par
\begin{description}[style=nextline,leftmargin=3.15cm,labelwidth=2.8cm,font=\normalfont\bfseries]
\item[Obligación] Las pantallas críticas deben alcanzar, en el protocolo UX definido, una valoración positiva o muy positiva en al menos 70\% de una muestra mínima de 20 usuarios representativos tanto a los 30 segundos como a los 5 minutos.
\item[Prioridad] Alta
\item[Procedencia] Consolidación de RAC-006; ID heredado: \texttt{RAC-006}.
\item[Verificación] \texttt{AC-QR-009}
\end{description}
\vspace{0.25em}\hrule\vspace{0.65em}
\noindent\textbf{QR-010 --- Jerarquía visual de información crítica}\par
\begin{description}[style=nextline,leftmargin=3.15cm,labelwidth=2.8cm,font=\normalfont\bfseries]
\item[Obligación] Las pantallas críticas deben diferenciar visualmente el estado actual, las acciones habilitadas y bloqueadas, la causa del bloqueo y, cuando apliquen, importe bruto, comisión, retención, devolución e importe neto.
\item[Prioridad] Alta
\item[Procedencia] Consolidación de RNF-009, RAC-007, RIE-015; ID heredado: \texttt{RNF-009, RAC-007, RIE-015}.
\item[Verificación] \texttt{AC-QR-010}
\end{description}
\vspace{0.25em}\hrule\vspace{0.65em}
\noindent\textbf{QR-011 --- Legibilidad del contenido crítico}\par
\begin{description}[style=nextline,leftmargin=3.15cm,labelwidth=2.8cm,font=\normalfont\bfseries]
\item[Obligación] El texto y los importes críticos deben cumplir WCAG 2.2 nivel AA y conservar legibilidad y funcionalidad en los viewports y niveles de zoom definidos por \texttt{ENV-UX-001}.
\item[Prioridad] Alta
\item[Procedencia] Consolidación de RNF-010, RAC-008; ID heredado: \texttt{RNF-010, RAC-008}.
\item[Verificación] \texttt{AC-QR-011}
\end{description}
\vspace{0.25em}\hrule\vspace{0.65em}
\noindent\textbf{QR-012 --- Contraste y codificación redundante de estados críticos}\par
\begin{description}[style=nextline,leftmargin=3.15cm,labelwidth=2.8cm,font=\normalfont\bfseries]
\item[Obligación] Los estados, acciones, bloqueos e importes críticos deben cumplir los criterios de contraste aplicables de WCAG 2.2 nivel AA y ninguna distinción crítica debe depender exclusivamente del color.
\item[Prioridad] Alta
\item[Procedencia] Consolidación de RNF-011, RAC-009; ID heredado: \texttt{RNF-011, RAC-009}.
\item[Verificación] \texttt{AC-QR-012}
\end{description}
\vspace{0.25em}\hrule\vspace{0.65em}
\noindent\textbf{QR-013 --- Consistencia visual y terminológica}\par
\begin{description}[style=nextline,leftmargin=3.15cm,labelwidth=2.8cm,font=\normalfont\bfseries]
\item[Obligación] Las pantallas que representen el mismo concepto del dominio deben usar la misma terminología canónica y patrones de interacción coherentes, salvo una excepción documentada.
\item[Prioridad] Alta
\item[Procedencia] Consolidación de RNF-012, RAC-010; ID heredado: \texttt{RNF-012, RAC-010}.
\item[Verificación] \texttt{AC-QR-013}
\end{description}
\vspace{0.25em}\hrule\vspace{0.65em}
\noindent\textbf{QR-014 --- Prevención de error en acciones sensibles}\par
\begin{description}[style=nextline,leftmargin=3.15cm,labelwidth=2.8cm,font=\normalfont\bfseries]
\item[Obligación] Antes de confirmar una acción sensible, la interfaz debe proporcionar contexto y frenos proporcionales a sus consecuencias y debe impedir confirmaciones ambiguas.
\item[Prioridad] Alta
\item[Procedencia] Consolidación de RNF-013, RAC-011; ID heredado: \texttt{RNF-013, RAC-011}.
\item[Verificación] \texttt{AC-QR-014}
\end{description}
\vspace{0.25em}\hrule\vspace{0.65em}
\noindent\textbf{QR-015 --- Claridad previa de consecuencias}\par
\begin{description}[style=nextline,leftmargin=3.15cm,labelwidth=2.8cm,font=\normalfont\bfseries]
\item[Obligación] Antes de confirmar una acción sensible, la interfaz debe informar las consecuencias operativas, económicas o jurídicas aplicables que el sistema conozca por sus reglas vigentes.
\item[Prioridad] Alta
\item[Procedencia] Consolidación de RNF-014, RAC-012; ID heredado: \texttt{RNF-014, RAC-012}.
\item[Verificación] \texttt{AC-QR-015}
\end{description}
\vspace{0.25em}\hrule\vspace{0.65em}
\section{Verificación y aceptación}
\subsection{Estrategia de verificación}
Los RF y RC-EXT se verifican principalmente mediante pruebas funcionales de caja negra y análisis de trazas autorizadas. Los CALC se verifican mediante conjuntos de datos de referencia y comparación numérica bajo NUM-001. Los UI se verifican mediante pruebas de interacción y control de acceso observable. Los QR se verifican mediante inspección estructurada, medición o pruebas con usuarios en ENV-UX-001 según corresponda.

Un criterio AC define el oráculo de aceptación. La automatización del caso de prueba puede residir en una suite separada; la SRS no prescribe framework de pruebas.

\subsection{Catálogo de criterios de aceptación}

\noindent\textbf{AC-RF-EVT-001 --- Crear evento}\par
\textbf{Método:} Prueba funcional de caja negra.\par
\textbf{Oráculo:} Cuando se registran los datos iniciales del evento y sus reglas publicadas, el sistema crea el evento y lo muestra en estado \texttt{creado}.\par\vspace{0.45em}
\noindent\textbf{AC-RF-EVT-002 --- Abrir inscripciones del evento}\par
\textbf{Método:} Prueba funcional de caja negra.\par
\textbf{Oráculo:} Si faltan datos públicos o reglas publicadas del evento, la apertura de inscripciones se rechaza; si ya están publicados, el estado cambia de \texttt{creado} a \texttt{inscripciones}.\par\vspace{0.45em}
\noindent\textbf{AC-RF-EVT-003 --- Vocabulario oficial de estados del evento}\par
\textbf{Método:} Prueba funcional de caja negra.\par
\textbf{Oráculo:} Crear o consultar eventos en cada fase y comprobar que ningún estado oficial distinto de los siete valores definidos puede registrarse.\par\vspace{0.45em}
\noindent\textbf{AC-RF-EVT-004 --- Bloqueo de transiciones no autorizadas}\par
\textbf{Método:} Prueba funcional de caja negra.\par
\textbf{Oráculo:} Intentar cada transición no definida en el catálogo de transiciones; la operación se rechaza y el estado previo permanece sin cambio.\par\vspace{0.45em}
\noindent\textbf{AC-RF-EVT-005 --- Cerrar inscripciones}\par
\textbf{Método:} Prueba funcional de caja negra.\par
\textbf{Oráculo:} Al ejecutarse la transición \texttt{inscripciones -> velocidad}, ya no se aceptan nuevas inscripciones.\par\vspace{0.45em}
\noindent\textbf{AC-RF-EVT-006 --- Precondición para la publicación oficial}\par
\textbf{Método:} Prueba funcional de caja negra.\par
\textbf{Oráculo:} Si faltan datos necesarios para determinar y publicar resultados oficiales, el evento no puede alcanzar el estado \texttt{publicado}.\par\vspace{0.45em}
\noindent\textbf{AC-RF-EVT-007 --- Publicar resultados oficiales del evento}\par
\textbf{Método:} Prueba funcional de caja negra.\par
\textbf{Oráculo:} Al publicar resultados, el evento queda en estado \texttt{publicado} e identifica las referencias a los datos de origen y a las reglas del evento.\par\vspace{0.45em}
\noindent\textbf{AC-RF-EVT-008 --- Inmutabilidad lógica de los resultados publicados}\par
\textbf{Método:} Prueba funcional de caja negra.\par
\textbf{Oráculo:} Los intentos de modificación directa sobre resultados oficiales de un evento en estado \texttt{publicado} se bloquean y quedan registrados.\par\vspace{0.45em}
\noindent\textbf{AC-RF-EVT-009 --- Registro histórico de hitos del evento}\par
\textbf{Método:} Prueba funcional de caja negra.\par
\textbf{Oráculo:} Para cada cambio de estado crítico existe un registro histórico consultable con marca temporal y actor responsable.\par\vspace{0.45em}
\noindent\textbf{AC-RF-EVT-010 --- Estado inicial obligatorio}\par
\textbf{Método:} Prueba funcional de caja negra.\par
\textbf{Oráculo:} Al crear un evento, su estado inicial registrado es siempre \texttt{creado}.\par\vspace{0.45em}
\noindent\textbf{AC-RF-EVT-011 --- Transición de apertura de inscripciones}\par
\textbf{Método:} Prueba funcional de caja negra.\par
\textbf{Oráculo:} Un evento en \texttt{creado} puede pasar a \texttt{inscripciones} y la transición queda registrada.\par\vspace{0.45em}
\noindent\textbf{AC-RF-EVT-012 --- Transición de inscripciones a velocidad}\par
\textbf{Método:} Prueba funcional de caja negra.\par
\textbf{Oráculo:} Con guardas satisfechas, ejecutar la transición y comprobar que el evento queda en \texttt{velocidad}; con una guarda incumplida, la transición se rechaza.\par\vspace{0.45em}
\noindent\textbf{AC-RF-EVT-013 --- Cierre automático de inscripciones al iniciar velocidad}\par
\textbf{Método:} Prueba funcional de caja negra.\par
\textbf{Oráculo:} Ejecutar la transición y comprobar que una nueva inscripción posterior es rechazada.\par\vspace{0.45em}
\noindent\textbf{AC-RF-EVT-014 --- Verificación de identidad previa al paso a velocidad}\par
\textbf{Método:} Prueba funcional de caja negra.\par
\textbf{Oráculo:} Si existe al menos una inscripción no cancelada asociada a un usuario no verificado, la transición a \texttt{velocidad} se bloquea; cuando esas identidades ya están verificadas, la transición puede ejecutarse conforme a las demás reglas aplicables.\par\vspace{0.45em}
\noindent\textbf{AC-RF-EVT-015 --- Transición a la fase de carrera}\par
\textbf{Método:} Prueba funcional de caja negra.\par
\textbf{Oráculo:} Un evento en \texttt{velocidad} puede pasar a \texttt{carrera} y la transición queda registrada.\par\vspace{0.45em}
\noindent\textbf{AC-RF-EVT-016 --- Transición a la fase de subasta}\par
\textbf{Método:} Prueba funcional de caja negra.\par
\textbf{Oráculo:} Si existe un vehículo que deba ir a subasta y todavía no ha sido entregado al promotor ni ha quedado fuera por cancelación o descalificación, la transición a \texttt{subasta} se bloquea; cuando esa condición ya está resuelta, la transición puede ejecutarse.\par\vspace{0.45em}
\noindent\textbf{AC-RF-EVT-017 --- Transición a publicado}\par
\textbf{Método:} Prueba funcional de caja negra.\par
\textbf{Oráculo:} Si falta información necesaria para determinar y publicar los resultados oficiales o la liquidación económica aplicable, la transición se bloquea; cuando esa información existe, el evento pasa a \texttt{publicado}.\par\vspace{0.45em}
\noindent\textbf{AC-RF-EVT-018 --- Registro del estado cancelado}\par
\textbf{Método:} Prueba funcional de caja negra.\par
\textbf{Oráculo:} Cancelar un evento válido; la consulta posterior devuelve estado \texttt{cancelado} y no existe publicación oficial de resultados del evento.\par\vspace{0.45em}
\noindent\textbf{AC-RF-EVT-019 --- Reversión económica inmediata por cancelación del evento}\par
\textbf{Método:} Prueba funcional de caja negra.\par
\textbf{Oráculo:} Preparar un evento con importes retenidos o cobrados sujetos a devolución; al cancelarlo, verificar las devoluciones y liberaciones aplicables y su registro temporal.\par\vspace{0.45em}
\noindent\textbf{AC-RF-EVT-020 --- Estado terminal publicado}\par
\textbf{Método:} Prueba funcional de caja negra.\par
\textbf{Oráculo:} Todo intento de transición desde \texttt{publicado} se bloquea.\par\vspace{0.45em}
\noindent\textbf{AC-RF-EVT-021 --- Estado terminal cancelado}\par
\textbf{Método:} Prueba funcional de caja negra.\par
\textbf{Oráculo:} Todo intento de transición desde \texttt{cancelado} se bloquea.\par\vspace{0.45em}
\noindent\textbf{AC-RF-EVT-022 --- Bloqueo de transiciones no definidas}\par
\textbf{Método:} Prueba funcional de caja negra.\par
\textbf{Oráculo:} Para cualquier par origen-destino no definido como válido, el sistema responde con un rechazo explícito.\par\vspace{0.45em}
\noindent\textbf{AC-RF-EVT-023 --- Registro de transiciones de estado}\par
\textbf{Método:} Prueba funcional de caja negra.\par
\textbf{Oráculo:} Al consultar el historial del evento se observan registros completos de transiciones válidas e intentos inválidos.\par\vspace{0.45em}
\noindent\textbf{AC-RF-EVT-024 --- Principio de operaciones habilitadas según el estado}\par
\textbf{Método:} Prueba funcional de caja negra.\par
\textbf{Oráculo:} Para cada estado oficial, solo se permiten sus operaciones asociadas; cualquier operación fuera de ventana se rechaza.\par\vspace{0.45em}
\noindent\textbf{AC-RF-EVT-025 --- Operaciones en estado inscripciones}\par
\textbf{Método:} Prueba funcional de caja negra.\par
\textbf{Oráculo:} Mientras el evento está en \texttt{inscripciones}, se aceptan altas y ediciones de inscripción, aportaciones directas a la bolsa, y solo cuando la función esté activa, aportaciones a participantes y operaciones de subasta de acceso; al pasar a \texttt{velocidad}, esas operaciones dejan de aceptarse salvo las consultas públicas que sigan vigentes.\par\vspace{0.45em}
\noindent\textbf{AC-RF-EVT-026 --- Operaciones en estado velocidad}\par
\textbf{Método:} Prueba funcional de caja negra.\par
\textbf{Oráculo:} Mientras el evento está en \texttt{velocidad}, se aceptan capturas y ediciones de resultados de velocidad y pueden ejecutarse las decisiones permitidas antes de \texttt{carrera}; fuera de ese estado, esas operaciones se bloquean.\par\vspace{0.45em}
\noindent\textbf{AC-RF-EVT-027 --- Ventana de inspección comunitaria y protestas antes de carrera}\par
\textbf{Método:} Prueba funcional de caja negra.\par
\textbf{Oráculo:} En \texttt{inscripciones} y \texttt{velocidad} el sistema permite esas operaciones; al pasar a \texttt{carrera}, deja de permitir nuevas inspecciones, objeciones, votaciones y decisiones de cancelación o descalificación por esa via.\par\vspace{0.45em}
\noindent\textbf{AC-RF-EVT-028 --- Operaciones en estado carrera}\par
\textbf{Método:} Prueba funcional de caja negra.\par
\textbf{Oráculo:} Mientras el evento está en \texttt{carrera}, se aceptan capturas y ediciones de resultados de carrera y se puede registrar la entrega del vehículo al promotor; fuera de ese estado, esas operaciones se bloquean.\par\vspace{0.45em}
\noindent\textbf{AC-RF-EVT-029 --- Operaciones en estado subasta}\par
\textbf{Método:} Prueba funcional de caja negra.\par
\textbf{Oráculo:} Mientras el evento está en \texttt{subasta}, se aceptan operaciones de subasta, de entrega y de disputa de entrega; fuera de ese estado, se bloquean.\par\vspace{0.45em}
\noindent\textbf{AC-RF-EVT-030 --- Operaciones en estado creado}\par
\textbf{Método:} Prueba funcional de caja negra.\par
\textbf{Oráculo:} Mientras el evento está en \texttt{creado}, se permite completar la información pública requerida para el evento; al pasar a \texttt{inscripciones}, esa operación queda sujeta a las reglas del nuevo estado.\par\vspace{0.45em}
\noindent\textbf{AC-RF-EVT-031 --- Consulta de resultados oficiales en estado publicado}\par
\textbf{Método:} Prueba funcional de caja negra.\par
\textbf{Oráculo:} Consultar un evento publicado y comprobar la disponibilidad de sus resultados oficiales.\par\vspace{0.45em}
\noindent\textbf{AC-RF-EVT-032 --- Bloqueo de nuevas capturas en estado publicado}\par
\textbf{Método:} Prueba funcional de caja negra.\par
\textbf{Oráculo:} Intentar capturas representativas después de la publicación; todas las no autorizadas se rechazan y los datos oficiales permanecen sin cambio.\par\vspace{0.45em}
\noindent\textbf{AC-RF-EVT-033 --- Operaciones en estado cancelado}\par
\textbf{Método:} Prueba funcional de caja negra.\par
\textbf{Oráculo:} En \texttt{cancelado} se bloquean las operaciones activas del evento y solo permanecen disponibles las consultas permitidas.\par\vspace{0.45em}
\noindent\textbf{AC-RF-EVT-034 --- Rechazo y mensaje ante operación inválida por estado}\par
\textbf{Método:} Prueba funcional de caja negra.\par
\textbf{Oráculo:} Ante cualquier intento de operación inválida por estado, el evento no se modifica y el sistema entrega un rechazo explícito con el estado actual y el motivo.\par\vspace{0.45em}
\noindent\textbf{AC-RF-EVT-035 --- Ventana para aportar a evento o participante del evento}\par
\textbf{Método:} Prueba funcional de caja negra.\par
\textbf{Oráculo:} En \texttt{inscripciones} se aceptan nuevas aportaciones al evento y a participantes; en cualquier otro estado se rechazan.\par\vspace{0.45em}
\noindent\textbf{AC-RF-EVT-036 --- Cierre de aportaciones del evento al entrar en velocidad}\par
\textbf{Método:} Prueba funcional de caja negra.\par
\textbf{Oráculo:} Tras la transición a \texttt{velocidad}, no existe ninguna operación válida para anadir nuevas aportaciones al evento o a sus participantes.\par\vspace{0.45em}
\noindent\textbf{AC-RF-EVT-037 --- Integracion de aportaciones acumuladas por participante}\par
\textbf{Método:} Prueba funcional de caja negra.\par
\textbf{Oráculo:} Al registrar una aportación válida y cobrarla, el total de contribución computable del participante aumenta exactamente en el monto aportado.\par\vspace{0.45em}
\noindent\textbf{AC-RF-EVT-038 --- Bloqueo y mensaje por aportación fuera de ventana}\par
\textbf{Método:} Prueba funcional de caja negra.\par
\textbf{Oráculo:} Ante un intento fuera de ventana, no cambia ningun total y se presenta un mensaje explicito de rechazo.\par\vspace{0.45em}
\noindent\textbf{AC-RF-EVT-039 --- Bloqueo preventivo de acciones no permitidas en la IU}\par
\textbf{Método:} Prueba funcional de caja negra.\par
\textbf{Oráculo:} En cualquier pantalla operativa, las acciones no permitidas para el estado actual no pueden ejecutarse desde la IU.\par\vspace{0.45em}
\noindent\textbf{AC-RF-EVT-040 --- Impedir la solicitud de acciones no permitidas}\par
\textbf{Método:} Prueba funcional de caja negra.\par
\textbf{Oráculo:} Para acciones no permitidas por estado, la IU no genera la solicitud operativa y el usuario no puede iniciarla.\par\vspace{0.45em}
\noindent\textbf{AC-RF-EVT-051 --- Configuración obligatoria del límite de carrera}\par
\textbf{Método:} Prueba funcional de caja negra.\par
\textbf{Oráculo:} Intentar abrir inscripciones sin tipo de límite, con valor no positivo y, para distancia, sin longitud positiva de vuelta; las operaciones se rechazan. Con \texttt{tiempo} y tiempo máximo positivo, o con \texttt{distancia}, distancia objetivo positiva y longitud de vuelta positiva, la configuración es aceptable.\par\vspace{0.45em}
\noindent\textbf{AC-RF-PAY-001 --- Consolidación del pozo económico oficial}\par
\textbf{Método:} Prueba funcional de caja negra.\par
\textbf{Oráculo:} Al cerrarse la ventana económica del evento, el sistema determina un único pozo económico oficial a partir del pozo inicial publicado y de las cantidades efectivamente cobradas y válidamente aceptadas.\par\vspace{0.45em}
\noindent\textbf{AC-RF-PAY-002 --- Conservación del registro económico oficial}\par
\textbf{Método:} Prueba funcional de caja negra.\par
\textbf{Oráculo:} En un evento con movimiento económico, el sistema permite consultar el total oficial consolidado y su correspondencia con el pozo inicial, con los cobros, con las retenciones, con los pagos y con las devoluciones que lo integran o afectan.\par\vspace{0.45em}
\noindent\textbf{AC-RF-PAY-003 --- Determinación del premio individual del evento}\par
\textbf{Método:} Prueba funcional de caja negra.\par
\textbf{Oráculo:} Al ejecutarse la determinación económica del evento, existe un premio individual definido para cada competidor computable con base en su desempeño, en las reglas publicadas y en los datos oficiales del evento.\par\vspace{0.45em}
\noindent\textbf{AC-RF-PAY-004 --- Registro trazable del premio individual del evento}\par
\textbf{Método:} Prueba funcional de caja negra.\par
\textbf{Oráculo:} Para cada competidor computable existe un registro consultable de su premio individual y de los datos oficiales aplicados.\par\vspace{0.45em}
\noindent\textbf{AC-RF-PAY-005 --- Distribución económica posterior del premio individual del evento}\par
\textbf{Método:} Prueba funcional de caja negra.\par
\textbf{Oráculo:} Cuando un competidor computable tenga aportantes asociados, el sistema genera la distribución económica correspondiente entre el competidor y dichos aportantes.\par\vspace{0.45em}
\noindent\textbf{AC-RF-PAY-006 --- Liquidación económica oficial del evento}\par
\textbf{Método:} Prueba funcional de caja negra.\par
\textbf{Oráculo:} Una vez consolidados el pozo económico oficial, los resultados del dominio y las devoluciones aplicables, el sistema genera una liquidación oficial del evento con premios individuales, distribución posterior y devoluciones correspondientes.\par\vspace{0.45em}
\noindent\textbf{AC-RF-PAY-007 --- Devolución inmediata por cancelación o descalificación previa a carrera}\par
\textbf{Método:} Prueba funcional de caja negra.\par
\textbf{Oráculo:} Si el evento se cancela o una inscripción queda fuera antes de \texttt{carrera}, el sistema ejecuta de inmediato la devolución de los cobros aplicables, incluidas las aportaciones de terceros asociadas a ese participante dentro del evento.\par\vspace{0.45em}
\noindent\textbf{AC-RF-PAY-008 --- Pago automático e inmediato a todos los beneficiarios del evento}\par
\textbf{Método:} Prueba funcional de caja negra.\par
\textbf{Oráculo:} Cuando la liquidación económica oficial del evento queda determinada, el sistema ejecuta de inmediato el pago a cada beneficiario según la distribución posterior aplicable.\par\vspace{0.45em}
\noindent\textbf{AC-RF-PAY-009 --- Publicación de resultados económicos oficiales}\par
\textbf{Método:} Prueba funcional de caja negra.\par
\textbf{Oráculo:} La publicación oficial incluye los premios individuales por competidor computable y la distribución económica posterior asociada, con anonimización de terceros cuando corresponda.\par\vspace{0.45em}
\noindent\textbf{AC-RF-AUC-001 --- Apertura de subasta de acceso de espectadores al evento}\par
\textbf{Método:} Prueba funcional de caja negra.\par
\textbf{Oráculo:} En un evento con subasta de acceso para espectadores prevista, el sistema registra la cantidad disponible de accesos y permite recibir pujas secretas mientras la subasta permanezca abierta.\par\vspace{0.45em}
\noindent\textbf{AC-RF-AUC-002 --- Registro de pujas secretas en subasta de acceso con retención suficiente}\par
\textbf{Método:} Prueba funcional de caja negra.\par
\textbf{Oráculo:} Durante la subasta de acceso, cada puja aceptada queda registrada como puja válida solo si existe una retención suficiente asociada a ella, y su monto individual no es visible antes del cierre.\par\vspace{0.45em}
\noindent\textbf{AC-RF-AUC-003 --- Una sola puja vigente por usuario en subasta de acceso}\par
\textbf{Método:} Prueba funcional de caja negra.\par
\textbf{Oráculo:} Si un mismo usuario participa en una subasta de acceso, el sistema conserva una sola puja vigente asociada a ese usuario para esa subasta.\par\vspace{0.45em}
\noindent\textbf{AC-RF-AUC-004 --- Sustitución solo al alza de la puja vigente de acceso}\par
\textbf{Método:} Prueba funcional de caja negra.\par
\textbf{Oráculo:} Intentar sustituir una puja por importes inferior, igual y superior; solo el superior se acepta y pasa a ser la única puja vigente.\par\vspace{0.45em}
\noindent\textbf{AC-RF-AUC-005 --- Ajuste de retención por actualización válida de la puja de acceso}\par
\textbf{Método:} Prueba funcional de caja negra.\par
\textbf{Oráculo:} Aumentar una puja conocida y comprobar que la retención vigente corresponde al nuevo importe y sustituye a la anterior.\par\vspace{0.45em}
\noindent\textbf{AC-RF-AUC-006 --- Publicación del conteo de pujas vigentes en subasta de acceso}\par
\textbf{Método:} Prueba funcional de caja negra.\par
\textbf{Oráculo:} Durante la subasta de acceso, el sistema publica el conteo actualizado de pujas vigentes y no muestra los montos individuales asociados.\par\vspace{0.45em}
\noindent\textbf{AC-RF-AUC-007 --- Ordenamiento de pujas válidas de acceso}\par
\textbf{Método:} Prueba funcional de caja negra.\par
\textbf{Oráculo:} Cerrar una subasta con un conjunto conocido de pujas válidas y comprobar que el orden resultante es descendente por importe.\par\vspace{0.45em}
\noindent\textbf{AC-RF-AUC-008 --- Adjudicación de accesos disponibles}\par
\textbf{Método:} Prueba funcional de caja negra.\par
\textbf{Oráculo:} Con $n$ accesos y más de $n$ pujas válidas, cerrar la subasta y comprobar que existen exactamente $n$ adjudicaciones a las posiciones superiores del orden oficial.\par\vspace{0.45em}
\noindent\textbf{AC-RF-AUC-009 --- Determinación del importe de cobro del acceso adjudicado}\par
\textbf{Método:} Prueba funcional de caja negra.\par
\textbf{Oráculo:} Usar una secuencia conocida de pujas y comprobar, para cada adjudicación, que el importe de cobro coincide con la siguiente puja válida inferior aplicable.\par\vspace{0.45em}
\noindent\textbf{AC-RF-AUC-010 --- Publicación del resultado de la subasta de acceso}\par
\textbf{Método:} Prueba funcional de caja negra.\par
\textbf{Oráculo:} Cerrar la subasta y comprobar que la publicación oficial identifica cada adjudicatario y el importe de cobro calculado.\par\vspace{0.45em}
\noindent\textbf{AC-RF-AUC-011 --- Notificación de acceso adjudicado con QR}\par
\textbf{Método:} Prueba funcional de caja negra.\par
\textbf{Oráculo:} Cuando una persona resulta ganadora de un acceso, el sistema deja disponible para ella una notificación asociada a ese acceso adjudicado que incluye un código QR y el importe correspondiente.\par\vspace{0.45em}
\noindent\textbf{AC-RF-AUC-012 --- Retencion, cobro y liberación en subasta de acceso}\par
\textbf{Método:} Prueba funcional de caja negra.\par
\textbf{Oráculo:} Al cerrarse la subasta de acceso, el sistema cobra automáticamente a las personas ganadoras el importe adjudicado y libera de inmediato las retenciones asociadas a las pujas no ganadoras.\par\vspace{0.45em}
\noindent\textbf{AC-RF-AUC-013 --- Devolución inmediata de accesos adjudicados por cancelación del evento}\par
\textbf{Método:} Prueba funcional de caja negra.\par
\textbf{Oráculo:} Si un evento con subasta de acceso ya cobrada pasa a \texttt{cancelado}, el sistema devuelve de inmediato los importes cobrados por accesos adjudicados y libera cualquier retención aún vigente.\par\vspace{0.45em}
\noindent\textbf{AC-RF-AUC-014 --- Apertura de subastas de vehículos después de la carrera}\par
\textbf{Método:} Prueba funcional de caja negra.\par
\textbf{Oráculo:} Al entrar el evento en estado \texttt{subasta}, solo los vehículos asociados a inscripciones computables que participaron en el evento y ya fueron recibidos por el promotor quedan disponibles para subasta.\par\vspace{0.45em}
\noindent\textbf{AC-RF-AUC-015 --- Registro de pujas selladas en subastas de vehículos}\par
\textbf{Método:} Prueba funcional de confidencialidad con dos usuarios y un dueño.\par
\textbf{Oráculo:} Registrar pujas con usuarios distintos y comprobar que cada usuario solo consulta su propia puja vigente antes del cierre y no obtiene importes de terceros.\par\vspace{0.45em}
\noindent\textbf{AC-RF-AUC-016 --- Verificación reforzada máxima obligatoria para pujar en vehículos}\par
\textbf{Método:} Prueba de autorización por niveles de verificación.\par
\textbf{Oráculo:} Un usuario con verificación reforzada activa, términos vigentes y sin exclusión puede pujar aunque no tenga otro papel funcional; usuarios por debajo del nivel máximo o excluidos son rechazados.\par\vspace{0.45em}
\noindent\textbf{AC-RF-AUC-017 --- Una sola puja vigente y sustitución únicamente al alza en cada vehículo}\par
\textbf{Método:} Prueba funcional con secuencia de sustituciones.\par
\textbf{Oráculo:} Para el mismo usuario y vehículo, una puja inferior o igual a la vigente es rechazada; una superior la sustituye y queda como única puja vigente.\par\vspace{0.45em}
\noindent\textbf{AC-RF-AUC-018 --- Publicación del resultado oficial de la subasta solo después del cierre}\par
\textbf{Método:} Prueba funcional antes y después del cierre.\par
\textbf{Oráculo:} Antes del cierre de la fase no se publica ningún resultado de vehículo; al cerrarla se publican para cada vehículo el precio oficial, el adjudicatario y la naturaleza de transferencia o autoadjudicación.\par\vspace{0.45em}
\noindent\textbf{AC-RF-AUC-019 --- Determinación por segundo precio de la subasta sellada de vehículos}\par
\textbf{Método:} Prueba funcional con pujas conocidas, incluida una puja del dueño.\par
\textbf{Oráculo:} Al cerrar la fase, para cada vehículo con al menos dos pujas válidas la mayor puja vigente determina al adjudicatario y el segundo importe vigente más alto determina el precio oficial, sin excluir la puja del dueño del ordenamiento.\par\vspace{0.45em}
\noindent\textbf{AC-RF-AUC-020 --- Entrega obligatoria del vehículo al promotor al terminar la carrera}\par
\textbf{Método:} Prueba funcional de caja negra.\par
\textbf{Oráculo:} Para un vehículo que deba ir a subasta, el sistema permite dejar constancia de su entrega al promotor al terminar la carrera.\par\vspace{0.45em}
\noindent\textbf{AC-RF-AUC-021 --- Descalificación, morosidad y exclusión del dueño por no entrega o por modificación antes de la entrega al promotor}\par
\textbf{Método:} Prueba funcional de caja negra.\par
\textbf{Oráculo:} Si el promotor registra la no entrega o la modificación indebida antes de recibir el vehículo, el sistema permite descalificar al competidor, declararlo moroso y aplicarle la exclusión correspondiente del evento, del sistema y de todos los campeonatos aplicables.\par\vspace{0.45em}
\noindent\textbf{AC-RF-AUC-022 --- Custodia del promotor y parc ferme del vehículo subastado}\par
\textbf{Método:} Prueba funcional de caja negra.\par
\textbf{Oráculo:} Para un vehículo recibido por el promotor para subasta, el sistema conserva su condición de control del evento y su sujeción a \texttt{parc ferme} y a transferencia obligatoria.\par\vspace{0.45em}
\noindent\textbf{AC-RF-AUC-023 --- Retención del principal de pujas de vehículos de usuarios distintos del dueño}\par
\textbf{Método:} Prueba económica con comprador externo y dueño.\par
\textbf{Oráculo:} La puja externa mantiene retenido su importe vigente; la puja del dueño sobre su propio vehículo no genera retención de principal.\par\vspace{0.45em}
\noindent\textbf{AC-RF-AUC-024 --- Cobro del principal al adjudicatario externo y liberación de retenciones no ganadoras}\par
\textbf{Método:} Prueba económica de cierre con adjudicatario externo.\par
\textbf{Oráculo:} Al cerrar, el sistema cobra al adjudicatario externo el precio oficial y libera las retenciones no ganadoras conforme a las reglas vigentes.\par\vspace{0.45em}
\noindent\textbf{AC-RF-AUC-025 --- Pago íntegro del precio oficial al dueño cuando gana un tercero}\par
\textbf{Método:} Prueba económica de cierre con venta a tercero.\par
\textbf{Oráculo:} Después de cobrar el precio oficial al comprador externo, el dueño recibe exactamente ese precio y la comisión del comprador no se descuenta del vendedor.\par\vspace{0.45em}
\noindent\textbf{AC-RF-AUC-026 --- Estado de promotor deudor únicamente en adjudicación a un tercero}\par
\textbf{Método:} Prueba de estados con resultado externo y autoadjudicación.\par
\textbf{Oráculo:} La adjudicación a tercero activa la condición de deudor hasta la confirmación de entrega; la autoadjudicación no la activa.\par\vspace{0.45em}
\noindent\textbf{AC-RF-AUC-027 --- Confirmación de entrega por adjudicatario distinto del dueño previo}\par
\textbf{Método:} Prueba de interacción con adjudicatario externo y dueño autoadjudicatario.\par
\textbf{Oráculo:} El adjudicatario externo puede confirmar recepción; el dueño que conserva su propio vehículo no recibe una acción de confirmación de entrega consigo mismo.\par\vspace{0.45em}
\noindent\textbf{AC-RF-AUC-028 --- Inicio de disputa de entrega por adjudicatario distinto del dueño previo}\par
\textbf{Método:} Prueba de disputa en ambos tipos de resultado.\par
\textbf{Oráculo:} La falta de entrega a un adjudicatario externo habilita disputa dentro del plazo; una autoadjudicación no habilita disputa de entrega del dueño contra sí mismo.\par\vspace{0.45em}
\noindent\textbf{AC-RF-AUC-029 --- Resolución administrativa de la disputa de entrega}\par
\textbf{Método:} Prueba funcional de caja negra.\par
\textbf{Oráculo:} Abrir una disputa y registrar cada decisión posible como \texttt{admin}; la resolución queda conservada y asociada a la disputa.\par\vspace{0.45em}
\noindent\textbf{AC-RF-AUC-030 --- Exclusión del promotor declarado moroso por disputa de entrega}\par
\textbf{Método:} Prueba funcional de caja negra.\par
\textbf{Oráculo:} Resolver una disputa manteniendo la morosidad y comprobar el estado de exclusión del promotor.\par\vspace{0.45em}
\noindent\textbf{AC-RF-AUC-031 --- Conservación de eventos previos del promotor excluido en campeonatos}\par
\textbf{Método:} Prueba funcional de caja negra.\par
\textbf{Oráculo:} Excluir a un promotor con eventos computables previos y comprobar que esos eventos conservan su participación en los cálculos de campeonato aplicables.\par\vspace{0.45em}
\noindent\textbf{AC-RF-AUC-032 --- Participación del dueño en la subasta de su propio vehículo}\par
\textbf{Método:} Prueba de elegibilidad y registro de puja.\par
\textbf{Oráculo:} Un dueño con verificación reforzada activa puede presentar una puja sobre su propio vehículo y esta queda vigente bajo las mismas reglas de ordenamiento que las demás.\par\vspace{0.45em}
\noindent\textbf{AC-RF-AUC-033 --- Exención de retención del principal para la puja propia del dueño}\par
\textbf{Método:} Prueba económica de retenciones.\par
\textbf{Oráculo:} Al registrar o aumentar la puja del dueño sobre su propio vehículo, el saldo no queda retenido por el principal de esa puja; una puja del mismo usuario sobre un vehículo ajeno sí aplica la retención normal.\par\vspace{0.45em}
\noindent\textbf{AC-RF-AUC-034 --- Autoadjudicación del vehículo al dueño previo}\par
\textbf{Método:} Prueba funcional y económica de cierre.\par
\textbf{Oráculo:} Cuando la puja del dueño es la mayor, el resultado se registra como autoadjudicación, no existe cobro/pago circular del principal ni flujo de entrega o disputa, y el vehículo permanece con el dueño.\par\vspace{0.45em}
\noindent\textbf{AC-RF-AUC-035 --- Advertencia de uso, estudio e inspección técnica posterior a la adjudicación}\par
\textbf{Método:} Prueba de interfaz y aceptación de términos.\par
\textbf{Oráculo:} El flujo de inscripción de vehículo y el flujo de puja de vehículo muestran una advertencia explícita sobre adquisición, inspección y estudio técnico posterior dentro de los límites legales y registran aceptación antes de continuar.\par\vspace{0.45em}
\noindent\textbf{AC-RF-CLS-001 --- Determinación de resultados del dominio del evento}\par
\textbf{Método:} Prueba funcional de caja negra.\par
\textbf{Oráculo:} Cuando el evento cuenta con los datos oficiales necesarios, el sistema genera los resultados del dominio aplicables al evento conforme a las reglas publicadas.\par\vspace{0.45em}
\noindent\textbf{AC-RF-CLS-002 --- Gestión de clasificaciones por modalidad}\par
\textbf{Método:} Prueba funcional de caja negra.\par
\textbf{Oráculo:} Para un evento con resultados oficiales, el sistema actualiza las clasificaciones aplicables por modalidad conforme a las reglas publicadas.\par\vspace{0.45em}
\noindent\textbf{AC-RF-CLS-003 --- Actualización de clasificaciones por alcance territorial y temporal}\par
\textbf{Método:} Prueba funcional de caja negra.\par
\textbf{Oráculo:} Al consolidarse los resultados del evento, el sistema actualiza las clasificaciones territoriales y temporales que correspondan a ese evento.\par\vspace{0.45em}
\noindent\textbf{AC-RF-CLS-004 --- Publicación de la matriz de pesos del evento desde inscripciones}\par
\textbf{Método:} Prueba funcional de caja negra.\par
\textbf{Oráculo:} A partir de la transición del evento a \texttt{inscripciones}, el sistema muestra la matriz de pesos correspondiente al evento junto con sus niveles territoriales y temporales.\par\vspace{0.45em}
\noindent\textbf{AC-RF-CLS-005 --- Publicación conjunta de resultados y clasificaciones del evento}\par
\textbf{Método:} Prueba funcional de caja negra.\par
\textbf{Oráculo:} Una vez publicados los resultados oficiales del evento, el sistema permite consultar en una misma publicación los resultados y las clasificaciones correspondientes, manteniendo visible la matriz de pesos del evento ya publicada desde \texttt{inscripciones}.\par\vspace{0.45em}
\noindent\textbf{AC-RF-CLS-006 --- Exclusión de inscripciones canceladas de los resultados del evento}\par
\textbf{Método:} Prueba funcional de caja negra.\par
\textbf{Oráculo:} Si una inscripción fue cancelada antes de la determinación de resultados, esa inscripción no aparece en los resultados oficiales, no altera las clasificaciones del evento y no influye en la publicación oficial del evento.\par\vspace{0.45em}
\noindent\textbf{AC-RF-CLS-007 --- Publicación del resultado final vigente de la subasta de vehículo}\par
\textbf{Método:} Prueba funcional de caja negra.\par
\textbf{Oráculo:} Si una subasta de vehículo sufrió descalificación del dueño, disputa de entrega o declaración de morosidad, los resultados oficiales del evento muestran solo el ganador y el precio finales vigentes de esa subasta.\par\vspace{0.45em}
\noindent\textbf{AC-RF-SAF-001 --- Lista estandarizada de comprobación de seguridad y cumplimiento}\par
\textbf{Método:} Prueba funcional de caja negra.\par
\textbf{Oráculo:} En un evento con comprobación de seguridad y cumplimiento, el sistema pone a disposición una lista estandarizada que incluye, como mínimo, jaula antivuelco, arneses, ausencia de humo visible y prohibición de combustibles con llama invisible.\par\vspace{0.45em}
\noindent\textbf{AC-RF-SAF-002 --- Soporte a la inspección comunitaria previa a carrera}\par
\textbf{Método:} Prueba funcional de caja negra.\par
\textbf{Oráculo:} Mientras el evento permanezca en \texttt{inscripciones} o \texttt{velocidad}, el sistema permite utilizar la lista de comprobación de seguridad y cumplimiento para la inspección comunitaria de las condiciones visibles de seguridad de los vehículos participantes, sin registrar en esa vía el pesaje externo.\par\vspace{0.45em}
\noindent\textbf{AC-RF-SAF-003 --- Registro de objeciones específicas de seguridad}\par
\textbf{Método:} Prueba funcional de caja negra.\par
\textbf{Oráculo:} Mientras el evento permanezca antes de \texttt{carrera}, un competidor puede registrar una objeción asociada a un vehículo e identificar el componente o la condición observada.\par\vspace{0.45em}
\noindent\textbf{AC-RF-SAF-004 --- Registro del presentante de la protesta de seguridad}\par
\textbf{Método:} Prueba funcional de caja negra.\par
\textbf{Oráculo:} Para cada objeción de seguridad registrada, el sistema conserva la identidad del usuario que la presentó como parte del registro de esa protesta.\par\vspace{0.45em}
\noindent\textbf{AC-RF-SAF-005 --- Registro del resultado de la comprobación de seguridad}\par
\textbf{Método:} Prueba funcional de caja negra.\par
\textbf{Oráculo:} Cuando se registra el resultado de una comprobación de seguridad, el sistema lo asocia al vehículo correspondiente y distingue si el vehículo queda aprobado o excluido del evento.\par\vspace{0.45em}
\noindent\textbf{AC-RF-SAF-006 --- Votación formal ante objeción no resuelta}\par
\textbf{Método:} Prueba funcional de caja negra.\par
\textbf{Oráculo:} Ante una objeción no resuelta, un competidor puede iniciar una votación formal; el sistema registra el total de votos emitidos y el voto de cada participante identificado por su nombre/usuario exacto, y conserva ese resultado como información para la decisión del promotor, sin cancelar automáticamente la inscripción afectada.\par\vspace{0.45em}
\noindent\textbf{AC-RF-SAF-007 --- Decisión final del promotor en disputas persistentes}\par
\textbf{Método:} Prueba funcional de caja negra.\par
\textbf{Oráculo:} Cuando una disputa de seguridad persiste, el sistema permite asociar al evento y al vehículo la decisión final registrada por el promotor y, si corresponde, cancelar la inscripción afectada conforme a las reglas aplicables.\par\vspace{0.45em}
\noindent\textbf{AC-RF-SAF-008 --- Registro opcional de datos de cumplimiento}\par
\textbf{Método:} Prueba funcional de caja negra.\par
\textbf{Oráculo:} En un evento con comprobación de seguridad y cumplimiento, el sistema permite registrar datos de seguridad e inspección, pero no exige completar ese registro como condición técnica obligatoria para continuar.\par\vspace{0.45em}
\noindent\textbf{AC-RF-SAF-009 --- Carácter informativo de la protesta de seguridad}\par
\textbf{Método:} Prueba funcional de caja negra.\par
\textbf{Oráculo:} Una objeción de seguridad y su votación formal pueden quedar registradas y consultables antes de \texttt{carrera}, con publicación del total de votos emitidos y del voto de cada participante identificado por su nombre/usuario exacto, pero la inscripción afectada solo se cancela si el promotor registra esa decisión.\par\vspace{0.45em}
\noindent\textbf{AC-RF-COM-001 --- Publicación en el foro comunitario}\par
\textbf{Método:} Prueba funcional de caja negra.\par
\textbf{Oráculo:} Cualquier usuario con identidad verificada puede registrar en el foro comunitario una discusión, una pregunta o una retroalimentación del ecosistema.\par\vspace{0.45em}
\noindent\textbf{AC-RF-COM-002 --- Consulta del foro comunitario}\par
\textbf{Método:} Prueba funcional de caja negra.\par
\textbf{Oráculo:} Un usuario puede consultar las discusiones, preguntas y retroalimentación publicadas en el foro comunitario del ecosistema.\par\vspace{0.45em}
\noindent\textbf{AC-RF-COM-003 --- Publicación en el canal de novedades del ecosistema}\par
\textbf{Método:} Prueba funcional de caja negra.\par
\textbf{Oráculo:} Cualquier usuario con identidad verificada puede registrar en el canal de novedades contenido correspondiente al ecosistema.\par\vspace{0.45em}
\noindent\textbf{AC-RF-COM-004 --- Consulta del canal de novedades del ecosistema}\par
\textbf{Método:} Prueba funcional de caja negra.\par
\textbf{Oráculo:} El canal de novedades del ecosistema permite consultar al menos anuncios de eventos, resultados de carrera, resultados de subasta o destacados de la comunidad.\par\vspace{0.45em}
\noindent\textbf{AC-RF-COM-005 --- Moderación del foro comunitario y del canal de novedades}\par
\textbf{Método:} Prueba funcional de caja negra.\par
\textbf{Oráculo:} \texttt{admin} dispone de funciones de moderación aplicables al foro comunitario y al canal de novedades del ecosistema.\par\vspace{0.45em}
\noindent\textbf{AC-RF-PUB-001 --- Consulta pública de la información públicada del evento}\par
\textbf{Método:} Prueba funcional de caja negra.\par
\textbf{Oráculo:} Un usuario puede consultar la información públicada del evento, incluidos resultados publicados, clasificaciones publicadas, matriz de pesos publicada y demás publicaciones públicas vigentes del evento.\par\vspace{0.45em}
\noindent\textbf{AC-RF-PUB-002 --- Consulta histórica de resultados oficiales del evento}\par
\textbf{Método:} Prueba funcional de caja negra.\par
\textbf{Oráculo:} Después de la publicación oficial de un evento, sus resultados oficiales permanecen consultables dentro del historial del evento.\par\vspace{0.45em}
\noindent\textbf{AC-RF-PUB-003 --- Trazabilidad de resultados oficiales a reglas y datos de origen}\par
\textbf{Método:} Prueba funcional de caja negra.\par
\textbf{Oráculo:} Para un resultado oficial publicado, el sistema permite consultar las reglas publicadas aplicadas y los datos oficiales del evento que lo sustentan.\par\vspace{0.45em}
\noindent\textbf{AC-RF-PUB-004 --- Registro histórico de resultados públicos de procesos críticos del evento}\par
\textbf{Método:} Prueba funcional de caja negra.\par
\textbf{Oráculo:} Cuando el sistema publica el resultado de una subasta, de una votación formal o de otro proceso crítico del evento, ese resultado permanece consultable dentro del historial del evento.\par\vspace{0.45em}
\noindent\textbf{AC-RF-PUB-005 --- Visibilidad pública de inscripciones canceladas o descalificadas}\par
\textbf{Método:} Prueba funcional de caja negra.\par
\textbf{Oráculo:} Si una inscripción fue cancelada o descalificada, sigue apareciendo en la consulta del evento con una indicación visible de su condición y con la nota asociada.\par\vspace{0.45em}
\noindent\textbf{AC-RF-PUB-006 --- Visibilidad pública de protestas de seguridad y sus resultados}\par
\textbf{Método:} Prueba funcional de caja negra.\par
\textbf{Oráculo:} Si una protesta de seguridad fue registrada en el evento, la consulta pública del evento permite ver esa protesta y el resultado de su tramitación; si hubo votación formal, también permite ver el total de votos emitidos y el voto de cada participante identificado por su nombre o usuario exacto, exista o no cancelación de la inscripción afectada.\par\vspace{0.45em}
\noindent\textbf{AC-RF-PUB-007 --- Reserva pública de la identidad de quien presenta la protesta}\par
\textbf{Método:} Prueba funcional de caja negra.\par
\textbf{Oráculo:} Cuando una protesta de seguridad es visible en la consulta pública del evento, esa consulta no muestra la identidad de quien la presentó.\par\vspace{0.45em}
\noindent\textbf{AC-RF-PUB-008 --- Visibilidad pública de disputas de entrega de vehículos adjudicados}\par
\textbf{Método:} Prueba funcional de caja negra.\par
\textbf{Oráculo:} Si una disputa de entrega fue registrada respecto de un vehículo adjudicado, la consulta pública del evento permite ver esa disputa con las identidades involucradas, los reportes registrados y la decisión final de \texttt{admin}.\par\vspace{0.45em}
\noindent\textbf{AC-RF-PUB-009 --- Visibilidad pública del incumplimiento del dueño o del promotor en la entrega del vehículo}\par
\textbf{Método:} Prueba funcional de caja negra.\par
\textbf{Oráculo:} Si el dueño o el promotor son declarados morosos dentro del flujo de subasta y entrega del vehículo, la consulta pública del evento muestra esa declaración y sus efectos públicos aplicables.\par\vspace{0.45em}
\noindent\textbf{AC-RF-PUB-010 --- Historial y trazabilidad temporal de disputas de entrega de vehículos adjudicados}\par
\textbf{Método:} Prueba funcional de caja negra.\par
\textbf{Oráculo:} Si existió una disputa de entrega de vehículo adjudicado, el historial del evento permite consultar cronológicamente su inicio, los reportes registrados, la decisión de \texttt{admin} y los efectos derivados de ella.\par\vspace{0.45em}
\noindent\textbf{AC-RF-IAM-001 --- Identificación del usuario mediante cuenta individual}\par
\textbf{Método:} Prueba funcional de caja negra.\par
\textbf{Oráculo:} Cuando un usuario accede a funciones que requieren participación en el sistema, el acceso queda asociado a una cuenta individual del usuario.\par\vspace{0.45em}
\noindent\textbf{AC-RF-IAM-002 --- Verificación de identidad previa a la participación}\par
\textbf{Método:} Prueba funcional de caja negra.\par
\textbf{Oráculo:} Si un usuario no tiene identidad verificada, el sistema no le habilita operaciones de inscripción, puja, aportación, promoción, protesta ni otras funciones de participación operativa.\par\vspace{0.45em}
\noindent\textbf{AC-RF-IAM-003 --- Aceptación versionada de términos y condiciones previa a la participación}\par
\textbf{Método:} Prueba funcional de caja negra.\par
\textbf{Oráculo:} Si un usuario no tiene registrada la aceptación de la versión vigente de los términos y condiciones, el sistema no le habilita funciones de inscripción, puja, aportación, promoción, protesta ni otras funciones de participación operativa.\par\vspace{0.45em}
\noindent\textbf{AC-RF-IAM-004 --- Acceso a funciones según papeles funcionales y nivel visible de acceso del usuario}\par
\textbf{Método:} Prueba funcional de caja negra.\par
\textbf{Oráculo:} Para un usuario dado, el sistema habilita el acceso a las funciones asociadas a sus papeles funcionales y a su nivel visible de acceso disponible, y no presenta como disponibles funciones que requieran un papel o un nivel no cumplido; además, el nivel \texttt{extendido} conserva las capacidades de \texttt{habilitado}, y \texttt{habilitado} conserva las capacidades de \texttt{base}.\par\vspace{0.45em}
\noindent\textbf{AC-RF-IAM-005 --- Acceso restringido a la identidad de quien presenta una protesta de seguridad}\par
\textbf{Método:} Prueba funcional de caja negra.\par
\textbf{Oráculo:} Cuando existe una protesta de seguridad registrada, la identidad de quien la presentó es visible para \texttt{admin} y para el promotor del evento, y no es visible para otros usuarios.\par\vspace{0.45em}
\noindent\textbf{AC-RF-IAM-006 --- Restricción de acceso por declaración válida de morosidad}\par
\textbf{Método:} Prueba funcional de caja negra.\par
\textbf{Oráculo:} Cuando un usuario queda declarado moroso por una decisión válida dentro del flujo de subasta y entrega de vehículos, el sistema deja de permitirle el acceso a las funciones del sistema.\par\vspace{0.45em}
\noindent\textbf{AC-RF-IAM-007 --- Habilitación reforzada de operaciones de mayor riesgo}\par
\textbf{Método:} Prueba funcional de caja negra.\par
\textbf{Oráculo:} Si un usuario no se encuentra en el nivel \texttt{extendido}, el sistema no le habilita las funciones de inscripción de vehículos, puja en subastas de vehículos ni promoción de eventos.\par\vspace{0.45em}
\noindent\textbf{AC-RF-IAM-008 --- Recordatorio al intentar una operación bloqueada por nivel visible de acceso}\par
\textbf{Método:} Prueba funcional de caja negra.\par
\textbf{Oráculo:} Si un usuario intenta una operación que requiere un nivel visible de acceso superior, el sistema le muestra una explicación concreta de su nivel actual, del nivel o condición faltante y de la acción necesaria para habilitarla.\par\vspace{0.45em}
\noindent\textbf{AC-RF-IAM-009 --- Acumulación de papeles funcionales en un mismo usuario}\par
\textbf{Método:} Prueba funcional de caja negra.\par
\textbf{Oráculo:} Un mismo usuario puede operar en más de un papel funcional dentro del sistema cuando esos papeles no correspondan a \texttt{admin}.\par\vspace{0.45em}
\noindent\textbf{AC-RF-IAM-010 --- Unicidad y exclusividad de admin}\par
\textbf{Método:} Prueba funcional de caja negra.\par
\textbf{Oráculo:} En el sistema existe un solo \texttt{admin}, y ese papel no aparece acumulado con otros papeles funcionales en un mismo usuario.\par\vspace{0.45em}
\noindent\textbf{AC-RF-EVT-041 --- Registro de datos mínimos de convocatoria del evento}\par
\textbf{Método:} Prueba funcional de caja negra.\par
\textbf{Oráculo:} Al registrar la convocatoria de un evento, el sistema permite capturar esos datos mínimos, incluido el costo fijo por pesaje, y asociarlos al evento.\par\vspace{0.45em}
\noindent\textbf{AC-RF-EVT-042 --- Conservación histórica del número esperado de participantes}\par
\textbf{Método:} Prueba funcional de caja negra.\par
\textbf{Oráculo:} Si se producen cancelaciones de inscripciones después de publicada la convocatoria, el número esperado de participantes permanece visible como dato histórico de la convocatoria original.\par\vspace{0.45em}
\noindent\textbf{AC-RF-EVT-043 --- Verificación reforzada obligatoria para promover eventos}\par
\textbf{Método:} Prueba funcional de caja negra.\par
\textbf{Oráculo:} Si un usuario no tiene verificación reforzada activa, el sistema no le habilita el registro ni la publicación de eventos.\par\vspace{0.45em}
\noindent\textbf{AC-RF-EVT-044 --- Publicación de la convocatoria del evento en inscripciones}\par
\textbf{Método:} Prueba funcional de caja negra.\par
\textbf{Oráculo:} Al entrar un evento en \texttt{inscripciones}, su convocatoria, sus condiciones públicas y sus reglas públicas quedan disponibles para consulta.\par\vspace{0.45em}
\noindent\textbf{AC-RF-EVT-045 --- Establecimiento obligatorio de la carga mínima del evento}\par
\textbf{Método:} Prueba funcional de caja negra.\par
\textbf{Oráculo:} En un evento con actividad competitiva pendiente, el sistema no permite continuar sin una carga mínima establecida y la conserva como parte de las condiciones publicadas del evento.\par\vspace{0.45em}
\noindent\textbf{AC-RF-EVT-046 --- Determinación de la clasificación combinada aplicable a la salida invertida}\par
\textbf{Método:} Prueba funcional de caja negra.\par
\textbf{Oráculo:} En un evento con salida invertida, el sistema conserva identificada la clasificación combinada aplicable a partir del nivel territorial y del nivel temporal del evento como parte de las condiciones publicadas del evento.\par\vspace{0.45em}
\noindent\textbf{AC-RF-EVT-047 --- Establecimiento obligatorio de la cantidad disponible de accesos para espectadores}\par
\textbf{Método:} Prueba funcional de caja negra.\par
\textbf{Oráculo:} En un evento con subasta de acceso para espectadores, el sistema no permite continuar sin una cantidad disponible de accesos establecida, y la conserva como parte de la convocatoria del evento.\par\vspace{0.45em}
\noindent\textbf{AC-RF-EVT-048 --- Establecimiento obligatorio del plazo máximo para iniciar disputas de entrega de vehículos}\par
\textbf{Método:} Prueba funcional de caja negra.\par
\textbf{Oráculo:} En un evento con subasta de vehículos, el sistema no permite continuar sin un plazo máximo establecido para iniciar disputas de entrega y lo conserva como parte de las condiciones publicadas del evento.\par\vspace{0.45em}
\noindent\textbf{AC-RF-EVT-049 --- Establecimiento obligatorio del plazo máximo para resolver disputas de entrega de vehículos}\par
\textbf{Método:} Prueba funcional de caja negra.\par
\textbf{Oráculo:} En un evento con subasta de vehículos, el sistema no permite continuar sin un plazo máximo establecido para resolver disputas de entrega y lo conserva como parte de las condiciones publicadas del evento.\par\vspace{0.45em}
\noindent\textbf{AC-RF-EVT-050 --- Publicación de la condición de entrega previa del vehículo al promotor}\par
\textbf{Método:} Prueba funcional de caja negra.\par
\textbf{Oráculo:} En un evento con subasta de vehículos, las condiciones publicadas del evento indican expresamente que la subasta solo puede iniciar respecto de vehículos ya recibidos por el promotor.\par\vspace{0.45em}
\noindent\textbf{AC-RF-VEH-001 --- Registro previo obligatorio del vehículo en el sistema}\par
\textbf{Método:} Prueba funcional de caja negra.\par
\textbf{Oráculo:} Si un vehículo no cuenta con registro previo en el sistema, no puede ser utilizado en una inscripción de evento.\par\vspace{0.45em}
\noindent\textbf{AC-RF-VEH-002 --- Asociación del vehículo registrado con su usuario dueño}\par
\textbf{Método:} Prueba funcional de caja negra.\par
\textbf{Oráculo:} Para cada vehículo registrado existe un usuario dueño asociado y consultable dentro del sistema.\par\vspace{0.45em}
\noindent\textbf{AC-RF-VEH-003 --- Datos mínimos obligatorios del vehículo registrado}\par
\textbf{Método:} Prueba funcional de caja negra.\par
\textbf{Oráculo:} Un vehículo no puede quedar registrado en el sistema si falta cualquiera de esos datos mínimos.\par\vspace{0.45em}
\noindent\textbf{AC-RF-VEH-004 --- Inscripción del vehículo en el evento por su usuario dueño}\par
\textbf{Método:} Prueba funcional de caja negra.\par
\textbf{Oráculo:} Si la inscripción del vehículo la intenta realizar un usuario distinto del usuario dueño registrado, la inscripción se rechaza.\par\vspace{0.45em}
\noindent\textbf{AC-RF-VEH-005 --- Verificación reforzada obligatoria del usuario dueño para inscribir un vehículo}\par
\textbf{Método:} Prueba funcional de caja negra.\par
\textbf{Oráculo:} Si el usuario dueño del vehículo no tiene verificación reforzada activa, el sistema no le habilita la inscripción de ese vehículo en el evento.\par\vspace{0.45em}
\noindent\textbf{AC-RF-VEH-006 --- Identificación del usuario piloto en la inscripción cuando exista pilotaje humano}\par
\textbf{Método:} Prueba funcional de caja negra.\par
\textbf{Oráculo:} Una inscripción con pilotaje humano directo no puede quedar registrada si no incluye el \texttt{id del usuario piloto}; una inscripción sin pilotaje humano directo puede quedar registrada sin ese dato.\par\vspace{0.45em}
\noindent\textbf{AC-RF-VEH-007 --- Posibilidad de que el usuario piloto sea distinto del usuario dueño cuando exista pilotaje humano}\par
\textbf{Método:} Prueba funcional de caja negra.\par
\textbf{Oráculo:} Cuando exista pilotaje humano directo, el sistema acepta una inscripción en la que el \texttt{id del usuario piloto} sea distinto del usuario dueño del vehículo, siempre que se cumplan las demás condiciones de la inscripción.\par\vspace{0.45em}
\noindent\textbf{AC-RF-VEH-008 --- Datos mínimos obligatorios de la inscripción del competidor}\par
\textbf{Método:} Prueba funcional de caja negra.\par
\textbf{Oráculo:} Una inscripción no puede quedar registrada si falta cualquiera de esos datos mínimos exigibles para su modalidad operativa, incluido el \texttt{id del usuario piloto} cuando exista pilotaje humano directo.\par\vspace{0.45em}
\noindent\textbf{AC-RF-VEH-009 --- Aportación mínima obligatoria en la inscripción}\par
\textbf{Método:} Prueba funcional de caja negra.\par
\textbf{Oráculo:} Intentar inscripciones con aportaciones inferior, igual y superior a 1 unidad; la inferior se rechaza y las demás siguen las guardas restantes.\par\vspace{0.45em}
\noindent\textbf{AC-RF-VEH-010 --- Cobro inmediato de importes asociados a la inscripción}\par
\textbf{Método:} Prueba funcional de caja negra.\par
\textbf{Oráculo:} Aceptar una inscripción con importes conocidos y comprobar los cargos aplicables y su relación con la inscripción.\par\vspace{0.45em}
\noindent\textbf{AC-RF-VEH-011 --- Cobro inmediato de la inscripción del competidor}\par
\textbf{Método:} Prueba funcional de caja negra.\par
\textbf{Oráculo:} Cuando una inscripción queda registrada válidamente, el sistema ejecuta de inmediato el cobro correspondiente a la inscripción y a los importes asociados que deban cobrarse con ella.\par\vspace{0.45em}
\noindent\textbf{AC-RF-VEH-012 --- Nota obligatoria en cancelación o descalificación decidida por el promotor}\par
\textbf{Método:} Prueba funcional de caja negra.\par
\textbf{Oráculo:} Toda cancelación o descalificación de inscripción decidida por el promotor queda acompañada por una nota escrita con la razón correspondiente.\par\vspace{0.45em}
\noindent\textbf{AC-RF-VEH-013 --- Clasificación válida para pasar a carrera}\par
\textbf{Método:} Prueba funcional de caja negra.\par
\textbf{Oráculo:} Preparar inscripciones no canceladas ni descalificadas con y sin tiempo válido de calificación; al pasar a \texttt{carrera}, solo las primeras quedan \texttt{calificadas} y las segundas quedan \texttt{no calificadas}.\par\vspace{0.45em}
\noindent\textbf{AC-RF-VEH-014 --- Efectos de la cancelación o descalificación de la inscripción}\par
\textbf{Método:} Prueba funcional de caja negra.\par
\textbf{Oráculo:} Si una inscripción queda cancelada o descalificada, el sistema conserva su registro visible con la nota correspondiente; además, esa inscripción no aporta resultados, no participa en clasificaciones, no influye en cálculos económicos y su vehículo no aparece en la subasta del evento.\par\vspace{0.45em}
\noindent\textbf{AC-RF-VEH-015 --- Devolución inmediata por cancelación o descalificación previa a carrera}\par
\textbf{Método:} Prueba funcional de caja negra.\par
\textbf{Oráculo:} Cuando una inscripción queda cancelada o descalificada antes de \texttt{carrera}, el sistema ejecuta de inmediato la devolución de todos los importes cobrados asociados a esa inscripción y de las aportaciones de terceros cobradas a favor de ese participante en el evento.\par\vspace{0.45em}
\noindent\textbf{AC-RF-VEH-016 --- Descalificación y morosidad del dueño por no entrega o modificación del vehículo}\par
\textbf{Método:} Prueba funcional de caja negra.\par
\textbf{Oráculo:} Si el promotor registra la no entrega o la modificación indebida antes de recibir el vehículo, el sistema permite descalificar al competidor, declararlo moroso y aplicarle la exclusión correspondiente del evento, del sistema y de los campeonatos aplicables.\par\vspace{0.45em}
\noindent\textbf{AC-RF-EPR-001 --- Consulta pública del premio individual del evento}\par
\textbf{Método:} Prueba funcional con conjunto de datos de referencia y análisis de derivación.\par
\textbf{Oráculo:} Una vez determinada la liquidación económica oficial del evento, el sistema permite consultar públicamente el premio individual correspondiente a cada participante computable y no publica un premio dependiente de carrera para una inscripción \texttt{no calificada}.\par\vspace{0.45em}
\noindent\textbf{AC-RF-EPR-002 --- Consulta pública de la distribución posterior asociada al premio individual del evento}\par
\textbf{Método:} Prueba funcional con conjunto de datos de referencia y análisis de derivación.\par
\textbf{Oráculo:} Una vez determinada la liquidación económica oficial del evento, el sistema permite consultar la distribución económica posterior asociada al premio individual de cada participante computable.\par\vspace{0.45em}
\noindent\textbf{AC-RF-EPR-003 --- Correspondencia trazable entre premio individual del evento y distribución posterior}\par
\textbf{Método:} Prueba funcional con conjunto de datos de referencia y análisis de derivación.\par
\textbf{Oráculo:} Para un participante computable del evento, el sistema permite consultar el premio individual y la distribución económica posterior asociada como partes vinculadas del mismo registro económico oficial.\par\vspace{0.45em}
\noindent\textbf{AC-RF-EPR-004 --- Publicación anonimizada de las aportaciones y de la parte de premio de terceros en evento}\par
\textbf{Método:} Prueba funcional con conjunto de datos de referencia y análisis de derivación.\par
\textbf{Oráculo:} Si un competidor del evento tiene terceros asociados a sus contribuciones, la consulta correspondiente muestra las cantidades aportadas y la parte del premio que corresponde a esos terceros, pero no muestra sus nombres ni su identidad.\par\vspace{0.45em}
\noindent\textbf{AC-RF-CPR-001 --- Registro y cobro inmediato de aportación directa a la bolsa del campeonato}\par
\textbf{Método:} Prueba funcional con conjunto de datos de referencia y análisis de derivación.\par
\textbf{Oráculo:} Una aportación directa a la bolsa de premio de campeonato queda asociada al campeonato seleccionado, se cobra de inmediato y no puede aceptarse fuera de la ventana que termina antes del cierre de inscripciones del primer evento de ese campeonato.\par\vspace{0.45em}
\noindent\textbf{AC-RF-CPR-002 --- Registro y cobro inmediato de aportación a la contribución de un participante del campeonato}\par
\textbf{Método:} Prueba funcional con conjunto de datos de referencia y análisis de derivación.\par
\textbf{Oráculo:} Cuando la función de aportaciones a participantes de campeonato está activa, una aportación a participante de campeonato queda asociada al campeonato y al participante seleccionados, se cobra de inmediato y no puede aceptarse fuera de la ventana que termina antes del cierre de inscripciones del primer evento de ese campeonato; cuando la función está inactiva, el sistema rechaza su registro.\par\vspace{0.45em}
\noindent\textbf{AC-RF-CPR-003 --- Contribución mínima por defecto de participantes al premio de campeonato}\par
\textbf{Método:} Prueba funcional con conjunto de datos de referencia y análisis de derivación.\par
\textbf{Oráculo:} En un campeonato con premio activo, cada participante del campeonato aparece con una contribución mínima por defecto de 1 unidad monetaria, sin perjuicio de aportaciones adicionales que las reglas permitan.\par\vspace{0.45em}
\noindent\textbf{AC-RF-CPR-004 --- Determinación del pozo de campeonato}\par
\textbf{Método:} Prueba funcional con conjunto de datos de referencia y análisis de derivación.\par
\textbf{Oráculo:} Para un campeonato con aportaciones cobradas dentro de la ventana permitida, el sistema genera un único pozo de campeonato integrado por ambas clases de aportación.\par\vspace{0.45em}
\noindent\textbf{AC-RF-CPR-005 --- Registro de la finalidad económica de la aportación a participante en campeonato}\par
\textbf{Método:} Prueba funcional con conjunto de datos de referencia y análisis de derivación.\par
\textbf{Oráculo:} Para una aportación a participante de campeonato, el sistema conserva la relación entre el aportante, el participante beneficiario y la distribución posterior del premio que corresponda.\par\vspace{0.45em}
\noindent\textbf{AC-RF-CPR-006 --- Registro de patrocinio asociado a aportación directa de campeonato}\par
\textbf{Método:} Prueba funcional con conjunto de datos de referencia y análisis de derivación.\par
\textbf{Oráculo:} Una aportación directa a la bolsa de campeonato puede quedar vinculada a un anuncio con imagen y esa vinculación se conserva como parte del registro del patrocinio.\par\vspace{0.45em}
\noindent\textbf{AC-RF-CPR-007 --- Publicación del patrocinio de campeonato}\par
\textbf{Método:} Prueba funcional con conjunto de datos de referencia y análisis de derivación.\par
\textbf{Oráculo:} Consultar un campeonato con patrocinio registrado y comprobar que el patrocinio aplicable está publicado.\par\vspace{0.45em}
\noindent\textbf{AC-RF-CPR-008 --- Trazabilidad económica del patrocinio de campeonato}\par
\textbf{Método:} Prueba funcional con conjunto de datos de referencia y análisis de derivación.\par
\textbf{Oráculo:} Seleccionar una operación de patrocinio y seguir sus referencias desde la aportación y el anuncio hasta las aportaciones a participantes y la distribución económica asociada, cuando existan.\par\vspace{0.45em}
\noindent\textbf{AC-RF-CPR-009 --- Consulta en la página principal de campeonatos patrocinados y patrocinadores}\par
\textbf{Método:} Prueba funcional con conjunto de datos de referencia y análisis de derivación.\par
\textbf{Oráculo:} Desde la página principal, un usuario puede consultar los campeonatos con patrocinio publicado y visualizar los patrocinadores asociados a cada uno.\par\vspace{0.45em}
\noindent\textbf{AC-RF-WGT-001 --- Determinación de campeonatos y clasificaciones aplicables al evento}\par
\textbf{Método:} Prueba funcional con conjunto de datos de referencia y análisis de derivación.\par
\textbf{Oráculo:} Cuando un evento tenga definidos su nivel territorial y su nivel temporal, el sistema identifica de forma expresa los campeonatos y las clasificaciones a los que ese evento aporta resultados.\par\vspace{0.45em}
\noindent\textbf{AC-RF-WGT-002 --- Determinación del puntaje de prioridad propio del evento}\par
\textbf{Método:} Prueba funcional con conjunto de datos de referencia y análisis de derivación.\par
\textbf{Oráculo:} Para un evento con pozo económico oficial y costo por participante determinados, el sistema produce un único valor de P/C propio del evento y lo mantiene distinguible de la matriz de pesos vigente.\par\vspace{0.45em}
\noindent\textbf{AC-RF-WGT-003 --- Determinación inicial de la matriz de pesos vigente del evento}\par
\textbf{Método:} Prueba funcional con conjunto de datos de referencia y análisis de derivación.\par
\textbf{Oráculo:} Para un evento recién considerado en la matriz de pesos, el sistema parte de peso 100 en cada posición de la matriz antes de aplicar comparaciones con otros eventos.\par\vspace{0.45em}
\noindent\textbf{AC-RF-WGT-004 --- Aplicación del peso fijo mínimo en los niveles más bajos}\par
\textbf{Método:} Prueba funcional con conjunto de datos de referencia y análisis de derivación.\par
\textbf{Oráculo:} En los niveles más bajos aplicables, la matriz de pesos del evento muestra peso 100 sin requerir comparación con otros eventos.\par\vspace{0.45em}
\noindent\textbf{AC-RF-WGT-005 --- Comparación operativa entre eventos para ajustar pesos}\par
\textbf{Método:} Prueba funcional con conjunto de datos de referencia y análisis de derivación.\par
\textbf{Oráculo:} Si dos eventos comparten el solapamiento exigido por las reglas del dominio para un nivel superior, el sistema los incluye en la comparación operativa de ese nivel.\par\vspace{0.45em}
\noindent\textbf{AC-RF-WGT-006 --- Conservación del peso del evento con mayor P/C}\par
\textbf{Método:} Prueba funcional con conjunto de datos de referencia y análisis de derivación.\par
\textbf{Oráculo:} Comparar dos eventos ajustables con P/C distintos y comprobar que el peso del evento con mayor P/C no desciende por esa comparación.\par\vspace{0.45em}
\noindent\textbf{AC-RF-WGT-007 --- Descenso del peso del evento con menor P/C}\par
\textbf{Método:} Prueba funcional con conjunto de datos de referencia y análisis de derivación.\par
\textbf{Oráculo:} Comparar dos eventos ajustables con P/C distintos y comprobar que el evento con menor P/C adopta el siguiente valor definido de la secuencia base.\par\vspace{0.45em}
\noindent\textbf{AC-RF-WGT-008 --- Desempate de P/C por antigüedad de registro}\par
\textbf{Método:} Prueba funcional con conjunto de datos de referencia y análisis de derivación.\par
\textbf{Oráculo:} Comparar dos eventos con P/C iguales y marcas temporales de registro distintas; comprobar que el evento más antiguo recibe el tratamiento favorable previsto para la comparación.\par\vspace{0.45em}
\noindent\textbf{AC-RF-WGT-009 --- Iteración del ajuste hasta estabilizar la matriz de pesos}\par
\textbf{Método:} Prueba funcional con conjunto de datos de referencia y análisis de derivación.\par
\textbf{Oráculo:} Tras aplicar las comparaciones pertinentes, el sistema conserva una matriz de pesos estable para el evento, sin descensos pendientes por resolver.\par\vspace{0.45em}
\noindent\textbf{AC-RF-WGT-010 --- Escalamiento de puntos por peso del evento}\par
\textbf{Método:} Prueba funcional con conjunto de datos de referencia y análisis de derivación.\par
\textbf{Oráculo:} Para una base de puntos y un peso vigente determinados, el sistema produce la puntuación escalada correspondiente aplicando la razón peso sobre 100.\par\vspace{0.45em}
\noindent\textbf{AC-RF-WGT-011 --- Publicación trazable del puntaje de prioridad y de la matriz de pesos vigente}\par
\textbf{Método:} Prueba funcional con conjunto de datos de referencia y análisis de derivación.\par
\textbf{Oráculo:} Un usuario puede consultar para un evento el P/C propio, la matriz de pesos vigente, la lógica de ajuste aplicada cuando corresponda y los campeonatos y clasificaciones aplicables como partes vinculadas de la misma publicación.\par\vspace{0.45em}
\noindent\textbf{AC-RF-EFF-001 --- Determinación del rendimiento observado antes de la subasta}\par
\textbf{Método:} Prueba funcional con tiempos y longitud de vuelta conocidos.\par
\textbf{Oráculo:} Para cada inscripción computable, el sistema toma el menor tiempo válido de calificación, calcula $S_i=d_e/T_i^*$ y conserva ese valor sin modificación durante y después de la subasta.\par\vspace{0.45em}
\noindent\textbf{AC-RF-EFF-002 --- Consolidación de la muestra oficial precio--rendimiento}\par
\textbf{Método:} Prueba funcional con resultados oficiales conocidos.\par
\textbf{Oráculo:} Después del cierre existe exactamente un par $(S_i,C_i)$ por inscripción computable y no se incluyen inscripciones no computables, precios no positivos ni duplicados técnicos.\par\vspace{0.45em}
\noindent\textbf{AC-RF-EFF-003 --- Ausencia de resultados provisionales de eficiencia durante la subasta}\par
\textbf{Método:} Prueba de interacción durante subasta abierta.\par
\textbf{Oráculo:} Modificar varias pujas válidas y comprobar que no se muestra clasificación, puntos ni premio provisional de eficiencia.\par\vspace{0.45em}
\noindent\textbf{AC-RF-EFF-004 --- Reserva de información derivada de las pujas antes del cierre}\par
\textbf{Método:} Prueba de confidencialidad con usuarios distintos.\par
\textbf{Oráculo:} Antes del cierre no se muestran conteos, posiciones, precios provisionales, adjudicatarios provisionales ni indicadores derivados que permitan inferir el orden de pujas.\par\vspace{0.45em}
\noindent\textbf{AC-RF-EFF-005 --- Creación de la señal oficial de mercado al cierre}\par
\textbf{Método:} Prueba funcional de secuencia temporal.\par
\textbf{Oráculo:} La muestra oficial se forma después del cierre utilizando rendimientos fijados antes de la subasta y precios oficiales finales.\par\vspace{0.45em}
\noindent\textbf{AC-RF-EFF-006 --- Línea de tendencia log-lineal global de precios en función del rendimiento}\par
\textbf{Método:} Prueba numérica con conjunto de referencia.\par
\textbf{Oráculo:} Para $n\ge3$ y rango suficiente, los parámetros coinciden con OLS para $\ln(C_i)=a+bS_i+\varepsilon_i$ y los precios de tendencia coinciden con $\widehat C(S_i)=\exp(a+bS_i)$; los demás casos siguen CALC-EFF-002.\par\vspace{0.45em}
\noindent\textbf{AC-RF-EFF-007 --- Estimación única de la línea de tendencia con todos los participantes computables}\par
\textbf{Método:} Prueba de instrumentación y reproducibilidad.\par
\textbf{Oráculo:} El evento conserva un único par de parámetros de tendencia estimado con el conjunto completo de participantes computables; no se ajusta una regresión distinta por participante.\par\vspace{0.45em}
\noindent\textbf{AC-RF-EFF-008 --- Conservación trazable de la muestra y parámetros de tendencia}\par
\textbf{Método:} Prueba de auditabilidad.\par
\textbf{Oráculo:} A partir del resultado oficial puede recuperarse la muestra válida y los parámetros usados para reproducir los precios de tendencia publicados.\par\vspace{0.45em}
\noindent\textbf{AC-RF-EFF-009 --- Tratamiento normativo de muestras de tamaño cero, uno y dos y de rango insuficiente}\par
\textbf{Método:} Prueba numérica de casos límite.\par
\textbf{Oráculo:} $n=0$ no produce clasificación; $n=1$ produce primer lugar único con $D=0$; $n=2$ produce dos diferencias normativas cero; y rango insuficiente con $n\ge3$ usa la media geométrica de los precios oficiales válidos como modelo de solo intercepto.\par\vspace{0.45em}
\noindent\textbf{AC-RF-EFF-010 --- Determinación de la diferencia de eficiencia precio--tendencia}\par
\textbf{Método:} Prueba numérica con precios de tendencia y oficiales conocidos.\par
\textbf{Oráculo:} Para cada participante se obtiene exactamente $D_i=\widehat C(S_i)-C_i$; un precio oficial menor que la tendencia produce $D_i>0$ y uno mayor produce $D_i<0$.\par\vspace{0.45em}
\noindent\textbf{AC-RF-EFF-011 --- Ordenamiento descendente por diferencia de eficiencia}\par
\textbf{Método:} Prueba funcional con diferencias conocidas.\par
\textbf{Oráculo:} El orden publicado es descendente por $D_i$.\par\vspace{0.45em}
\noindent\textbf{AC-RF-EFF-012 --- Asignación de base de puntos de eficiencia}\par
\textbf{Método:} Prueba numérica sobre posiciones conocidas.\par
\textbf{Oráculo:} Las posiciones reciben la secuencia de CALC-PTS-001 sin alteración.\par\vspace{0.45em}
\noindent\textbf{AC-RF-EFF-013 --- Desempate de eficiencia por antigüedad de inscripción}\par
\textbf{Método:} Prueba funcional con diferencias iguales.\par
\textbf{Oráculo:} Entre diferencias normativamente iguales, la inscripción válida más antigua aparece primero.\par\vspace{0.45em}
\noindent\textbf{AC-RF-EFF-014 --- Publicación de la cadena de trazabilidad de eficiencia después del cierre}\par
\textbf{Método:} Prueba de navegación y reproducibilidad.\par
\textbf{Oráculo:} Después del cierre, para un participante se puede recorrer desde su mejor vuelta y rendimiento hasta precio oficial, precio de tendencia, diferencia, posición y base de puntos.\par\vspace{0.45em}
\noindent\textbf{AC-RF-EFF-015 --- Publicación exclusiva de resultados oficiales de eficiencia}\par
\textbf{Método:} Prueba antes y después del cierre.\par
\textbf{Oráculo:} No existen resultados provisionales visibles; después del cierre se publican resultados basados en la muestra oficial.\par\vspace{0.45em}
\noindent\textbf{AC-RF-EFF-016 --- Interpretación del signo de la diferencia de eficiencia}\par
\textbf{Método:} Prueba de presentación con $D>0$, $D=0$ y $D<0$.\par
\textbf{Oráculo:} La interfaz explica respectivamente valoración inferior a tendencia, coincidencia con tendencia y valoración superior a tendencia.\par\vspace{0.45em}
\noindent\textbf{AC-RF-EFF-017 --- Uso del precio oficial de autoadjudicación en la línea de tendencia}\par
\textbf{Método:} Prueba numérica con una autoadjudicación.\par
\textbf{Oráculo:} El precio oficial de segundo precio del vehículo autoadjudicado aparece como $C_i$ en la misma muestra usada para ajustar la línea global.\par\vspace{0.45em}
\noindent\textbf{AC-RF-EFF-018 --- Explicación posterior del incentivo de eficiencia}\par
\textbf{Método:} Prueba de comprensión del contenido posterior al cierre.\par
\textbf{Oráculo:} La explicación permite identificar que, a rendimiento fijo, un precio por encima de la tendencia reduce $D_i$ y que un vehículo competitivo por debajo de la tendencia obtiene una diferencia favorable y puede mejorar su premio.\par\vspace{0.45em}
\noindent\textbf{AC-RF-EPR-005 --- Determinación de las aportaciones directas computables a la bolsa del evento}\par
\textbf{Método:} Prueba funcional con conjunto de datos de referencia y análisis de derivación.\par
\textbf{Oráculo:} Concluida la ventana permitida del evento, el sistema identifica las aportaciones directas cobradas a la bolsa y las distingue de las contribuciones asociadas a participantes.\par\vspace{0.45em}
\noindent\textbf{AC-RF-EPR-006 --- Determinación de las contribuciones de eficiencia computables por participante en el evento}\par
\textbf{Método:} Prueba funcional con conjunto de datos de referencia y análisis de derivación.\par
\textbf{Oráculo:} Para cada participante computable, el sistema conserva una contribución total de eficiencia integrada por su mínimo obligatorio cobrado y por las aportaciones de eficiencia cobradas a su favor.\par\vspace{0.45em}
\noindent\textbf{AC-RF-EPR-007 --- Determinación de las contribuciones de terminación computables por participante en el evento}\par
\textbf{Método:} Prueba funcional con conjunto de datos de referencia y análisis de derivación.\par
\textbf{Oráculo:} Para cada inscripción computable, el sistema conserva una contribución total de terminación integrada por la cantidad registrada con la inscripción y por las aportaciones de terminación cobradas a su favor.\par\vspace{0.45em}
\noindent\textbf{AC-RF-EPR-008 --- Asignación base de puntos para el premio del evento}\par
\textbf{Método:} Prueba funcional con conjunto de datos de referencia y análisis de derivación.\par
\textbf{Oráculo:} Una vez determinado el ordenamiento de eficiencia del evento, el sistema asigna la base de puntos siguiendo la secuencia base del dominio.\par\vspace{0.45em}
\noindent\textbf{AC-RF-EPR-009 --- Determinación del coeficiente de premio del evento por participante}\par
\textbf{Método:} Prueba funcional con conjunto de datos de referencia y análisis de derivación.\par
\textbf{Oráculo:} Determinadas la base de puntos por eficiencia y la contribución total de eficiencia computable de un participante, el sistema produce para ese participante un coeficiente de premio igual al producto de ambas magnitudes.\par\vspace{0.45em}
\noindent\textbf{AC-RF-EPR-010 --- Determinación de la suma de coeficientes de premio del evento}\par
\textbf{Método:} Prueba funcional con conjunto de datos de referencia y análisis de derivación.\par
\textbf{Oráculo:} Una vez calculados los coeficientes individuales del evento, el sistema produce una única suma de coeficientes para el evento.\par\vspace{0.45em}
\noindent\textbf{AC-RF-EPR-011 --- Determinación operativa del premio del evento por participante}\par
\textbf{Método:} Prueba funcional con conjunto de datos de referencia y análisis de derivación.\par
\textbf{Oráculo:} Con el pozo económico oficial, el coeficiente individual y la suma de coeficientes ya determinados, el sistema calcula las cuotas ideales y asigna importes monetarios enteros mediante \texttt{CALC-MON-001}; la suma de premios coincide exactamente con el pozo. Para $K=0$, no se adjudican premios y se restituyen las cantidades al origen.\par\vspace{0.45em}
\noindent\textbf{AC-RF-EPR-012 --- Elegibilidad para el premio de terminación}\par
\textbf{Método:} Prueba funcional con conjunto de datos de referencia y análisis de derivación.\par
\textbf{Oráculo:} Preparar resultados con participantes terminados y no terminados; únicamente los terminados aparecen en el conjunto de beneficiarios computables.\par\vspace{0.45em}
\noindent\textbf{AC-RF-EPR-013 --- Determinación del pozo computable de terminación}\par
\textbf{Método:} Prueba funcional con conjunto de datos de referencia y análisis de derivación.\par
\textbf{Oráculo:} Con importes conocidos, comparar el pozo computable con la suma de sus componentes.\par\vspace{0.45em}
\noindent\textbf{AC-RF-EPR-014 --- Determinación de vueltas esperadas}\par
\textbf{Método:} Prueba funcional con conjunto de datos de referencia y análisis de derivación.\par
\textbf{Oráculo:} En una carrera por tiempo, comprobar $E=T/L$ con resultado fraccionario; en una carrera por distancia, comprobar $E=D/d$ para todas las inscripciones computables. Una inscripción sin tiempo válido de calificación no entra en el cálculo.\par\vspace{0.45em}
\noindent\textbf{AC-RF-EPR-015 --- Determinación del peso de terminación}\par
\textbf{Método:} Prueba funcional con conjunto de datos de referencia y análisis de derivación.\par
\textbf{Oráculo:} Con contribución, vueltas completadas y vueltas esperadas conocidas, comparar el peso de terminación con el cálculo de referencia.\par\vspace{0.45em}
\noindent\textbf{AC-RF-EPR-016 --- Asignación proporcional del premio de terminación}\par
\textbf{Método:} Prueba funcional con conjunto de datos de referencia y análisis de derivación.\par
\textbf{Oráculo:} Con un conjunto de pesos positivos conocido, comprobar que cada premio individual es proporcional a su peso y que la suma de participaciones coincide con el pozo bajo la política monetaria vigente.\par\vspace{0.45em}
\noindent\textbf{AC-RF-EPR-017 --- Objetivo de incentivo del premio de terminación}\par
\textbf{Método:} Prueba funcional con conjunto de datos de referencia y análisis de derivación.\par
\textbf{Oráculo:} Inspeccionar las definiciones \texttt{CALC-EPR-005} y \texttt{CALC-EPR-006} y verificar que ambas variables intervienen en el peso y, por tanto, en la proporción asignada.\par\vspace{0.45em}
\noindent\textbf{AC-RF-EPR-018 --- Determinación operativa de la distribución económica posterior del premio del evento}\par
\textbf{Método:} Prueba funcional con conjunto de datos de referencia y análisis de derivación.\par
\textbf{Oráculo:} Si un participante computable tiene personas asociadas a sus contribuciones de eficiencia, el sistema distribuye el premio del evento entre ellas y el propio participante según la proporción aplicable y \texttt{CALC-MON-001}; la suma coincide exactamente con el premio individual. Si $Q=0$, el participante recibe el 100\% del premio.\par\vspace{0.45em}
\noindent\textbf{AC-RF-EPR-019 --- Determinación operativa de la distribución económica posterior del premio de terminación}\par
\textbf{Método:} Prueba funcional con conjunto de datos de referencia y análisis de derivación.\par
\textbf{Oráculo:} Si una inscripción computable tiene personas asociadas a sus contribuciones de terminación, el sistema distribuye el premio de terminación entre ellas y el propio participante según la proporción aplicable y \texttt{CALC-MON-001}; la suma coincide exactamente con el premio individual. Si $Q=0$, el participante recibe el 100\% del premio.\par\vspace{0.45em}
\noindent\textbf{AC-RF-EPR-020 --- Publicación trazable del patrocinio y de los premios del evento}\par
\textbf{Método:} Prueba funcional con conjunto de datos de referencia y análisis de derivación.\par
\textbf{Oráculo:} Un usuario puede consultar la relación entre pozo económico oficial, aportaciones directas cobradas a la bolsa, patrocinio publicado del evento, contribuciones de eficiencia cobradas a participantes, base de puntos por eficiencia, coeficientes, premio individual del evento, pozo del premio de terminación, contribuciones de terminación, premio de terminación y distribuciones posteriores para cada participante computable.\par\vspace{0.45em}
\noindent\textbf{AC-RF-EPR-021 --- Pago automático e inmediato de los premios del evento a todos los beneficiarios}\par
\textbf{Método:} Prueba funcional con conjunto de datos de referencia y análisis de derivación.\par
\textbf{Oráculo:} Cuando las distribuciones económicas posteriores válidas del evento quedan determinadas, el sistema ejecuta de inmediato el pago correspondiente a cada beneficiario.\par\vspace{0.45em}
\noindent\textbf{AC-RF-EPR-022 --- Reserva del pozo consolidado y de las aportaciones antes del cierre de la subasta de vehículos}\par
\textbf{Método:} Prueba de publicación antes y después del cierre con pozo inicial y aportaciones conocidas.\par
\textbf{Oráculo:} Durante la fase de subasta se consulta públicamente el evento y solo es visible el pozo inicial publicado como referencia del premio de eficiencia; no se muestran el pozo oficial consolidado ni importes de aportaciones o contribuciones. Después del cierre, la publicación oficial aplicable puede mostrar esos datos conforme a las reglas vigentes.\par\vspace{0.45em}
\noindent\textbf{AC-RF-CPR-010 --- Determinación de las aportaciones directas computables a la bolsa del campeonato}\par
\textbf{Método:} Prueba funcional con conjunto de datos de referencia y análisis de derivación.\par
\textbf{Oráculo:} Para un campeonato con aportaciones directas aceptadas, el sistema identifica el conjunto de aportaciones directas computables cobradas que integran su bolsa patrocinada.\par\vspace{0.45em}
\noindent\textbf{AC-RF-CPR-011 --- Determinación de las contribuciones computables por participante en el campeonato}\par
\textbf{Método:} Prueba funcional con conjunto de datos de referencia y análisis de derivación.\par
\textbf{Oráculo:} Para un campeonato con premio activo, el sistema conserva por cada participante una contribución total integrada por su mínimo por defecto y por las aportaciones cobradas a su favor.\par\vspace{0.45em}
\noindent\textbf{AC-RF-CPR-012 --- Resultado computable de campeonato de eficiencia}\par
\textbf{Método:} Prueba funcional con conjunto de datos de referencia y análisis de derivación.\par
\textbf{Oráculo:} Con una clasificación oficial y secuencia base conocidas, comparar el resultado computable de cada participante con la asignación de referencia.\par\vspace{0.45em}
\noindent\textbf{AC-RF-CPR-013 --- Resultado computable de campeonato de terminación}\par
\textbf{Método:} Prueba funcional con conjunto de datos de referencia y análisis de derivación.\par
\textbf{Oráculo:} Con relaciones por evento conocidas, comprobar que el resultado de campeonato es su suma para cada participante.\par\vspace{0.45em}
\noindent\textbf{AC-RF-CPR-014 --- Desempate del resultado computable de campeonato}\par
\textbf{Método:} Prueba funcional con conjunto de datos de referencia y análisis de derivación.\par
\textbf{Oráculo:} Crear un empate y marcas temporales distintas en el primer evento de participación; comprobar que se favorece la inscripción más antigua aplicable.\par\vspace{0.45em}
\noindent\textbf{AC-RF-CPR-015 --- Determinación del coeficiente de premio de campeonato por participante}\par
\textbf{Método:} Prueba funcional con conjunto de datos de referencia y análisis de derivación.\par
\textbf{Oráculo:} Determinados el resultado computable aplicable y la contribución total computable de un participante, el sistema produce para ese participante un coeficiente de premio de campeonato igual al producto de ambas magnitudes.\par\vspace{0.45em}
\noindent\textbf{AC-RF-CPR-016 --- Determinación de la suma de coeficientes de premio de campeonato}\par
\textbf{Método:} Prueba funcional con conjunto de datos de referencia y análisis de derivación.\par
\textbf{Oráculo:} Una vez calculados los coeficientes individuales del campeonato, el sistema produce una única suma de coeficientes para ese campeonato.\par\vspace{0.45em}
\noindent\textbf{AC-RF-CPR-017 --- Determinación operativa del premio de campeonato por participante}\par
\textbf{Método:} Prueba funcional con conjunto de datos de referencia y análisis de derivación.\par
\textbf{Oráculo:} Con el pozo de campeonato, el coeficiente individual y la suma de coeficientes ya determinados, el sistema calcula las cuotas ideales y asigna importes monetarios enteros mediante \texttt{CALC-MON-001}; la suma coincide exactamente con el pozo de campeonato. Para $K=0$, no se adjudican premios y se restituyen las cantidades al origen.\par\vspace{0.45em}
\noindent\textbf{AC-RF-CPR-018 --- Determinación operativa de la distribución económica posterior del premio de campeonato}\par
\textbf{Método:} Prueba funcional con conjunto de datos de referencia y análisis de derivación.\par
\textbf{Oráculo:} Si un participante del campeonato tiene personas asociadas a sus contribuciones, el sistema distribuye el premio del participante entre ellas y el propio participante mediante \texttt{CALC-DST-001} y \texttt{CALC-MON-001}; la suma coincide exactamente con el premio individual.\par\vspace{0.45em}
\noindent\textbf{AC-RF-CPR-019 --- Publicación y trazabilidad del patrocinio y del premio de campeonato}\par
\textbf{Método:} Prueba funcional con conjunto de datos de referencia y análisis de derivación.\par
\textbf{Oráculo:} Para un campeonato patrocinado, el sistema permite seguir la relación entre bolsa patrocinada, patrocinadores, anuncios publicados, contribuciones a participantes, resultado computable aplicable, coeficientes, premio individual del campeonato y distribución posterior correspondiente.\par\vspace{0.45em}
\noindent\textbf{AC-RF-CPR-020 --- Pago automático e inmediato del premio de campeonato a todos los beneficiarios}\par
\textbf{Método:} Prueba funcional con conjunto de datos de referencia y análisis de derivación.\par
\textbf{Oráculo:} Cuando la distribución económica posterior del premio de campeonato queda determinada, el sistema ejecuta de inmediato el pago correspondiente a cada beneficiario.\par\vspace{0.45em}
\noindent\textbf{AC-RF-SPN-001 --- Registro y cobro inmediato de aportación directa a la bolsa del evento}\par
\textbf{Método:} Prueba funcional de caja negra.\par
\textbf{Oráculo:} Una aportación directa a la bolsa del evento queda asociada al evento seleccionado, se cobra de inmediato y no puede aceptarse fuera de la ventana que termina al pasar el evento de \texttt{inscripciones} a \texttt{velocidad}.\par\vspace{0.45em}
\noindent\textbf{AC-RF-SPN-002 --- Registro y cobro inmediato de aportación a la contribución de eficiencia de un participante del evento}\par
\textbf{Método:} Prueba funcional de caja negra.\par
\textbf{Oráculo:} Cuando la función de aportaciones a participantes del evento está activa, una aportación a la contribución de eficiencia de un participante del evento queda asociada al evento y al participante seleccionados, se cobra de inmediato y no puede aceptarse fuera de la ventana que termina al pasar el evento de \texttt{inscripciones} a \texttt{velocidad}; cuando la función está inactiva, el sistema rechaza su registro.\par\vspace{0.45em}
\noindent\textbf{AC-RF-SPN-003 --- Registro y cobro inmediato de aportación a la contribución de terminación de un participante del evento}\par
\textbf{Método:} Prueba funcional de caja negra.\par
\textbf{Oráculo:} Cuando la función de aportaciones a participantes del evento está activa, una aportación a la contribución de terminación de una inscripción del evento queda asociada al evento y a la inscripción seleccionados, se cobra de inmediato y no puede aceptarse fuera de la ventana que termina al pasar el evento de \texttt{inscripciones} a \texttt{velocidad}; cuando la función está inactiva, el sistema rechaza su registro.\par\vspace{0.45em}
\noindent\textbf{AC-RF-SPN-004 --- Registro de patrocinio asociado a aportación directa de evento}\par
\textbf{Método:} Prueba funcional de caja negra.\par
\textbf{Oráculo:} Una aportación directa a la bolsa del evento puede quedar vinculada a un anuncio con imagen y esa vinculación se conserva como parte del registro del patrocinio.\par\vspace{0.45em}
\noindent\textbf{AC-RF-SPN-005 --- Publicación del patrocinio directo del evento}\par
\textbf{Método:} Prueba funcional de caja negra.\par
\textbf{Oráculo:} Consultar un evento con patrocinio directo y comprobar que el patrocinio está publicado.\par\vspace{0.45em}
\noindent\textbf{AC-RF-SPN-006 --- Publicación distinguida del patrocinio de campeonato aplicable}\par
\textbf{Método:} Prueba funcional de caja negra.\par
\textbf{Oráculo:} Consultar un evento con patrocinio directo y de campeonato y comprobar que la procedencia de cada uno se distingue.\par\vspace{0.45em}
\noindent\textbf{AC-RF-SPN-007 --- Trazabilidad económica del patrocinio y aportaciones del evento}\par
\textbf{Método:} Prueba funcional de caja negra.\par
\textbf{Oráculo:} Seleccionar operaciones económicas representativas y seguir sus referencias desde el cobro y anuncio hasta la distribución asociada, cuando exista.\par\vspace{0.45em}
\noindent\textbf{AC-RF-SPN-008 --- Devolución inmediata de aportaciones del evento por cancelación del evento}\par
\textbf{Método:} Prueba funcional de caja negra.\par
\textbf{Oráculo:} Cuando el evento pasa a \texttt{cancelado}, el sistema devuelve de inmediato las aportaciones directas a la bolsa y las aportaciones de eficiencia y de terminación cobradas dentro de ese evento.\par\vspace{0.45em}
\noindent\textbf{AC-RF-SPN-009 --- Devolución inmediata de aportaciones a participante por cancelación o descalificación previa a carrera}\par
\textbf{Método:} Prueba funcional de caja negra.\par
\textbf{Oráculo:} Cuando un participante queda fuera antes de \texttt{carrera}, el sistema devuelve de inmediato las aportaciones de terceros cobradas a su favor dentro del evento para eficiencia y para terminación.\par\vspace{0.45em}
\noindent\textbf{AC-RF-PTS-001 --- Base de puntos de velocidad}\par
\textbf{Método:} Prueba funcional con conjunto de datos de referencia y análisis de derivación.\par
\textbf{Oráculo:} Con un orden oficial conocido, comprobar la asignación posición--base de puntos.\par\vspace{0.45em}
\noindent\textbf{AC-RF-PTS-002 --- Desempate del orden de velocidad para puntos}\par
\textbf{Método:} Prueba funcional con conjunto de datos de referencia y análisis de derivación.\par
\textbf{Oráculo:} Crear un empate de velocidad y comprobar que la inscripción más antigua recibe la posición favorable.\par\vspace{0.45em}
\noindent\textbf{AC-RF-PTS-003 --- Base de puntos de carrera}\par
\textbf{Método:} Prueba funcional con conjunto de datos de referencia y análisis de derivación.\par
\textbf{Oráculo:} Con un orden oficial conocido, comprobar la asignación posición--base de puntos.\par\vspace{0.45em}
\noindent\textbf{AC-RF-PTS-004 --- Desempate del orden de carrera para puntos}\par
\textbf{Método:} Prueba funcional con conjunto de datos de referencia y análisis de derivación.\par
\textbf{Oráculo:} Crear un empate de carrera y comprobar que la inscripción más antigua recibe la posición favorable.\par\vspace{0.45em}
\noindent\textbf{AC-RF-PTS-005 --- Base de puntos de eficiencia para campeonato}\par
\textbf{Método:} Prueba funcional con conjunto de datos de referencia y análisis de derivación.\par
\textbf{Oráculo:} Con un orden oficial conocido, comprobar la asignación posición--base de puntos.\par\vspace{0.45em}
\noindent\textbf{AC-RF-PTS-006 --- Desempate del orden de eficiencia para campeonato}\par
\textbf{Método:} Prueba funcional con conjunto de datos de referencia y análisis de derivación.\par
\textbf{Oráculo:} Crear un empate de eficiencia y comprobar que la inscripción más antigua recibe la posición favorable.\par\vspace{0.45em}
\noindent\textbf{AC-RF-PTS-007 --- Escalamiento de puntos por modalidad en cada campeonato aplicable}\par
\textbf{Método:} Prueba funcional con conjunto de datos de referencia y análisis de derivación.\par
\textbf{Oráculo:} Determinadas la base de puntos por modalidad y el peso vigente del evento, el sistema produce la puntuación escalada del evento para cada campeonato aplicable sin aplicar redondeo de dominio intermedio no prescrito.\par\vspace{0.45em}
\noindent\textbf{AC-RF-PTS-008 --- Actualización de clasificaciones de velocidad, carrera y eficiencia}\par
\textbf{Método:} Prueba funcional con conjunto de datos de referencia y análisis de derivación.\par
\textbf{Oráculo:} Una vez determinada la puntuación escalada del evento por modalidad, el sistema la incorpora a las clasificaciones correspondientes de cada campeonato aplicable.\par\vspace{0.45em}
\noindent\textbf{AC-RF-PTS-009 --- Actualización de la clasificación combinada del campeonato}\par
\textbf{Método:} Prueba funcional con conjunto de datos de referencia y análisis de derivación.\par
\textbf{Oráculo:} Determinadas las puntuaciones escaladas de velocidad y carrera del evento, el sistema actualiza la clasificación combinada del campeonato sumando ambas aportaciones conforme a las reglas del dominio.\par\vspace{0.45em}
\noindent\textbf{AC-RF-PTS-010 --- Exclusión de inscripciones canceladas de la actualización de puntos y clasificaciones}\par
\textbf{Método:} Prueba funcional con conjunto de datos de referencia y análisis de derivación.\par
\textbf{Oráculo:} Si una inscripción fue cancelada, esa inscripción no recibe puntos, no altera las clasificaciones por modalidad y no influye en la actualización de ningún campeonato aplicable.\par\vspace{0.45em}
\noindent\textbf{AC-RF-PTS-011 --- Publicación trazable de la actualización de puntos y clasificaciones por modalidad}\par
\textbf{Método:} Prueba funcional con conjunto de datos de referencia y análisis de derivación.\par
\textbf{Oráculo:} Un usuario puede consultar para un evento y un campeonato aplicable la relación entre resultados oficiales, base de puntos por modalidad, peso vigente, puntuación escalada y actualización de las clasificaciones correspondientes.\par\vspace{0.45em}
\noindent\textbf{AC-RF-CPR-021 --- Consulta pública del premio individual del campeonato}\par
\textbf{Método:} Prueba funcional con conjunto de datos de referencia y análisis de derivación.\par
\textbf{Oráculo:} Una vez determinada la distribución económica oficial del premio de campeonato, el sistema permite consultar públicamente el premio individual correspondiente a cada competidor de ese campeonato.\par\vspace{0.45em}
\noindent\textbf{AC-RF-CPR-022 --- Consulta pública de la distribución posterior del premio de campeonato}\par
\textbf{Método:} Prueba funcional con conjunto de datos de referencia y análisis de derivación.\par
\textbf{Oráculo:} Una vez determinada la distribución económica oficial del premio de campeonato, el sistema permite consultar la distribución económica posterior asociada al premio individual de cada competidor.\par\vspace{0.45em}
\noindent\textbf{AC-RF-CPR-023 --- Publicación anonimizada de las aportaciones y de la parte de premio de terceros en campeonato}\par
\textbf{Método:} Prueba funcional con conjunto de datos de referencia y análisis de derivación.\par
\textbf{Oráculo:} Si un competidor de un campeonato tiene terceros asociados a sus contribuciones, la consulta correspondiente muestra las cantidades aportadas y la parte del premio que corresponde a esos terceros, pero no muestra sus nombres ni su identidad.\par\vspace{0.45em}
\noindent\textbf{AC-RF-CPR-024 --- Correspondencia trazable entre premio individual del campeonato y distribución posterior}\par
\textbf{Método:} Prueba funcional con conjunto de datos de referencia y análisis de derivación.\par
\textbf{Oráculo:} Para un competidor de campeonato con premio activo, el sistema permite consultar el premio individual y la distribución económica posterior asociada como partes vinculadas del mismo registro económico oficial.\par\vspace{0.45em}
\noindent\textbf{AC-RF-GOV-001 --- Publicación transversal de la información declarada pública}\par
\textbf{Método:} Prueba funcional de caja negra.\par
\textbf{Oráculo:} Cuando una regla del dominio declare pública una cancelación, un patrocinio, una protesta, una disputa, una decisión del promotor o un resultado, el sistema la muestra en la consulta pública aplicable con el alcance exigido por esa regla.\par\vspace{0.45em}
\noindent\textbf{AC-RF-GOV-002 --- Reserva selectiva de identidades sujetas a restricción pública}\par
\textbf{Método:} Prueba funcional de caja negra.\par
\textbf{Oráculo:} Cuando una regla del dominio reserve una identidad para la consulta pública, esa identidad no se muestra públicamente y solo queda visible para los papeles funcionales autorizados por la propia regla.\par\vspace{0.45em}
\noindent\textbf{AC-RF-GOV-003 --- Publicación anonimizada de terceros en distribuciones económicas}\par
\textbf{Método:} Prueba funcional de caja negra.\par
\textbf{Oráculo:} Cuando un premio individual de evento o de campeonato incluya distribución posterior hacia terceros, la consulta correspondiente muestra las cantidades y partes de premio de esos terceros sin exponer sus nombres ni su identidad.\par\vspace{0.45em}
\noindent\textbf{AC-RF-GOV-004 --- Moderación administrativa de espacios comunitarios}\par
\textbf{Método:} Prueba funcional de caja negra.\par
\textbf{Oráculo:} \texttt{admin} dispone de funciones de moderación aplicables al foro comunitario y al canal de novedades, y el ejercicio de esas funciones no cambia por sí mismo las reglas generales del dominio sobre qué información debe permanecer pública, reservada o trazable.\par\vspace{0.45em}
\noindent\textbf{AC-RF-GOV-005 --- Permanencia histórica de publicaciones y procesos críticos visibles}\par
\textbf{Método:} Prueba funcional de caja negra.\par
\textbf{Oráculo:} Cuando una publicación o un proceso crítico del dominio deba permanecer visible con trazabilidad histórica, el sistema lo mantiene consultable en el historial aplicable junto con los datos públicos que las reglas correspondientes exijan.\par\vspace{0.45em}
\noindent\textbf{AC-RF-IAM-011 --- Registro lógico del estado de identidad verificada del usuario}\par
\textbf{Método:} Prueba funcional de caja negra.\par
\textbf{Oráculo:} Para cada cuenta individual, el sistema conserva un valor lógico verificable que expresa si la identidad del usuario está marcada como verificada o no verificada.\par\vspace{0.45em}
\noindent\textbf{AC-RF-IAM-012 --- Registro lógico del estado de verificación reforzada del usuario}\par
\textbf{Método:} Prueba funcional de caja negra.\par
\textbf{Oráculo:} Para cada cuenta individual, el sistema conserva un valor lógico verificable que expresa si la verificación reforzada del usuario está activa o inactiva.\par\vspace{0.45em}
\noindent\textbf{AC-RF-IAM-013 --- Registro lógico del estado de exclusión del usuario}\par
\textbf{Método:} Prueba funcional de caja negra.\par
\textbf{Oráculo:} Para cada cuenta individual, el sistema conserva un valor lógico verificable que expresa si el usuario está marcado como excluido o no excluido.\par\vspace{0.45em}
\noindent\textbf{AC-RF-IAM-014 --- Facultad global de admin sobre verificación y exclusión}\par
\textbf{Método:} Prueba funcional de caja negra.\par
\textbf{Oráculo:} \texttt{admin} dispone de funciones con las que puede cambiar, para cualquier usuario, tanto el estado booleano de identidad verificada como el estado booleano de exclusión.\par\vspace{0.45em}
\noindent\textbf{AC-RF-IAM-015 --- Facultad global de admin sobre verificación reforzada}\par
\textbf{Método:} Prueba funcional de caja negra.\par
\textbf{Oráculo:} \texttt{admin} dispone de funciones con las que puede activar o desactivar la verificación reforzada de cualquier usuario.\par\vspace{0.45em}
\noindent\textbf{AC-RF-IAM-016 --- Facultad limitada del promotor sobre verificación y exclusión}\par
\textbf{Método:} Prueba funcional de caja negra.\par
\textbf{Oráculo:} El promotor puede marcar como verificada la identidad de un usuario o excluirlo solo si ese usuario está vinculado operativamente a un evento suyo y solo mientras el evento correspondiente permanezca vigente dentro de su ciclo de vida; fuera de ese alcance, el sistema no le habilita esas funciones.\par\vspace{0.45em}
\noindent\textbf{AC-RF-IAM-017 --- Imposibilidad del promotor para modificar la verificación reforzada}\par
\textbf{Método:} Prueba funcional de caja negra.\par
\textbf{Oráculo:} El promotor no dispone de funciones para cambiar la verificación reforzada de usuarios vinculados o no vinculados a su evento.\par\vspace{0.45em}
\noindent\textbf{AC-RF-IAM-018 --- Efecto combinado de los estados de verificación, exclusión y aceptación vigente}\par
\textbf{Método:} Prueba funcional de caja negra.\par
\textbf{Oráculo:} Un usuario solo puede ejecutar funciones operativas generales cuando su estado de identidad verificada sea verdadero, su estado de exclusión sea falso y tenga aceptación vigente de términos y condiciones; si el estado de exclusión es verdadero, no puede participar aunque la identidad está verificada.\par\vspace{0.45em}
\noindent\textbf{AC-RF-IAM-019 --- Efecto de la verificación reforzada sobre operaciones de mayor riesgo}\par
\textbf{Método:} Prueba funcional de caja negra.\par
\textbf{Oráculo:} Si un usuario no tiene verificación reforzada activa, el sistema no le habilita las funciones de inscripción de vehículos, puja en subastas de vehículos ni promoción de eventos, aunque cumpla las demás condiciones generales de participación.\par\vspace{0.45em}
\noindent\textbf{AC-RF-IAM-020 --- Cancelación de sesión en curso por exclusión}\par
\textbf{Método:} Prueba funcional de caja negra.\par
\textbf{Oráculo:} Activar la exclusión de un usuario autenticado y comprobar que la sesión deja de autorizar operaciones.\par\vspace{0.45em}
\noindent\textbf{AC-RF-IAM-021 --- Rechazo de inicio de sesión durante exclusión activa}\par
\textbf{Método:} Prueba funcional de caja negra.\par
\textbf{Oráculo:} Intentar iniciar sesión con un usuario excluido y comprobar que el acceso se rechaza.\par\vspace{0.45em}
\noindent\textbf{AC-RF-IAM-022 --- Mensaje de exclusión al rechazar el inicio de sesión}\par
\textbf{Método:} Prueba funcional de caja negra.\par
\textbf{Oráculo:} Intentar iniciar sesión con un usuario excluido y comprobar la presencia del mensaje aplicable.\par\vspace{0.45em}
\noindent\textbf{AC-RF-IAM-023 --- Conservación trazable de cambios en verificación, verificación reforzada y exclusión}\par
\textbf{Método:} Prueba funcional de caja negra.\par
\textbf{Oráculo:} Cuando \texttt{admin} o un promotor competente cambian alguno de esos estados de un usuario, el sistema conserva un registro consultable del actor, la marca temporal, el usuario afectado y el cambio efectuado, sin alterar por ello las reglas de reserva o visibilidad aplicables a esos estados.\par\vspace{0.45em}
\noindent\textbf{AC-RF-IAM-024 --- Exclusión de nueva identidad detectada como correspondiente a un usuario excluido}\par
\textbf{Método:} Prueba funcional de caja negra.\par
\textbf{Oráculo:} Cuando el sistema registra una detección válida de correspondencia entre una nueva identidad y un usuario ya excluido, la nueva identidad queda marcada también como excluida.\par\vspace{0.45em}
\noindent\textbf{AC-RF-IAM-025 --- Consulta personal de los estados booleanos del usuario}\par
\textbf{Método:} Prueba funcional de caja negra.\par
\textbf{Oráculo:} Un usuario puede consultar en el sistema si su identidad está marcada como verificada o no verificada, si su verificación reforzada está activa o inactiva y si su estado de exclusión está activo o no.\par\vspace{0.45em}
\noindent\textbf{AC-RF-IAM-026 --- Consulta personal del nivel visible de acceso}\par
\textbf{Método:} Prueba funcional de caja negra.\par
\textbf{Oráculo:} Consultar usuarios con combinaciones representativas de estados y comprobar que el nivel mostrado coincide con la regla de derivación vigente.\par\vspace{0.45em}
\noindent\textbf{AC-RF-IAM-027 --- Explicación de condiciones faltantes para el siguiente nivel de acceso}\par
\textbf{Método:} Prueba funcional de caja negra.\par
\textbf{Oráculo:} Consultar un usuario con condiciones faltantes y comparar la explicación con los estados y reglas de habilitación vigentes.\par\vspace{0.45em}
\noindent\textbf{AC-RF-IAM-028 --- Recordatorio al intentar una operación bloqueada por nivel visible de acceso}\par
\textbf{Método:} Prueba funcional de caja negra.\par
\textbf{Oráculo:} Cuando un usuario intenta una operación que el sistema no le habilita por su nivel visible de acceso, el sistema le muestra una explicación que incluye su nivel actual, la condición faltante y la acción necesaria para habilitar esa operación.\par\vspace{0.45em}
\noindent\textbf{AC-RF-IAM-029 --- Consulta administrativa de estados de verificación y exclusión}\par
\textbf{Método:} Prueba funcional de caja negra.\par
\textbf{Oráculo:} \texttt{admin} puede consultar para cualquier usuario si su identidad está verificada, si su verificación reforzada está activa y si su estado de exclusión está activo.\par\vspace{0.45em}
\noindent\textbf{AC-RF-IAM-030 --- Consulta limitada del promotor sobre estados de usuarios de su evento}\par
\textbf{Método:} Prueba funcional de caja negra.\par
\textbf{Oráculo:} El promotor puede consultar los booleanos de verificación y exclusión de un usuario solo si ese usuario está vinculado operativamente a uno de sus eventos y solo mientras ese evento permanezca vigente; fuera de ese alcance, el sistema no le muestra esos estados.\par\vspace{0.45em}
\noindent\textbf{AC-RF-IAM-031 --- Reserva de estados de otros usuarios frente a usuarios no autorizados}\par
\textbf{Método:} Prueba funcional de caja negra.\par
\textbf{Oráculo:} Un usuario distinto de \texttt{admin}, del propio usuario o de un promotor autorizado dentro de su alcance no puede consultar los estados de identidad verificada, verificación reforzada o exclusión de otro usuario, salvo cuando una regla específica del dominio lo autorice expresamente.\par\vspace{0.45em}
\noindent\textbf{AC-RF-IAM-032 --- Reserva de los estados y del nivel del usuario en la consulta pública del evento}\par
\textbf{Método:} Prueba funcional de caja negra.\par
\textbf{Oráculo:} En la consulta pública del evento se muestran los datos públicos del evento y de sus procesos visibles, pero no se muestran los estados de identidad verificada, verificación reforzada o exclusión de los usuarios vinculados a ese evento, ni el nivel derivado de ellos.\par\vspace{0.45em}
\noindent\textbf{AC-RF-IAM-033 --- Visibilidad general de usuarios no verificados y reserva de usuarios excluidos}\par
\textbf{Método:} Prueba funcional de caja negra.\par
\textbf{Oráculo:} Un usuario no verificado puede seguir apareciendo en las consultas del sistema o del evento que le correspondan por las demás reglas del dominio; un usuario con exclusión activa no aparece en consultas generales de usuarios para otros usuarios y sí aparece para \texttt{admin}, mientras el propio usuario puede seguir consultando sus estados y su nivel.\par\vspace{0.45em}
\noindent\textbf{AC-RF-REP-001 --- Registro de calificación del promotor por participantes}\par
\textbf{Método:} Prueba funcional de caja negra.\par
\textbf{Oráculo:} Después del evento, un participante vinculado a ese evento puede registrar en el sistema una calificación del promotor correspondiente.\par\vspace{0.45em}
\noindent\textbf{AC-RF-REP-002 --- Registro de calificación del promotor por espectadores}\par
\textbf{Método:} Prueba funcional de caja negra.\par
\textbf{Oráculo:} Después del evento, un espectador vinculado a ese evento puede registrar en el sistema una calificación del promotor correspondiente.\par\vspace{0.45em}
\noindent\textbf{AC-RF-REP-003 --- Determinación de la reputación del promotor}\par
\textbf{Método:} Prueba funcional de caja negra.\par
\textbf{Oráculo:} Cuando existan calificaciones registradas de eventos anteriores, el sistema conserva una reputación del promotor derivada de esas calificaciones.\par\vspace{0.45em}
\noindent\textbf{AC-RF-REP-004 --- Aplicación de la reputación del promotor a eventos futuros}\par
\textbf{Método:} Prueba funcional de caja negra.\par
\textbf{Oráculo:} Para un promotor con eventos futuros registrados, el sistema conserva la relación entre la reputación vigente de ese promotor y esos eventos futuros.\par\vspace{0.45em}
\noindent\textbf{AC-RF-BUS-001 --- Configuración administrativa de parámetros de comisión y cuota fija}\par
\textbf{Método:} Prueba funcional con cuenta admin y cuenta ordinaria.\par
\textbf{Oráculo:} Admin puede configurar cada parámetro, incluida la comisión del comprador de vehículo adjudicado; un usuario ordinario no puede modificar esos parámetros.\par\vspace{0.45em}
\noindent\textbf{AC-RF-BUS-002 --- Base positiva de comisión del competidor por evento}\par
\textbf{Método:} Prueba funcional con conjunto de datos de referencia y análisis de derivación.\par
\textbf{Oráculo:} Preparar casos con $R-C$ negativo, cero y positivo; la base de comisión es cero en los dos primeros y $R-C$ en el tercero, sin incluir compras de vehículos en $C$.\par\vspace{0.45em}
\noindent\textbf{AC-RF-BUS-003 --- Base positiva de comisión del competidor por campeonato}\par
\textbf{Método:} Prueba funcional con conjunto de datos de referencia y análisis de derivación.\par
\textbf{Oráculo:} Preparar casos con $R-C$ negativo, cero y positivo; la base de comisión es cero en los dos primeros y $R-C$ en el tercero.\par\vspace{0.45em}
\noindent\textbf{AC-RF-BUS-004 --- Base positiva de comisión del aportante asociado a un participante}\par
\textbf{Método:} Prueba funcional con conjunto de datos de referencia y análisis de derivación.\par
\textbf{Oráculo:} Preparar, para evento y campeonato, casos con $R-C$ negativo, cero y positivo; la base de comisión es cero en los dos primeros y $R-C$ en el tercero.\par\vspace{0.45em}
\noindent\textbf{AC-RF-BUS-005 --- Determinación de la comisión del comprador de vehículo adjudicado}\par
\textbf{Método:} Prueba económica con venta a tercero y autoadjudicación.\par
\textbf{Oráculo:} En ambos resultados, la comisión se calcula sobre el precio oficial y se carga al comprador/adjudicatario; el vendedor no soporta esa comisión por vender.\par\vspace{0.45em}
\noindent\textbf{AC-RF-BUS-006 --- Exención de comisión para patrocinadores directos}\par
\textbf{Método:} Prueba funcional con conjunto de datos de referencia y análisis de derivación.\par
\textbf{Oráculo:} Cuando un usuario realiza una aportación directa a una bolsa de evento o campeonato, el sistema registra el patrocinio sin calcular ni cobrar comisión al patrocinador por esa aportación.\par\vspace{0.45em}
\noindent\textbf{AC-RF-BUS-007 --- Cobro de cuota fija al promotor a partir del umbral anual configurado}\par
\textbf{Método:} Prueba funcional con conjunto de datos de referencia y análisis de derivación.\par
\textbf{Oráculo:} Cuando un promotor crea un evento nuevo y su cantidad anual de eventos creados ya alcanzó el umbral de gratuidad vigente, el sistema cobra la cuota fija correspondiente al crear ese evento; cuando no lo ha alcanzado, no la cobra.\par\vspace{0.45em}
\noindent\textbf{AC-RF-BUS-008 --- Cobro automático de comisiones y cuotas del modelo de negocio}\par
\textbf{Método:} Prueba funcional y económica con operaciones gravadas.\par
\textbf{Oráculo:} Cuando se produzca un premio positivo gravable, una adjudicación de vehículo --incluida autoadjudicación-- o la creación gravable de un evento, el sistema calcula y cobra automáticamente la comisión o cuota correspondiente.\par\vspace{0.45em}
\noindent\textbf{AC-RF-BUS-009 --- Información previa de comisión o cuota aplicable}\par
\textbf{Método:} Prueba de interacción antes y después de una subasta.\par
\textbf{Oráculo:} Antes de pujar se informa la regla y tasa vigente de comisión del comprador; después del cierre se muestra el importe efectivo derivado del precio oficial.\par\vspace{0.45em}
\noindent\textbf{AC-RF-BUS-010 --- Publicación posterior del detalle de comisión o cuota cobrada}\par
\textbf{Método:} Prueba funcional con conjunto de datos de referencia y análisis de derivación.\par
\textbf{Oráculo:} Cuando el sistema cobra una comisión o cuota del modelo de negocio, el usuario afectado puede consultar el detalle del cálculo, del importe cobrado y del importe neto resultante.\par\vspace{0.45em}
\noindent\textbf{AC-RF-BUS-011 --- Información expresa de exención de comisión cuando corresponda}\par
\textbf{Método:} Prueba funcional con conjunto de datos de referencia y análisis de derivación.\par
\textbf{Oráculo:} En operaciones de venta de vehículo por el vendedor o de patrocinio directo a bolsa, el sistema informa expresamente que la operación se encuentra exenta de comisión para ese actor.\par\vspace{0.45em}
\noindent\textbf{AC-UI-001 --- Acceso a funciones según habilitación del usuario}\par
\textbf{Método:} Prueba de interacción y autorización observable.\par
\textbf{Oráculo:} Comparar usuarios con niveles y papeles distintos; cada uno ve las funciones públicas y solo puede iniciar las operaciones para las que está habilitado.\par\vspace{0.45em}
\noindent\textbf{AC-UI-002 --- Consulta personal de estados y nivel visible de acceso}\par
\textbf{Método:} Prueba de interacción y autorización observable.\par
\textbf{Oráculo:} Iniciar sesión con usuarios representativos de cada combinación relevante y comprobar la correspondencia entre estados y nivel visible.\par\vspace{0.45em}
\noindent\textbf{AC-UI-003 --- Explicación de condiciones faltantes de acceso}\par
\textbf{Método:} Prueba de interacción y autorización observable.\par
\textbf{Oráculo:} Con un usuario que no cumpla una condición, consultar su nivel e intentar una operación bloqueada; la interfaz identifica la condición faltante en ambos contextos.\par\vspace{0.45em}
\noindent\textbf{AC-UI-004 --- Consulta de campeonatos patrocinados desde la interfaz principal}\par
\textbf{Método:} Prueba de interacción y autorización observable.\par
\textbf{Oráculo:} Desde la interfaz principal, consultar un campeonato patrocinado y comprobar la presencia de sus patrocinadores publicados.\par\vspace{0.45em}
\noindent\textbf{AC-UI-005 --- Cancelación inmediata de sesión por exclusión}\par
\textbf{Método:} Prueba de interacción y autorización observable.\par
\textbf{Oráculo:} Activar la exclusión de un usuario autenticado y comprobar que la sesión deja de autorizar operaciones.\par\vspace{0.45em}
\noindent\textbf{AC-UI-006 --- Bloqueo de inicio de sesión de usuario excluido}\par
\textbf{Método:} Prueba de interacción y autorización observable.\par
\textbf{Oráculo:} Intentar iniciar sesión con un usuario excluido; el acceso se rechaza y se muestra el mensaje definido.\par\vspace{0.45em}
\noindent\textbf{AC-UI-007 --- Interfaz para convocar y registrar eventos}\par
\textbf{Método:} Prueba de interacción y autorización observable.\par
\textbf{Oráculo:} Al registrar un evento, la interfaz solicita y publica esos datos, incluido el costo fijo por pesaje, al pasar el evento al estado de \texttt{inscripciones}.\par\vspace{0.45em}
\noindent\textbf{AC-UI-008 --- Captura de datos mínimos de inscripción}\par
\textbf{Método:} Prueba de interacción y autorización observable.\par
\textbf{Oráculo:} Registrar inscripciones con y sin pilotaje humano directo y comprobar que los datos mínimos aplicables quedan asociados a la inscripción correspondiente.\par\vspace{0.45em}
\noindent\textbf{AC-UI-009 --- Visualización del cobro inmediato de la inscripción}\par
\textbf{Método:} Prueba de interacción y autorización observable.\par
\textbf{Oráculo:} Aceptar una inscripción con importes conocidos y comprobar que la interfaz muestra el cobro ejecutado y su relación con la inscripción.\par\vspace{0.45em}
\noindent\textbf{AC-UI-010 --- Publicación de la cantidad de accesos disponibles}\par
\textbf{Método:} Prueba de interacción y autorización observable.\par
\textbf{Oráculo:} Abrir una subasta de acceso y comparar la cantidad publicada con la cantidad oficial configurada para el evento.\par\vspace{0.45em}
\noindent\textbf{AC-UI-011 --- Captura de pujas secretas de acceso con respaldo suficiente}\par
\textbf{Método:} Prueba de interacción y autorización observable.\par
\textbf{Oráculo:} Intentar pujar con saldo suficiente e insuficiente; solo la puja respaldada se acepta y queda asociada a su retención.\par\vspace{0.45em}
\noindent\textbf{AC-UI-012 --- Unicidad de la puja vigente de acceso por usuario}\par
\textbf{Método:} Prueba de interacción y autorización observable.\par
\textbf{Oráculo:} Registrar y actualizar pujas del mismo usuario y comprobar que solo una permanece vigente.\par\vspace{0.45em}
\noindent\textbf{AC-UI-013 --- Actualización solo al alza de la puja de acceso}\par
\textbf{Método:} Prueba de interacción y autorización observable.\par
\textbf{Oráculo:} Intentar sustituir una puja por importes inferior, igual y superior; solo el superior se acepta.\par\vspace{0.45em}
\noindent\textbf{AC-UI-014 --- Conteo público de pujas de acceso sin importes individuales}\par
\textbf{Método:} Prueba de interacción y autorización observable.\par
\textbf{Oráculo:} Comparar el conteo visible con las pujas vigentes y verificar con usuarios distintos que no se muestran importes individuales ajenos.\par\vspace{0.45em}
\noindent\textbf{AC-UI-015 --- Publicación del resultado de acceso al cierre}\par
\textbf{Método:} Prueba de interacción y autorización observable.\par
\textbf{Oráculo:} Cerrar una subasta conocida y comparar la publicación con el resultado oficial de adjudicación.\par\vspace{0.45em}
\noindent\textbf{AC-UI-016 --- Cobro automático de accesos adjudicados}\par
\textbf{Método:} Prueba de interacción y autorización observable.\par
\textbf{Oráculo:} Cerrar una subasta con ganadores y comprobar el cargo registrado para cada adjudicación.\par\vspace{0.45em}
\noindent\textbf{AC-UI-017 --- Liberación de retenciones de pujas de acceso no ganadoras}\par
\textbf{Método:} Prueba de interacción y autorización observable.\par
\textbf{Oráculo:} Cerrar una subasta con pujas perdedoras y comprobar que sus retenciones dejan de estar vigentes.\par\vspace{0.45em}
\noindent\textbf{AC-UI-018 --- Disponibilidad de la notificación QR del acceso ganado}\par
\textbf{Método:} Prueba de interacción y autorización observable.\par
\textbf{Oráculo:} Acceder como persona ganadora después del cierre y comprobar que la notificación y su QR están disponibles.\par\vspace{0.45em}
\noindent\textbf{AC-UI-019 --- Listado de vehículos computables en la subasta}\par
\textbf{Método:} Prueba de interacción y autorización observable.\par
\textbf{Oráculo:} Al terminar la carrera, comparar la lista de subasta con las inscripciones computables y registros de custodia; solo aparecen los vehículos elegibles.\par\vspace{0.45em}
\noindent\textbf{AC-UI-020 --- Captura de pujas selladas con la regla de retención aplicable}\par
\textbf{Método:} Prueba de interacción, autorización y reglas económicas observables.\par
\textbf{Oráculo:} Una puja de un usuario distinto del dueño se rechaza cuando carece de retención suficiente; la puja del dueño sobre su propio vehículo se acepta sin retención del principal cuando cumple las demás condiciones de elegibilidad; una puja de ese mismo dueño sobre un vehículo ajeno vuelve a exigir la retención ordinaria.\par\vspace{0.45em}
\noindent\textbf{AC-UI-021 --- Reserva total de información agregada de pujas antes del cierre}\par
\textbf{Método:} Prueba de interacción y autorización observable.\par
\textbf{Oráculo:} Durante una subasta abierta, ningún usuario visualiza conteos, orden, precios, adjudicatarios provisionales ni indicadores derivados de terceros.\par\vspace{0.45em}
\noindent\textbf{AC-UI-022 --- Consulta y aumento de la puja propia vigente}\par
\textbf{Método:} Prueba de interacción y autorización observable.\par
\textbf{Oráculo:} Cada usuario visualiza su propia puja vigente; intentos de reducirla o mantenerla igual son rechazados y solo un aumento la sustituye.\par\vspace{0.45em}
\noindent\textbf{AC-UI-023 --- Ausencia de eficiencia, puntos y premios provisionales durante la subasta}\par
\textbf{Método:} Prueba de interacción y autorización observable.\par
\textbf{Oráculo:} Modificar pujas durante una subasta abierta no produce vistas visibles de clasificación, puntos ni premio de eficiencia.\par\vspace{0.45em}
\noindent\textbf{AC-UI-024 --- Publicación posterior del resultado oficial y su naturaleza}\par
\textbf{Método:} Prueba de interacción y autorización observable.\par
\textbf{Oráculo:} Tras el cierre de la fase se muestran, para cada vehículo, precio oficial, adjudicatario y la etiqueta de transferencia a tercero o autoadjudicación; antes del cierre esos datos no aparecen.\par\vspace{0.45em}
\noindent\textbf{AC-UI-025 --- Publicación del precio final y ganador al cierre}\par
\textbf{Método:} Prueba de interacción y autorización observable.\par
\textbf{Oráculo:} Cerrar la subasta y comparar los datos publicados con el resultado oficial.\par\vspace{0.45em}
\noindent\textbf{AC-UI-026 --- Información sobre transferencia obligatoria}\par
\textbf{Método:} Prueba de interacción y autorización observable.\par
\textbf{Oráculo:} Consultar un vehículo subastado y comprobar que la condición de transferencia obligatoria es visible.\par\vspace{0.45em}
\noindent\textbf{AC-UI-027 --- Gestión de entrega y disputa del vehículo adjudicado}\par
\textbf{Método:} Prueba de interacción y autorización observable.\par
\textbf{Oráculo:} Ejecutar cada operación dentro y fuera de plazo; dentro de plazo se registra y consulta, fuera de plazo se rechaza conforme a las restricciones externas.\par\vspace{0.45em}
\noindent\textbf{AC-UI-028 --- Consulta pública de resultados y clasificaciones del evento}\par
\textbf{Método:} Prueba de interacción y autorización observable.\par
\textbf{Oráculo:} Consultar un evento publicado y comparar los resultados y clasificaciones mostrados con sus registros oficiales.\par\vspace{0.45em}
\noindent\textbf{AC-UI-029 --- Consulta pública de la matriz de pesos del evento}\par
\textbf{Método:} Prueba de interacción y autorización observable.\par
\textbf{Oráculo:} Consultar un evento y comparar la matriz publicada con la matriz oficial aplicable.\par\vspace{0.45em}
\noindent\textbf{AC-UI-030 --- Selección del corte territorial y temporal de standings}\par
\textbf{Método:} Prueba de interacción y autorización observable.\par
\textbf{Oráculo:} Seleccionar combinaciones válidas de niveles y comprobar que el corte consultado cambia conforme a la selección.\par\vspace{0.45em}
\noindent\textbf{AC-UI-031 --- Publicación de campeonatos y puntajes del corte seleccionado}\par
\textbf{Método:} Prueba de interacción y autorización observable.\par
\textbf{Oráculo:} Consultar un corte conocido y comparar campeonatos y puntajes publicados con los resultados oficiales correspondientes.\par\vspace{0.45em}
\noindent\textbf{AC-UI-032 --- Lista estandarizada de comprobación de seguridad}\par
\textbf{Método:} Prueba de interacción y autorización observable.\par
\textbf{Oráculo:} Abrir la comprobación de un evento y comparar la lista mostrada con la versión vigente de la lista estandarizada.\par\vspace{0.45em}
\noindent\textbf{AC-UI-033 --- Registro del proceso formal de objeciones de seguridad}\par
\textbf{Método:} Prueba de interacción y autorización observable.\par
\textbf{Oráculo:} Registrar una objeción, una votación y una decisión; comprobar su asociación y secuencia trazable en el evento.\par\vspace{0.45em}
\noindent\textbf{AC-UI-034 --- Captura opcional de datos de seguridad e inspección}\par
\textbf{Método:} Prueba de interacción y autorización observable.\par
\textbf{Oráculo:} Continuar un caso permitido sin datos opcionales y otro con datos capturados; ambos siguen las mismas guardas normativas restantes.\par\vspace{0.45em}
\noindent\textbf{AC-UI-035 --- Exclusión del pesaje externo de la interfaz de seguridad}\par
\textbf{Método:} Prueba de interacción y autorización observable.\par
\textbf{Oráculo:} Inspeccionar y probar la interfaz de seguridad; no existe operación de captura de pesajes externos.\par\vspace{0.45em}
\noindent\textbf{AC-UI-036 --- Interfaz comunitaria de discusión y novedades}\par
\textbf{Método:} Prueba de interacción y autorización observable.\par
\textbf{Oráculo:} La interfaz permite consultar el foro comunitario y el canal de novedades, permite que cualquier usuario con identidad verificada publique en ambos y permite que \texttt{admin} ejerza funciones de moderación sobre ellos.\par\vspace{0.45em}
\noindent\textbf{AC-UI-037 --- Interfaz de premio individual del evento}\par
\textbf{Método:} Prueba de interacción y autorización observable.\par
\textbf{Oráculo:} Después de la determinación económica del evento, la interfaz permite consultar públicamente el premio individual de cada participante computable y la distribución económica posterior asociada, con anonimización de terceros cuando corresponda.\par\vspace{0.45em}
\noindent\textbf{AC-UI-038 --- Selección del campeonato para operaciones económicas}\par
\textbf{Método:} Prueba de interacción y autorización observable.\par
\textbf{Oráculo:} Seleccionar campeonatos de ambos tipos y comprobar que las operaciones posteriores quedan asociadas al campeonato elegido.\par\vspace{0.45em}
\noindent\textbf{AC-UI-039 --- Registro de aportaciones directas a la bolsa del campeonato}\par
\textbf{Método:} Prueba de interacción y autorización observable.\par
\textbf{Oráculo:} Registrar una aportación directa dentro y fuera de la ventana; solo la operación temporalmente válida se acepta.\par\vspace{0.45em}
\noindent\textbf{AC-UI-040 --- Registro condicionado de aportaciones a participantes de campeonato}\par
\textbf{Método:} Prueba de interacción y autorización observable.\par
\textbf{Oráculo:} Probar la misma aportación con la función activa e inactiva; solo el caso activo se acepta.\par\vspace{0.45em}
\noindent\textbf{AC-UI-041 --- Cobro inmediato de aportaciones de campeonato aceptadas}\par
\textbf{Método:} Prueba de interacción y autorización observable.\par
\textbf{Oráculo:} Registrar una aportación aceptada y comprobar el cargo y su marca temporal.\par\vspace{0.45em}
\noindent\textbf{AC-UI-042 --- Publicación del anuncio de patrocinio directo del campeonato}\par
\textbf{Método:} Prueba de interacción y autorización observable.\par
\textbf{Oráculo:} Registrar un patrocinio directo con imagen y comprobar su publicación en el campeonato.\par\vspace{0.45em}
\noindent\textbf{AC-UI-043 --- Consulta pública del premio individual de campeonato}\par
\textbf{Método:} Prueba de interacción y autorización observable.\par
\textbf{Oráculo:} Consultar un campeonato con premios oficiales y comparar los importes publicados con la liquidación oficial.\par\vspace{0.45em}
\noindent\textbf{AC-UI-044 --- Consulta pública de la distribución posterior del premio de campeonato}\par
\textbf{Método:} Prueba de interacción y autorización observable.\par
\textbf{Oráculo:} Consultar una distribución con aportantes terceros y comprobar importes públicos y anonimización de identidades cuando corresponda.\par\vspace{0.45em}
\noindent\textbf{AC-UI-045 --- Registro de aportaciones directas a la bolsa del evento}\par
\textbf{Método:} Prueba de interacción y autorización observable.\par
\textbf{Oráculo:} Registrar una aportación directa durante \texttt{inscripciones} y comprobar su asociación con la bolsa del evento.\par\vspace{0.45em}
\noindent\textbf{AC-UI-046 --- Registro condicionado de aportaciones de eficiencia a participantes del evento}\par
\textbf{Método:} Prueba de interacción y autorización observable.\par
\textbf{Oráculo:} Probar una aportación de eficiencia con la función activa e inactiva; solo el caso activo se acepta.\par\vspace{0.45em}
\noindent\textbf{AC-UI-047 --- Registro condicionado de aportaciones de terminación a inscripciones}\par
\textbf{Método:} Prueba de interacción y autorización observable.\par
\textbf{Oráculo:} Probar una aportación de terminación con la función activa e inactiva; solo el caso activo se acepta.\par\vspace{0.45em}
\noindent\textbf{AC-UI-048 --- Cobro inmediato de aportaciones del evento aceptadas}\par
\textbf{Método:} Prueba de interacción y autorización observable.\par
\textbf{Oráculo:} Registrar una aportación aceptada y comprobar el cargo y su marca temporal.\par\vspace{0.45em}
\noindent\textbf{AC-UI-049 --- Publicación del anuncio de patrocinio directo del evento}\par
\textbf{Método:} Prueba de interacción y autorización observable.\par
\textbf{Oráculo:} Registrar un patrocinio directo con imagen y comprobar su publicación en el evento.\par\vspace{0.45em}
\noindent\textbf{AC-UI-050 --- Distinción del patrocinio de campeonatos aplicable al evento}\par
\textbf{Método:} Prueba de interacción y autorización observable.\par
\textbf{Oráculo:} Consultar un evento con ambos tipos de patrocinio y comprobar que su procedencia se distingue inequívocamente.\par\vspace{0.45em}
\noindent\textbf{AC-UI-051 --- Interfaz de calificación del promotor}\par
\textbf{Método:} Prueba de interacción y autorización observable.\par
\textbf{Oráculo:} Después de un evento, un participante o espectador vinculado a ese evento puede acceder a una interfaz en la que registra una calificación del promotor correspondiente.\par\vspace{0.45em}
\noindent\textbf{AC-UI-052 --- Configuración administrativa de parámetros del modelo de negocio}\par
\textbf{Método:} Prueba de interacción y autorización observable.\par
\textbf{Oráculo:} Acceder como \texttt{admin}, actualizar cada familia de parámetros y comprobar la vigencia del nuevo valor en su consulta administrativa.\par\vspace{0.45em}
\noindent\textbf{AC-UI-053 --- Información previa sobre comisión o cuota aplicable}\par
\textbf{Método:} Prueba de interacción y autorización observable.\par
\textbf{Oráculo:} Abrir operaciones gravadas y exentas y comparar la información previa con la configuración administrativa vigente.\par\vspace{0.45em}
\noindent\textbf{AC-UI-054 --- Consulta posterior del cobro o exención del modelo de negocio}\par
\textbf{Método:} Prueba de interacción y autorización observable.\par
\textbf{Oráculo:} Ejecutar casos gravados y exentos y comprobar que el detalle posterior identifica el resultado y su base normativa aplicable.\par\vspace{0.45em}
\noindent\textbf{AC-UI-055 --- Conjunto de pantallas críticas para evaluación UX}\par
\textbf{Método:} Prueba de interacción y autorización observable.\par
\textbf{Oráculo:} El protocolo de evaluación UX de percepción inicial positiva debe incluir como mínimo esas interfaces y debe medir en ellas la reacción a los 30 segundos y a los 5 minutos conforme al criterio definido en RAC-006.\par\vspace{0.45em}
\noindent\textbf{AC-UI-056 --- Contexto común de inscripciones en el detalle del evento}\par
\textbf{Método:} Prueba de interacción y autorización observable.\par
\textbf{Oráculo:} Abrir el detalle de un evento y comprobar que las inscripciones y sus datos asociados son consultables en un mismo contexto operativo.\par\vspace{0.45em}
\noindent\textbf{AC-UI-057 --- Captura de mejor vuelta válida durante velocidad}\par
\textbf{Método:} Prueba de interacción y autorización observable.\par
\textbf{Oráculo:} Con el evento en \texttt{velocidad}, registrar un tiempo válido para una inscripción y comprobar su consulta posterior.\par\vspace{0.45em}
\noindent\textbf{AC-UI-058 --- Ausencia de velocidad válida por incumplimiento de pesaje decidido fuera del sistema}\par
\textbf{Método:} Prueba de interacción y autorización observable.\par
\textbf{Oráculo:} Marcar el caso aplicable y comprobar que la inscripción queda sin velocidad válida capturada y conserva la trazabilidad de la decisión operativa correspondiente.\par\vspace{0.45em}
\noindent\textbf{AC-UI-059 --- Captura de lugar oficial durante carrera}\par
\textbf{Método:} Prueba de interacción y autorización observable.\par
\textbf{Oráculo:} Con el evento en \texttt{carrera}, registrar lugares oficiales y comprobar su persistencia y consulta.\par\vspace{0.45em}
\noindent\textbf{AC-UI-060 --- Captura de vueltas durante carrera}\par
\textbf{Método:} Prueba de interacción y autorización observable.\par
\textbf{Oráculo:} Registrar vueltas en carrera y comprobar que el cálculo de terminación puede referenciar esos valores.\par\vspace{0.45em}
\noindent\textbf{AC-UI-061 --- Descalificación con causa textual}\par
\textbf{Método:} Prueba de interacción y autorización observable.\par
\textbf{Oráculo:} Intentar descalificar sin causa y con causa; la primera operación se rechaza y la segunda conserva la causa.\par\vspace{0.45em}
\noindent\textbf{AC-UI-062 --- Consulta económica de terminación en el detalle del evento}\par
\textbf{Método:} Prueba de interacción y autorización observable.\par
\textbf{Oráculo:} Después del cálculo, abrir el detalle y comprobar la presencia de los tres elementos para cada inscripción aplicable.\par\vspace{0.45em}
\noindent\textbf{AC-UI-063 --- Reserva de datos privados en el detalle operativo}\par
\textbf{Método:} Prueba de interacción y autorización observable.\par
\textbf{Oráculo:} Comparar el detalle con dos usuarios; cada uno accede solo a sus datos privados autorizados.\par\vspace{0.45em}
\noindent\textbf{AC-UI-064 --- Bloqueo de edición fuera de condiciones habilitantes}\par
\textbf{Método:} Prueba de interacción y autorización observable.\par
\textbf{Oráculo:} Intentar editar la misma información en estados y con papeles no habilitados; la lectura permitida permanece y la edición se rechaza.\par\vspace{0.45em}
\noindent\textbf{AC-UI-065 --- Selección del corte territorial y temporal de prioridades}\par
\textbf{Método:} Prueba de interacción y autorización observable.\par
\textbf{Oráculo:} Seleccionar combinaciones de niveles y comparar la lista resultante con los eventos pertenecientes al corte.\par\vspace{0.45em}
\noindent\textbf{AC-UI-066 --- Publicación de discriminadores y P/C por evento}\par
\textbf{Método:} Prueba de interacción y autorización observable.\par
\textbf{Oráculo:} Consultar un corte con eventos conocidos y comparar cada dato publicado con el registro oficial correspondiente.\par\vspace{0.45em}
\noindent\textbf{AC-UI-067 --- Ordenamiento consultable por P/C}\par
\textbf{Método:} Prueba de interacción y autorización observable.\par
\textbf{Oráculo:} Solicitar el orden por P/C y comprobar que la secuencia corresponde a los valores publicados.\par\vspace{0.45em}
\noindent\textbf{AC-UI-068 --- Navegación al detalle público del evento}\par
\textbf{Método:} Prueba de interacción y autorización observable.\par
\textbf{Oráculo:} Seleccionar un evento de la lista y comprobar que se abre su detalle público correspondiente.\par\vspace{0.45em}
\noindent\textbf{AC-UI-069 --- Información pública del pozo durante la subasta de vehículos}\par
\textbf{Método:} Prueba de interacción durante la fase de subasta con aportaciones previamente aceptadas.\par
\textbf{Oráculo:} El detalle público muestra el pozo inicial publicado y no muestra el pozo consolidado ni importes de aportaciones o contribuciones antes del cierre de la fase.\par\vspace{0.45em}
\noindent\textbf{AC-UI-070 --- Declaración de categoría tecnológica y documentación abierta obligatoria}\par
\textbf{Método:} Prueba de interfaz de publicación e inscripción.\par
\textbf{Oráculo:} Al publicar un evento se registra y muestra la categoría tecnológica admitida; al inscribir un vehículo el sistema exige la documentación abierta obligatoria del vehículo conforme a la definición de tecnología abierta del vehículo, sin que esa documentación sustituya las reglas de seguridad, medición, participación ni el cálculo de eficiencia.\par\vspace{0.45em}
\noindent\textbf{AC-RC-EXT-001 --- Procesamiento obligatorio de retenciones, cobros, pagos y devoluciones reales}\par
\textbf{Método:} Prueba negativa/positiva de restricción externa.\par
\textbf{Oráculo:} En los procesos de inscripción, aportación, patrocinio, subasta y premio, el sistema ejecuta las operaciones económicas reales que el dominio exige para cada caso.\par\vspace{0.45em}
\noindent\textbf{AC-RC-EXT-002 --- Sujeción a plazos publicados del evento en disputas de entrega de vehículos}\par
\textbf{Método:} Prueba negativa/positiva de restricción externa.\par
\textbf{Oráculo:} En un evento con subasta de vehículos, no puede registrarse válidamente una disputa de entrega ni una resolución administrativa fuera de los plazos publicados del evento.\par\vspace{0.45em}
\noindent\textbf{AC-RC-EXT-003 --- Custodia previa obligatoria del vehículo antes de su subasta}\par
\textbf{Método:} Prueba negativa/positiva de restricción externa.\par
\textbf{Oráculo:} Si un vehículo no ha sido entregado al promotor, la subasta correspondiente no puede iniciarse válidamente respecto de ese vehículo.\par\vspace{0.45em}
\noindent\textbf{AC-RC-EXT-004 --- Régimen obligatorio de publicidad del dominio}\par
\textbf{Método:} Prueba negativa/positiva de restricción externa.\par
\textbf{Oráculo:} En la consulta pública del evento, de los campeonatos patrocinados y del historial del evento, la información declarada pública por el dominio permanece visible con el alcance exigido por sus reglas correspondientes.\par\vspace{0.45em}
\noindent\textbf{AC-RC-EXT-005 --- Régimen obligatorio de reserva selectiva de identidad}\par
\textbf{Método:} Prueba negativa/positiva de restricción externa.\par
\textbf{Oráculo:} Cuando exista una protesta de seguridad visible en la consulta pública del evento, la identidad de quien la presentó no se muestra públicamente, pero sí está disponible para \texttt{admin} y para el promotor.\par\vspace{0.45em}
\noindent\textbf{AC-RC-EXT-006 --- Verificación de identidad y no exclusión antes de la participación}\par
\textbf{Método:} Prueba negativa/positiva de restricción externa.\par
\textbf{Oráculo:} Un usuario puede consultar la información pública disponible, pero no puede inscribirse, pujar, aportar, promover, presentar protestas ni ejecutar otras operaciones de participación dentro del sistema si su identidad no está verificada o si su estado de exclusión está activo.\par\vspace{0.45em}
\noindent\textbf{AC-RC-EXT-007 --- Verificación reforzada obligatoria para operaciones de mayor riesgo}\par
\textbf{Método:} Prueba negativa/positiva de restricción externa.\par
\textbf{Oráculo:} Un usuario que no tenga verificación reforzada activa no puede inscribir vehículos, pujar en subastas de vehículos ni promover eventos, aunque tenga identidad verificada, no está excluido y haya aceptado la versión vigente de términos y condiciones.\par\vspace{0.45em}
\noindent\textbf{AC-RC-EXT-008 --- Aceptación de la versión vigente de términos y condiciones antes de la participación}\par
\textbf{Método:} Prueba negativa/positiva de restricción externa.\par
\textbf{Oráculo:} Un usuario puede consultar la información pública disponible, pero no puede inscribirse, pujar, aportar, promover, presentar protestas ni ejecutar otras operaciones de participación dentro del sistema si no tiene aceptación registrada de la versión vigente de términos y condiciones.\par\vspace{0.45em}
\noindent\textbf{AC-RC-EXT-009 --- Sujeción del vehículo subastado a transferencia obligatoria y prohibición de retiro no autorizado}\par
\textbf{Método:} Prueba negativa/positiva de restricción externa.\par
\textbf{Oráculo:} En un evento con subasta de vehículos, las reglas visibles del dominio establecen que el vehículo subastado no puede retirarse del control del evento contra el resultado oficial vigente de la subasta y que el retiro no autorizado se conserva como incumplimiento grave.\par\vspace{0.45em}
\noindent\textbf{AC-QR-001 --- Claridad operativa de estados y acciones}\par
\textbf{Método:} Inspección, medición o prueba UX según el atributo.\par
\textbf{Oráculo:} En cada estado oficial, revisar el conjunto de operaciones representativas: el estado es visible, las operaciones habilitadas pueden iniciarse y las bloqueadas muestran una causa.\par\vspace{0.45em}
\noindent\textbf{AC-QR-002 --- Auditabilidad y trazabilidad de operaciones críticas}\par
\textbf{Método:} Inspección, medición o prueba UX según el atributo.\par
\textbf{Oráculo:} Seleccionar cambios de estado, operaciones económicas y resultados oficiales; reconstruir cada caso mediante los registros publicados o autorizados y comprobar actor, tiempo, datos y reglas aplicables.\par\vspace{0.45em}
\noindent\textbf{AC-QR-003 --- Integridad de derivaciones del dominio}\par
\textbf{Método:} Inspección, medición o prueba UX según el atributo.\par
\textbf{Oráculo:} Ejecutar casos de invalidación y comparar todas las derivaciones afectadas antes y después; los registros invalidados dejan de computar sin perder su trazabilidad histórica.\par\vspace{0.45em}
\noindent\textbf{AC-QR-004 --- Transparencia pública controlada}\par
\textbf{Método:} Inspección, medición o prueba UX según el atributo.\par
\textbf{Oráculo:} Ejecutar una matriz de pruebas por tipo de dato y papel de usuario; cada dato público es consultable y cada dato reservado es inaccesible para actores no autorizados.\par\vspace{0.45em}
\noindent\textbf{AC-QR-005 --- Permanencia y consultabilidad histórica}\par
\textbf{Método:} Inspección, medición o prueba UX según el atributo.\par
\textbf{Oráculo:} Publicar condiciones y resultados; recorrer el ciclo de vida y el historial y comprobar que la información permanece disponible en los periodos definidos.\par\vspace{0.45em}
\noindent\textbf{AC-QR-006 --- Reproducibilidad numérica sin redondeo intermedio no prescrito}\par
\textbf{Método:} Inspección, medición o prueba UX según el atributo.\par
\textbf{Oráculo:} Repetir un conjunto de cálculos analíticos de referencia bajo punto flotante binario de 64 bits y comparar sus resultados dentro de la tolerancia definida por el caso de prueba; para distribuciones monetarias, comprobar que todos los importes son enteros y que su suma coincide exactamente con el total mediante \texttt{CALC-MON-001}.\par\vspace{0.45em}
\noindent\textbf{AC-QR-007 --- Retroalimentación perceptible al inicio de acciones interactivas}\par
\textbf{Método:} Inspección, medición o prueba UX según el atributo.\par
\textbf{Oráculo:} Medir en \texttt{ENV-UX-001} desde la activación hasta la primera señal perceptible de recepción en una muestra representativa; el percentil 95 no excede 100 ms y ninguna medición individual excede 250 ms.\par\vspace{0.45em}
\noindent\textbf{AC-QR-008 --- Visibilidad de espera y bloqueo por incompatibilidad operativa}\par
\textbf{Método:} Inspección, medición o prueba UX según el atributo.\par
\textbf{Oráculo:} Ejecutar operaciones representativas de Q, M, T, F y A; si una operación supera 500 ms, el estado de espera permanece visible. El conjunto de acciones bloqueadas coincide con \texttt{OPS-001} y una consulta independiente permanece habilitada.\par\vspace{0.45em}
\noindent\textbf{AC-QR-009 --- Percepción inicial positiva en pantallas críticas}\par
\textbf{Método:} Inspección, medición o prueba UX según el atributo.\par
\textbf{Oráculo:} Aplicar el protocolo a las pantallas críticas definidas; registrar tamaño y criterios de selección de la muestra, instrumento y resultados en ambos momentos.\par\vspace{0.45em}
\noindent\textbf{AC-QR-010 --- Jerarquía visual de información crítica}\par
\textbf{Método:} Inspección, medición o prueba UX según el atributo.\par
\textbf{Oráculo:} Realizar inspección estructurada de cada pantalla crítica y una prueba de identificación con usuarios representativos; cada elemento aplicable puede localizarse y relacionarse con su función.\par\vspace{0.45em}
\noindent\textbf{AC-QR-011 --- Legibilidad del contenido crítico}\par
\textbf{Método:} Inspección, medición o prueba UX según el atributo.\par
\textbf{Oráculo:} Verificar WCAG 2.2 AA en los cuatro viewports de \texttt{ENV-UX-001}, a 100\% y 200\% de zoom, sin pérdida de información o funcionalidad crítica; complementar con una prueba de campo en pista al aire libre.\par\vspace{0.45em}
\noindent\textbf{AC-QR-012 --- Contraste y codificación redundante de estados críticos}\par
\textbf{Método:} Inspección, medición o prueba UX según el atributo.\par
\textbf{Oráculo:} Inspeccionar las pantallas críticas sin información cromática y comprobar que las distinciones críticas conservan una señal adicional; medir los contrastes aplicables y comprobar conformidad con WCAG 2.2 nivel AA.\par\vspace{0.45em}
\noindent\textbf{AC-QR-013 --- Consistencia visual y terminológica}\par
\textbf{Método:} Inspección, medición o prueba UX según el atributo.\par
\textbf{Oráculo:} Comparar pantallas equivalentes mediante un inventario de conceptos, etiquetas y patrones; toda divergencia debe estar justificada y registrada.\par\vspace{0.45em}
\noindent\textbf{AC-QR-014 --- Prevención de error en acciones sensibles}\par
\textbf{Método:} Inspección, medición o prueba UX según el atributo.\par
\textbf{Oráculo:} Probar pagos, pujas, cancelaciones, exclusiones, entrega y aceptación legal; cada flujo presenta contexto suficiente y no permite una confirmación sin identificar inequívocamente la acción y su objeto.\par\vspace{0.45em}
\noindent\textbf{AC-QR-015 --- Claridad previa de consecuencias}\par
\textbf{Método:} Inspección, medición o prueba UX según el atributo.\par
\textbf{Oráculo:} En cada acción sensible representativa, comparar la información previa con los efectos ejecutados: los cambios de estado, cargos, devoluciones, comisiones, bloqueos, exclusiones y efectos legales aplicables se informan antes de confirmar.\par\vspace{0.45em}
\section{Trazabilidad}
\subsection{Mapa de identificadores heredados a requisitos canónicos}
La tabla sustituye las listas masivas de trazabilidad inversa del documento heredado. Un ID heredado puede mapearse a varios requisitos canónicos cuando se atomizó una obligación compuesta.
\begin{longtable}{P{3.0cm}P{10.6cm}}
\toprule
\textbf{ID heredado} & \textbf{ID canónico v2}\\
\midrule

\texttt{RAC-001} & \texttt{QR-006}\\
\texttt{RAC-002} & \texttt{QR-002}\\
\texttt{RAC-003} & \texttt{QR-003}\\
\texttt{RAC-004} & \texttt{QR-005}\\
\texttt{RAC-005} & \texttt{QR-004}\\
\texttt{RAC-006} & \texttt{QR-009}\\
\texttt{RAC-007} & \texttt{QR-010}\\
\texttt{RAC-008} & \texttt{QR-011}\\
\texttt{RAC-009} & \texttt{QR-012}\\
\texttt{RAC-010} & \texttt{QR-013}\\
\texttt{RAC-011} & \texttt{QR-014}\\
\texttt{RAC-012} & \texttt{QR-015}\\
\texttt{REX-001} & \texttt{RC-EXT-001}\\
\texttt{REX-002} & \texttt{RC-EXT-002}\\
\texttt{REX-003} & \texttt{RC-EXT-003}\\
\texttt{REX-004} & \texttt{RC-EXT-004}\\
\texttt{REX-005} & \texttt{RC-EXT-005}\\
\texttt{REX-006} & \texttt{RC-EXT-006}\\
\texttt{REX-006A} & \texttt{RC-EXT-007}\\
\texttt{REX-007} & \texttt{RC-EXT-008}\\
\texttt{REX-008} & \texttt{RC-EXT-009}\\
\texttt{RF-001} & \texttt{RF-EVT-001}\\
\texttt{RF-002} & \texttt{RF-EVT-002}\\
\texttt{RF-003} & \texttt{RF-EVT-003}, \texttt{RF-EVT-004}\\
\texttt{RF-004} & \texttt{RF-EVT-005}\\
\texttt{RF-005} & \texttt{RF-EVT-006}\\
\texttt{RF-006} & \texttt{RF-EVT-007}\\
\texttt{RF-007} & \texttt{RF-EVT-008}\\
\texttt{RF-008} & \texttt{RF-EVT-009}\\
\texttt{RF-009} & \texttt{RF-EVT-010}\\
\texttt{RF-010} & \texttt{RF-EVT-011}\\
\texttt{RF-011} & \texttt{RF-EVT-012}, \texttt{RF-EVT-013}\\
\texttt{RF-011A} & \texttt{RF-EVT-014}\\
\texttt{RF-012} & \texttt{RF-EVT-015}\\
\texttt{RF-013} & \texttt{RF-EVT-016}\\
\texttt{RF-014} & \texttt{RF-EVT-017}\\
\texttt{RF-015} & \texttt{RF-EVT-018}, \texttt{RF-EVT-019}\\
\texttt{RF-016} & \texttt{RF-EVT-020}\\
\texttt{RF-017} & \texttt{RF-EVT-021}\\
\texttt{RF-018} & \texttt{RF-EVT-022}\\
\texttt{RF-019} & \texttt{RF-EVT-023}\\
\texttt{RF-020} & \texttt{RF-EVT-024}\\
\texttt{RF-021} & \texttt{RF-EVT-025}\\
\texttt{RF-022} & \texttt{RF-EVT-026}\\
\texttt{RF-022A} & \texttt{RF-EVT-027}\\
\texttt{RF-023} & \texttt{RF-EVT-028}\\
\texttt{RF-024} & \texttt{RF-EVT-029}\\
\texttt{RF-025} & \texttt{RF-EVT-030}\\
\texttt{RF-026} & \texttt{RF-EVT-031}, \texttt{RF-EVT-032}\\
\texttt{RF-027} & \texttt{RF-EVT-033}\\
\texttt{RF-028} & \texttt{RF-EVT-034}\\
\texttt{RF-029} & \texttt{RF-EVT-035}\\
\texttt{RF-030} & \texttt{RF-EVT-036}\\
\texttt{RF-031} & \texttt{RF-EVT-037}\\
\texttt{RF-032} & \texttt{RF-EVT-038}\\
\texttt{RF-033} & \texttt{RF-EVT-039}\\
\texttt{RF-034} & \texttt{RF-EVT-040}\\
\texttt{RF-035} & \texttt{RF-PAY-001}\\
\texttt{RF-036} & \texttt{RF-PAY-002}\\
\texttt{RF-037} & \texttt{RF-PAY-003}\\
\texttt{RF-038} & \texttt{RF-PAY-004}\\
\texttt{RF-039} & \texttt{RF-PAY-005}\\
\texttt{RF-040} & \texttt{RF-PAY-006}\\
\texttt{RF-040A} & \texttt{RF-PAY-007}\\
\texttt{RF-040B} & \texttt{RF-PAY-008}\\
\texttt{RF-041} & \texttt{RF-PAY-009}\\
\texttt{RF-042} & \texttt{RF-AUC-001}\\
\texttt{RF-043} & \texttt{RF-AUC-002}\\
\texttt{RF-043A} & \texttt{RF-AUC-003}\\
\texttt{RF-043B} & \texttt{RF-AUC-004}, \texttt{RF-AUC-005}\\
\texttt{RF-044} & \texttt{RF-AUC-006}\\
\texttt{RF-045} & \texttt{RF-AUC-007}, \texttt{RF-AUC-008}, \texttt{RF-AUC-009}, \texttt{RF-AUC-010}\\
\texttt{RF-045A} & \texttt{RF-AUC-011}\\
\texttt{RF-045B} & \texttt{RF-AUC-012}\\
\texttt{RF-045C} & \texttt{RF-AUC-013}\\
\texttt{RF-046} & \texttt{RF-AUC-014}\\
\texttt{RF-047} & \texttt{RF-AUC-015}\\
\texttt{RF-047A} & \texttt{RF-AUC-016}\\
\texttt{RF-048} & \texttt{RF-AUC-017}\\
\texttt{RF-049} & \texttt{RF-AUC-018}\\
\texttt{RF-050} & \texttt{RF-AUC-019}\\
\texttt{RF-050A} & \texttt{RF-AUC-020}\\
\texttt{RF-050B} & \texttt{RF-AUC-021}\\
\texttt{RF-050C} & \texttt{RF-AUC-022}\\
\texttt{RF-050D} & \texttt{RF-AUC-023}\\
\texttt{RF-050E} & \texttt{RF-AUC-024}\\
\texttt{RF-050F} & \texttt{RF-AUC-025}\\
\texttt{RF-050G} & \texttt{RF-AUC-026}\\
\texttt{RF-050H} & \texttt{RF-AUC-027}\\
\texttt{RF-050I} & \texttt{RF-AUC-028}\\
\texttt{RF-050J} & \texttt{RF-AUC-029}, \texttt{RF-AUC-030}, \texttt{RF-AUC-031}\\
\texttt{RF-051} & \texttt{RF-CLS-001}\\
\texttt{RF-052} & \texttt{RF-CLS-002}\\
\texttt{RF-053} & \texttt{RF-CLS-003}\\
\texttt{RF-054} & \texttt{RF-CLS-004}\\
\texttt{RF-055} & \texttt{RF-CLS-005}\\
\texttt{RF-055A} & \texttt{RF-CLS-006}\\
\texttt{RF-055B} & \texttt{RF-CLS-007}\\
\texttt{RF-056} & \texttt{RF-SAF-001}\\
\texttt{RF-057} & \texttt{RF-SAF-002}\\
\texttt{RF-058} & \texttt{RF-SAF-003}\\
\texttt{RF-058A} & \texttt{RF-SAF-004}\\
\texttt{RF-059} & \texttt{RF-SAF-005}\\
\texttt{RF-060} & \texttt{RF-SAF-006}\\
\texttt{RF-061} & \texttt{RF-SAF-007}\\
\texttt{RF-062} & \texttt{RF-SAF-008}\\
\texttt{RF-062A} & \texttt{RF-SAF-009}\\
\texttt{RF-063} & \texttt{RF-COM-001}\\
\texttt{RF-064} & \texttt{RF-COM-002}\\
\texttt{RF-065} & \texttt{RF-COM-003}\\
\texttt{RF-066} & \texttt{RF-COM-004}\\
\texttt{RF-066A} & \texttt{RF-COM-005}\\
\texttt{RF-067} & \texttt{RF-PUB-001}\\
\texttt{RF-067A} & \texttt{RF-PUB-005}\\
\texttt{RF-067B} & \texttt{RF-PUB-006}\\
\texttt{RF-067C} & \texttt{RF-PUB-007}\\
\texttt{RF-067D} & \texttt{RF-PUB-008}\\
\texttt{RF-067E} & \texttt{RF-PUB-009}\\
\texttt{RF-068} & \texttt{RF-PUB-002}\\
\texttt{RF-069} & \texttt{RF-PUB-003}\\
\texttt{RF-070} & \texttt{RF-PUB-004}\\
\texttt{RF-070A} & \texttt{RF-PUB-010}\\
\texttt{RF-071} & \texttt{RF-IAM-001}\\
\texttt{RF-071A} & \texttt{RF-IAM-002}\\
\texttt{RF-071B} & \texttt{RF-IAM-003}\\
\texttt{RF-072} & \texttt{RF-IAM-004}\\
\texttt{RF-072A} & \texttt{RF-IAM-005}\\
\texttt{RF-072B} & \texttt{RF-IAM-006}\\
\texttt{RF-072C} & \texttt{RF-IAM-007}\\
\texttt{RF-072D} & \texttt{RF-IAM-008}\\
\texttt{RF-073} & \texttt{RF-IAM-009}\\
\texttt{RF-074} & \texttt{RF-IAM-010}\\
\texttt{RF-075} & \texttt{RF-EVT-041}\\
\texttt{RF-075A} & \texttt{RF-EVT-042}\\
\texttt{RF-075B} & \texttt{RF-EVT-043}\\
\texttt{RF-076} & \texttt{RF-EVT-044}\\
\texttt{RF-077} & \texttt{RF-EVT-045}\\
\texttt{RF-078} & \texttt{RF-EVT-046}\\
\texttt{RF-079} & \texttt{RF-EVT-047}\\
\texttt{RF-079A} & \texttt{RF-EVT-048}\\
\texttt{RF-079B} & \texttt{RF-EVT-049}\\
\texttt{RF-079C} & \texttt{RF-EVT-050}\\
\texttt{RF-080} & \texttt{RF-VEH-001}\\
\texttt{RF-081} & \texttt{RF-VEH-002}\\
\texttt{RF-082} & \texttt{RF-VEH-003}\\
\texttt{RF-083} & \texttt{RF-VEH-004}\\
\texttt{RF-083A} & \texttt{RF-VEH-005}\\
\texttt{RF-084} & \texttt{RF-VEH-006}\\
\texttt{RF-085} & \texttt{RF-VEH-007}\\
\texttt{RF-086} & \texttt{RF-VEH-008}\\
\texttt{RF-086A} & \texttt{RF-VEH-009}, \texttt{RF-VEH-010}\\
\texttt{RF-087} & \texttt{RF-VEH-011}\\
\texttt{RF-088} & \texttt{RF-VEH-012}\\
\texttt{RF-089} & \texttt{RF-VEH-013}\\
\texttt{RF-090} & \texttt{RF-VEH-014}\\
\texttt{RF-090A} & \texttt{RF-VEH-015}\\
\texttt{RF-090B} & \texttt{RF-VEH-016}\\
\texttt{RF-091} & \texttt{RF-EPR-001}\\
\texttt{RF-092} & \texttt{RF-EPR-002}\\
\texttt{RF-093} & \texttt{RF-EPR-003}\\
\texttt{RF-094} & \texttt{RF-EPR-004}\\
\texttt{RF-095} & \texttt{RF-CPR-001}\\
\texttt{RF-096} & \texttt{RF-CPR-002}\\
\texttt{RF-097} & \texttt{RF-CPR-003}\\
\texttt{RF-098} & \texttt{RF-CPR-004}\\
\texttt{RF-099} & \texttt{RF-CPR-005}\\
\texttt{RF-100} & \texttt{RF-CPR-006}\\
\texttt{RF-101} & \texttt{RF-CPR-007}, \texttt{RF-CPR-008}\\
\texttt{RF-102} & \texttt{RF-CPR-009}\\
\texttt{RF-103} & \texttt{RF-WGT-001}\\
\texttt{RF-104} & \texttt{RF-WGT-002}\\
\texttt{RF-105} & \texttt{RF-WGT-003}\\
\texttt{RF-106} & \texttt{RF-WGT-004}\\
\texttt{RF-107} & \texttt{RF-WGT-005}\\
\texttt{RF-108} & \texttt{RF-WGT-006}, \texttt{RF-WGT-007}, \texttt{RF-WGT-008}\\
\texttt{RF-108A} & \texttt{RF-WGT-009}\\
\texttt{RF-108B} & \texttt{RF-WGT-010}\\
\texttt{RF-109} & \texttt{RF-WGT-011}\\
\texttt{RF-110} & \texttt{RF-EFF-001}\\
\texttt{RF-111} & \texttt{RF-EFF-002}, \texttt{RF-EFF-003}, \texttt{RF-EFF-004}, \texttt{RF-EFF-005}\\
\texttt{RF-112} & \texttt{RF-EFF-006}, \texttt{RF-EFF-007}, \texttt{RF-EFF-008}, \texttt{RF-EFF-009}\\
\texttt{RF-113} & \texttt{RF-EFF-010}\\
\texttt{RF-113A} & \texttt{RF-EFF-011}, \texttt{RF-EFF-012}, \texttt{RF-EFF-013}\\
\texttt{RF-114} & \texttt{RF-EFF-014}, \texttt{RF-EFF-015}, \texttt{RF-EFF-016}, \texttt{RF-EFF-017}, \texttt{RF-EFF-018}\\
\texttt{RF-115} & \texttt{RF-EPR-005}\\
\texttt{RF-116} & \texttt{RF-EPR-006}\\
\texttt{RF-116A} & \texttt{RF-EPR-007}\\
\texttt{RF-117} & \texttt{RF-EPR-008}\\
\texttt{RF-118} & \texttt{RF-EPR-009}\\
\texttt{RF-118A} & \texttt{RF-EPR-010}\\
\texttt{RF-118B} & \texttt{RF-EPR-011}\\
\texttt{RF-118C} & \texttt{RF-EPR-012}, \texttt{RF-EPR-013}, \texttt{RF-EPR-014}, \texttt{RF-EPR-015}, \texttt{RF-EPR-016}, \texttt{RF-EPR-017}\\
\texttt{RF-119} & \texttt{RF-EPR-018}\\
\texttt{RF-119A} & \texttt{RF-EPR-019}\\
\texttt{RF-120} & \texttt{RF-EPR-020}\\
\texttt{RF-120A} & \texttt{RF-EPR-021}\\
\texttt{RF-121} & \texttt{RF-CPR-010}\\
\texttt{RF-122} & \texttt{RF-CPR-011}\\
\texttt{RF-123} & \texttt{RF-CPR-012}, \texttt{RF-CPR-013}, \texttt{RF-CPR-014}\\
\texttt{RF-123A} & \texttt{RF-CPR-015}\\
\texttt{RF-123B} & \texttt{RF-CPR-016}\\
\texttt{RF-123C} & \texttt{RF-CPR-017}\\
\texttt{RF-124} & \texttt{RF-CPR-018}\\
\texttt{RF-125} & \texttt{RF-CPR-019}\\
\texttt{RF-125A} & \texttt{RF-CPR-020}\\
\texttt{RF-126} & \texttt{RF-SPN-001}\\
\texttt{RF-127} & \texttt{RF-SPN-002}\\
\texttt{RF-127A} & \texttt{RF-SPN-003}\\
\texttt{RF-128} & \texttt{RF-SPN-004}\\
\texttt{RF-129} & \texttt{RF-SPN-005}, \texttt{RF-SPN-006}, \texttt{RF-SPN-007}\\
\texttt{RF-129A} & \texttt{RF-SPN-008}\\
\texttt{RF-129B} & \texttt{RF-SPN-009}\\
\texttt{RF-130} & \texttt{RF-PTS-001}, \texttt{RF-PTS-002}\\
\texttt{RF-131} & \texttt{RF-PTS-003}, \texttt{RF-PTS-004}\\
\texttt{RF-132} & \texttt{RF-PTS-005}, \texttt{RF-PTS-006}\\
\texttt{RF-133} & \texttt{RF-PTS-007}\\
\texttt{RF-134} & \texttt{RF-PTS-008}\\
\texttt{RF-135} & \texttt{RF-PTS-009}\\
\texttt{RF-136} & \texttt{RF-PTS-010}\\
\texttt{RF-137} & \texttt{RF-PTS-011}\\
\texttt{RF-138} & \texttt{RF-CPR-021}\\
\texttt{RF-139} & \texttt{RF-CPR-022}\\
\texttt{RF-140} & \texttt{RF-CPR-023}\\
\texttt{RF-141} & \texttt{RF-CPR-024}\\
\texttt{RF-142} & \texttt{RF-GOV-001}\\
\texttt{RF-143} & \texttt{RF-GOV-002}\\
\texttt{RF-144} & \texttt{RF-GOV-003}\\
\texttt{RF-145} & \texttt{RF-GOV-004}\\
\texttt{RF-146} & \texttt{RF-GOV-005}\\
\texttt{RF-147} & \texttt{RF-IAM-011}\\
\texttt{RF-147A} & \texttt{RF-IAM-012}\\
\texttt{RF-148} & \texttt{RF-IAM-013}\\
\texttt{RF-149} & \texttt{RF-IAM-014}\\
\texttt{RF-149A} & \texttt{RF-IAM-015}\\
\texttt{RF-150} & \texttt{RF-IAM-016}\\
\texttt{RF-150A} & \texttt{RF-IAM-017}\\
\texttt{RF-151} & \texttt{RF-IAM-018}\\
\texttt{RF-151A} & \texttt{RF-IAM-019}\\
\texttt{RF-151B} & \texttt{RF-IAM-020}, \texttt{RF-IAM-021}, \texttt{RF-IAM-022}\\
\texttt{RF-152} & \texttt{RF-IAM-023}\\
\texttt{RF-152A} & \texttt{RF-IAM-024}\\
\texttt{RF-153} & \texttt{RF-IAM-025}\\
\texttt{RF-153A} & \texttt{RF-IAM-026}, \texttt{RF-IAM-027}\\
\texttt{RF-153B} & \texttt{RF-IAM-028}\\
\texttt{RF-154} & \texttt{RF-IAM-029}\\
\texttt{RF-154A} & \texttt{RF-IAM-030}\\
\texttt{RF-155} & \texttt{RF-IAM-031}\\
\texttt{RF-156} & \texttt{RF-IAM-032}\\
\texttt{RF-157} & \texttt{RF-IAM-033}\\
\texttt{RF-158} & \texttt{RF-REP-001}\\
\texttt{RF-159} & \texttt{RF-REP-002}\\
\texttt{RF-160} & \texttt{RF-REP-003}\\
\texttt{RF-161} & \texttt{RF-REP-004}\\
\texttt{RF-162} & \texttt{RF-BUS-001}\\
\texttt{RF-163} & \texttt{RF-BUS-002}\\
\texttt{RF-164} & \texttt{RF-BUS-003}\\
\texttt{RF-165} & \texttt{RF-BUS-004}\\
\texttt{RF-166} & \texttt{RF-BUS-005}\\
\texttt{RF-167} & \texttt{RF-BUS-006}\\
\texttt{RF-168} & \texttt{RF-BUS-007}\\
\texttt{RF-169} & \texttt{RF-BUS-008}\\
\texttt{RF-170} & \texttt{RF-BUS-009}\\
\texttt{RF-171} & \texttt{RF-BUS-010}\\
\texttt{RF-172} & \texttt{RF-BUS-011}\\
\texttt{RIE-001} & \texttt{UI-001}, \texttt{UI-002}, \texttt{UI-003}, \texttt{UI-004}, \texttt{UI-005}, \texttt{UI-006}\\
\texttt{RIE-002} & \texttt{UI-007}\\
\texttt{RIE-003} & \texttt{UI-008}, \texttt{UI-009}\\
\texttt{RIE-004} & \texttt{UI-010}, \texttt{UI-011}, \texttt{UI-012}, \texttt{UI-013}, \texttt{UI-014}, \texttt{UI-015}, \texttt{UI-016}, \texttt{UI-017}, \texttt{UI-018}\\
\texttt{RIE-005} & \texttt{UI-019}, \texttt{UI-020}, \texttt{UI-021}, \texttt{UI-022}, \texttt{UI-023}, \texttt{UI-024}, \texttt{UI-025}, \texttt{UI-026}, \texttt{UI-027}\\
\texttt{RIE-006} & \texttt{UI-028}, \texttt{UI-029}, \texttt{UI-030}, \texttt{UI-031}\\
\texttt{RIE-007} & \texttt{UI-032}, \texttt{UI-033}, \texttt{UI-034}, \texttt{UI-035}\\
\texttt{RIE-008} & \texttt{UI-036}\\
\texttt{RIE-009} & \texttt{UI-037}\\
\texttt{RIE-010} & \texttt{UI-038}, \texttt{UI-039}, \texttt{UI-040}, \texttt{UI-041}, \texttt{UI-042}, \texttt{UI-043}, \texttt{UI-044}\\
\texttt{RIE-011} & \texttt{UI-045}, \texttt{UI-046}, \texttt{UI-047}, \texttt{UI-048}, \texttt{UI-049}, \texttt{UI-050}\\
\texttt{RIE-012} & \texttt{UI-051}\\
\texttt{RIE-013} & \texttt{UI-052}, \texttt{UI-053}, \texttt{UI-054}\\
\texttt{RIE-014} & \texttt{UI-055}\\
\texttt{RIE-015} & \texttt{QR-010}\\
\texttt{RIE-016} & \texttt{UI-056}, \texttt{UI-057}, \texttt{UI-058}, \texttt{UI-059}, \texttt{UI-060}, \texttt{UI-061}, \texttt{UI-062}, \texttt{UI-063}, \texttt{UI-064}\\
\texttt{RIE-017} & \texttt{UI-065}, \texttt{UI-066}, \texttt{UI-067}, \texttt{UI-068}\\
\texttt{RNF-001} & \texttt{QR-001}\\
\texttt{RNF-002} & \texttt{QR-002}\\
\texttt{RNF-003} & \texttt{QR-003}\\
\texttt{RNF-004} & \texttt{QR-002}\\
\texttt{RNF-005} & \texttt{QR-005}\\
\texttt{RNF-006} & \texttt{QR-006}\\
\texttt{RNF-007} & \texttt{QR-007}\\
\texttt{RNF-008} & \texttt{QR-008}\\
\texttt{RNF-009} & \texttt{QR-010}\\
\texttt{RNF-010} & \texttt{QR-011}\\
\texttt{RNF-011} & \texttt{QR-012}\\
\texttt{RNF-012} & \texttt{QR-013}\\
\texttt{RNF-013} & \texttt{QR-014}\\
\texttt{RNF-014} & \texttt{QR-015}\\
\bottomrule
\end{longtable}

\subsection{Trazabilidad canónica hacia procedencia}
\begin{landscape}
\begin{longtable}{P{2.7cm}P{2.6cm}P{15.7cm}}
\toprule
\textbf{ID canónico} & \textbf{ID heredado} & \textbf{Procedencia declarada en la SRS heredada}\\
\midrule

\texttt{RF-EVT-001} & \texttt{RF-001} & Secciones 1.2, 2.2 y 2.6.1.\\
\texttt{RF-EVT-002} & \texttt{RF-002} & Secciones 2.5, 2.6 y 2.6.1.\\
\texttt{RF-EVT-003} & \texttt{RF-003} & Secciones 2.2, 2.5.1 y 2.7.\\
\texttt{RF-EVT-004} & \texttt{RF-003} & Secciones 2.2, 2.5.1 y 2.7.\\
\texttt{RF-EVT-005} & \texttt{RF-004} & Secciones 1.2, 2.5.1 y RF-011.\\
\texttt{RF-EVT-006} & \texttt{RF-005} & Secciones 2.2, 2.6.1 y 2.7.\\
\texttt{RF-EVT-007} & \texttt{RF-006} & Secciones 1.2, 2.5, 2.6.1 y 2.7.\\
\texttt{RF-EVT-008} & \texttt{RF-007} & Secciones 2.5 y 2.7.\\
\texttt{RF-EVT-009} & \texttt{RF-008} & Secciones 2.4 y 2.7.\\
\texttt{RF-EVT-010} & \texttt{RF-009} & Sección 2.5.1.\\
\texttt{RF-EVT-011} & \texttt{RF-010} & Secciones 2.2 y 2.5.1.\\
\texttt{RF-EVT-012} & \texttt{RF-011} & Secciones 2.2 y 2.5.1.\\
\texttt{RF-EVT-013} & \texttt{RF-011} & Secciones 2.2 y 2.5.1.\\
\texttt{RF-EVT-014} & \texttt{RF-011A} & RF-011, RF-071A y decisiones prioritarias del cliente sobre paso a \texttt{velocidad}.\\
\texttt{RF-EVT-015} & \texttt{RF-012} & Secciones 2.2 y 2.5.1.\\
\texttt{RF-EVT-016} & \texttt{RF-013} & Secciones 2.2, 2.4, 2.5.1, RF-046 y decisiones prioritarias del cliente sobre entrega previa del vehículo al promotor.\\
\texttt{RF-EVT-017} & \texttt{RF-014} & Secciones 2.5, 2.6.4, 2.7, RF-005 y RF-006.\\
\texttt{RF-EVT-018} & \texttt{RF-015} & Secciones 2.5.1, 2.6.4 y decisiones prioritarias del cliente sobre devolución inmediata en cancelación del evento.\\
\texttt{RF-EVT-019} & \texttt{RF-015} & Secciones 2.5.1, 2.6.4 y decisiones prioritarias del cliente sobre devolución inmediata en cancelación del evento.\\
\texttt{RF-EVT-020} & \texttt{RF-016} & Secciones 2.5 y 2.7.\\
\texttt{RF-EVT-021} & \texttt{RF-017} & Secciones 2.5.1 y 2.7.\\
\texttt{RF-EVT-022} & \texttt{RF-018} & Secciones 2.5.1 y 2.7.\\
\texttt{RF-EVT-023} & \texttt{RF-019} & Secciones 2.4, 2.7 y RF-008.\\
\texttt{RF-EVT-024} & \texttt{RF-020} & Secciones 2.5.1 y 2.7.\\
\texttt{RF-EVT-025} & \texttt{RF-021} & Reglas operativas del ciclo del evento (2.2, 2.5.1), RF-042 a RF-045B, RF-086A y decisiones prioritarias del cliente sobre cobro inmediato.\\
\texttt{RF-EVT-026} & \texttt{RF-022} & Reglas operativas del ciclo del evento (2.2, 2.5.1) y decisiones prioritarias del cliente sobre vigencia hasta antes de \texttt{carrera}.\\
\texttt{RF-EVT-027} & \texttt{RF-022A} & RF-057 a RF-061 y decisiones prioritarias del cliente sobre vigencia hasta antes de \texttt{carrera}.\\
\texttt{RF-EVT-028} & \texttt{RF-023} & Reglas operativas del ciclo del evento (2.2, 2.4, 2.5.1) y decisiones prioritarias del cliente sobre entrega del vehículo al promotor al terminar la carrera.\\
\texttt{RF-EVT-029} & \texttt{RF-024} & Reglas operativas del ciclo del evento (2.2, 2.4, 2.5.1) y RF-046 a RF-050J.\\
\texttt{RF-EVT-030} & \texttt{RF-025} & Secciones 2.2, 2.5.1 y 2.6.1.\\
\texttt{RF-EVT-031} & \texttt{RF-026} & Secciones 2.5, 2.7 y RF-007.\\
\texttt{RF-EVT-032} & \texttt{RF-026} & Secciones 2.5, 2.7 y RF-007.\\
\texttt{RF-EVT-033} & \texttt{RF-027} & Secciones 2.5.1 y 2.7.\\
\texttt{RF-EVT-034} & \texttt{RF-028} & Sección 2.7 y RF-020.\\
\texttt{RF-EVT-035} & \texttt{RF-029} & RF-021, sección 2.6.2 y decisiones prioritarias del cliente sobre cierre al pasar a \texttt{velocidad}.\\
\texttt{RF-EVT-036} & \texttt{RF-030} & Integridad del proceso económico previo a la actividad competitiva (2.5.1, 2.7) y decisiones prioritarias del cliente.\\
\texttt{RF-EVT-037} & \texttt{RF-031} & Reglas de pozo y reparto (1.2, 2.2, 2.6.2).\\
\texttt{RF-EVT-038} & \texttt{RF-032} & Sección 2.7.\\
\texttt{RF-EVT-039} & \texttt{RF-033} & Sección 2.7 y RF-020.\\
\texttt{RF-EVT-040} & \texttt{RF-034} & Sección 2.7 y RF-020.\\
\texttt{RF-PAY-001} & \texttt{RF-035} & Secciones 2.6.2, 2.6.3 y 2.6.4.\\
\texttt{RF-PAY-002} & \texttt{RF-036} & Secciones 2.6.2, 2.6.3 y 2.7.\\
\texttt{RF-PAY-003} & \texttt{RF-037} & Sección 2.6.4.\\
\texttt{RF-PAY-004} & \texttt{RF-038} & Secciones 2.6.4 y 2.7.\\
\texttt{RF-PAY-005} & \texttt{RF-039} & Sección 2.6.4 y fuente prioritaria confirmada por el usuario sobre contribuciones de terceros.\\
\texttt{RF-PAY-006} & \texttt{RF-040} & Secciones 2.6.3 y 2.6.4.\\
\texttt{RF-PAY-007} & \texttt{RF-040A} & Secciones 2.6.4 y decisiones prioritarias del cliente sobre devolución inmediata.\\
\texttt{RF-PAY-008} & \texttt{RF-040B} & Sección 2.6.4 y decisiones prioritarias del cliente sobre pago automático a todos los beneficiarios.\\
\texttt{RF-PAY-009} & \texttt{RF-041} & Secciones 2.6.3, 2.6.4 y 2.7.\\
\texttt{RF-AUC-001} & \texttt{RF-042} & Secciones 2.2, 2.4, 2.6.1 y fuente documental \texttt{TOP\_Racing\_Grok\_Clarified.tex} (líneas 146 y 219).\\
\texttt{RF-AUC-002} & \texttt{RF-043} & Secciones 2.4, 2.6.4 y decisiones prioritarias del cliente sobre pagos reales y retención en subasta de acceso.\\
\texttt{RF-AUC-003} & \texttt{RF-043A} & RF-043 y decisiones prioritarias del cliente sobre puja vigente única en subasta de acceso.\\
\texttt{RF-AUC-004} & \texttt{RF-043B} & RF-043A y decisiones prioritarias del cliente sobre actualización solo al alza en subasta de acceso.\\
\texttt{RF-AUC-005} & \texttt{RF-043B} & RF-043A y decisiones prioritarias del cliente sobre actualización solo al alza en subasta de acceso.\\
\texttt{RF-AUC-006} & \texttt{RF-044} & Secciones 2.4, 2.7 y fuente documental \texttt{TOP\_Racing\_Grok\_Clarified.tex} (línea 219).\\
\texttt{RF-AUC-007} & \texttt{RF-045} & Secciones 2.4, 2.7 y fuente documental \texttt{TOP\_Racing\_Grok\_Clarified.tex} (línea 219).\\
\texttt{RF-AUC-008} & \texttt{RF-045} & Secciones 2.4, 2.7 y fuente documental \texttt{TOP\_Racing\_Grok\_Clarified.tex} (línea 219).\\
\texttt{RF-AUC-009} & \texttt{RF-045} & Secciones 2.4, 2.7 y fuente documental \texttt{TOP\_Racing\_Grok\_Clarified.tex} (línea 219).\\
\texttt{RF-AUC-010} & \texttt{RF-045} & Secciones 2.4, 2.7 y fuente documental \texttt{TOP\_Racing\_Grok\_Clarified.tex} (línea 219).\\
\texttt{RF-AUC-011} & \texttt{RF-045A} & RF-045, fuente documental \texttt{TOP\_Racing\_Grok\_Clarified.tex} (líneas 219 y 223) y decisiones prioritarias del cliente sobre QR en lugar de ticket.\\
\texttt{RF-AUC-012} & \texttt{RF-045B} & Secciones 2.6.4 y decisiones prioritarias del cliente sobre pagos reales en subasta de acceso.\\
\texttt{RF-AUC-013} & \texttt{RF-045C} & RF-045B y decisiones prioritarias del cliente sobre devolución inmediata por cancelación del evento.\\
\texttt{RF-AUC-014} & \texttt{RF-046} & Secciones 2.2, 2.4, 2.5.1, 2.7, fuente documental \texttt{TOP\_Racing\_Grok\_Clarified.tex} (líneas 146, 221 y 223) y decisiones prioritarias del cliente sobre entrega del vehículo al promotor.\\
\texttt{RF-AUC-015} & \texttt{RF-047} & Secciones 2.4, 2.6.4 y decisiones prioritarias del cliente sobre garantía real de pago en subasta de vehículos.\\
\texttt{RF-AUC-016} & \texttt{RF-047A} & Secciones 2.3, 2.5, 2.6, 3.1.24 y decisiones prioritarias del cliente sobre puja de vehículos con verificación reforzada máxima.\\
\texttt{RF-AUC-017} & \texttt{RF-048} & Secciones 2.4, 2.7 y fuente documental \texttt{TOP\_Racing\_Grok\_Clarified.tex} (líneas 221 y 223).\\
\texttt{RF-AUC-018} & \texttt{RF-049} & Secciones 2.4, 2.7, RF-050C y fuente documental \texttt{TOP\_Racing\_Grok\_Clarified.tex} (líneas 221 y 223).\\
\texttt{RF-AUC-019} & \texttt{RF-050} & Sección 1.3, fuente documental \texttt{TOP\_Racing\_Grok\_Clarified.tex} y conversación fuente prioritaria consolidada.\\
\texttt{RF-AUC-020} & \texttt{RF-050A} & Decisiones prioritarias del cliente sobre entrega del vehículo al promotor antes de subasta.\\
\texttt{RF-AUC-021} & \texttt{RF-050B} & Decisiones prioritarias del cliente sobre no entrega o modificación del vehículo, descalificación y morosidad del dueño.\\
\texttt{RF-AUC-022} & \texttt{RF-050C} & Secciones 2.4, 2.5, RF-050A y decisiones prioritarias del cliente sobre \texttt{parc ferme} y control del evento.\\
\texttt{RF-AUC-023} & \texttt{RF-050D} & Secciones 2.6.4 y decisiones prioritarias del cliente sobre garantía real de pago en subasta de vehículos.\\
\texttt{RF-AUC-024} & \texttt{RF-050E} & Secciones 2.6.4 y decisiones prioritarias del cliente sobre cobro automático al ganador.\\
\texttt{RF-AUC-025} & \texttt{RF-050F} & Secciones 2.6.4 y decisiones prioritarias del cliente sobre pago automático al dueño.\\
\texttt{RF-AUC-026} & \texttt{RF-050G} & Decisiones prioritarias del cliente sobre condición de deudor del promotor hasta la entrega.\\
\texttt{RF-AUC-027} & \texttt{RF-050H} & Decisiones prioritarias del cliente sobre confirmación de entrega en la aplicación.\\
\texttt{RF-AUC-028} & \texttt{RF-050I} & RF-075, decisiones prioritarias del cliente sobre disputa de entrega y plazos del evento.\\
\texttt{RF-AUC-029} & \texttt{RF-050J} & Decisiones prioritarias del cliente sobre resolución por \texttt{admin} y efectos de la morosidad del promotor.\\
\texttt{RF-AUC-030} & \texttt{RF-050J} & Decisiones prioritarias del cliente sobre resolución por \texttt{admin} y efectos de la morosidad del promotor.\\
\texttt{RF-AUC-031} & \texttt{RF-050J} & Decisiones prioritarias del cliente sobre resolución por \texttt{admin} y efectos de la morosidad del promotor.\\
\texttt{RF-AUC-032} & -- & \texttt{DEC-164}; participación del dueño en la subasta de su propio vehículo.\\
\texttt{RF-AUC-033} & -- & \texttt{DEC-165}; excepción de retención del principal para puja propia.\\
\texttt{RF-AUC-034} & -- & \texttt{DEC-166}; autoadjudicación sin flujo circular de principal ni entrega a tercero.\\
\texttt{RF-AUC-035} & -- & \texttt{DEC-172}, \texttt{DEC-173}; advertencia sobre inspección y estudio técnico posterior a la adjudicación.\\
\texttt{RF-CLS-001} & \texttt{RF-051} & Secciones 2.2, 2.6.4, 2.7 y fuente documental \texttt{TOP\_Racing\_Grok\_Clarified.tex} (líneas 150 y 167).\\
\texttt{RF-CLS-002} & \texttt{RF-052} & Secciones 1.2, 2.2 y fuente documental \texttt{TOP\_Racing\_Grok\_Clarified.tex} (líneas 150, 163 y 167).\\
\texttt{RF-CLS-003} & \texttt{RF-053} & Secciones 1.2, 2.2, 2.4 y fuente documental \texttt{TOP\_Racing\_Grok\_Clarified.tex} (líneas 150, 163 y 171).\\
\texttt{RF-CLS-004} & \texttt{RF-054} & Secciones 2.2, 2.4, 2.7 y fuente documental \texttt{TOP\_Racing\_Grok\_Clarified.tex} (líneas 150 y 171).\\
\texttt{RF-CLS-005} & \texttt{RF-055} & Secciones 1.2, 2.2, 2.4, 2.7 y fuente documental \texttt{TOP\_Racing\_Grok\_Clarified.tex} (líneas 150 y 171).\\
\texttt{RF-CLS-006} & \texttt{RF-055A} & RF-051 a RF-055 y fuente prioritaria confirmada por el usuario sobre efectos de la cancelación de la inscripción.\\
\texttt{RF-CLS-007} & \texttt{RF-055B} & RF-050B, RF-050J, RF-067E y fuentes prioritarias confirmadas por el usuario sobre resultado final vigente de la subasta.\\
\texttt{RF-SAF-001} & \texttt{RF-056} & Secciones 2.2, 2.3 y fuente documental \texttt{TOP\_Racing\_Grok\_Clarified.tex} (líneas 131 y 152).\\
\texttt{RF-SAF-002} & \texttt{RF-057} & Secciones 2.2 y 2.3, y fuente documental \texttt{TOP\_Racing\_Grok\_Clarified.tex} (líneas 95, 97 y 152), y fuente prioritaria confirmada por el usuario sobre vigencia de la inspección hasta antes de \texttt{carrera}.\\
\texttt{RF-SAF-003} & \texttt{RF-058} & Secciones 2.2, 2.3 y fuente documental \texttt{TOP\_Racing\_Grok\_Clarified.tex} (línea 99), y fuente prioritaria confirmada por el usuario sobre vigencia de las protestas hasta antes de \texttt{carrera}.\\
\texttt{RF-SAF-004} & \texttt{RF-058A} & RF-058 y fuente prioritaria confirmada por el usuario sobre visibilidad restringida de quien presenta la protesta.\\
\texttt{RF-SAF-005} & \texttt{RF-059} & Secciones 2.2 y 2.3, y fuente documental \texttt{TOP\_Racing\_Grok\_Clarified.tex} (líneas 99, 101 y 131).\\
\texttt{RF-SAF-006} & \texttt{RF-060} & Secciones 2.2, 2.3 y fuente documental \texttt{TOP\_Racing\_Grok\_Clarified.tex} (líneas 108 a 112), y fuentes prioritarias confirmadas por el usuario sobre carácter no vinculante de la protesta y publicidad del detalle de votos.\\
\texttt{RF-SAF-007} & \texttt{RF-061} & Secciones 2.2, 2.3 y fuente documental \texttt{TOP\_Racing\_Grok\_Clarified.tex} (líneas 99 y 131), y fuentes prioritarias confirmadas por el usuario sobre decisión final del promotor y cancelación derivada de protesta.\\
\texttt{RF-SAF-008} & \texttt{RF-062} & Secciones 2.2, 2.3, 2.7 y fuente documental \texttt{TOP\_Racing\_Grok\_Clarified.tex} (líneas 101 y 152).\\
\texttt{RF-SAF-009} & \texttt{RF-062A} & RF-058 a RF-061 y fuente prioritaria confirmada por el usuario sobre carácter no vinculante de la protesta.\\
\texttt{RF-COM-001} & \texttt{RF-063} & Secciones 2.2, 2.3, 2.5, 3.4 y decisiones prioritarias del cliente sobre publicación comunitaria.\\
\texttt{RF-COM-002} & \texttt{RF-064} & Secciones 2.2, 2.3 y fuente documental \texttt{TOP\_Racing\_Grok\_Clarified.tex} (línea 154).\\
\texttt{RF-COM-003} & \texttt{RF-065} & Secciones 2.1, 2.2, 2.3, 2.5, 3.4 y decisiones prioritarias del cliente sobre publicación de novedades.\\
\texttt{RF-COM-004} & \texttt{RF-066} & Secciones 1.2, 2.2 y fuente documental \texttt{TOP\_Racing\_Grok\_Clarified.tex} (línea 154).\\
\texttt{RF-COM-005} & \texttt{RF-066A} & Secciones 2.2, 2.3 y decisiones prioritarias del cliente sobre moderación.\\
\texttt{RF-PUB-001} & \texttt{RF-067} & Secciones 1.2, 2.2, 2.7 y fuente documental \texttt{TOP\_Racing\_Grok\_Clarified.tex} (líneas 146, 150, 154, 171 y 219).\\
\texttt{RF-PUB-002} & \texttt{RF-068} & Secciones 1.2, 2.2, 2.7 y fuente documental \texttt{TOP\_Racing\_Grok\_Clarified.tex} (líneas 113, 146, 150 y 171).\\
\texttt{RF-PUB-003} & \texttt{RF-069} & Secciones 1.2, 2.7 y fuente documental \texttt{TOP\_Racing\_Grok\_Clarified.tex} (líneas 150, 171, 210, 268 y 272).\\
\texttt{RF-PUB-004} & \texttt{RF-070} & Secciones 1.2, 2.2, 2.7 y fuente documental \texttt{TOP\_Racing\_Grok\_Clarified.tex} (líneas 113, 146, 171 y 219).\\
\texttt{RF-PUB-005} & \texttt{RF-067A} & RF-067, RF-088, RF-090 y fuentes prioritarias confirmadas por el usuario sobre visibilidad de inscripciones canceladas o descalificadas y de sus notas.\\
\texttt{RF-PUB-006} & \texttt{RF-067B} & RF-060, RF-061, RF-062A y fuente prioritaria confirmada por el usuario sobre visibilidad pública de protestas y detalle de votos.\\
\texttt{RF-PUB-007} & \texttt{RF-067C} & RF-067B y fuente prioritaria confirmada por el usuario sobre reserva pública de la identidad de quien presenta la protesta.\\
\texttt{RF-PUB-008} & \texttt{RF-067D} & RF-050G a RF-050J y fuentes prioritarias confirmadas por el usuario sobre visibilidad pública de disputas de entrega.\\
\texttt{RF-PUB-009} & \texttt{RF-067E} & RF-050B, RF-050J y fuentes prioritarias confirmadas por el usuario sobre visibilidad de la morosidad del dueño y del promotor.\\
\texttt{RF-PUB-010} & \texttt{RF-070A} & RF-050G a RF-050J, RF-067D, RF-067E y fuentes prioritarias confirmadas por el usuario sobre historial y trazabilidad temporal de la disputa de entrega.\\
\texttt{RF-IAM-001} & \texttt{RF-071} & Secciones 2.1, 2.3 y fuente documental \texttt{TOP\_Racing\_Grok\_Clarified.tex} (líneas 142 y 154).\\
\texttt{RF-IAM-002} & \texttt{RF-071A} & Secciones 2.3, 2.5, REX-006 y fuente prioritaria confirmada por el usuario sobre verificación previa de identidad.\\
\texttt{RF-IAM-003} & \texttt{RF-071B} & Secciones 2.2, 2.5, REX-007 y decisiones prioritarias del cliente sobre aceptación legal versionada.\\
\texttt{RF-IAM-004} & \texttt{RF-072} & Secciones 2.3, 2.5, 2.7, RIE-001 y fuente documental \texttt{TOP\_Racing\_Grok\_Clarified.tex} (líneas 142, 144, 146 y 154), además de decisiones prioritarias del cliente sobre verificación reforzada.\\
\texttt{RF-IAM-005} & \texttt{RF-072A} & RF-058A y fuente prioritaria confirmada por el usuario sobre visibilidad restringida de quien presenta la protesta.\\
\texttt{RF-IAM-006} & \texttt{RF-072B} & RF-050B, RF-050J y fuentes prioritarias confirmadas por el usuario sobre exclusión del moroso del sistema.\\
\texttt{RF-IAM-007} & \texttt{RF-072C} & Secciones 2.3, 2.5, 2.6, 3.1.24 y decisiones prioritarias del cliente sobre verificación reforzada.\\
\texttt{RF-IAM-008} & \texttt{RF-072D} & Secciones 2.3, 2.5, RIE-001, 3.1.25 y decisiones prioritarias del cliente sobre recordatorio al usuario.\\
\texttt{RF-IAM-009} & \texttt{RF-073} & Secciones 2.3 y 2.7, y fuente prioritaria confirmada por el usuario sobre papeles funcionales acumulables.\\
\texttt{RF-IAM-010} & \texttt{RF-074} & Secciones 2.3 y fuente prioritaria confirmada por el usuario sobre \texttt{admin} único y exclusivo.\\
\texttt{RF-EVT-041} & \texttt{RF-075} & Secciones 2.4, 2.6.1, RIE-002, fuente documental \texttt{TOP\_Racing\_Grok\_Clarified.tex} y decisiones prioritarias del cliente sobre nivel temporal y plazos de disputa de entrega.\\
\texttt{RF-EVT-042} & \texttt{RF-075A} & RF-075 y fuente prioritaria confirmada por el usuario sobre conservación histórica del número esperado de participantes.\\
\texttt{RF-EVT-043} & \texttt{RF-075B} & Secciones 2.3, 2.5, 2.6, 3.1.24 y decisiones prioritarias del cliente sobre promoción de eventos con verificación reforzada.\\
\texttt{RF-EVT-044} & \texttt{RF-076} & Secciones 2.5, 2.6.1 y decisiones prioritarias del cliente sobre convocatoria en \texttt{inscripciones}.\\
\texttt{RF-EVT-045} & \texttt{RF-077} & Secciones 2.6.1 y fuente documental \texttt{TOP\_Racing\_Grok\_Clarified.tex} (línea 126).\\
\texttt{RF-EVT-046} & \texttt{RF-078} & Secciones 2.6.1, fuente documental \texttt{TOP\_Racing\_Grok\_Clarified.tex} y decisiones prioritarias del cliente sobre determinación de la clasificación combinada por nivel territorial y temporal.\\
\texttt{RF-EVT-047} & \texttt{RF-079} & Secciones 2.6.1, RF-042 y fuente documental \texttt{TOP\_Racing\_Grok\_Clarified.tex} (línea 219).\\
\texttt{RF-EVT-048} & \texttt{RF-079A} & RF-075 y decisiones prioritarias del cliente sobre plazo máximo para iniciar disputas como dato del evento.\\
\texttt{RF-EVT-049} & \texttt{RF-079B} & RF-075 y decisiones prioritarias del cliente sobre plazo máximo para resolver disputas como dato del evento.\\
\texttt{RF-EVT-050} & \texttt{RF-079C} & RF-046 y decisiones prioritarias del cliente sobre entrega previa del vehículo al promotor.\\
\texttt{RF-EVT-051} & --- & \texttt{DEC-155}, decisión confirmada sobre carreras limitadas por tiempo o distancia.\\
\texttt{RF-VEH-001} & \texttt{RF-080} & Fuente prioritaria confirmada por el usuario sobre registro previo del vehículo.\\
\texttt{RF-VEH-002} & \texttt{RF-081} & Fuente documental \texttt{TOP\_Racing\_Grok\_Clarified.tex} (líneas 99, 146 y 221) y fuente prioritaria confirmada por el usuario sobre propiedad del vehículo.\\
\texttt{RF-VEH-003} & \texttt{RF-082} & Fuente prioritaria confirmada por el usuario sobre datos mínimos obligatorios del vehículo registrado.\\
\texttt{RF-VEH-004} & \texttt{RF-083} & Fuente prioritaria confirmada por el usuario sobre inscripción del vehículo por su dueño.\\
\texttt{RF-VEH-005} & \texttt{RF-083A} & Secciones 2.3, 2.5, 2.6, 3.1.24 y decisiones prioritarias del cliente sobre inscripción de vehículos con verificación reforzada.\\
\texttt{RF-VEH-006} & \texttt{RF-084} & Fuente prioritaria confirmada por el usuario sobre datos mínimos obligatorios de la inscripción.\\
\texttt{RF-VEH-007} & \texttt{RF-085} & Conversación fuente prioritaria consolidada y fuente prioritaria confirmada por el usuario sobre piloto distinto del dueño.\\
\texttt{RF-VEH-008} & \texttt{RF-086} & Fuente documental \texttt{TOP\_Racing\_Grok\_Clarified.tex} (líneas 133, 144, 204, 206 y 212) y fuente prioritaria confirmada por el usuario sobre datos mínimos obligatorios de la inscripción.\\
\texttt{RF-VEH-009} & \texttt{RF-086A} & Fuentes prioritarias confirmadas por el usuario sobre contribución mínima obligatoria en la inscripción del evento y sobre cobro real inmediato.\\
\texttt{RF-VEH-010} & \texttt{RF-086A} & Fuentes prioritarias confirmadas por el usuario sobre contribución mínima obligatoria en la inscripción del evento y sobre cobro real inmediato.\\
\texttt{RF-VEH-011} & \texttt{RF-087} & Fuentes prioritarias confirmadas por el usuario sobre cobro inmediato de la inscripción.\\
\texttt{RF-VEH-012} & \texttt{RF-088} & Fuentes prioritarias confirmadas por el usuario sobre cancelaciones y descalificaciones decididas por el promotor.\\
\texttt{RF-VEH-013} & \texttt{RF-089} & Fuente heredada RF-089, precisada por \texttt{DEC-155} sobre calificación mediante tiempo de vuelta válido.\\
\texttt{RF-VEH-014} & \texttt{RF-090} & Fuentes prioritarias confirmadas por el usuario sobre efectos de la cancelación de la inscripción, descalificación del competidor y visibilidad de la nota.\\
\texttt{RF-VEH-015} & \texttt{RF-090A} & Fuentes prioritarias confirmadas por el usuario sobre devoluciones inmediatas.\\
\texttt{RF-VEH-016} & \texttt{RF-090B} & Fuentes prioritarias confirmadas por el usuario sobre no entrega, modificación, descalificación y morosidad del dueño.\\
\texttt{RF-EPR-001} & \texttt{RF-091} & RF-037, RF-038, RF-040, RF-041, RIE-009 y decisiones prioritarias del cliente sobre publicidad del premio individual del evento.\\
\texttt{RF-EPR-002} & \texttt{RF-092} & RF-039, RF-040, RF-041, RIE-009 y decisiones prioritarias del cliente sobre consulta económica del evento.\\
\texttt{RF-EPR-003} & \texttt{RF-093} & RF-038, RF-039, RF-040, sección 2.7 y decisiones prioritarias del cliente sobre publicación del premio individual del evento.\\
\texttt{RF-EPR-004} & \texttt{RF-094} & RF-039, RF-040, sección 2.6.4 y decisiones prioritarias del cliente sobre publicación sin identidad de terceros en evento.\\
\texttt{RF-CPR-001} & \texttt{RF-095} & Secciones 2.2, 2.6.2 y 2.6.4.3, además de decisiones prioritarias del cliente sobre patrocinio directo y ventana de aportación al campeonato.\\
\texttt{RF-CPR-002} & \texttt{RF-096} & Secciones 2.6.2 y 2.6.4.3 y decisiones prioritarias del cliente sobre aportaciones a participantes en campeonatos y ventana de aportación al campeonato.\\
\texttt{RF-CPR-003} & \texttt{RF-097} & Secciones 2.6.2 y 2.6.4.3 y fuente prioritaria confirmada por el usuario sobre contribución mínima por defecto en premio de campeonato.\\
\texttt{RF-CPR-004} & \texttt{RF-098} & Secciones 2.6.2 y 2.6.4.3, RF-095, RF-096 y decisiones prioritarias del cliente.\\
\texttt{RF-CPR-005} & \texttt{RF-099} & Secciones 2.6.2 y 2.6.4.3 y decisiones prioritarias del cliente sobre finalidad de las aportaciones a participantes.\\
\texttt{RF-CPR-006} & \texttt{RF-100} & Sección 2.6.4.3 y fuente prioritaria confirmada por el usuario sobre patrocinio público con anuncio e imagen.\\
\texttt{RF-CPR-007} & \texttt{RF-101} & Secciones 2.6.2 y 2.6.4.3 y decisiones prioritarias del cliente sobre patrocinio y aportaciones a participantes.\\
\texttt{RF-CPR-008} & \texttt{RF-101} & Secciones 2.6.2 y 2.6.4.3 y decisiones prioritarias del cliente sobre patrocinio y aportaciones a participantes.\\
\texttt{RF-CPR-009} & \texttt{RF-102} & Secciones 2.1, 2.2, 2.6.4.3 y fuente prioritaria confirmada por el usuario sobre consulta en la página principal.\\
\texttt{RF-WGT-001} & \texttt{RF-103} & Secciones 2.2, 2.6.3.1 y 2.6.4.\\
\texttt{RF-WGT-002} & \texttt{RF-104} & Secciones 2.6.3, 2.6.3.1 y fuente documental TOP\_Racing\_Grok\_Clarified.tex (líneas 202 a 208).\\
\texttt{RF-WGT-003} & \texttt{RF-105} & Secciones 2.6.3.1 y fuente documental TOP\_Racing\_Grok\_Clarified.tex (líneas 171 y 210).\\
\texttt{RF-WGT-004} & \texttt{RF-106} & Secciones 2.6.3.1 y fuente documental TOP\_Racing\_Grok\_Clarified.tex (líneas 171 y 210).\\
\texttt{RF-WGT-005} & \texttt{RF-107} & Secciones 2.6.3.1 y fuente documental TOP\_Racing\_Grok\_Clarified.tex (líneas 171 y 210).\\
\texttt{RF-WGT-006} & \texttt{RF-108} & Secciones 2.6.3.1 y decisiones prioritarias del cliente confirmadas sobre secuencia y desempate, además de la fuente documental TOP\_Racing\_Grok\_Clarified.tex (líneas 171 y 210).\\
\texttt{RF-WGT-007} & \texttt{RF-108} & Secciones 2.6.3.1 y decisiones prioritarias del cliente confirmadas sobre secuencia y desempate, además de la fuente documental TOP\_Racing\_Grok\_Clarified.tex (líneas 171 y 210).\\
\texttt{RF-WGT-008} & \texttt{RF-108} & Secciones 2.6.3.1 y decisiones prioritarias del cliente confirmadas sobre secuencia y desempate, además de la fuente documental TOP\_Racing\_Grok\_Clarified.tex (líneas 171 y 210).\\
\texttt{RF-WGT-009} & \texttt{RF-108A} & Secciones 2.6.3.1 y fuente documental TOP\_Racing\_Grok\_Clarified.tex (líneas 171 y 210).\\
\texttt{RF-WGT-010} & \texttt{RF-108B} & Secciones 2.6.3.1 y fuente documental TOP\_Racing\_Grok\_Clarified.tex (línea 171).\\
\texttt{RF-WGT-011} & \texttt{RF-109} & Secciones 2.6.3, 2.6.3.1, RF-054 y RF-055.\\
\texttt{RF-EFF-001} & \texttt{RF-110} & Secciones 2.6.4, 2.6.4.1, RF-024 a RF-026 y fuente documental TOP\_Racing\_Grok\_Clarified.tex (líneas 234 y 250).\\
\texttt{RF-EFF-002} & \texttt{RF-111} & Secciones 2.6.4.1, RF-046 a RF-050 y RF-090.\\
\texttt{RF-EFF-003} & \texttt{RF-111} & Secciones 2.6.4.1, RF-046 a RF-050 y RF-090.\\
\texttt{RF-EFF-004} & \texttt{RF-111} & Secciones 2.6.4.1, RF-046 a RF-050 y RF-090.\\
\texttt{RF-EFF-005} & \texttt{RF-111} & Secciones 2.6.4.1, RF-046 a RF-050 y RF-090.\\
\texttt{RF-EFF-006} & \texttt{RF-112} & Secciones 2.6.4, 2.6.4.1, fuente documental TOP\_Racing\_Grok\_Clarified.tex (línea 238) y decisiones prioritarias del cliente sobre ausencia de redondeo.\\
\texttt{RF-EFF-007} & \texttt{RF-112} & Secciones 2.6.4, 2.6.4.1, fuente documental TOP\_Racing\_Grok\_Clarified.tex (línea 238) y decisiones prioritarias del cliente sobre ausencia de redondeo.\\
\texttt{RF-EFF-008} & \texttt{RF-112} & Secciones 2.6.4, 2.6.4.1, fuente documental TOP\_Racing\_Grok\_Clarified.tex (línea 238) y decisiones prioritarias del cliente sobre ausencia de redondeo.\\
\texttt{RF-EFF-009} & \texttt{RF-112} & Secciones 2.6.4, 2.6.4.1, fuente documental TOP\_Racing\_Grok\_Clarified.tex (línea 238) y decisiones prioritarias del cliente sobre ausencia de redondeo.\\
\texttt{RF-EFF-010} & \texttt{RF-113} & Secciones 2.6.4, 2.6.4.1, RF-110, RF-112 y fuente documental TOP\_Racing\_Grok\_Clarified.tex (líneas 240 y 242).\\
\texttt{RF-EFF-011} & \texttt{RF-113A} & Secciones 2.6.4.1 y fuente documental TOP\_Racing\_Grok\_Clarified.tex (líneas 242 a 247), además de la confirmación prioritaria del usuario y de las decisiones prioritarias del cliente sobre desempate.\\
\texttt{RF-EFF-012} & \texttt{RF-113A} & Secciones 2.6.4.1 y fuente documental TOP\_Racing\_Grok\_Clarified.tex (líneas 242 a 247), además de la confirmación prioritaria del usuario y de las decisiones prioritarias del cliente sobre desempate.\\
\texttt{RF-EFF-013} & \texttt{RF-113A} & Secciones 2.6.4.1 y fuente documental TOP\_Racing\_Grok\_Clarified.tex (líneas 242 a 247), además de la confirmación prioritaria del usuario y de las decisiones prioritarias del cliente sobre desempate.\\
\texttt{RF-EFF-014} & \texttt{RF-114} & Secciones 2.6.4.1, RF-051 a RF-055, RF-112, RF-113 y RF-113A.\\
\texttt{RF-EFF-015} & \texttt{RF-114} & Secciones 2.6.4.1, RF-051 a RF-055, RF-112, RF-113 y RF-113A.\\
\texttt{RF-EFF-016} & \texttt{RF-114} & Secciones 2.6.4.1, RF-051 a RF-055, RF-112, RF-113 y RF-113A.\\
\texttt{RF-EFF-017} & \texttt{RF-114} & Secciones 2.6.4.1, RF-051 a RF-055, RF-112, RF-113 y RF-113A.\\
\texttt{RF-EFF-018} & \texttt{RF-114} & Secciones 2.6.4.1, RF-051 a RF-055, RF-112, RF-113 y RF-113A.\\
\texttt{RF-EPR-005} & \texttt{RF-115} & Secciones 2.6.2 y 2.6.4.2 y decisiones prioritarias del cliente sobre patrocinio directo de eventos.\\
\texttt{RF-EPR-006} & \texttt{RF-116} & Secciones 2.6.2 y 2.6.4.2, RF-086A y decisiones prioritarias del cliente sobre aportaciones de eficiencia a participantes.\\
\texttt{RF-EPR-007} & \texttt{RF-116A} & Secciones 2.6.2, 2.6.4.2, 3.1.13 y decisiones prioritarias del cliente sobre separación entre contribuciones de eficiencia y de terminación.\\
\texttt{RF-EPR-008} & \texttt{RF-117} & Secciones 2.6.4.2 y fuente documental \texttt{TOP\_Racing\_Grok\_Clarified.tex}, además de la confirmación prioritaria del usuario.\\
\texttt{RF-EPR-009} & \texttt{RF-118} & Secciones 2.6.4.2, RF-116, RF-117 y fuente documental \texttt{TOP\_Racing\_Grok\_Clarified.tex}.\\
\texttt{RF-EPR-010} & \texttt{RF-118A} & Secciones 2.6.4.2, RF-118 y fuente documental \texttt{TOP\_Racing\_Grok\_Clarified.tex}.\\
\texttt{RF-EPR-011} & \texttt{RF-118B} & Secciones 2.6.4.2, RF-115, RF-118, RF-118A, fuente documental \texttt{TOP\_Racing\_Grok\_Clarified.tex} y decisiones prioritarias del cliente sobre ausencia de redondeo.\\
\texttt{RF-EPR-012} & \texttt{RF-118C} & Secciones 2.6.2, 2.6.4.2, 3.1.12, 3.1.13, 3.1.17 y decisiones prioritarias del cliente confirmadas sobre separación entre contribuciones de eficiencia y de terminación.\\
\texttt{RF-EPR-013} & \texttt{RF-118C} & Secciones 2.6.2, 2.6.4.2, 3.1.12, 3.1.13, 3.1.17 y decisiones prioritarias del cliente confirmadas sobre separación entre contribuciones de eficiencia y de terminación.\\
\texttt{RF-EPR-014} & \texttt{RF-118C} & Secciones 2.6.2, 2.6.4.2, 3.1.12, 3.1.13, 3.1.17 y decisiones prioritarias del cliente confirmadas sobre separación entre contribuciones de eficiencia y de terminación.\\
\texttt{RF-EPR-015} & \texttt{RF-118C} & Secciones 2.6.2, 2.6.4.2, 3.1.12, 3.1.13, 3.1.17 y decisiones prioritarias del cliente confirmadas sobre separación entre contribuciones de eficiencia y de terminación.\\
\texttt{RF-EPR-016} & \texttt{RF-118C} & Secciones 2.6.2, 2.6.4.2, 3.1.12, 3.1.13, 3.1.17 y decisiones prioritarias del cliente confirmadas sobre separación entre contribuciones de eficiencia y de terminación.\\
\texttt{RF-EPR-017} & \texttt{RF-118C} & Secciones 2.6.2, 2.6.4.2, 3.1.12, 3.1.13, 3.1.17 y decisiones prioritarias del cliente confirmadas sobre separación entre contribuciones de eficiencia y de terminación.\\
\texttt{RF-EPR-018} & \texttt{RF-119} & Secciones 2.6.4.2, RF-116, RF-118B y decisiones prioritarias del cliente sobre aportantes asociados, reparto del premio del participante y ausencia de redondeo.\\
\texttt{RF-EPR-019} & \texttt{RF-119A} & Secciones 2.6.4.2, RF-116A, RF-118C y decisiones prioritarias del cliente sobre separación entre contribuciones de eficiencia y de terminación.\\
\texttt{RF-EPR-020} & \texttt{RF-120} & Secciones 2.6.4.2, RF-041, RF-091 a RF-094 y RF-115 a RF-119A.\\
\texttt{RF-EPR-021} & \texttt{RF-120A} & Secciones 2.6.4.2 y decisiones prioritarias del cliente sobre pago automático a todos.\\
\texttt{RF-EPR-022} & -- & \texttt{DEC-170}; reserva del pozo consolidado y de las aportaciones durante la subasta de vehículos.\\
\texttt{RF-CPR-010} & \texttt{RF-121} & Secciones 2.6.2, 2.6.4.3, RF-095 y RF-098.\\
\texttt{RF-CPR-011} & \texttt{RF-122} & Secciones 2.6.4.3, RF-096, RF-097 y RF-099.\\
\texttt{RF-CPR-012} & \texttt{RF-123} & Secciones 2.6.4.3 y decisiones prioritarias del cliente sobre la lógica del premio de campeonato y su desempate.\\
\texttt{RF-CPR-013} & \texttt{RF-123} & Secciones 2.6.4.3 y decisiones prioritarias del cliente sobre la lógica del premio de campeonato y su desempate.\\
\texttt{RF-CPR-014} & \texttt{RF-123} & Secciones 2.6.4.3 y decisiones prioritarias del cliente sobre la lógica del premio de campeonato y su desempate.\\
\texttt{RF-CPR-015} & \texttt{RF-123A} & Secciones 2.6.4.3, RF-122, RF-123 y decisiones prioritarias del cliente sobre la lógica del premio de campeonato.\\
\texttt{RF-CPR-016} & \texttt{RF-123B} & Secciones 2.6.4.3, RF-123A y decisiones prioritarias del cliente sobre la lógica del premio de campeonato.\\
\texttt{RF-CPR-017} & \texttt{RF-123C} & Secciones 2.6.4.3, RF-121, RF-123A, RF-123B y decisiones prioritarias del cliente sobre la lógica del premio de campeonato y ausencia de redondeo.\\
\texttt{RF-CPR-018} & \texttt{RF-124} & Secciones 2.6.4.3, RF-122, RF-123C y decisiones prioritarias del cliente sobre aportaciones a participantes, reparto del premio del participante y ausencia de redondeo.\\
\texttt{RF-CPR-019} & \texttt{RF-125} & Secciones 2.6.4.3, RF-101, RF-102 y RF-121 a RF-124.\\
\texttt{RF-CPR-020} & \texttt{RF-125A} & Secciones 2.6.4.3 y decisiones prioritarias del cliente sobre pago automático a todos.\\
\texttt{RF-SPN-001} & \texttt{RF-126} & Secciones 2.2, 2.6.2, 2.6.4.2 y decisiones prioritarias del cliente sobre patrocinio directo de eventos.\\
\texttt{RF-SPN-002} & \texttt{RF-127} & Secciones 2.6.2, 2.6.4.2 y decisiones prioritarias del cliente sobre aportaciones de eficiencia a participantes del evento.\\
\texttt{RF-SPN-003} & \texttt{RF-127A} & Secciones 2.6.2, 2.6.4.2 y decisiones prioritarias del cliente sobre aportaciones de terminación a participantes del evento.\\
\texttt{RF-SPN-004} & \texttt{RF-128} & Sección 2.6.4.2 y decisiones prioritarias del cliente sobre publicidad asociada al patrocinio directo.\\
\texttt{RF-SPN-005} & \texttt{RF-129} & Secciones 2.6.2, 2.6.4.2, RF-126, RF-127, RF-127A y RF-128.\\
\texttt{RF-SPN-006} & \texttt{RF-129} & Secciones 2.6.2, 2.6.4.2, RF-126, RF-127, RF-127A y RF-128.\\
\texttt{RF-SPN-007} & \texttt{RF-129} & Secciones 2.6.2, 2.6.4.2, RF-126, RF-127, RF-127A y RF-128.\\
\texttt{RF-SPN-008} & \texttt{RF-129A} & Secciones 2.6.4 y decisiones prioritarias del cliente sobre devolución inmediata en cancelación del evento.\\
\texttt{RF-SPN-009} & \texttt{RF-129B} & Secciones 2.6.4 y decisiones prioritarias del cliente sobre devolución inmediata al cancelarse o descalificarse un participante antes de \texttt{carrera}.\\
\texttt{RF-PTS-001} & \texttt{RF-130} & Secciones 2.6.3.1, 3.1.7, 3.1.16, fuente documental TOP\_Racing\_Grok\_Clarified.tex (líneas 167 a 171) y decisiones prioritarias del cliente sobre desempate.\\
\texttt{RF-PTS-002} & \texttt{RF-130} & Secciones 2.6.3.1, 3.1.7, 3.1.16, fuente documental TOP\_Racing\_Grok\_Clarified.tex (líneas 167 a 171) y decisiones prioritarias del cliente sobre desempate.\\
\texttt{RF-PTS-003} & \texttt{RF-131} & Secciones 2.6.3.1, 3.1.7, 3.1.16, fuente documental TOP\_Racing\_Grok\_Clarified.tex (líneas 167 a 171) y decisiones prioritarias del cliente sobre desempate.\\
\texttt{RF-PTS-004} & \texttt{RF-131} & Secciones 2.6.3.1, 3.1.7, 3.1.16, fuente documental TOP\_Racing\_Grok\_Clarified.tex (líneas 167 a 171) y decisiones prioritarias del cliente sobre desempate.\\
\texttt{RF-PTS-005} & \texttt{RF-132} & Secciones 2.6.3.1, 3.1.7, 3.1.17, fuente documental TOP\_Racing\_Grok\_Clarified.tex (líneas 167 a 171 y 242) y decisiones prioritarias del cliente sobre desempate.\\
\texttt{RF-PTS-006} & \texttt{RF-132} & Secciones 2.6.3.1, 3.1.7, 3.1.17, fuente documental TOP\_Racing\_Grok\_Clarified.tex (líneas 167 a 171 y 242) y decisiones prioritarias del cliente sobre desempate.\\
\texttt{RF-PTS-007} & \texttt{RF-133} & Secciones 2.6.3.1, RF-108B, RF-130, RF-131, RF-132 y decisiones prioritarias del cliente sobre ausencia de redondeo.\\
\texttt{RF-PTS-008} & \texttt{RF-134} & Secciones 2.6.3.1, 3.1.7, RF-133 y fuente documental TOP\_Racing\_Grok\_Clarified.tex (líneas 150, 167 y 171).\\
\texttt{RF-PTS-009} & \texttt{RF-135} & Secciones 2.6.3.1, RF-133, fuente documental TOP\_Racing\_Grok\_Clarified.tex (líneas 167 a 171) y decisiones prioritarias del cliente sobre clasificación combinada.\\
\texttt{RF-PTS-010} & \texttt{RF-136} & RF-055A, RF-090, RF-130 a RF-135 y decisiones prioritarias del cliente sobre efectos de la cancelación.\\
\texttt{RF-PTS-011} & \texttt{RF-137} & Secciones 2.6.3.1, 3.1.7, RF-108B y RF-130 a RF-136.\\
\texttt{RF-CPR-021} & \texttt{RF-138} & RF-123C, RF-125, RIE-010 y decisiones prioritarias del cliente sobre publicidad del premio individual del campeonato.\\
\texttt{RF-CPR-022} & \texttt{RF-139} & RF-124, RF-125, RIE-010 y decisiones prioritarias del cliente sobre consulta económica de campeonato.\\
\texttt{RF-CPR-023} & \texttt{RF-140} & RF-124, RF-125 y decisiones prioritarias del cliente sobre publicación sin identidad de terceros en premio de campeonato.\\
\texttt{RF-CPR-024} & \texttt{RF-141} & RF-123C, RF-124, RF-125 y sección 2.7.\\
\texttt{RF-GOV-001} & \texttt{RF-142} & Secciones 2.2, 2.5, 2.7, 3.4 y decisiones prioritarias del cliente sobre publicidad del dominio.\\
\texttt{RF-GOV-002} & \texttt{RF-143} & Secciones 2.5, 2.7, 3.4, RF-058A, RF-067C, RF-153 a RF-157 y decisiones prioritarias del cliente sobre reserva de identidad.\\
\texttt{RF-GOV-003} & \texttt{RF-144} & Secciones 2.6.4.2, 2.6.4.3, 2.7, RIE-009, RIE-010 y decisiones prioritarias del cliente sobre publicación de premios.\\
\texttt{RF-GOV-004} & \texttt{RF-145} & Secciones 2.2, 2.3, 2.7, RF-063 a RF-066A y decisiones prioritarias del cliente sobre moderación.\\
\texttt{RF-GOV-005} & \texttt{RF-146} & Secciones 2.2, 2.7, RF-067 a RF-070A, RF-101, RF-102, RF-129 y decisiones prioritarias del cliente sobre permanencia histórica.\\
\texttt{RF-IAM-011} & \texttt{RF-147} & Secciones 1.3, 2.3, 2.5 y decisiones prioritarias del cliente sobre estados booleanos del usuario.\\
\texttt{RF-IAM-012} & \texttt{RF-147A} & Secciones 1.3, 2.3, 2.5 y decisiones prioritarias del cliente sobre verificación reforzada.\\
\texttt{RF-IAM-013} & \texttt{RF-148} & Secciones 1.3, 2.3, 2.5, RF-072B y decisiones prioritarias del cliente sobre estados booleanos del usuario.\\
\texttt{RF-IAM-014} & \texttt{RF-149} & Secciones 2.3, 2.5, RF-074 y decisiones prioritarias del cliente sobre facultades de \texttt{admin}.\\
\texttt{RF-IAM-015} & \texttt{RF-149A} & Secciones 2.3, 2.5 y decisiones prioritarias del cliente sobre facultades de \texttt{admin} respecto de la verificación reforzada.\\
\texttt{RF-IAM-016} & \texttt{RF-150} & Secciones 2.3, 2.5, 2.6 y decisiones prioritarias del cliente sobre alcance del promotor.\\
\texttt{RF-IAM-017} & \texttt{RF-150A} & Secciones 2.3, 2.5 y decisiones prioritarias del cliente sobre gestión de la verificación reforzada.\\
\texttt{RF-IAM-018} & \texttt{RF-151} & Secciones 2.5, 2.6, RF-071A, RF-071B, RF-072B, REX-006, REX-007 y decisiones prioritarias del cliente sobre efecto de exclusión.\\
\texttt{RF-IAM-019} & \texttt{RF-151A} & Secciones 2.3, 2.5, 2.6 y decisiones prioritarias del cliente sobre operaciones sujetas a verificación reforzada.\\
\texttt{RF-IAM-020} & \texttt{RF-151B} & Secciones 2.5, RF-148, RF-151, 3.1.25, RIE-001 y decisión prioritaria del cliente confirmada sobre cancelación de sesión y rechazo de inicio de sesión por exclusión activa.\\
\texttt{RF-IAM-021} & \texttt{RF-151B} & Secciones 2.5, RF-148, RF-151, 3.1.25, RIE-001 y decisión prioritaria del cliente confirmada sobre cancelación de sesión y rechazo de inicio de sesión por exclusión activa.\\
\texttt{RF-IAM-022} & \texttt{RF-151B} & Secciones 2.5, RF-148, RF-151, 3.1.25, RIE-001 y decisión prioritaria del cliente confirmada sobre cancelación de sesión y rechazo de inicio de sesión por exclusión activa.\\
\texttt{RF-IAM-023} & \texttt{RF-152} & Secciones 2.7, RNF-002 y decisiones prioritarias del cliente sobre gestión de esos estados.\\
\texttt{RF-IAM-024} & \texttt{RF-152A} & Decisiones prioritarias del cliente sobre exclusión de nuevas identidades detectadas.\\
\texttt{RF-IAM-025} & \texttt{RF-153} & Secciones 1.3, 2.3, 2.5, 3.1.24 y decisiones prioritarias del cliente sobre consulta personal de estados.\\
\texttt{RF-IAM-026} & \texttt{RF-153A} & Secciones 1.3, 2.3, 2.5, 3.1.24 y decisiones prioritarias del cliente sobre nivel actual visible.\\
\texttt{RF-IAM-027} & \texttt{RF-153A} & Secciones 1.3, 2.3, 2.5, 3.1.24 y decisiones prioritarias del cliente sobre nivel actual visible.\\
\texttt{RF-IAM-028} & \texttt{RF-153B} & Secciones 2.3, 2.5, 3.1.11, 3.1.24 y decisiones prioritarias del cliente sobre recordatorio ante operaciones bloqueadas.\\
\texttt{RF-IAM-029} & \texttt{RF-154} & Secciones 2.3, 2.5, 3.1.24, RF-149 y RF-149A.\\
\texttt{RF-IAM-030} & \texttt{RF-154A} & Secciones 2.3, 2.5, RF-150 y decisiones prioritarias del cliente sobre consulta del promotor.\\
\texttt{RF-IAM-031} & \texttt{RF-155} & Secciones 2.5, 3.4, RF-143 y decisiones prioritarias del cliente sobre reserva de estados.\\
\texttt{RF-IAM-032} & \texttt{RF-156} & Secciones 1.3, 2.5, RF-067, RF-143 y decisiones prioritarias del cliente sobre consulta pública del evento.\\
\texttt{RF-IAM-033} & \texttt{RF-157} & Secciones 2.3, 2.5, RF-148, RF-154 y decisiones prioritarias del cliente sobre usuarios no verificados y usuarios excluidos.\\
\texttt{RF-REP-001} & \texttt{RF-158} & Secciones 1.2, 2.2 y fuente documental \texttt{TOP\_Racing\_Grok\_Clarified.tex} (línea 133).\\
\texttt{RF-REP-002} & \texttt{RF-159} & Secciones 1.2, 2.2 y fuente documental \texttt{TOP\_Racing\_Grok\_Clarified.tex} (línea 133).\\
\texttt{RF-REP-003} & \texttt{RF-160} & Secciones 2.2 y fuente documental \texttt{TOP\_Racing\_Grok\_Clarified.tex} (línea 133).\\
\texttt{RF-REP-004} & \texttt{RF-161} & Secciones 1.2, 2.2 y fuente documental \texttt{TOP\_Racing\_Grok\_Clarified.tex} (línea 133).\\
\texttt{RF-BUS-001} & \texttt{RF-162} & Secciones 2.2, 2.6.2.2 y decisiones prioritarias del cliente sobre configuración administrativa del modelo de negocio.\\
\texttt{RF-BUS-002} & \texttt{RF-163} & Secciones 2.6.2.2, 2.6.4.2 y decisiones prioritarias del cliente sobre comisión del competidor por evento.\\
\texttt{RF-BUS-003} & \texttt{RF-164} & Secciones 2.6.2.2, 2.6.4.3 y decisiones prioritarias del cliente sobre comisión del competidor por campeonato.\\
\texttt{RF-BUS-004} & \texttt{RF-165} & Secciones 2.6.2.2, 2.6.4.2, 2.6.4.3 y decisiones prioritarias del cliente sobre comisión del aportante.\\
\texttt{RF-BUS-005} & \texttt{RF-166} & Secciones 2.6.2.2, 3.1.6 y decisiones prioritarias del cliente sobre comisión a cargo del comprador o adjudicatario y vendedor sin comisión por vender.\\
\texttt{RF-BUS-006} & \texttt{RF-167} & Secciones 2.6.2.2, 3.1.15, 3.1.20 y decisiones prioritarias del cliente sobre patrocinadores sin comisión.\\
\texttt{RF-BUS-007} & \texttt{RF-168} & Secciones 2.6.2.2 y decisiones prioritarias del cliente sobre cuota fija del promotor.\\
\texttt{RF-BUS-008} & \texttt{RF-169} & Secciones 2.6.2.2, 2.6.4 y decisiones prioritarias del cliente sobre pagos reales, comisiones pequeñas y cuota del promotor.\\
\texttt{RF-BUS-009} & \texttt{RF-170} & Secciones 2.6.2.2, 3.2 y decisiones prioritarias del cliente sobre publicación de comisiones.\\
\texttt{RF-BUS-010} & \texttt{RF-171} & Secciones 2.6.2.2, 2.6.4 y decisiones prioritarias del cliente sobre detalle posterior de comisiones.\\
\texttt{RF-BUS-011} & \texttt{RF-172} & Secciones 2.6.2.2 y decisiones prioritarias del cliente sobre vendedor y patrocinador sin comisión.\\
\texttt{UI-001} & \texttt{RIE-001} & Secciones 1.2, 2.1, 2.2, 2.3, 2.4, 3.1.11, 3.1.24, 3.1.25 y decisiones prioritarias del cliente sobre página principal, niveles de verificación y aceptación legal previa y versionada.\\
\texttt{UI-002} & \texttt{RIE-001} & Secciones 1.2, 2.1, 2.2, 2.3, 2.4, 3.1.11, 3.1.24, 3.1.25 y decisiones prioritarias del cliente sobre página principal, niveles de verificación y aceptación legal previa y versionada.\\
\texttt{UI-003} & \texttt{RIE-001} & Secciones 1.2, 2.1, 2.2, 2.3, 2.4, 3.1.11, 3.1.24, 3.1.25 y decisiones prioritarias del cliente sobre página principal, niveles de verificación y aceptación legal previa y versionada.\\
\texttt{UI-004} & \texttt{RIE-001} & Secciones 1.2, 2.1, 2.2, 2.3, 2.4, 3.1.11, 3.1.24, 3.1.25 y decisiones prioritarias del cliente sobre página principal, niveles de verificación y aceptación legal previa y versionada.\\
\texttt{UI-005} & \texttt{RIE-001} & Secciones 1.2, 2.1, 2.2, 2.3, 2.4, 3.1.11, 3.1.24, 3.1.25 y decisiones prioritarias del cliente sobre página principal, niveles de verificación y aceptación legal previa y versionada.\\
\texttt{UI-006} & \texttt{RIE-001} & Secciones 1.2, 2.1, 2.2, 2.3, 2.4, 3.1.11, 3.1.24, 3.1.25 y decisiones prioritarias del cliente sobre página principal, niveles de verificación y aceptación legal previa y versionada.\\
\texttt{UI-007} & \texttt{RIE-002} & Secciones 1.2, 2.2, 2.3, 2.6.1, fuente documental \texttt{TOP\_Racing\_Grok\_Clarified.tex} (línea 144) y decisiones prioritarias del cliente sobre nivel temporal, convocatoria en \texttt{inscripciones} y plazos de disputa de entrega.\\
\texttt{UI-008} & \texttt{RIE-003} & Secciones 1.2, 2.2, 2.3, 2.6.2, 3.1.13 y fuente documental \texttt{TOP\_Racing\_Grok\_Clarified.tex} (línea 144), además de decisiones prioritarias del cliente sobre cobro inmediato de la inscripción.\\
\texttt{UI-009} & \texttt{RIE-003} & Secciones 1.2, 2.2, 2.3, 2.6.2, 3.1.13 y fuente documental \texttt{TOP\_Racing\_Grok\_Clarified.tex} (línea 144), además de decisiones prioritarias del cliente sobre cobro inmediato de la inscripción.\\
\texttt{UI-010} & \texttt{RIE-004} & Secciones 1.2, 2.2, 2.3, 3.1.6 y fuente documental \texttt{TOP\_Racing\_Grok\_Clarified.tex} (líneas 146 y 219), además de decisiones prioritarias del cliente sobre QR, actualización al alza y pagos reales con retención.\\
\texttt{UI-011} & \texttt{RIE-004} & Secciones 1.2, 2.2, 2.3, 3.1.6 y fuente documental \texttt{TOP\_Racing\_Grok\_Clarified.tex} (líneas 146 y 219), además de decisiones prioritarias del cliente sobre QR, actualización al alza y pagos reales con retención.\\
\texttt{UI-012} & \texttt{RIE-004} & Secciones 1.2, 2.2, 2.3, 3.1.6 y fuente documental \texttt{TOP\_Racing\_Grok\_Clarified.tex} (líneas 146 y 219), además de decisiones prioritarias del cliente sobre QR, actualización al alza y pagos reales con retención.\\
\texttt{UI-013} & \texttt{RIE-004} & Secciones 1.2, 2.2, 2.3, 3.1.6 y fuente documental \texttt{TOP\_Racing\_Grok\_Clarified.tex} (líneas 146 y 219), además de decisiones prioritarias del cliente sobre QR, actualización al alza y pagos reales con retención.\\
\texttt{UI-014} & \texttt{RIE-004} & Secciones 1.2, 2.2, 2.3, 3.1.6 y fuente documental \texttt{TOP\_Racing\_Grok\_Clarified.tex} (líneas 146 y 219), además de decisiones prioritarias del cliente sobre QR, actualización al alza y pagos reales con retención.\\
\texttt{UI-015} & \texttt{RIE-004} & Secciones 1.2, 2.2, 2.3, 3.1.6 y fuente documental \texttt{TOP\_Racing\_Grok\_Clarified.tex} (líneas 146 y 219), además de decisiones prioritarias del cliente sobre QR, actualización al alza y pagos reales con retención.\\
\texttt{UI-016} & \texttt{RIE-004} & Secciones 1.2, 2.2, 2.3, 3.1.6 y fuente documental \texttt{TOP\_Racing\_Grok\_Clarified.tex} (líneas 146 y 219), además de decisiones prioritarias del cliente sobre QR, actualización al alza y pagos reales con retención.\\
\texttt{UI-017} & \texttt{RIE-004} & Secciones 1.2, 2.2, 2.3, 3.1.6 y fuente documental \texttt{TOP\_Racing\_Grok\_Clarified.tex} (líneas 146 y 219), además de decisiones prioritarias del cliente sobre QR, actualización al alza y pagos reales con retención.\\
\texttt{UI-018} & \texttt{RIE-004} & Secciones 1.2, 2.2, 2.3, 3.1.6 y fuente documental \texttt{TOP\_Racing\_Grok\_Clarified.tex} (líneas 146 y 219), además de decisiones prioritarias del cliente sobre QR, actualización al alza y pagos reales con retención.\\
\texttt{UI-019} & \texttt{RIE-005} & Secciones 1.2, 2.2, 2.3, 2.4, 3.1.6 y fuente documental \texttt{TOP\_Racing\_Grok\_Clarified.tex} (líneas 146, 221 y 223), además de decisiones prioritarias del cliente sobre garantía real de pago, custodia del promotor y transferencia obligatoria del vehículo subastado.\\
\texttt{UI-020} & \texttt{RIE-005} & Secciones 1.2, 2.2, 2.3, 2.4, 3.1.6 y fuente documental \texttt{TOP\_Racing\_Grok\_Clarified.tex} (líneas 146, 221 y 223), además de decisiones prioritarias del cliente sobre garantía real de pago, custodia del promotor y transferencia obligatoria del vehículo subastado.\\
\texttt{UI-021} & \texttt{RIE-005} & Secciones 1.2, 2.2, 2.3, 2.4, 3.1.6 y fuente documental \texttt{TOP\_Racing\_Grok\_Clarified.tex} (líneas 146, 221 y 223), además de decisiones prioritarias del cliente sobre garantía real de pago, custodia del promotor y transferencia obligatoria del vehículo subastado.\\
\texttt{UI-022} & \texttt{RIE-005} & Secciones 1.2, 2.2, 2.3, 2.4, 3.1.6 y fuente documental \texttt{TOP\_Racing\_Grok\_Clarified.tex} (líneas 146, 221 y 223), además de decisiones prioritarias del cliente sobre garantía real de pago, custodia del promotor y transferencia obligatoria del vehículo subastado.\\
\texttt{UI-023} & \texttt{RIE-005} & Secciones 1.2, 2.2, 2.3, 2.4, 3.1.6 y fuente documental \texttt{TOP\_Racing\_Grok\_Clarified.tex} (líneas 146, 221 y 223), además de decisiones prioritarias del cliente sobre garantía real de pago, custodia del promotor y transferencia obligatoria del vehículo subastado.\\
\texttt{UI-024} & \texttt{RIE-005} & Secciones 1.2, 2.2, 2.3, 2.4, 3.1.6 y fuente documental \texttt{TOP\_Racing\_Grok\_Clarified.tex} (líneas 146, 221 y 223), además de decisiones prioritarias del cliente sobre garantía real de pago, custodia del promotor y transferencia obligatoria del vehículo subastado.\\
\texttt{UI-025} & \texttt{RIE-005} & Secciones 1.2, 2.2, 2.3, 2.4, 3.1.6 y fuente documental \texttt{TOP\_Racing\_Grok\_Clarified.tex} (líneas 146, 221 y 223), además de decisiones prioritarias del cliente sobre garantía real de pago, custodia del promotor y transferencia obligatoria del vehículo subastado.\\
\texttt{UI-026} & \texttt{RIE-005} & Secciones 1.2, 2.2, 2.3, 2.4, 3.1.6 y fuente documental \texttt{TOP\_Racing\_Grok\_Clarified.tex} (líneas 146, 221 y 223), además de decisiones prioritarias del cliente sobre garantía real de pago, custodia del promotor y transferencia obligatoria del vehículo subastado.\\
\texttt{UI-027} & \texttt{RIE-005} & Secciones 1.2, 2.2, 2.3, 2.4, 3.1.6 y fuente documental \texttt{TOP\_Racing\_Grok\_Clarified.tex} (líneas 146, 221 y 223), además de decisiones prioritarias del cliente sobre garantía real de pago, custodia del promotor y transferencia obligatoria del vehículo subastado.\\
\texttt{UI-028} & \texttt{RIE-006} & Secciones 1.2, 2.2, 2.3, 2.4, 3.1.16, 3.1.21, fuente documental TOP\_Racing\_Grok\_Clarified.tex (líneas 150 y 171) y análisis funcional confirmado de la app legacy sobre la consulta de standings por selección de nivel temporal y territorial.\\
\texttt{UI-029} & \texttt{RIE-006} & Secciones 1.2, 2.2, 2.3, 2.4, 3.1.16, 3.1.21, fuente documental TOP\_Racing\_Grok\_Clarified.tex (líneas 150 y 171) y análisis funcional confirmado de la app legacy sobre la consulta de standings por selección de nivel temporal y territorial.\\
\texttt{UI-030} & \texttt{RIE-006} & Secciones 1.2, 2.2, 2.3, 2.4, 3.1.16, 3.1.21, fuente documental TOP\_Racing\_Grok\_Clarified.tex (líneas 150 y 171) y análisis funcional confirmado de la app legacy sobre la consulta de standings por selección de nivel temporal y territorial.\\
\texttt{UI-031} & \texttt{RIE-006} & Secciones 1.2, 2.2, 2.3, 2.4, 3.1.16, 3.1.21, fuente documental TOP\_Racing\_Grok\_Clarified.tex (líneas 150 y 171) y análisis funcional confirmado de la app legacy sobre la consulta de standings por selección de nivel temporal y territorial.\\
\texttt{UI-032} & \texttt{RIE-007} & Secciones 2.2, 2.3 y fuente documental \texttt{TOP\_Racing\_Grok\_Clarified.tex} (línea 152).\\
\texttt{UI-033} & \texttt{RIE-007} & Secciones 2.2, 2.3 y fuente documental \texttt{TOP\_Racing\_Grok\_Clarified.tex} (línea 152).\\
\texttt{UI-034} & \texttt{RIE-007} & Secciones 2.2, 2.3 y fuente documental \texttt{TOP\_Racing\_Grok\_Clarified.tex} (línea 152).\\
\texttt{UI-035} & \texttt{RIE-007} & Secciones 2.2, 2.3 y fuente documental \texttt{TOP\_Racing\_Grok\_Clarified.tex} (línea 152).\\
\texttt{UI-036} & \texttt{RIE-008} & Secciones 1.2, 2.1, 2.2, 2.3, fuente documental \texttt{TOP\_Racing\_Grok\_Clarified.tex} (línea 154) y decisiones prioritarias del cliente sobre publicación por cualquier usuario con identidad verificada y moderación por \texttt{admin}.\\
\texttt{UI-037} & \texttt{RIE-009} & Secciones 1.2, 2.2, 2.3, 2.6.3, 2.6.4 y fuente documental \texttt{TOP\_Racing\_Grok\_Clarified.tex} (línea 148).\\
\texttt{UI-038} & \texttt{RIE-010} & Secciones 2.2, 2.6.4.3 y decisiones prioritarias del cliente sobre premio de campeonato, patrocinio, aportaciones a participantes y cobro inmediato.\\
\texttt{UI-039} & \texttt{RIE-010} & Secciones 2.2, 2.6.4.3 y decisiones prioritarias del cliente sobre premio de campeonato, patrocinio, aportaciones a participantes y cobro inmediato.\\
\texttt{UI-040} & \texttt{RIE-010} & Secciones 2.2, 2.6.4.3 y decisiones prioritarias del cliente sobre premio de campeonato, patrocinio, aportaciones a participantes y cobro inmediato.\\
\texttt{UI-041} & \texttt{RIE-010} & Secciones 2.2, 2.6.4.3 y decisiones prioritarias del cliente sobre premio de campeonato, patrocinio, aportaciones a participantes y cobro inmediato.\\
\texttt{UI-042} & \texttt{RIE-010} & Secciones 2.2, 2.6.4.3 y decisiones prioritarias del cliente sobre premio de campeonato, patrocinio, aportaciones a participantes y cobro inmediato.\\
\texttt{UI-043} & \texttt{RIE-010} & Secciones 2.2, 2.6.4.3 y decisiones prioritarias del cliente sobre premio de campeonato, patrocinio, aportaciones a participantes y cobro inmediato.\\
\texttt{UI-044} & \texttt{RIE-010} & Secciones 2.2, 2.6.4.3 y decisiones prioritarias del cliente sobre premio de campeonato, patrocinio, aportaciones a participantes y cobro inmediato.\\
\texttt{UI-045} & \texttt{RIE-011} & Secciones 1.2, 2.2, 2.6.2, 2.6.4.2 y decisiones prioritarias del cliente sobre patrocinio directo de eventos, aportaciones a participantes y cobro inmediato.\\
\texttt{UI-046} & \texttt{RIE-011} & Secciones 1.2, 2.2, 2.6.2, 2.6.4.2 y decisiones prioritarias del cliente sobre patrocinio directo de eventos, aportaciones a participantes y cobro inmediato.\\
\texttt{UI-047} & \texttt{RIE-011} & Secciones 1.2, 2.2, 2.6.2, 2.6.4.2 y decisiones prioritarias del cliente sobre patrocinio directo de eventos, aportaciones a participantes y cobro inmediato.\\
\texttt{UI-048} & \texttt{RIE-011} & Secciones 1.2, 2.2, 2.6.2, 2.6.4.2 y decisiones prioritarias del cliente sobre patrocinio directo de eventos, aportaciones a participantes y cobro inmediato.\\
\texttt{UI-049} & \texttt{RIE-011} & Secciones 1.2, 2.2, 2.6.2, 2.6.4.2 y decisiones prioritarias del cliente sobre patrocinio directo de eventos, aportaciones a participantes y cobro inmediato.\\
\texttt{UI-050} & \texttt{RIE-011} & Secciones 1.2, 2.2, 2.6.2, 2.6.4.2 y decisiones prioritarias del cliente sobre patrocinio directo de eventos, aportaciones a participantes y cobro inmediato.\\
\texttt{UI-051} & \texttt{RIE-012} & Secciones 1.2, 2.2, 3.1.26 y fuente documental \texttt{TOP\_Racing\_Grok\_Clarified.tex} (línea 133).\\
\texttt{UI-052} & \texttt{RIE-013} & Secciones 1.2, 2.2, 2.6.2.2, 3.1.27 y decisiones prioritarias del cliente sobre configuración y publicación de comisiones.\\
\texttt{UI-053} & \texttt{RIE-013} & Secciones 1.2, 2.2, 2.6.2.2, 3.1.27 y decisiones prioritarias del cliente sobre configuración y publicación de comisiones.\\
\texttt{UI-054} & \texttt{RIE-013} & Secciones 1.2, 2.2, 2.6.2.2, 3.1.27 y decisiones prioritarias del cliente sobre configuración y publicación de comisiones.\\
\texttt{UI-055} & \texttt{RIE-014} & Secciones 2.2, 2.7, 3.5, RIE-001, RIE-003, RIE-004, RIE-005, RIE-009, RIE-010, RIE-011 y decisiones prioritarias del cliente sobre importancia de la estética y la UX para la adopción.\\
\texttt{UI-056} & \texttt{RIE-016} & Secciones 1.2, 2.2, 2.3, 3.1.4, 3.1.6, 3.1.13, 3.1.17, 3.1.18, 3.1.21, 3.2 RIE-003, RIE-005 y decisiones prioritarias del cliente confirmadas a partir del análisis funcional de la app legacy sobre la lista común de inscripciones con captura condicionada de resultados, premio de terminación y pujas.\\
\texttt{UI-057} & \texttt{RIE-016} & Secciones 1.2, 2.2, 2.3, 3.1.4, 3.1.6, 3.1.13, 3.1.17, 3.1.18, 3.1.21, 3.2 RIE-003, RIE-005 y decisiones prioritarias del cliente confirmadas a partir del análisis funcional de la app legacy sobre la lista común de inscripciones con captura condicionada de resultados, premio de terminación y pujas.\\
\texttt{UI-058} & \texttt{RIE-016} & Secciones 1.2, 2.2, 2.3, 3.1.4, 3.1.6, 3.1.13, 3.1.17, 3.1.18, 3.1.21, 3.2 RIE-003, RIE-005 y decisiones prioritarias del cliente confirmadas a partir del análisis funcional de la app legacy sobre la lista común de inscripciones con captura condicionada de resultados, premio de terminación y pujas.\\
\texttt{UI-059} & \texttt{RIE-016} & Secciones 1.2, 2.2, 2.3, 3.1.4, 3.1.6, 3.1.13, 3.1.17, 3.1.18, 3.1.21, 3.2 RIE-003, RIE-005 y decisiones prioritarias del cliente confirmadas a partir del análisis funcional de la app legacy sobre la lista común de inscripciones con captura condicionada de resultados, premio de terminación y pujas.\\
\texttt{UI-060} & \texttt{RIE-016} & Secciones 1.2, 2.2, 2.3, 3.1.4, 3.1.6, 3.1.13, 3.1.17, 3.1.18, 3.1.21, 3.2 RIE-003, RIE-005 y decisiones prioritarias del cliente confirmadas a partir del análisis funcional de la app legacy sobre la lista común de inscripciones con captura condicionada de resultados, premio de terminación y pujas.\\
\texttt{UI-061} & \texttt{RIE-016} & Secciones 1.2, 2.2, 2.3, 3.1.4, 3.1.6, 3.1.13, 3.1.17, 3.1.18, 3.1.21, 3.2 RIE-003, RIE-005 y decisiones prioritarias del cliente confirmadas a partir del análisis funcional de la app legacy sobre la lista común de inscripciones con captura condicionada de resultados, premio de terminación y pujas.\\
\texttt{UI-062} & \texttt{RIE-016} & Secciones 1.2, 2.2, 2.3, 3.1.4, 3.1.6, 3.1.13, 3.1.17, 3.1.18, 3.1.21, 3.2 RIE-003, RIE-005 y decisiones prioritarias del cliente confirmadas a partir del análisis funcional de la app legacy sobre la lista común de inscripciones con captura condicionada de resultados, premio de terminación y pujas.\\
\texttt{UI-063} & \texttt{RIE-016} & Secciones 1.2, 2.2, 2.3, 3.1.4, 3.1.6, 3.1.13, 3.1.17, 3.1.18, 3.1.21, 3.2 RIE-003, RIE-005 y decisiones prioritarias del cliente confirmadas a partir del análisis funcional de la app legacy sobre la lista común de inscripciones con captura condicionada de resultados, premio de terminación y pujas.\\
\texttt{UI-064} & \texttt{RIE-016} & Secciones 1.2, 2.2, 2.3, 3.1.4, 3.1.6, 3.1.13, 3.1.17, 3.1.18, 3.1.21, 3.2 RIE-003, RIE-005 y decisiones prioritarias del cliente confirmadas a partir del análisis funcional de la app legacy sobre la lista común de inscripciones con captura condicionada de resultados, premio de terminación y pujas.\\
\texttt{UI-065} & \texttt{RIE-017} & Secciones 1.2, 2.2, 2.4, 3.1.16, RF-104, RF-109 y análisis funcional confirmado de la app legacy sobre la consulta de prioridades de eventos por selección de nivel temporal y territorial.\\
\texttt{UI-066} & \texttt{RIE-017} & Secciones 1.2, 2.2, 2.4, 3.1.16, RF-104, RF-109 y análisis funcional confirmado de la app legacy sobre la consulta de prioridades de eventos por selección de nivel temporal y territorial.\\
\texttt{UI-067} & \texttt{RIE-017} & Secciones 1.2, 2.2, 2.4, 3.1.16, RF-104, RF-109 y análisis funcional confirmado de la app legacy sobre la consulta de prioridades de eventos por selección de nivel temporal y territorial.\\
\texttt{UI-068} & \texttt{RIE-017} & Secciones 1.2, 2.2, 2.4, 3.1.16, RF-104, RF-109 y análisis funcional confirmado de la app legacy sobre la consulta de prioridades de eventos por selección de nivel temporal y territorial.\\
\texttt{UI-069} & -- & \texttt{DEC-170}, \texttt{RF-EPR-022}; reserva de información económica del premio antes del cierre.\\
\texttt{UI-070} & -- & \texttt{DEC-173}, \texttt{DEC-174}, \texttt{DEC-175}, \texttt{DEC-180}; categoría tecnológica admitida y documentación abierta obligatoria.\\
\texttt{RC-EXT-001} & \texttt{REX-001} & Sección 2.5 y decisiones prioritarias del cliente sobre pagos reales.\\
\texttt{RC-EXT-002} & \texttt{REX-002} & Sección 2.5, RF-075, RF-079A, RF-079B y decisiones prioritarias del cliente.\\
\texttt{RC-EXT-003} & \texttt{REX-003} & Sección 2.5, RF-046 y decisiones prioritarias del cliente sobre entrega previa del vehículo al promotor.\\
\texttt{RC-EXT-004} & \texttt{REX-004} & Sección 2.5, RF-067A a RF-067E, RF-070A, RF-101, RF-102, RF-129 y decisiones prioritarias del cliente.\\
\texttt{RC-EXT-005} & \texttt{REX-005} & Sección 2.5, RF-058A, RF-067C y DEC-027.\\
\texttt{RC-EXT-006} & \texttt{REX-006} & Sección 2.5, sección 2.3 y DEC-069.\\
\texttt{RC-EXT-007} & \texttt{REX-006A} & Sección 2.5, sección 2.6, RF-151A y decisiones prioritarias del cliente sobre verificación reforzada.\\
\texttt{RC-EXT-008} & \texttt{REX-007} & Sección 2.5, RF-071B y decisiones prioritarias del cliente sobre aceptación versionada.\\
\texttt{RC-EXT-009} & \texttt{REX-008} & Sección 2.5, RF-050C y DEC-088 a DEC-090.\\
\texttt{QR-001} & \texttt{RNF-001} & Consolidación de RNF-001\\
\texttt{QR-002} & \texttt{RNF-002, RNF-004, RAC-002} & Consolidación de RNF-002, RNF-004, RAC-002\\
\texttt{QR-003} & \texttt{RNF-003, RAC-003} & Consolidación de RNF-003, RAC-003\\
\texttt{QR-004} & \texttt{RAC-005} & Consolidación de RAC-005\\
\texttt{QR-005} & \texttt{RNF-005, RAC-004} & Consolidación de RNF-005, RAC-004\\
\texttt{QR-006} & \texttt{RNF-006, RAC-001} & Consolidación de RNF-006, RAC-001\\
\texttt{QR-007} & \texttt{RNF-007} & Consolidación de RNF-007\\
\texttt{QR-008} & \texttt{RNF-008} & Consolidación de RNF-008\\
\texttt{QR-009} & \texttt{RAC-006} & Consolidación de RAC-006\\
\texttt{QR-010} & \texttt{RNF-009, RAC-007, RIE-015} & Consolidación de RNF-009, RAC-007, RIE-015\\
\texttt{QR-011} & \texttt{RNF-010, RAC-008} & Consolidación de RNF-010, RAC-008\\
\texttt{QR-012} & \texttt{RNF-011, RAC-009} & Consolidación de RNF-011, RAC-009\\
\texttt{QR-013} & \texttt{RNF-012, RAC-010} & Consolidación de RNF-012, RAC-010\\
\texttt{QR-014} & \texttt{RNF-013, RAC-011} & Consolidación de RNF-013, RAC-011\\
\texttt{QR-015} & \texttt{RNF-014, RAC-012} & Consolidación de RNF-014, RAC-012\\
\bottomrule
\end{longtable}
\end{landscape}

\subsection{Rationales e hipótesis separadas de las obligaciones}

\noindent\textbf{RAT-RF-AUC-032.} Rationale confirmado sobre la estrategia de puja del dueño: sea $v$ el precio mínimo real al que el dueño está dispuesto a vender, $H$ la mayor puja externa y $M$ la segunda mayor puja externa. Pujar $b=v$ implementa directamente el umbral real de venta: si $H>v$, un tercero puede ganar y el precio oficial no queda por debajo de $v$; si $H\leq v$, el dueño conserva el vehículo. Frente a esta referencia, una puja $b>v$ solo aumenta directamente el precio recibido por una venta cuando ocurre $M<b<H$: la puja propia se inserta entre las dos mayores pujas externas y sustituye a $M$ como segundo precio. Si $b\leq M$, la elevación no mejora el precio; si $v<H\leq b$, la elevación impide una venta que el dueño habría aceptado y produce autoadjudicación. Por tanto, $b>v$ constituye una apuesta por ocupar un intervalo oculto $(M,H)$ y cerrar parte de la brecha $H-M$. La subasta no revela $H$, $M$ ni indicadores provisionales de su posición; en consecuencia tampoco revela $H-M$ ni ofrece una vía independiente para inferir esa brecha. Estimarla exige un modelo externo de la distribución conjunta de las pujas. Bajo la información disponible dentro de TOP Racing, pujar $v$ es la estrategia robusta recomendada; pujar por encima de $v$ es una estrategia especulativa que solo resulta razonable con información externa suficientemente sólida y no se declara estrategia dominante ni óptima universal.\par
\noindent\textbf{RAT-RF-EFF-018.} Rationale confirmado: el dueño del vehículo subastado es siempre participante y, por ello, está expuesto al premio de eficiencia de su propia inscripción. El precio oficial de subasta es simultáneamente resultado de adjudicación y señal económica para la línea de tendencia precio--rendimiento. Un vehículo competitivo que el mercado valore por debajo de la tendencia obtiene una diferencia de eficiencia positiva y puede recibir mayor incentivo económico. Una valoración superior a la tendencia reduce la diferencia de eficiencia y puede disminuir el premio. Antes del cierre, el participante conoce el pozo inicial publicado, pero no el pozo económico oficial consolidado ni los importes de aportaciones y contribuciones que lo integran; tampoco conoce los precios oficiales ajenos ni la clasificación final. Por ello, el efecto monetario esperado de una sobrepuja sobre el premio de eficiencia es difícil de cuantificar de forma fiable. Este acoplamiento refuerza el carácter robusto de pujar la disposición mínima real de venta y añade riesgo a una puja superior: una estimación excesiva puede producir autoadjudicación, comisión del comprador y una peor posición de eficiencia cuyo costo económico exacto no es observable antes del cierre.\par
\noindent\textbf{RAT-BR-TEC-001.} Rationale confirmado: el objetivo institucional de TOP Racing es incentivar mérito tecnodeportivo, no gasto absoluto. El mecanismo premia a participantes que logran rendimiento observado alto con una valoración económica inferior a la tendencia del evento, y simultáneamente reconoce la destreza deportiva necesaria para extraer rendimiento del vehículo. La plataforma puede operar con categorías locales, caseras, artesanales, reutilizadas o de bajo costo sujetas a tecnología abierta obligatoria, porque su referencia de eficiencia es endógena al evento y no exige que la región compita contra una frontera tecnológica global. Este rationale no convierte por sí mismo una hipótesis de política pública en obligación de software ni afirma causalidad empírica sin validación posterior.\par
\noindent\textbf{RAT-BR-TEC-002.} Rationale confirmado: cuando el dueño puja $v$, toda puja externa superior puede adquirir el vehículo. Esa adquisición puede tener valor no solo de uso o reventa, sino también de aprendizaje, inspección e ingeniería inversa física permitida por las reglas aplicables. En consecuencia, la subasta puede funcionar como canal de comercio y difusión tecnológica regional sobre vehículos cuya documentación técnica abierta ya forma parte de las obligaciones del dominio. La SRS no presupone que la plataforma pueda sustituir límites legales de propiedad intelectual, seguridad, derechos de terceros, control físico del evento o restricciones externas de acceso.\par
\noindent\textbf{RAT-BR-TEC-003.} Rationale de política pública: un patrocinio público que aumente bolsas o infraestructura común puede amplificar incentivos privados hacia desarrollo tecnológico costo--eficiente y destreza deportiva en regiones con distinto nivel de desarrollo. La conveniencia de esa política depende de mantener bajas las barreras de entrada, evitar selección discrecional de ganadores, preservar seguridad y medir si el esfuerzo competitivo genera capacidades reutilizables fuera de la competición.\par
\noindent\textbf{RAT-BR-WGT-001.} Rationale confirmado: la matriz de pesos introduce competencia por relevancia entre promotores cuyos eventos se solapan en ámbitos territoriales y temporales comparables. La entrada de un evento con P/C suficiente puede desplazar a otros hacia valores inferiores de la secuencia de pesos en los campeonatos superiores compartidos. Este efecto se interpreta como congestión competitiva de peso y no como invalidación de los eventos afectados.\par
\noindent\textbf{RAT-BR-WGT-002.} Rationale confirmado: la competencia por peso crea dos respuestas estratégicas para el promotor. Puede mejorar el P/C aumentando el pozo o reduciendo el costo obligatorio por participante, o puede diferenciar territorio o periodo para disminuir el conjunto de eventos con los que compite por relevancia superior. La diferencia entre peso 100 y el peso vigente de rango inferior constituye una prima potencial de diferenciación antes de considerar demanda, costos de sede, traslado, organización y economías de red.\par
\noindent\textbf{RAT-BR-WGT-003.} Rationale académico y de política pública: la presión de diferenciación territorial y temporal puede inducir expansión descentralizada hacia regiones o periodos con menor densidad de eventos. Un nuevo nodo competitivo puede activar localmente el ciclo de participación, desarrollo tecnodeportivo, subasta, transferencia y difusión tecnológica. La SRS no afirma que la expansión sea automática ni socialmente óptima; deben evaluarse fragmentación, demanda mínima, calidad organizativa, seguridad, manipulación de delimitaciones y persistencia de economías de red.\par
\noindent\textbf{RAT-RF-EPR-017.} Rationale heredado: favorecer la terminación pretende incentivar durabilidad y reducir abandonos que disminuyan el espectáculo. La efectividad causal de ese incentivo requiere validación empírica.\par
\section{Control documental}
\begin{longtable}{P{2.0cm}P{2.6cm}P{8.8cm}}
\toprule
\textbf{Versión} & \textbf{Fecha} & \textbf{Cambio principal}\\
\midrule
1.x heredada & hasta 2026-03-08 & SRS integrada con requisitos RF/RIE/RNF/REX/RAC, 150 decisiones prioritarias y control de iteraciones detallado.\\
2.0 candidato & 2026-07-05 & Refactorización integral: jerarquía LaTeX corregida, metamodelo DEC/BR/CALC/RF/UI/RC-EXT/QR/AC, atomización de requisitos compuestos críticos, consolidación de calidad, separación de rationale, trazabilidad tabular y brechas normativas explícitas.\\
2.1 normativa & 2026-07-05 & Cierre de PEND-001 a PEND-009 mediante DEC-151 a DEC-159; costo por participante, secuencia recurrente, política numérica, calificación, carrera por tiempo o distancia, UX, accesibilidad y denominadores cero.\\
2.2 normativa consolidada & 2026-07-06 & Revisión acoplada de subasta de vehículos y eficiencia: línea global de tendencia de precios en función del rendimiento, diferencia precio--tendencia, confidencialidad total hasta el cierre, pujas únicamente crecientes, participación del dueño, excepción de retención del principal, autoadjudicación y comisión a cargo del comprador.\\
2.3 normativa consolidada & 2026-07-06 & Consolidación de la estrategia de puja del dueño: el dueño inscrito es el participante, solo vehículos de inscripciones computables participantes entran en subasta, el pozo consolidado y las aportaciones permanecen reservados durante la fase y el rationale recomienda $b=v$ como estrategia robusta; una sobrepuja solo mejora directamente el precio cuando $M<b<H$ y la brecha $H-M$ no es observable desde el sistema.\\
2.4 normativa consolidada & 2026-07-06 & Consolidación histórica de la dimensión tecnodeportiva: bajo costo, tecnología casera u abierta según categoría, adquisición para uso o estudio técnico, difusión tecnológica por transferencia de vehículos, política pública como rationale no normativo y generalización del vehículo de competición más allá de vehículos motorizados caros.\\
2.5 normativa consolidada & 2026-07-06 & Explicitación de la matriz de pesos como mecanismo de competencia por relevancia entre eventos comparables, congestión competitiva de peso, prima potencial de diferenciación territorial--temporal y posible expansión descentralizada hacia nuevos territorios o periodos.\\
2.6 normativa consolidada & 2026-07-21 & Corrección del estatuto de tecnología abierta del vehículo como obligación transversal del dominio TOP Racing; separación entre apertura tecnológica obligatoria y acceso irrestricto al evento físico, que puede quedar limitado por reglas publicadas del promotor dentro del alcance operativo del sistema.\\
\bottomrule
\end{longtable}

El historial narrativo completo de iteraciones de la SRS heredada no forma parte del cuerpo normativo de esta versión. Se conserva en el artefacto fuente heredado para auditoría histórica.

\appendix
\section{Vocabulario controlado completo}
Las definiciones siguientes preservan el glosario completo de la SRS heredada, con correcciones editoriales menores. En caso de conflicto con una definición CALC o BR más específica, prevalece la definición normativa más específica.

\noindent\textbf{Admin.} único usuario administrativo del sistema.\par
\noindent\textbf{Aportación directa a bolsa.} aportación económica no asociada a un participante específico, destinada directamente a la bolsa de un evento o de un campeonato. En el dominio TOP Racing, toda aportación directa a bolsa constituye patrocinio.\par
\noindent\textbf{Aportante.} usuario o tercero que realiza una o varias aportaciones económicas, ya sea directamente a la bolsa de un evento o campeonato como patrocinio, o a favor de un participante específico para compartir el premio que a ese participante corresponda.\par
\noindent\textbf{Clasificación.} ordenación de resultados según una modalidad, un ámbito territorial y un periodo temporal.\par
\noindent\textbf{Clasificación combinada.} clasificación del dominio que integra resultados de velocidad y de carrera cuando así corresponda; para salida invertida, la aplicable al evento queda determinada por su nivel territorial y su nivel temporal.\par
\noindent\textbf{Competidor.} papel funcional del usuario que accede a un evento como competidor y compite en él.\par
\noindent\textbf{Cuenta individual.} cuenta propia de un usuario dentro del sistema, mediante la cual accede a las funciones que le correspondan.\par
\noindent\textbf{Aceptación de términos y condiciones.} manifestación registrada y versionada por el sistema mediante la cual el usuario acepta una versión concreta de las condiciones legales aplicables al uso operativo del ecosistema.\par
\noindent\textbf{Identidad verificada del usuario.} condición por la cual el sistema reconoce que la identidad de un usuario ha sido verificada y, por tanto, puede habilitarle funciones que impliquen participación dentro del ecosistema.\par
\noindent\textbf{Verificación reforzada del usuario.} máximo nivel de verificación operativa del dominio, gestionado por \texttt{admin}, requerido para operaciones de mayor riesgo como inscribir vehículos, pujar en subastas de vehículos y promover eventos.\par
\noindent\textbf{Estado booleano de identidad verificada.} valor lógico que expresa si la identidad del usuario está marcada como verificada o no verificada dentro del sistema.\par
\noindent\textbf{Estado booleano de verificación reforzada.} valor lógico que expresa si la verificación reforzada del usuario está activa o inactiva dentro del sistema.\par
\noindent\textbf{Estado booleano de exclusión del usuario.} valor lógico que expresa si el usuario está marcado como excluido o no excluido dentro del sistema.\par
\noindent\textbf{Nivel de verificación del usuario.} nivel visible derivado de los estados booleanos de identidad verificada, verificación reforzada y exclusión del usuario, utilizado para informar qué operaciones puede realizar.\par
\noindent\textbf{Operación bloqueada por nivel de verificación.} operación que el sistema no habilita porque el usuario no cumple el nivel o la condición de verificación requerida por el dominio.\par
\noindent\textbf{Recordatorio de habilitación operativa.} mensaje emitido por el sistema cuando una operación queda bloqueada por nivel de verificación, indicando el nivel actual del usuario, la condición faltante y la acción necesaria para habilitarla.\par
\noindent\textbf{Usuario excluido.} usuario cuyo estado booleano de exclusión está activo y que, por ello, no puede participar operativamente en el sistema.\par
\noindent\textbf{Contribución económica asociada a participante.} aportación económica registrada a favor de un participante específico. En eventos puede pertenecer a la familia de eficiencia o a la de terminación; en campeonatos puede pertenecer a la familia de eficiencia o a la de terminación según la modalidad del premio correspondiente.\par
\noindent\textbf{Contribución de eficiencia del participante.} aportación económica asociada a un participante específico que sirve para compartir y aumentar el premio basado en eficiencia que a ese participante corresponda en un evento o en un campeonato de eficiencia.\par
\noindent\textbf{Contribución total de eficiencia del participante.} suma de todas las contribuciones de eficiencia válidamente aceptadas a favor de un mismo participante dentro de un evento o de un campeonato de eficiencia.\par
\noindent\textbf{Contribución de terminación del campeonato.} aportación económica asociada a un participante específico de un campeonato de terminación que sirve para compartir y aumentar el premio de terminación que a ese participante corresponda en ese campeonato.\par
\noindent\textbf{Contribución total de terminación del campeonato.} suma de todas las contribuciones de terminación válidamente aceptadas a favor de un mismo participante dentro de un campeonato de terminación.\par
\noindent\textbf{Campeonato de terminación.} campeonato cuya determinación económica usa resultados acumulados de terminación en lugar de resultados de eficiencia.\par
\noindent\textbf{Criterio de verificación.} condición objetiva que permite comprobar si un requerimiento se cumple.\par
\noindent\textbf{Datos oficiales del evento.} datos del evento que el sistema reconoce como válidos para producir resultados, clasificaciones y resultados económicos oficiales.\par
\noindent\textbf{Desempeño participante.} actuación competitiva de un participante dentro de un evento o dentro de un campeonato, considerada para efectos competitivos y económicos.\par
\noindent\textbf{Distribución económica posterior.} reparto posterior al resultado económico individual de un participante entre el propio participante y las personas asociadas a las contribuciones hechas a su favor, según las reglas publicadas.\par
\noindent\textbf{Dueño de vehículo.} papel funcional del usuario al que pertenece un vehículo registrado en el sistema y que administra su información dentro del contexto del sistema y del evento.\par
\noindent\textbf{Vehículo de competición.} artefacto físico admitido por una categoría de evento para producir rendimiento medible en pista; puede ser motorizado o no motorizado, casero, artesanal, reutilizado o de bajo costo si satisface las reglas publicadas, las restricciones de seguridad del evento y la condición transversal de tecnología abierta del vehículo.\par
\noindent\textbf{Tecnología abierta del vehículo.} condición obligatoria del vehículo de competición dentro del dominio TOP Racing por la cual la solución técnica del vehículo debe quedar documentada con información suficiente para ser consultada, estudiada, replicada o reutilizada conforme a los términos publicados aplicables. La apertura tecnológica no implica ausencia de reglas de seguridad, restricciones legales externas, derechos de terceros, control físico del evento ni admisión irrestricta de participantes por parte del promotor.\par
\noindent\textbf{Mérito tecnodeportivo.} mérito integrado que combina desarrollo tecnológico costo--eficiente, capacidad deportiva para extraer rendimiento del vehículo y validación práctica de esa combinación mediante rendimiento medido en pista.\par
\noindent\textbf{Rendimiento observado.} velocidad media correspondiente a la vuelta válida más rápida de calificación de un participante computable; queda fijada antes de la subasta de vehículos.\par
\noindent\textbf{Eficiencia (del dominio).} indicador económico relativo que compara el precio oficial de subasta de un vehículo con el precio de tendencia asociado a su rendimiento por la relación precio--rendimiento del evento.\par
\noindent\textbf{Eficiencia relativa al mercado.} diferencia entre el precio indicado por la línea de tendencia del evento para el rendimiento observado y el precio oficial de subasta del vehículo; es positiva cuando el mercado valoró el vehículo por debajo de la tendencia y negativa cuando lo valoró por encima.\par
\noindent\textbf{Espectador.} papel funcional del usuario que consulta información pública del ecosistema y puede participar en las funciones públicas disponibles para un evento.\par
\noindent\textbf{Estado del evento.} situación oficial del evento dentro de su ciclo de vida. En esta SRS los estados oficiales son \texttt{creado}, \texttt{inscripciones}, \texttt{velocidad}, \texttt{carrera}, \texttt{subasta}, \texttt{publicado} y \texttt{cancelado}.\par
\noindent\textbf{Evento.} instancia competitiva con un ciclo de vida definido.\par
\noindent\textbf{IEEE 830.} estándar de referencia para estructurar especificaciones de requerimientos.\par
\noindent\textbf{Interfaz.} medio de interacción del usuario con las funciones del sistema.\par
\noindent\textbf{Interacción comunitaria.} intercambio de discusiones, preguntas o retroalimentación dentro del ecosistema TOP Racing.\par
\noindent\textbf{Canal de novedades del ecosistema.} medio disponible para difundir novedades del ecosistema sin fijar en esta especificación su implementación concreta.\par
\noindent\textbf{Foro comunitario.} espacio comunitario disponible para publicar y consultar discusiones, preguntas y retroalimentación del ecosistema sin fijar en esta especificación su implementación concreta.\par
\noindent\textbf{Moderación del espacio comunitario y de novedades.} función atribuida a \texttt{admin} para moderar el foro comunitario y el canal de novedades del ecosistema.\par
\noindent\textbf{Función de aportaciones a participantes.} capacidad configurable del sistema para aceptar aportaciones a favor de participantes específicos y vincularlas a la distribución posterior del premio correspondiente.\par
\noindent\textbf{Estado de activación de aportaciones a participantes.} estado booleano global, gestionado solo por \texttt{admin}, que determina si el sistema acepta o rechaza aportaciones a participantes en eventos o en campeonatos, según el ámbito configurado.\par
\noindent\textbf{Novedad del ecosistema.} anuncio o actualización relevante del ecosistema TOP Racing, incluidos anuncios de eventos, resultados competitivos, resultados de subasta y destacados de la comunidad.\par
\noindent\textbf{Nota de cancelación.} texto registrado por el promotor para dejar constancia de la razón de una cancelación de inscripción por una causa permitida.\par
\noindent\textbf{Inspección comunitaria.} comprobación previa al inicio de la actividad competitiva del evento, realizada por competidores registrados sobre los vehículos participantes para verificar condiciones visibles de seguridad antes de \texttt{carrera}.\par
\noindent\textbf{Lista de comprobación de seguridad y cumplimiento.} conjunto estandarizado de verificaciones visibles de seguridad aplicables al evento.\par
\noindent\textbf{Objeción de seguridad.} observación específica formulada respecto de un vehículo durante la inspección comunitaria por una posible condición de riesgo o incumplimiento.\par
\noindent\textbf{Protesta de seguridad.} tramitación formal de una objeción de seguridad dentro de la ventana permitida para inspección y decisión previa a \texttt{carrera}.\par
\noindent\textbf{Resultado de la protesta de seguridad.} constancia pública de la votación formal, de la decisión del promotor o de ambas, relativa a una protesta de seguridad tramitada dentro del evento.\par
\noindent\textbf{Resultado detallado de la votación formal.} publicación del total de votos emitidos y del voto de cada participante en una votación formal del evento, identificado por su nombre/usuario exacto.\par
\noindent\textbf{Presentante de la protesta de seguridad.} usuario que registra una objeción de seguridad o inicia su tramitación formal dentro del evento.\par
\noindent\textbf{Inscripción del competidor.} registro mediante el cual el dueño de un vehículo previamente registrado lo inscribe en un evento, identifica al usuario piloto cuando exista pilotaje humano directo y deja asociados los datos económicos mínimos de esa inscripción.\par
\noindent\textbf{Inscripción cancelada.} inscripción cuyo registro se conserva, pero que queda invalidada a todos los efectos del evento por una causa permitida por las reglas aplicables.\par
\noindent\textbf{Vehículo registrado.} vehículo previamente registrado en el sistema con sus datos mínimos obligatorios y asociado a un usuario dueño.\par
\noindent\textbf{Liquidación económica oficial.} cierre económico oficial del evento una vez consolidados el pozo económico oficial, los resultados del dominio que correspondan y la distribución económica posterior.\par
\noindent\textbf{Historial del evento.} conjunto de resultados oficiales, publicaciones y decisiones registradas del evento que permanecen consultables con el tiempo.\par
\noindent\textbf{Consulta pública.} acceso de los usuarios a la información pública que el sistema pone a disposición del ecosistema.\par
\noindent\textbf{Consulta pública del evento.} vista pública del evento en la que el sistema muestra la información del evento que las reglas del dominio declaren visible.\par
\noindent\textbf{Matriz de campeonatos.} estructura del dominio que resulta de combinar seis niveles temporales, ocho niveles territoriales y cuatro modalidades del sistema, dando lugar a 48 campeonatos y a sus clasificaciones aplicables.\par
\noindent\textbf{Matriz de pesos.} estructura publicada por el sistema que expresa el peso relativo con el que un evento aporta puntos a los distintos niveles territoriales y temporales que le resultan aplicables.\par
\noindent\textbf{Matriz de pesos vigente.} matriz de pesos actualmente aplicable a un evento, una vez consideradas las recalculaciones válidas por afectación entre eventos.\par
\noindent\textbf{Ámbito comparativo de peso.} combinación territorial y temporal en la que dos o más eventos resultan comparables para determinar su contribución a un campeonato superior conforme a las reglas de la matriz de pesos.\par
\noindent\textbf{Congestión competitiva de peso.} efecto por el cual la presencia o entrada de eventos comparables puede reducir el peso de otros eventos en ámbitos superiores compartidos al modificar su rango por P/C.\par
\noindent\textbf{Prima de diferenciación territorial-temporal.} incremento potencial de peso que un evento de rango inferior puede obtener al operar en un ámbito menos congestionado donde alcance un rango superior, antes de considerar costos de entrada y demanda.\par
\noindent\textbf{Campeonatos y clasificaciones aplicables al evento.} subconjunto de la estructura de campeonatos y clasificaciones al que un evento aporta resultados según su nivel territorial, su nivel temporal y las modalidades que le correspondan.\par
\noindent\textbf{Precio de tendencia aplicable.} precio indicado por la línea de tendencia global precio--rendimiento del evento para el rendimiento observado de un participante computable.\par
\noindent\textbf{Base económica computable del participante.} relación transparente entre la eficiencia relativa al mercado determinada para un participante o su posición oficial en un campeonato, y las contribuciones computables asociadas a ese participante.\par
\noindent\textbf{Aportaciones computables al pozo de campeonato.} conjunto de aportaciones directas a la bolsa del campeonato correspondiente y de contribuciones válidamente aceptadas a favor de los participantes de ese campeonato que integran el pozo correspondiente.\par
\noindent\textbf{Mejor oferta.} puja válida más alta dentro de una subasta.\par
\noindent\textbf{Papel funcional.} forma de participación que un usuario puede asumir dentro del ecosistema TOP Racing. Un mismo usuario puede acumular varios papeles funcionales, salvo \texttt{admin}, que es único y exclusivo.\par
\noindent\textbf{Modalidad.} dimensión de clasificación del dominio. En las fuentes del proyecto comprende velocidad, carrera, velocidad + carrera y eficiencia.\par
\noindent\textbf{Patrocinador de campeonato.} usuario que realiza una aportación directa a la bolsa de un campeonato y asocia a esa aportación un anuncio para su publicación.\par
\noindent\textbf{Patrocinador de evento.} usuario que realiza una aportación directa a la bolsa de un evento y asocia a esa aportación un anuncio para su publicación.\par
\noindent\textbf{Patrocinio aplicable por campeonato al evento.} patrocinio de un campeonato al que un evento pertenezca y que debe mostrarse en la consulta de ese evento, de forma distinguida del patrocinio directo del propio evento.\par
\noindent\textbf{Anuncio de patrocinio de campeonato.} contenido con imagen asociado a una aportación directa a la bolsa de un campeonato para su publicación como patrocinio.\par
\noindent\textbf{Anuncio de patrocinio de evento.} contenido con imagen asociado a una aportación directa a la bolsa de un evento para su publicación como patrocinio.\par
\noindent\textbf{Participante.} usuario dueño del vehículo inscrito en un evento para competir; cuando existe pilotaje humano directo, el piloto puede ser un usuario distinto del participante dueño.\par
\noindent\textbf{Pozo económico oficial.} monto económico consolidado del evento a partir del pozo inicial publicado, de las aportaciones directas a la bolsa del evento y de las contribuciones válidamente aceptadas a favor de participantes dentro de la ventana permitida.\par
\noindent\textbf{Pozo de campeonato.} monto económico asociado a un campeonato, integrado por las aportaciones directas a su bolsa y por las contribuciones válidamente aceptadas a favor de sus participantes para distribuirse según los resultados oficiales correspondientes.\par
\noindent\textbf{Pozo inicial.} monto económico inicial publicado para un evento antes de la consolidación del pozo económico oficial.\par
\noindent\textbf{Pozo del premio de terminación.} monto económico específico del evento, publicado como condición económica del mismo, destinado exclusivamente al premio de terminación.\par
\noindent\textbf{Contribución de terminación del participante.} aportación económica asociada a una inscripción específica del evento que interviene exclusivamente en la determinación del premio de terminación y sirve para compartir y aumentar ese premio.\par
\noindent\textbf{Contribución total de terminación del participante.} suma de todas las contribuciones de terminación válidamente aceptadas a favor de una misma inscripción computable dentro del evento.\par
\noindent\textbf{Premio de terminación.} premio económico del evento, separado del premio del evento basado en eficiencia, orientado a favorecer que los competidores terminen la carrera.\par
\noindent\textbf{Vueltas esperadas del competidor.} referencia operativa fraccionaria de vueltas que el sistema deriva conforme al tipo de límite de carrera: tiempo máximo dividido entre el mejor tiempo de vuelta válido de calificación, o distancia objetivo dividida entre la longitud oficial de una vuelta.\par
\noindent\textbf{Postor.} papel funcional del usuario que participa en subastas del ecosistema con arreglo a las reglas publicadas.\par
\noindent\textbf{Precio de mercado.} valoración económica que el mercado del evento revela a partir de las pujas válidas observadas bajo las reglas aplicables.\par
\noindent\textbf{Precio oficial de subasta.} segundo precio válido determinado para un vehículo al cierre de la fase de subasta del evento; se usa como precio de adjudicación y como señal económica $C_i$ para la línea de tendencia de eficiencia.\par
\noindent\textbf{Autoadjudicación.} resultado en el que el dueño previo presenta la mayor puja vigente y conserva su propio vehículo; no existe transferencia circular del principal, pero el dueño actúa como comprador adjudicatario para la comisión aplicable.\par
\noindent\textbf{Precio de subasta.} precio resultante de la subasta conforme al mecanismo definido para ella.\par
\noindent\textbf{Precio final.} precio que el sistema publica al cierre de una subasta.\par
\noindent\textbf{Promotor.} usuario que convoca y administra un evento específico conforme a las reglas publicadas.\par
\noindent\textbf{Puntaje de prioridad (P/C).} indicador publicado propio de un evento que interviene en la determinación de su matriz de pesos conforme a las reglas del dominio.\par
\noindent\textbf{Coeficiente de premio del evento.} Producto entre la base de puntos por eficiencia del participante y su contribución total de eficiencia computable dentro del evento.\par
\noindent\textbf{Coeficiente de premio de campeonato.} magnitud operativa que combina el resultado acumulado computable del participante en el campeonato correspondiente con su contribución total computable dentro de ese mismo campeonato para determinar la parte del pozo que le corresponda.\par
\noindent\textbf{Distribución económica posterior proporcional.} Reparto del premio de un participante entre ese participante y las personas asociadas a sus contribuciones de la misma familia económica, según la proporción que cada cantidad represente dentro de la contribución total computable correspondiente.\par
\noindent\textbf{Secuencia base de pesos y puntos del dominio.} secuencia descendente utilizada para pesos y bases de puntos, definida por la repetición del bloque 100, 63, 40, 25, 16 en órdenes de magnitud sucesivamente menores: 10, 6.3, 4, 2.5, 1.6; 1, 0.63, 0.4, 0.25, 0.16; y así sucesivamente conforme a \texttt{CALC-PTS-001}.\par
\noindent\textbf{Regresión log-lineal de mercado.} línea de tendencia global de los precios oficiales de subasta en función del rendimiento observado, ajustada una sola vez sobre todos los participantes computables mediante $\ln(\widehat C(S))=a+bS$, equivalente a $\widehat C(S)=\exp(a+bS)$.\par
\noindent\textbf{Puja.} oferta económica presentada por un postor dentro de una subasta.\par
\noindent\textbf{Puja secreta.} puja cuyo monto y posición relativa permanecen reservados mientras la subasta esté abierta; el postor solo puede consultar su propia puja vigente.\par
\noindent\textbf{Puja válida.} puja que cumple las reglas publicadas de la subasta correspondiente.\par
\noindent\textbf{Reglas publicadas.} reglas del evento que el sistema pone a disposición antes de las fases en las que deban ser conocidas por quienes participan en él.\par
\noindent\textbf{Convocatoria del evento.} publicación del evento en estado \texttt{inscripciones}, junto con sus condiciones, reglas y datos públicos.\par
\noindent\textbf{Carga mínima.} carga total mínima obligatoria para competir en un evento, integrada por el peso del piloto y de su equipo protector personal cuando exista pilotaje humano, más el lastre verificable que corresponda como condición publicada del evento.\par
\noindent\textbf{Costo fijo por pesaje del evento.} importe fijo que el promotor publica para cada pesaje realizado fuera del sistema como condición económica del evento.\par
\noindent\textbf{Resultado competitivo.} resultado derivado del desempeño competitivo de un participante en el evento.\par
\noindent\textbf{Resultado económico individual.} resultado económico derivado del desempeño de un participante en el evento o en un campeonato, determinado con base en las reglas publicadas y en los datos oficiales aplicables.\par
\noindent\textbf{Resultado oficial.} resultado que el sistema publica con carácter oficial y que puede trazarse hasta los datos de origen y las reglas vigentes del evento.\par
\noindent\textbf{Resultados del dominio.} resultados calculados por el sistema conforme a las reglas del dominio aplicables al evento o al campeonato correspondiente, incluidos los resultados competitivos, las clasificaciones y la eficiencia cuando corresponda.\par
\noindent\textbf{SRS.} Especificación de Requerimientos de Software.\par
\noindent\textbf{Señal de mercado.} información económica observable dentro del evento que el dominio usa para formar expectativas de desempeño relativo.\par
\noindent\textbf{Segundo precio.} segundo precio más alto válido dentro de una subasta Vickrey.\par
\noindent\textbf{Subasta de acceso.} subasta vinculada al acceso de espectadores a un evento cuando las reglas publicadas del evento así lo establecen.\par
\noindent\textbf{Adjudicación de accesos por orden de pujas.} regla de la subasta de acceso por la cual los accesos disponibles se asignan a las mejores pujas válidas hasta completar la cantidad publicada y cada ganador paga la siguiente puja válida inferior a la suya.\par
\noindent\textbf{Transferencia obligatoria de vehículo subastado.} efecto del dominio por el cual el vehículo sometido a subasta debe quedar asignado al adjudicatario que resulte del resultado oficial vigente de la subasta correspondiente.\par
\noindent\textbf{Retiro no autorizado de vehículo subastado.} extracción del vehículo sometido a subasta fuera del control del evento sin que el resultado oficial vigente permita al dueño conservarlo o sin que se haya cumplido la transferencia obligatoria correspondiente.\par
\noindent\textbf{Subasta de vehículos.} subasta sellada de segundo precio realizada después de la carrera sobre los vehículos computables del evento. Pueden participar usuarios con verificación reforzada activa; el dueño puede pujar por su propio vehículo; antes del cierre no se publican pujas, conteos ni consecuencias derivadas.\par
\noindent\textbf{Subasta Vickrey.} mecanismo de puja sellada en el que gana la mejor oferta y se paga el segundo precio más alto válido.\par
\noindent\textbf{TOP.} sigla de \texttt{Technology Open Professional}, utilizada para expresar que la identidad del sistema se centra en tecnología, apertura de diseño y orientación profesional, entendida tanto como estándar serio de operación como posibilidad de profesionalización competitiva dentro del ecosistema.\par
\noindent\textbf{TOP Racing.} denominación completa de \texttt{Technology Open Professional Racing}; ecosistema de competencia técnica y de habilidad basado en desempeño competitivo, señales de mercado y exposición a venta de soluciones dentro del propio ecosistema.\par
\noindent\textbf{Professionalización competitiva dentro del ecosistema.} posibilidad de que los protagonistas de la competencia lleguen a obtener ingresos legítimos derivados de premios, patrocinios u otras consecuencias económicas regladas de su participación dentro de TOP Racing.\par
\noindent\textbf{Trazabilidad.} capacidad de relacionar cada requerimiento con su origen y con sus artefactos de validación.\par
\noindent\textbf{Trazabilidad histórica de resultados.} capacidad de consultar un resultado oficial junto con las reglas publicadas y los datos oficiales que lo sustentan, dentro del historial del evento.\par
\noindent\textbf{Usuario.} persona con acceso al sistema que puede asumir uno o varios papeles funcionales dentro del ecosistema TOP Racing.\par
\noindent\textbf{Consulta personal de estados del usuario.} consulta en la que el propio usuario puede ver sus estados booleanos de identidad verificada y de exclusión.\par
\noindent\textbf{Consulta limitada del promotor sobre estados del usuario.} consulta restringida por la cual el promotor puede ver los estados booleanos de verificación y exclusión solo de usuarios vinculados operativamente a su evento y solo durante el ciclo de vida de ese evento.\par
\noindent\textbf{Vehículo.} unidad técnica con la que se compite y que puede ser objeto de subasta según las reglas del evento.\par
\noindent\textbf{Ventana de contribuciones.} periodo del ciclo de vida del evento durante el cual el sistema admite nuevas contribuciones a favor de participantes o aportaciones directas a la bolsa. En los eventos, esa ventana termina al pasar el evento de \texttt{inscripciones} a \texttt{velocidad}.\par
\noindent\textbf{Ventana permitida.} intervalo del ciclo de vida del evento en el que una operación específica puede ejecutarse válidamente.\par
\noindent\textbf{Verificabilidad.} propiedad por la cual un requerimiento puede comprobarse de manera objetiva.\par
\noindent\textbf{Diferencia de eficiencia relativa al mercado.} resultado de restar el precio oficial de subasta al precio de tendencia correspondiente al rendimiento observado: $D_i=\widehat C(S_i)-C_i$.\par
\noindent\textbf{Puntuación escalada por peso.} Resultado de multiplicar la base de puntos de una modalidad por el peso vigente del evento y dividir el producto entre 100.\par
\noindent\textbf{Base de puntos por modalidad.} Valor de puntos asignado dentro de una modalidad según el ordenamiento oficial correspondiente y la secuencia base del dominio.\par
\noindent\textbf{Puntuación escalada del evento.} Resultado de aplicar el peso vigente del evento a la base de puntos de una modalidad en un campeonato aplicable.\par
\noindent\textbf{Actualización de clasificación combinada.} Incorporación separada de las puntuaciones escaladas de velocidad y carrera a la clasificación combinada del campeonato correspondiente.\par
\noindent\textbf{Desempate por inscripción al evento.} Regla según la cual, ante un empate en un ordenamiento del evento que deba producir puntos, base de puntos o premio, se favorece a quien se haya inscrito antes en ese evento.\par
\noindent\textbf{Cálculo sin redondeo.} término heredado para la ausencia de redondeos de dominio intermedios no prescritos. En la versión 2.1 los cálculos analíticos usan punto flotante binario de 64 bits conforme a \texttt{NUM-001}; los importes monetarios son enteros y toda distribución proporcional se reconcilia mediante \texttt{CALC-MON-001}.\par
\noindent\textbf{Resultado acumulado de terminación del campeonato.} suma de las relaciones entre vueltas completadas y vueltas esperadas de un participante en los eventos computables del campeonato de terminación correspondiente.\par
\noindent\textbf{Desempate en premio de campeonato por primera inscripción al campeonato.} Regla según la cual, ante un empate en el resultado computable aplicable del campeonato para efectos del premio de campeonato, se favorece a quien se haya inscrito antes en el primer evento de ese campeonato en el que haya participado.\par
\noindent\textbf{Ventana de aportaciones al premio de campeonato.} Intervalo que termina antes de que el primer evento del campeonato correspondiente haya cerrado inscripciones, dentro del cual pueden aceptarse aportaciones directas a la bolsa del campeonato y, solo cuando la función respectiva esté activa, aportaciones a participantes de ese campeonato.\par
\noindent\textbf{Ventana de aportaciones al evento.} intervalo que transcurre mientras el evento permanece en \texttt{inscripciones}, dentro del cual pueden aceptarse aportaciones directas a la bolsa del evento y, solo cuando la función respectiva esté activa, aportaciones a participantes de ese evento.\par
\noindent\textbf{Puja vigente de acceso.} única puja válida activa que un usuario mantiene en una subasta de acceso determinada.\par
\noindent\textbf{Actualización al alza de puja de acceso.} reemplazo de la puja vigente de acceso por otra de monto mayor presentada por el mismo usuario mientras la subasta permanezca abierta.\par
\noindent\textbf{Notificación de acceso adjudicado con QR.} notificación emitida al cierre de la subasta de acceso para la persona ganadora, que incluye un código QR asociado al acceso adjudicado y al importe correspondiente.\par
\noindent\textbf{Retención económica suficiente.} respaldo económico real mantenido por el sistema para cubrir el importe de una puja válida hasta que corresponda cobrarlo o liberarlo.\par
\noindent\textbf{Liberación de retención.} operación económica real mediante la cual el sistema deja sin efecto la retención asociada a una puja no ganadora o ya no vigente.\par
\noindent\textbf{Cobro automático al ganador.} operación económica real mediante la cual el sistema ejecuta el cargo correspondiente a la puja o adjudicación ganadora conforme al resultado oficial vigente.\par
\noindent\textbf{Pago automático al beneficiario.} operación económica real mediante la cual el sistema entrega de inmediato a un beneficiario la cantidad que le corresponde según el dominio.\par
\noindent\textbf{Devolución inmediata.} operación económica real mediante la cual el sistema reintegra de inmediato un importe previamente cobrado cuando las reglas del dominio así lo exigen.\par
\noindent\textbf{Promotor deudor de entrega.} condición temporal del promotor respecto del vehículo adjudicado mientras la persona ganadora no haya confirmado la recepción física en la app.\par
\noindent\textbf{Disputa de entrega de vehículo adjudicado.} controversia iniciada por la persona ganadora contra el promotor por falta de entrega del vehículo adjudicado dentro del plazo del evento.\par
\noindent\textbf{Reacción visible inmediata de la interfaz.} respuesta perceptible de la IU que confirma de forma inmediata que una acción solicitada por el usuario ya fue recibida por el sistema.\par
\noindent\textbf{Estado de espera operativa.} estado visible de la interfaz que indica que una operación está pendiente de un resultado o procesamiento para poder continuar.\par
\noindent\textbf{Bloqueo de acciones problemáticas durante el procesamiento.} deshabilitación transitoria de las demás acciones del usuario mientras una operación se encuentra en espera o procesamiento, siempre que realizarlas en ese momento genere problemas.\par
\noindent\textbf{Percepción inicial positiva de la interfaz.} respuesta favorable temprana de usuarios representativos durante sus primeros segundos de contacto con la interfaz, evaluable mediante pruebas de campo y técnicas de investigación UX.\par
\noindent\textbf{Jerarquía visual de información crítica del dominio.} organización visual clara de los elementos operativos y económicos más importantes de una pantalla, de modo que el usuario pueda distinguir el estado actual, las acciones habilitadas o bloqueadas, la razón del bloqueo y, cuando apliquen, los importes brutos, las comisiones, las retenciones, las devoluciones y los importes netos.\par
\noindent\textbf{Importes económicos críticos de la interfaz.} importes brutos, comisiones, retenciones, devoluciones e importes netos que una pantalla debe mostrar de forma diferenciada cuando resulten aplicables al caso de uso correspondiente.\par
\noindent\textbf{Legibilidad del contenido crítico.} cualidad por la cual el texto y los importes relevantes de una pantalla pueden leerse con rapidez y sin ambigüedad en condiciones normales de uso.\par
\noindent\textbf{Contraste visual suficiente.} diferenciación visual clara entre elementos relevantes de la interfaz que permite distinguir estados, acciones, bloqueos e información económica crítica.\par
\noindent\textbf{Consistencia visual y terminológica.} coherencia de patrones visuales, estructura informativa y terminología entre pantallas que representan conceptos equivalentes del dominio.\par
\noindent\textbf{Acción sensible del dominio.} operación cuya confirmación puede producir efectos operativos, económicos o jurídicos relevantes para el usuario o para terceros dentro del sistema.\par
\noindent\textbf{Prevención de error en acciones sensibles.} cualidad de la interfaz por la cual una acción sensible ofrece contexto, advertencias y frenos suficientes para reducir errores evitables antes de su confirmación.\par
\noindent\textbf{Claridad previa de consecuencias.} cualidad de la interfaz por la cual, antes de confirmar una acción sensible, informa con claridad las consecuencias operativas, económicas o jurídicas relevantes que esa acción produce cuando correspondan.\par
\section{Registro de decisiones de origen}
Las decisiones DEC preservan el contenido del registro prioritario heredado y reciben identificadores estables. Se utilizan como procedencia; la aceptación del sistema se realiza contra los elementos normativos canónicos derivados.

\noindent\textbf{DEC-001 --- Papeles funcionales acumulables y \texttt{admin} único}\par
Un mismo usuario puede asumir varios papeles funcionales dentro del ecosistema TOP Racing. El único papel exclusivo es \texttt{admin}, y existe un solo \texttt{admin}.\par\vspace{0.45em}
\noindent\textbf{DEC-002 --- Contribuciones de terceros al participante}\par
Las contribuciones voluntarias de terceros están disponibles como parte del modelo económico del sistema.\par\vspace{0.45em}
\noindent\textbf{DEC-003 --- Subasta de acceso para espectadores}\par
La subasta de acceso corresponde al acceso de espectadores al evento. No corresponde al acceso del competidor.\par\vspace{0.45em}
\noindent\textbf{DEC-004 --- Acceso del competidor por inscripción}\par
El acceso del competidor al evento se produce por inscripción y queda comprendido en esa inscripción.\par\vspace{0.45em}
\noindent\textbf{DEC-005 --- Publicación de la matriz de pesos desde \texttt{inscripciones}}\par
La matriz de pesos del evento debe estar publicada desde el momento en que el evento entra en estado \texttt{inscripciones} y debe permanecer visible durante el ciclo del evento.\par\vspace{0.45em}
\noindent\textbf{DEC-006 --- La convocatoria corresponde a la publicación de \texttt{inscripciones}}\par
La convocatoria del evento corresponde a la publicación del evento en estado \texttt{inscripciones}.\par\vspace{0.45em}
\noindent\textbf{DEC-007 --- La convocatoria incluye nivel temporal}\par
La información mínima de la convocatoria del evento debe incluir también su nivel temporal.\par\vspace{0.45em}
\noindent\textbf{DEC-008 --- La carga mínima es obligatoria en el evento}\par
La carga mínima del evento no solo debe poder registrarse; debe quedar establecida como condición obligatoria del evento.\par\vspace{0.45em}
\noindent\textbf{DEC-009 --- La salida invertida usa siempre la clasificación combinada definida por el nivel territorial y el nivel temporal}\par
La salida invertida usa siempre la clasificación combinada correspondiente al nivel territorial y al nivel temporal definidos para el evento.\par\vspace{0.45em}
\noindent\textbf{DEC-010 --- La cantidad disponible de accesos para espectadores es obligatoria cuando exista subasta de acceso}\par
Cuando un evento contemple subasta de acceso para espectadores, la cantidad disponible de accesos debe quedar establecida como condición obligatoria del evento.\par\vspace{0.45em}
\noindent\textbf{DEC-011 --- La inscripción del vehículo en un evento corresponde al dueño del vehículo}\par
El usuario que registra una inscripción para competir con un vehículo en un evento debe ser el dueño de ese vehículo.\par\vspace{0.45em}
\noindent\textbf{DEC-012 --- El vehículo debe estar previamente registrado en el sistema como propiedad de su dueño}\par
Un vehículo solo puede inscribirse en un evento si ya se encuentra previamente registrado en el sistema como perteneciente al usuario dueño que lo inscribe.\par\vspace{0.45em}
\noindent\textbf{DEC-013 --- Datos mínimos obligatorios del vehículo registrado}\par
Todo vehículo registrado en el sistema debe contar, como mínimo, con: \texttt{id}, \texttt{número de chasis}, \texttt{fabricante}, \texttt{número de motor} y \texttt{usuario dueño}.\par\vspace{0.45em}
\noindent\textbf{DEC-014 --- Datos mínimos obligatorios de la inscripción del competidor}\par
Toda inscripción debe registrar, como mínimo: \texttt{id del vehículo}, \texttt{id del usuario piloto} cuando exista pilotaje humano directo, \texttt{aportación de eficiencia}, \texttt{contribución de terminación}, \texttt{pago de inscripción} y \texttt{pago de renta de pista}. Cuando exista usuario piloto, este puede ser distinto del usuario dueño que realiza la inscripción.\par\vspace{0.45em}
\noindent\textbf{DEC-015 --- Esta versión sí procesa pagos reales}\par
El sistema sí procesa retenciones, cobros, liberaciones, pagos y devoluciones reales conforme a las reglas del dominio.\par\vspace{0.45em}
\noindent\textbf{DEC-016 --- La inscripción del competidor se cobra de inmediato}\par
La inscripción válida a un evento debe cobrarse de inmediato, incluida su aportación mínima obligatoria y la renta de pista cuando corresponda.\par\vspace{0.45em}
\noindent\textbf{DEC-017 --- Validez de la inscripción al pasar a \texttt{carrera} [superada por DEC-155]}\par
Esta decisión heredada consideraba válida para la carrera toda inscripción no cancelada ni descalificada. \texttt{DEC-155} la precisa y sustituye: además se requiere al menos un tiempo de vuelta válido en calificación.\par\vspace{0.45em}
\noindent\textbf{DEC-018 --- No existe un estado separado de pago en la inscripción y toda cancelación o descalificación decidida por el promotor exige nota}\par
El sistema no debe registrar un estado específico de pago para la inscripción. Si el promotor cancela o descalifica una inscripción por una causa permitida, debe dejar una nota escrita con la razón correspondiente.\par\vspace{0.45em}
\noindent\textbf{DEC-019 --- La cancelación invalida toda la inscripción pero conserva su registro}\par
Cuando una inscripción se cancela por una causa permitida antes de carrera, queda invalidada a todos los efectos del evento y no influye en ningún resultado, clasificación ni cálculo económico, pero el sistema conserva el dato de esa inscripción cancelada.\par\vspace{0.45em}
\noindent\textbf{DEC-020 --- La inscripción cancelada sigue visible como cancelada}\par
Una inscripción cancelada debe seguir mostrándose en las consultas del evento con indicación de su condición de cancelada.\par\vspace{0.45em}
\noindent\textbf{DEC-021 --- La inscripción cancelada queda excluida también de la subasta}\par
Si una inscripción queda cancelada, el vehículo asociado a esa inscripción no debe aparecer en la subasta posterior del evento.\par\vspace{0.45em}
\noindent\textbf{DEC-022 --- El número esperado de participantes se conserva como dato histórico de convocatoria}\par
El número esperado de participantes publicado en la convocatoria debe conservarse como dato histórico de esa convocatoria y no debe ajustarse por cancelaciones posteriores.\par\vspace{0.45em}
\noindent\textbf{DEC-023 --- La inspección comunitaria, las objeciones y las votaciones formales son válidas hasta antes de \texttt{carrera}}\par
La inspección comunitaria de vehículos, las objeciones de seguridad y las votaciones formales derivadas de esas objeciones pueden realizarse mientras el evento no haya entrado en \texttt{carrera}.\par\vspace{0.45em}
\noindent\textbf{DEC-024 --- Las protestas y votaciones de seguridad informan la decisión del promotor y no cancelan por sí mismas}\par
Las objeciones de seguridad y sus votaciones formales se manejan aparte y no producen por sí mismas la cancelación de una inscripción; constituyen información para que el promotor tome su decisión antes de \texttt{carrera}.\par\vspace{0.45em}
\noindent\textbf{DEC-025 --- La decisión del promotor tras una protesta de seguridad puede terminar en cancelación de la inscripción}\par
Si, con base en una objeción de seguridad, su votación formal y demás elementos disponibles, el promotor decide excluir al vehículo antes de \texttt{carrera}, el sistema debe permitir la cancelación de la inscripción afectada.\par\vspace{0.45em}
\noindent\textbf{DEC-026 --- Las protestas de seguridad y sus resultados siguen siendo visibles públicamente}\par
Aunque el promotor decida no cancelar la inscripción afectada, la protesta de seguridad, su votación formal cuando exista y el resultado final de esa tramitación deben seguir visibles en la consulta del evento.\par\vspace{0.45em}
\noindent\textbf{DEC-027 --- La identidad de quien presenta la protesta no se publica, pero sí es visible para \texttt{admin} y para el promotor}\par
La consulta pública del evento no debe mostrar quién presentó una protesta de seguridad. Esa identidad sí debe quedar visible para \texttt{admin} y para el promotor del evento.\par\vspace{0.45em}
\noindent\textbf{DEC-028 --- La votación formal de protestas de seguridad se publica con detalle por participante}\par
La publicación de una votación formal de seguridad debe mostrar el total de votos emitidos y el voto de cada participante.\par\vspace{0.45em}
\noindent\textbf{DEC-029 --- El detalle público de la votación formal identifica cada voto por nombre/usuario exacto}\par
Cuando se publique el detalle de una votación formal de seguridad, el voto de cada participante debe mostrarse asociado al nombre/usuario exacto de ese participante.\par\vspace{0.45em}
\noindent\textbf{DEC-030 --- La nota de cancelación debe quedar visible en la consulta del evento}\par
Cuando una inscripción sea cancelada por decisión del promotor, la nota escrita de esa cancelación debe quedar visible en la consulta del evento junto con la indicación de cancelación.\par\vspace{0.45em}
\noindent\textbf{DEC-031 --- La matriz de pesos puede recalcularse por afectación entre eventos, pero el P/C sigue siendo propio del evento}\par
La matriz de pesos de un evento no queda necesariamente fija desde inscripciones. Si surgen nuevos eventos que afecten los mismos campeonatos aplicables, el sistema debe recalcular los eventos afectados y reflejar ese cambio en sus matrices de pesos, sin que ello cambie por sí mismo el valor propio de P/C del evento.\par\vspace{0.45em}
\noindent\textbf{DEC-032 --- Debe existir premio de campeonato para eficiencia y puede existir también para terminación}\par
El sistema debe permitir que cualquier usuario aporte una cantidad de dinero destinada a un campeonato específico de eficiencia o de terminación para su distribución según los resultados oficiales aplicables de ese campeonato.\par\vspace{0.45em}
\noindent\textbf{DEC-033 --- La aportación directa a bolsa de campeonato da derecho a patrocinio público con anuncio e imagen}\par
Quien aporte directamente a la bolsa de un campeonato específico tiene derecho a asociar a esa aportación un anuncio con imagen para su publicación como patrocinador en una interfaz específica del campeonato patrocinado.\par\vspace{0.45em}
\noindent\textbf{DEC-034 --- Los participantes del campeonato también pueden aportar al premio de campeonato correspondiente}\par
Los participantes del campeonato correspondiente también pueden realizar aportaciones para aumentar el premio de ese campeonato, sea de eficiencia o de terminación.\par\vspace{0.45em}
\noindent\textbf{DEC-035 --- Toda inscripción a evento incluye una contribución mínima de 1 unidad monetaria}\par
El pago de inscripción a un evento incluye una contribución mínima obligatoria de 1 unidad de la moneda del evento, definida por el promotor, de modo que ninguna inscripción válida quede con contribución de evento igual a cero.\par\vspace{0.45em}
\noindent\textbf{DEC-036 --- Cada participante del campeonato tiene una contribución mínima por defecto de 1 unidad monetaria al premio de campeonato correspondiente}\par
Cuando exista premio de campeonato, cada participante del campeonato correspondiente debe contar por defecto con una contribución mínima de 1 unidad monetaria al pozo de ese campeonato, según la familia aplicable.\par\vspace{0.45em}
\noindent\textbf{DEC-037 --- La página principal debe mostrar consultas de campeonatos patrocinados y de sus patrocinadores}\par
La página principal del sistema debe poner a disposición consultas de los campeonatos patrocinados y de sus patrocinadores asociados.\par\vspace{0.45em}
\noindent\textbf{DEC-038 --- La secuencia base de pesos y puntos queda fijada en 100, 63, 40, 25, 16 y sus valores sucesivos del dominio}\par
La SRS debe fijar explícitamente la secuencia base del dominio para pesos y puntos, comenzando por 100, 63, 40, 25 y 16, con sus valores sucesivos conforme a las fuentes del proyecto.\par\vspace{0.45em}
\noindent\textbf{DEC-039 --- Regresión logarítmica de mercado [superada por DEC-160]}\par
La formulación heredada establecía una regresión con rendimiento como variable dependiente del precio. Esta dirección queda superada por \texttt{DEC-160}, que define una línea de tendencia de precios en función del rendimiento.\par\vspace{0.45em}
\noindent\textbf{DEC-040 --- El desempate en la matriz de pesos favorece al evento registrado más temprano}\par
Cuando dos eventos comparables tengan el mismo P/C en la determinación de la matriz de pesos, el desempate debe favorecer al evento registrado antes en el sistema.\par\vspace{0.45em}
\noindent\textbf{DEC-041 --- Toda aportación directa a la bolsa de un evento o de un campeonato constituye patrocinio}\par
Toda aportación que se registre directamente a la bolsa de un evento o de un campeonato debe considerarse patrocinio de ese evento o de ese campeonato.\par\vspace{0.45em}
\noindent\textbf{DEC-042 --- Las aportaciones a participantes son una función configurable y apagada por defecto}\par
El sistema puede admitir aportaciones a la contribución de un participante específico solo cuando esa función haya sido activada administrativamente; por defecto debe permanecer inactiva tanto para eventos como para campeonatos.\par\vspace{0.45em}
\noindent\textbf{DEC-043 --- Cuando la función de aportaciones a participantes esté activa, esas aportaciones sirven para compartir y aumentar el premio del participante}\par
Toda aportación válidamente aceptada a favor de un participante específico debe integrarse a la contribución total de ese participante y debe tener como finalidad compartir el premio que le corresponda y aumentar su magnitud.\par 43A. \textbf{La distinción entre patrocinio directo a bolsa y aportación a participante aplica tanto a eventos como a campeonatos cuando la función esté activa}\par La separación entre aportaciones directas a la bolsa como patrocinio y aportaciones a la contribución de un participante debe aplicarse de la misma manera en eventos individuales y en premios de campeonato, siempre que la función de aportaciones a participantes se encuentre activa.\par 43B. \textbf{Solo \texttt{admin} puede activar o desactivar la función de aportaciones a participantes}\par El promotor no puede activar ni desactivar la aceptación de aportaciones a participantes del evento ni del campeonato; esa facultad corresponde únicamente a \texttt{admin}.\par\vspace{0.45em}
\noindent\textbf{DEC-044 --- Debe ser posible patrocinar directamente un evento}\par
El sistema debe permitir aportaciones directas a la bolsa de un evento como forma de patrocinio del propio evento.\par\vspace{0.45em}
\noindent\textbf{DEC-045 --- Toda aportación directa a la bolsa da derecho a publicidad con anuncio e imagen}\par
Toda aportación directa a la bolsa de un evento o de un campeonato debe dar derecho a la publicación de un anuncio con imagen asociado a ese patrocinio.\par\vspace{0.45em}
\noindent\textbf{DEC-046 --- El premio de campeonato usa la misma lógica general del premio correspondiente, pero con resultados acumulados del campeonato}\par
El premio de campeonato de eficiencia debe seguir la misma lógica general del premio de eficiencia del evento, usando los lugares oficiales del campeonato de eficiencia en lugar del ordenamiento de eficiencia del evento. El premio de campeonato de terminación debe seguir una lógica análoga a la del premio de terminación del evento, usando como resultado acumulado de cada participante la suma de las relaciones entre vueltas completadas y vueltas esperadas en los eventos computables del campeonato correspondiente.\par\vspace{0.45em}
\noindent\textbf{DEC-047 --- La página del evento debe mostrar también los patrocinadores de los campeonatos patrocinados a los que el evento pertenezca, distinguiéndolos de los patrocinadores del evento}\par
Cuando un evento pertenezca a campeonatos patrocinados, la consulta del evento debe mostrar también esos patrocinadores y su publicidad, diferenciados de los patrocinadores directos del evento.\par\vspace{0.45em}
\noindent\textbf{DEC-048 --- Los empates en resultados del evento se deciden por fecha y hora de inscripción al evento}\par
Cuando dos participantes queden empatados en un ordenamiento del evento que deba producir puntos, base de puntos o premio, el desempate debe resolverse favoreciendo a quien se haya inscrito antes en ese evento.\par\vspace{0.45em}
\noindent\textbf{DEC-049 --- No hay redondeo de dominio intermedio [precisada por DEC-154]}\par
La intención heredada de no introducir redondeos arbitrarios se conserva, pero la expresión «todos los decimales» queda sustituida por \texttt{DEC-154} y \texttt{NUM-001}: punto flotante binario de 64 bits para cálculos analíticos, dinero entero y reconciliación monetaria exacta mediante mayor residuo.\par\vspace{0.45em}
\noindent\textbf{DEC-050 --- Los empates en el resultado computable del campeonato para efectos del premio se deciden por la fecha y hora de inscripción al primer evento de ese campeonato}\par
Cuando dos participantes queden empatados en los lugares oficiales de un campeonato de eficiencia o en el resultado acumulado de terminación de un campeonato de terminación para efectos del premio de campeonato, el desempate debe resolverse favoreciendo a quien se haya inscrito antes en el primer evento de ese campeonato en el que haya participado.\par\vspace{0.45em}
\noindent\textbf{DEC-051 --- La ventana de aportaciones al premio de campeonato cierra antes de que el primer evento del campeonato cierre inscripciones}\par
Las aportaciones directas a la bolsa del campeonato y, cuando la función de aportaciones a participantes esté activa, las aportaciones a la contribución de un participante de ese campeonato, sea de eficiencia o de terminación, solo pueden aceptarse hasta antes de que el primer evento perteneciente a ese campeonato haya cerrado inscripciones.\par\vspace{0.45em}
\noindent\textbf{DEC-052 --- La subasta de acceso usa retención suficiente y QR de acceso}\par
En la subasta de acceso para espectadores, cada puja válida debe quedar respaldada por retención suficiente; al cierre se cobra automáticamente a las personas ganadoras, se liberan de inmediato las retenciones no ganadoras y cada ganadora recibe una notificación con QR asociada a su acceso adjudicado.\par\vspace{0.45em}
\noindent\textbf{DEC-053 --- Garantía real de pago en pujas de vehículos [precisada por DEC-165]}\par
Las pujas de usuarios distintos del dueño deben quedar respaldadas por retención suficiente. La puja del dueño sobre su propio vehículo queda exenta de retención del principal conforme a \texttt{DEC-165}.\par\vspace{0.45em}
\noindent\textbf{DEC-054 --- La subasta de vehículos no inicia sin entrega previa del vehículo al promotor}\par
La subasta de un vehículo no puede iniciar si el promotor no tiene ya ese vehículo bajo su control al terminar la carrera.\par\vspace{0.45em}
\noindent\textbf{DEC-055 --- El dueño debe entregar el vehículo al promotor al terminar la carrera}\par
Al terminar la carrera, el dueño debe entregar el vehículo al promotor para su custodia, \texttt{parc fermé} y eventual transferencia al adjudicatario.\par\vspace{0.45em}
\noindent\textbf{DEC-056 --- La falta de entrega o la modificación del vehículo antes de entregarlo permiten descalificación y declaración de morosidad del dueño}\par
Si el dueño no entrega el vehículo al promotor o lo modifica antes de entregarlo, el promotor puede descalificar al competidor, declararlo moroso, excluirlo del sistema y excluirlo de los campeonatos aplicables.\par\vspace{0.45em}
\noindent\textbf{DEC-057 --- El promotor debe conservar el vehículo en \texttt{parc fermé} y entregarlo al ganador en el estado en que compitió}\par
El promotor debe custodiar el vehículo bajo \texttt{parc fermé}, impedir modificaciones entre calificación y carrera y garantizar su entrega al ganador en el estado en que compitió.\par\vspace{0.45em}
\noindent\textbf{DEC-058 --- El promotor queda como deudor hasta la confirmación de entrega del ganador}\par
Una vez cobrado el ganador y pagado el dueño, el promotor queda como deudor respecto de la entrega del vehículo hasta que la persona ganadora confirme la recepción en la app.\par\vspace{0.45em}
\noindent\textbf{DEC-059 --- El ganador puede iniciar disputa de entrega contra el promotor}\par
Si el promotor no entrega el vehículo al ganador, este puede iniciar una disputa de entrega dentro del plazo máximo publicado del evento.\par\vspace{0.45em}
\noindent\textbf{DEC-060 --- El plazo máximo para iniciar disputas de entrega es un dato del evento}\par
Todo evento con subasta de vehículos debe publicar el plazo máximo para iniciar disputas de entrega del vehículo adjudicado.\par\vspace{0.45em}
\noindent\textbf{DEC-061 --- El plazo máximo para resolver disputas de entrega es un dato del evento}\par
Todo evento con subasta de vehículos debe publicar el plazo máximo para resolver disputas de entrega del vehículo adjudicado.\par\vspace{0.45em}
\noindent\textbf{DEC-062 --- La disputa de entrega se publica}\par
Toda disputa de entrega iniciada respecto de un vehículo adjudicado debe quedar visible públicamente.\par\vspace{0.45em}
\noindent\textbf{DEC-063 --- La disputa de entrega publica las identidades del promotor y del ganador involucrados}\par
La publicación de la disputa de entrega debe mostrar la identidad del promotor y la de la persona ganadora involucrados en ella.\par\vspace{0.45em}
\noindent\textbf{DEC-064 --- La disputa de entrega publica los reportes de las partes y la decisión final de \texttt{admin}}\par
La publicación de la disputa de entrega debe mostrar los reportes presentados por las partes y la decisión final de \texttt{admin}.\par\vspace{0.45em}
\noindent\textbf{DEC-065 --- La declaración de morosidad del dueño o del promotor debe quedar visible en la consulta del evento}\par
Cuando se declare moroso al dueño por no entregar el vehículo al promotor o al promotor por no entregarlo al ganador, esa declaración debe quedar visible en la consulta del evento.\par\vspace{0.45em}
\noindent\textbf{DEC-066 --- El resultado final de la subasta de vehículo muestra solo el ganador y el precio finales vigentes}\par
Cuando existan descalificaciones, morosidad del dueño o disputas de entrega, el resultado oficial final del evento debe mostrar solo el ganador y el precio finales vigentes de la subasta correspondiente.\par\vspace{0.45em}
\noindent\textbf{DEC-067 --- La cancelación o descalificación del dueño por incumplimiento de entrega permanece visible en la consulta del evento}\par
Si el dueño incumple la entrega del vehículo al promotor y ello produce cancelación, descalificación o morosidad, esa situación debe permanecer visible en la consulta del evento junto con la información pública aplicable.\par\vspace{0.45em}
\noindent\textbf{DEC-068 --- La disputa de entrega de vehículo adjudicado debe quedar en el historial del evento con su propia secuencia temporal}\par
Toda disputa de entrega de vehículo adjudicado debe permanecer consultable en el historial del evento con su secuencia temporal propia, incluyendo su inicio, los reportes registrados, la decisión de \texttt{admin} y los efectos que de ella resulten.\par\vspace{0.45em}
\noindent\textbf{DEC-069 --- La identidad del usuario debe estar verificada antes de cualquier participación}\par
El sistema no debe permitir ninguna participación operativa en el ecosistema a un usuario cuya identidad no haya quedado verificada. La consulta pública no participativa no requiere esa verificación.\par\vspace{0.45em}
\noindent\textbf{DEC-070 --- La ventana de aportaciones al evento cierra al pasar de \texttt{inscripciones} a \texttt{velocidad}}\par
Las aportaciones directas a la bolsa del evento y las aportaciones a la contribución de un participante del evento solo pueden aceptarse mientras el evento permanezca en \texttt{inscripciones} y quedan cerradas al pasar el evento a \texttt{velocidad}.\par\vspace{0.45em}
\noindent\textbf{DEC-071 --- Todos los usuarios pueden publicar en el foro y en el canal de novedades}\par
Todos los usuarios del sistema pueden publicar tanto en el foro comunitario como en el canal de novedades del ecosistema.\par\vspace{0.45em}
\noindent\textbf{DEC-072 --- \texttt{admin} es el moderador del foro y del canal de novedades}\par
La moderación del foro comunitario y del canal de novedades del ecosistema corresponde a \texttt{admin}.\par\vspace{0.45em}
\noindent\textbf{DEC-073 --- El premio individual del campeonato es público}\par
La consulta del premio de campeonato debe mostrar públicamente el premio individual que corresponda a cada competidor en ese campeonato, sea de eficiencia o de terminación.\par\vspace{0.45em}
\noindent\textbf{DEC-074 --- Las aportaciones y la parte de premio de terceros en premio de campeonato se publican sin identidad}\par
Cuando existan aportaciones de terceros a la contribución de un competidor en un premio de campeonato, la consulta correspondiente debe publicar las cantidades aportadas y la parte del premio que les corresponda sin mostrar nombres ni identidad de esas personas.\par\vspace{0.45em}
\noindent\textbf{DEC-075 --- El premio individual del evento es público y las cantidades de terceros se publican sin identidad}\par
La consulta del evento debe mostrar públicamente el premio individual que corresponda a cada competidor, y cuando existan terceros asociados a sus contribuciones debe publicar las cantidades aportadas y la parte del premio que les corresponda sin mostrar nombres ni identidad de esas personas.\par\vspace{0.45em}
\noindent\textbf{DEC-076 --- El premio individual del campeonato sigue la misma regla de publicación que el premio individual del evento}\par
La consulta del premio de campeonato debe aplicar la misma regla de publicación: premio individual público del competidor y publicación sin identidad de las cantidades y de la parte de premio correspondientes a terceros asociados a sus contribuciones.\par\vspace{0.45em}
\noindent\textbf{DEC-077 --- El estado de identidad verificada y el estado de exclusión del usuario se representan como dos booleanos independientes}\par
El estado del usuario respecto de identidad y exclusión debe representarse mediante dos valores booleanos independientes: verificada/no verificada y excluida/no excluida.\par\vspace{0.45em}
\noindent\textbf{DEC-078 --- El promotor solo puede verificar identidad y excluir a usuarios vinculados a su evento y solo durante ese evento}\par
El promotor puede marcar como verificada la identidad de un usuario y puede excluirlo únicamente si ese usuario está vinculado operativamente a un evento que él promueve, y solo mientras ese evento se encuentre dentro de su propio ciclo de vida.\par\vspace{0.45em}
\noindent\textbf{DEC-079 --- \texttt{admin} puede verificar, desverificar, excluir y desexcluir a cualquier usuario}\par
\texttt{admin} puede modificar globalmente ambos booleanos del usuario, tanto para verificar o desverificar identidad como para excluir o desexcluir.\par\vspace{0.45em}
\noindent\textbf{DEC-080 --- El usuario excluido no puede participar aunque tenga identidad verificada}\par
Si un usuario está marcado como excluido, no puede participar operativamente en el sistema aunque su identidad permanezca verificada.\par\vspace{0.45em}
\noindent\textbf{DEC-081 --- El propio usuario ve sus dos booleanos}\par
El propio usuario puede consultar en el sistema si su identidad está verificada o no y si su estado de exclusión está activo o no.\par\vspace{0.45em}
\noindent\textbf{DEC-082 --- Otros usuarios no ven los estados de otros usuarios}\par
Ningún usuario distinto de \texttt{admin} o del propio usuario puede consultar públicamente o por simple navegación los booleanos de verificación o exclusión de otro usuario.\par\vspace{0.45em}
\noindent\textbf{DEC-083 --- Los usuarios con exclusión activa solo son visibles para \texttt{admin} en las consultas generales de usuarios}\par
En las consultas generales de usuarios del sistema, los usuarios con exclusión activa solo deben ser visibles para \texttt{admin}, sin perjuicio de que el propio usuario pueda consultar sus propios estados.\par\vspace{0.45em}
\noindent\textbf{DEC-084 --- Los usuarios sin verificación pueden permanecer visibles}\par
La falta de verificación de identidad no elimina por sí misma la visibilidad del usuario dentro de las consultas del sistema o del evento que sí le correspondan según las demás reglas del dominio.\par\vspace{0.45em}
\noindent\textbf{DEC-085 --- El evento no puede pasar a \texttt{velocidad} mientras existan inscripciones no canceladas asociadas a usuarios no verificados}\par
Mientras exista alguna inscripción no cancelada asociada a un usuario cuya identidad no está verificada, el evento no puede pasar de \texttt{inscripciones} a \texttt{velocidad}; para hacerlo, el promotor debe verificar antes la identidad de esos usuarios dentro del alcance permitido.\par\vspace{0.45em}
\noindent\textbf{DEC-086 --- El promotor solo puede consultar los booleanos de usuarios vinculados a su evento y solo durante ese evento}\par
El promotor puede consultar los estados de identidad verificada y exclusión únicamente respecto de usuarios vinculados operativamente a un evento que él promueve, y solo mientras ese evento permanezca dentro de su ciclo de vida.\par\vspace{0.45em}
\noindent\textbf{DEC-087 --- La aceptación legal de términos y condiciones es obligatoria antes de participar}\par
El sistema debe obtener, conservar y versionar la aceptación de términos y condiciones por parte de cada usuario antes de habilitarle cualquier participación operativa en el ecosistema, y debe exigir la aceptación de la versión vigente cuando esos términos cambien.\par\vspace{0.45em}
\noindent\textbf{DEC-088 --- El vehículo subastado queda sujeto a transferencia obligatoria según el resultado oficial}\par
Cuando un vehículo participe en un evento y quede sometido a la subasta de vehículos correspondiente, su dueño acepta que la transferencia del vehículo queda determinada por el resultado oficial vigente de esa subasta.\par\vspace{0.45em}
\noindent\textbf{DEC-089 --- El dueño no puede retirar del control del evento un vehículo subastado antes de resolver su adjudicación oficial}\par
El dueño de un vehículo sometido a la subasta del evento no puede retirarlo del control del evento mientras no exista resultado oficial vigente que le permita conservarlo o mientras no se haya cumplido la transferencia obligatoria correspondiente.\par\vspace{0.45em}
\noindent\textbf{DEC-090 --- El retiro no autorizado de un vehículo subastado constituye incumplimiento grave con consecuencias legales aplicables}\par
Si el dueño retira del control del evento un vehículo sometido a subasta sin que el resultado oficial vigente le permita conservarlo o sin cumplir la transferencia obligatoria correspondiente, el sistema debe tratar ese hecho como incumplimiento grave y conservarlo para las consecuencias del dominio y las acciones legales aplicables.\par\vspace{0.45em}
\noindent\textbf{DEC-091 --- La persona ganadora de un acceso recibe una notificación con QR en lugar de ticket}\par
Al cerrarse la subasta de acceso, cada persona ganadora debe recibir una notificación con un código QR asociado a su acceso adjudicado y al importe correspondiente, en lugar de un ticket independiente.\par\vspace{0.45em}
\noindent\textbf{DEC-092 --- Cada usuario tiene una sola puja vigente por subasta de acceso}\par
En cada subasta de acceso de espectadores, cada usuario solo puede mantener una puja vigente a la vez.\par\vspace{0.45em}
\noindent\textbf{DEC-093 --- La puja de acceso solo puede aumentar y nunca disminuir}\par
Mientras la subasta de acceso permanezca abierta, el usuario puede aumentar el monto de su puja vigente, pero no puede reducirlo; cada aumento reemplaza la puja vigente anterior.\par\vspace{0.45em}
\noindent\textbf{DEC-094 --- Las aportaciones y patrocinios se cobran de inmediato}\par
Toda aportación a participante y todo patrocinio directo a bolsa de evento o campeonato deben cobrarse de inmediato al quedar registrados válidamente.\par\vspace{0.45em}
\noindent\textbf{DEC-095 --- Las devoluciones por cancelación se ejecutan de inmediato}\par
Si un evento se cancela o si una inscripción se cancela o descalifica antes de \texttt{carrera}, el sistema debe devolver de inmediato los importes cobrados que correspondan.\par\vspace{0.45em}
\noindent\textbf{DEC-096 --- El sistema paga de inmediato premios y distribuciones posteriores a todos los beneficiarios}\par
Cuando el sistema determine una distribución económica posterior válida, debe pagar de inmediato a todos los beneficiarios la parte que les corresponda.\par\vspace{0.45em}
\noindent\textbf{DEC-097 --- El promotor resuelve la entrega física al ganador y \texttt{admin} resuelve la disputa de entrega}\par
El promotor realiza la entrega física del vehículo adjudicado y la persona ganadora confirma esa entrega en la app; si existe disputa de entrega, \texttt{admin} decide si se resuelve o si el promotor permanece moroso.\par\vspace{0.45em}
\noindent\textbf{DEC-098 --- Si el promotor permanece moroso, queda excluido del sistema pero sus eventos ya realizados siguen valiendo para campeonatos}\par
Cuando \texttt{admin} mantenga la morosidad del promotor por no entregar un vehículo adjudicado, el promotor queda excluido del sistema, sin invalidar por ello los eventos ya realizados para efectos de campeonatos.\par\vspace{0.45em}
\noindent\textbf{DEC-099 --- Las subastas de acceso usan retención, cobro al ganar y liberación al perder}\par
En la subasta de acceso para espectadores, cada puja válida debe quedar respaldada por retención suficiente; al cierre se cobra automáticamente a las personas ganadoras el importe adjudicado y se liberan de inmediato las retenciones de las pujas no ganadoras.\par\vspace{0.45em}
\noindent\textbf{DEC-100 --- La cancelación del evento devuelve de inmediato los accesos cobrados y libera retenciones vigentes}\par
Si un evento se cancela después de haberse adjudicado accesos de espectadores, el sistema debe devolver de inmediato los importes cobrados por esos accesos y liberar cualquier retención que permanezca vigente.\par\vspace{0.45em}
\noindent\textbf{DEC-101 --- La verificación reforzada es un estado lógico adicional del usuario}\par
El estado del usuario respecto de habilitación operativa de mayor riesgo debe incluir, además de identidad verificada y exclusión, un booleano independiente de verificación reforzada.\par\vspace{0.45em}
\noindent\textbf{DEC-102 --- Solo \texttt{admin} puede activar o desactivar la verificación reforzada}\par
El promotor no puede modificar la verificación reforzada de ningún usuario; esa facultad corresponde solo a \texttt{admin}.\par\vspace{0.45em}
\noindent\textbf{DEC-103 --- La verificación reforzada es requisito para inscribir vehículos en eventos}\par
La inscripción de un vehículo en un evento solo puede habilitarse si el usuario dueño que lo inscribe tiene verificación reforzada activa.\par\vspace{0.45em}
\noindent\textbf{DEC-104 --- El máximo nivel de verificación es requisito para pujar en subastas de vehículos}\par
La presentación de pujas en subastas de vehículos solo puede habilitarse a usuarios con verificación reforzada activa, que constituye el máximo nivel de verificación operativa. Esta elegibilidad no exige otro papel funcional.\par\vspace{0.45em}
\noindent\textbf{DEC-105 --- La verificación reforzada es requisito para promover eventos}\par
El registro y la publicación de eventos solo pueden habilitarse a usuarios con verificación reforzada activa.\par\vspace{0.45em}
\noindent\textbf{DEC-106 --- El sistema debe mostrar al usuario su nivel actual de verificación y recordarle qué le falta para habilitar operaciones bloqueadas}\par
El propio usuario debe ver sus estados aplicables, el nivel de verificación que de ellos resulte y, cuando intente una operación bloqueada por ese nivel, el sistema debe indicarle qué condición le falta cumplir para habilitarla.\par\vspace{0.45em}
\noindent\textbf{DEC-107 --- Si el sistema detecta que una nueva identidad corresponde a una persona ya excluida, debe excluir también esa nueva identidad}\par
La detección válida de correspondencia entre una nueva identidad y una persona ya excluida debe producir la exclusión de la nueva identidad.\par\vspace{0.45em}
\noindent\textbf{DEC-108 --- Al competidor se le cobra comisión solo sobre la parte positiva de su premio individual}\par
La comisión del competidor se calcula por evento y por campeonato por separado, y solo sobre la parte positiva del premio individual correspondiente respecto de los pagos que el dominio le atribuya en ese mismo caso; en el evento no se cuentan compras de vehículos.\par\vspace{0.45em}
\noindent\textbf{DEC-109 --- Al aportante asociado a un participante se le cobra comisión solo sobre la parte positiva de su premio}\par
La comisión del aportante se calcula por evento y por campeonato por separado, y solo sobre la parte positiva de la parte de premio que reciba respecto de la cantidad que aportó en ese mismo caso.\par\vspace{0.45em}
\noindent\textbf{DEC-110 --- La comisión de vehículo corresponde al comprador [precisada por DEC-166]}\par
Toda adjudicación oficial de un vehículo debe causar una comisión del sistema a cargo del comprador o adjudicatario, incluso cuando el adjudicatario sea el dueño previo. El vendedor no paga esa comisión en su condición de vendedor.\par\vspace{0.45em}
\noindent\textbf{DEC-111 --- Al vendedor de vehículo no se le cobra comisión}\par
El vendedor debe recibir integra la suma acordada de la venta del vehículo sin comisión a su cargo.\par\vspace{0.45em}
\noindent\textbf{DEC-112 --- Al patrocinador directo no se le cobra comisión}\par
Las aportaciones directas a la bolsa de eventos o campeonatos, en cuanto patrocinios, no causan comisión del sistema al patrocinador.\par\vspace{0.45em}
\noindent\textbf{DEC-113 --- Al promotor se le cobra una cuota fija por evento a partir del sexto evento del anio}\par
El promotor no paga cuota fija en sus primeros cinco eventos del anio; desde el sexto evento del mismo anio, el sistema cobra una cuota fija por cada evento nuevo al momento de crearlo.\par\vspace{0.45em}
\noindent\textbf{DEC-114 --- \texttt{admin} puede configurar parámetros de comisión y de cuota fija del modelo de negocio}\par
\texttt{admin} puede configurar los porcentajes de comisión aplicables, la cuota fija por evento del promotor y el umbral anual a partir del cuál esa cuota se activa.\par\vspace{0.45em}
\noindent\textbf{DEC-115 --- En eventos existen dos familias distintas de aportaciones asociadas a un competidor}\par
En los eventos, la aportación dirigida al premio basado en eficiencia y la aportación dirigida al premio de terminación son conceptos distintos y no deben mezclarse en la especificación ni en los cálculos del dominio.\par\vspace{0.45em}
\noindent\textbf{DEC-116 --- Cada familia de aportación comparte solo el premio de su misma familia}\par
Las aportaciones de eficiencia comparten y aumentan solo el premio del evento basado en eficiencia; las aportaciones de terminación comparten y aumentan solo el premio de terminación del mismo competidor.\par\vspace{0.45em}
\noindent\textbf{DEC-117 --- El pesaje de entrada y de salida se realiza fuera del sistema con báscula provista por el promotor}\par
El pesaje que sirve para comprobar la carga mínima al entrar y al salir de pista se realiza fuera del sistema y con báscula provista por el promotor del evento.\par\vspace{0.45em}
\noindent\textbf{DEC-118 --- El sistema solo publica el costo fijo por pesaje del evento}\par
El sistema no gestiona el pesaje ni registra cada medición; solo debe conservar y publicar el costo fijo que el promotor cobrará por cada pesaje del evento.\par\vspace{0.45em}
\noindent\textbf{DEC-119 --- El incumplimiento de pesaje de entrada impide entrar a pista y no forma parte de la inspección comunitaria}\par
Si un vehículo no cumple el pesaje de entrada, no puede entrar a pista; esa comprobación no se tramita por inspección comunitaria ni por objeción de seguridad dentro del sistema.\par\vspace{0.45em}
\noindent\textbf{DEC-120 --- El incumplimiento de pesaje de salida permite al promotor descalificar o dejar sin captura válida la velocidad}\par
Si un vehículo sale de pista sin cumplir el pesaje de salida, el promotor decide fuera del sistema si descalifica a la inscripción afectada o si simplemente no captura una velocidad válida para ella.\par\vspace{0.45em}
\noindent\textbf{DEC-121 --- La IU debe reaccionar de inmediato a toda acción iniciada por el usuario}\par
Toda acción iniciada por el usuario debe producir una reacción visible inmediata de la interfaz que confirme que el sistema ya recibio la solicitud, sin esperar a que termine la carga de la siguiente interfaz o el resultado final.\par\vspace{0.45em}
\noindent\textbf{DEC-122 --- La IU debe indicar la espera operativa y debe deshabilitar acciones problemáticas durante el procesamiento}\par
Cuando una operación quede pendiente de un resultado o procesamiento para poder continuar, la interfaz debe indicarlo expresamente y debe deshabilitar las demás acciones del usuario mientras se procesa la anterior siempre que esas acciones generen problemas si se realizan en ese momento.\par\vspace{0.45em}
\noindent\textbf{DEC-123 --- La estética, el diseño gráfico, la teoría psicológica aplicable y la UX son factores críticos de negocio}\par
La calidad estética y de experiencia de usuario de la interfaz es casi tan importante como lo funcional porque influye de forma directa en la adopción de la plataforma por el público.\par\vspace{0.45em}
\noindent\textbf{DEC-124 --- La interfaz debe poder producir una reacción positiva temprana en usuarios reales}\par
La interfaz debe poder producir una reacción positiva en usuarios reales a los 30 segundos de contacto y mantenerla a los 5 minutos, y esta cualidad debe evaluarse mediante pruebas de campo con una muestra mínima de 20 usuarios representativos, de los cuales al menos el 70\% debe reportar una valoración positiva o muy positiva en ambos momentos.\par\vspace{0.45em}
\noindent\textbf{DEC-125 --- La prueba UX de percepción inicial positiva debe aplicarse al menos a las pantallas críticas del dominio}\par
La evaluación de percepción inicial positiva debe cubrir como mínimo la interfaz principal del sistema, la inscripción de competidores, la subasta de acceso, la subasta de vehículos, el premio individual del evento, el premio de campeonato y patrocinio, y la interfaz del evento para patrocinio y aportaciones a participantes.\par\vspace{0.45em}
\noindent\textbf{DEC-126 --- Las pantallas críticas del dominio deben mostrar con jerarquía visual clara la información operativa y económica relevante}\par
Las pantallas críticas del dominio deben mostrar con jerarquía visual clara el estado actual, las acciones habilitadas y bloqueadas, la razón del bloqueo y, cuando corresponda, los importes brutos, las comisiones, las retenciones, las devoluciones y los importes netos.\par\vspace{0.45em}
\noindent\textbf{DEC-127 --- Las pantallas críticas del dominio deben presentar texto e importes con legibilidad suficiente}\par
Las pantallas críticas del dominio deben presentar texto e importes con legibilidad suficiente para una lectura rápida y sin ambigüedad de la información operativa y económica relevante.\par\vspace{0.45em}
\noindent\textbf{DEC-128 --- Las pantallas críticas del dominio deben mantener contraste visual suficiente}\par
Las pantallas críticas del dominio deben mantener contraste visual suficiente entre los elementos relevantes para distinguir estados, acciones, bloqueos e información económica crítica.\par\vspace{0.45em}
\noindent\textbf{DEC-129 --- Las pantallas equivalentes del dominio deben conservar consistencia visual y terminológica}\par
Las pantallas que representen conceptos equivalentes del dominio deben conservar patrones visuales y terminología coherentes entre sí.\par\vspace{0.45em}
\noindent\textbf{DEC-130 --- La interfaz debe prevenir errores evitables antes de confirmar acciones sensibles}\par
Antes de confirmar una acción sensible del dominio, la interfaz debe ayudar a prevenir errores mediante advertencias, contexto suficiente y mecanismos que eviten confirmaciones ambiguas o precipitadas cuando correspondan.\par\vspace{0.45em}
\noindent\textbf{DEC-131 --- La interfaz debe mostrar con claridad las consecuencias relevantes antes de confirmar acciones sensibles}\par
Antes de confirmar una acción sensible del dominio, la interfaz debe mostrar con claridad las consecuencias operativas, económicas o jurídicas relevantes que esa acción produzca cuando correspondan.\par\vspace{0.45em}
\noindent\textbf{DEC-132 --- La naturaleza de TOP no descansa en la competencia física entre deportistas humanos}\par
La naturaleza de TOP no descansa en la competencia física entre deportistas humanos, sino en la evaluación competitiva de soluciones de movilidad o desplazamiento, con control humano directo, control remoto o control autónomo, bajo reglas comunes y trazabilidad verificable.\par\vspace{0.45em}
\noindent\textbf{DEC-133 --- Las tres variantes de control compiten dentro del mismo marco disciplinario}\par
El sistema debe admitir dentro del mismo marco disciplinario las tres variantes de control computable: pilotaje humano directo, operación de control remoto y sistema autónomo.\par\vspace{0.45em}
\noindent\textbf{DEC-134 --- La inscripción de un sistema autónomo identifica a una persona responsable}\par
Cuando la inscripción corresponda a un sistema autónomo, el sistema debe registrar a la persona responsable de ese sistema para efectos de identificación, cumplimiento, seguridad y trazabilidad.\par\vspace{0.45em}
\noindent\textbf{DEC-135 --- La carrera cumple una función principalmente recreativa, espectacular y de validación operativa}\par
La carrera del evento cumple una función principalmente recreativa, espectacular y de validación operativa; puede generar puntos y clasificaciones, pero no determina por sí misma el premio principal del evento.\par\vspace{0.45em}
\noindent\textbf{DEC-136 --- La consecuencia económica directa propia de la carrera es el premio de terminación}\par
La consecuencia económica directa propia de la carrera es el premio de terminación, orientado a incentivar la completitud, la durabilidad y la preservación del espectáculo en pista.\par\vspace{0.45em}
\noindent\textbf{DEC-137 --- TOP debe conservar amplia libertad de diseño vehicular dentro de restricciones comunes}\par
TOP no debe cerrarse a una familia vehicular predefinida y debe conservar amplia libertad de diseño de vehículos dentro de las restricciones comunes de seguridad, control, carga mínima y medición.\par\vspace{0.45em}
\noindent\textbf{DEC-138 --- TOP significa Technology Open Professional}\par
La sigla \texttt{TOP} dentro del nombre institucional \texttt{TOP Racing} significa \texttt{Technology Open Professional} y debe interpretarse conforme a la identidad del proyecto como referencia a tecnología, apertura de diseño y orientación profesional del sistema.\par\vspace{0.45em}
\noindent\textbf{DEC-139 --- Professional expresa la posibilidad de profesionalización competitiva dentro del ecosistema}\par
El componente \texttt{Professional} de \texttt{TOP} no debe interpretarse solo como seriedad organizativa, sino también como posibilidad de que los protagonistas lleguen a profesionalizarse mediante ingresos derivados de su participación competitiva conforme a las reglas del sistema.\par\vspace{0.45em}
\noindent\textbf{DEC-140 --- El público pagador del evento es la fuente principal esperada de recursos del ecosistema}\par
El modelo del ecosistema debe apoyarse principalmente en que el público pague por divertirse en los eventos. De esa base deben surgir los recursos del evento que sostienen a promotores, facilitan mejorar el P/C y hacen posible la profesionalización progresiva de los protagonistas.\par\vspace{0.45em}
\noindent\textbf{DEC-141 --- El interés del público debe atraer patrocinadores hacia las bolsas de premios}\par
El interés real del público por los eventos debe servir para atraer patrocinadores a las bolsas de premios. Esa bolsa fortalece la profesionalización de los protagonistas, pero no debe confundirse con recursos libres del evento, porque su destino es el premio.\par\vspace{0.45em}
\noindent\textbf{DEC-142 --- La competencia entre promotores incentiva directamente menores costos de inscripción y mejores sedes}\par
La competencia entre promotores dentro del ecosistema debe incentivar de forma directa la oferta de costos de inscripción más bajos y de sedes o condiciones más atractivas para participantes y público.\par\vspace{0.45em}
\noindent\textbf{DEC-143 --- La bolsa mejora la atracción de participantes, no los recursos libres del evento}\par
El incremento de la bolsa del evento no debe tratarse como aumento de recursos disponibles del propio evento, porque su destino es el premio; su efecto principal dentro del círculo es aumentar la atracción de participantes.\par\vspace{0.45em}
\noindent\textbf{DEC-144 --- TOP busca crear un círculo virtuoso retroalimentado}\par
El ecosistema TOP debe entenderse como un círculo virtuoso retroalimentado en el que menores costos, mayor peso competitivo, mayor atracción, más participantes, mayor interés del público y mayor patrocinio se fortalecen mutuamente, distinguiendo entre recursos del evento y bolsas destinadas a premios.\par\vspace{0.45em}
\noindent\textbf{DEC-145 --- Los eventos de mayor peso en campeonatos de mayor nivel atraen participantes desde mayor distancia}\par
Los eventos con mayor peso dentro de campeonatos de mayor nivel deben entenderse como más buscados por los protagonistas del ecosistema y capaces de atraer participantes desde mayor distancia.\par\vspace{0.45em}
\noindent\textbf{DEC-146 --- Más participantes y mejores participantes deben tratarse como efectos distinguibles}\par
El modelo no debe fundir en una sola variable el aumento en número de participantes y la mejora en la calidad competitiva de los participantes; ambos efectos pueden coexistir, pero deben entenderse por separado dentro del círculo del ecosistema.\par\vspace{0.45em}
\noindent\textbf{DEC-147 --- Los recursos del evento aumentan los ingresos de promotores}\par
Los recursos económicos generados por el evento, especialmente a partir del público pagador, deben entenderse como base para mayores ingresos de promotores.\par\vspace{0.45em}
\noindent\textbf{DEC-148 --- Mayores ingresos de promotores aumentan el interés de otros promotores y facilitan mejorar el P/C}\par
Cuando un evento produce mayores ingresos a promotores, ello debe aumentar el interés de otros promotores por participar en el ecosistema y facilitar nuevas mejoras del P/C.\par\vspace{0.45em}
\noindent\textbf{DEC-149 --- Mejores sedes aumentan público y participantes}\par
La mejora de sedes y condiciones del evento debe entenderse como un factor capaz de aumentar tanto el interés del público como la atracción de participantes.\par\vspace{0.45em}
\noindent\textbf{DEC-150 --- Publicación provisional durante la subasta [superada por DEC-162]}\par
La decisión heredada de publicar consecuencias provisionales de eficiencia y premio queda superada. Conforme a \texttt{DEC-162}, antes del cierre no se publican resultados de subasta ni consecuencias derivadas de las pujas.\par\vspace{0.45em}
\noindent\textbf{DEC-151 --- Definición operacional de costo por participante para P/C}\par
El costo por participante es el costo mínimo obligatorio publicado de una inscripción ordinaria: inscripción, renta de pista imputable y pesaje obligatorio imputable. Se fija al abrir inscripciones, excluye aportaciones voluntarias y operaciones posteriores, y usa el número esperado publicado cuando la renta solo se conoce como total del evento.\par\vspace{0.45em}
\noindent\textbf{DEC-152 --- Secuencia base por bloques de cinco y órdenes de magnitud}\par
La secuencia base repite el bloque 100, 63, 40, 25 y 16 en órdenes de magnitud sucesivamente menores: 10, 6.3, 4, 2.5 y 1.6; después 1, 0.63, 0.4, 0.25 y 0.16; y así sucesivamente.\par\vspace{0.45em}
\noindent\textbf{DEC-153 --- Política de regresión y casos de muestra pequeña [superada en dirección por DEC-160]}\par
Se conserva la política de muestra pequeña: con un participante, este ocupa el primer lugar; con dos, ambos empatan normativamente en eficiencia y se desempatan por antigüedad de inscripción. Para tres o más participantes, la dirección y forma del modelo quedan sustituidas por \texttt{DEC-160}.\par\vspace{0.45em}
\noindent\textbf{DEC-154 --- Dinero entero y cálculo analítico en punto flotante}\par
Los importes monetarios de TOP Racing usan unidades enteras completas de la moneda del evento o campeonato, sin centavos ni otras fracciones. Los cálculos analíticos usan punto flotante binario de 64 bits. Las distribuciones monetarias proporcionales deben reconciliar exactamente el total mediante el método de mayor residuo.\par\vspace{0.45em}
\noindent\textbf{DEC-155 --- Calificación válida y carreras limitadas por tiempo o distancia}\par
Una inscripción sin al menos un tiempo de vuelta válido en calificación no califica para la carrera. La carrera puede estar limitada por tiempo o por distancia. Las vueltas esperadas se calculan sin redondeo entero: tiempo máximo sobre mejor vuelta válida para carreras por tiempo, o distancia objetivo sobre longitud oficial de vuelta para carreras por distancia.\par\vspace{0.45em}
\noindent\textbf{DEC-156 --- Umbral de retroalimentación perceptible}\par
La primera señal perceptible de recepción de una acción interactiva debe alcanzar un percentil 95 no mayor que 100 ms y ninguna medición individual mayor que 250 ms en el entorno UX normativo. Si una operación aceptada no termina dentro de 500 ms, la interfaz mantiene un estado de espera explícito.\par\vspace{0.45em}
\noindent\textbf{DEC-157 --- Bloqueo únicamente por incompatibilidad operativa}\par
La interfaz bloquea únicamente operaciones que entren en conflicto por agregado, estado dependiente, fondos, puja vigente, precondición o alcance de idempotencia. Las consultas permanecen disponibles por defecto.\par\vspace{0.45em}
\noindent\textbf{DEC-158 --- WCAG 2.2 AA y entorno UX de referencia}\par
La interfaz web debe cumplir WCAG 2.2 nivel AA y verificarse en los viewports, niveles de zoom y medios de entrada definidos por ENV-UX-001. Las pruebas de campo en pista son complementarias y no sustituyen los criterios cuantitativos de accesibilidad.\par\vspace{0.45em}
\noindent\textbf{DEC-159 --- Política para denominadores económicos cero}\par
Si la suma de coeficientes de eficiencia o de pesos de terminación es cero, no se adjudica el premio correspondiente y las cantidades computables se restituyen a su origen económico. Si un premio individual ya existe y la suma de contribuciones para su distribución posterior es cero, el participante recibe el 100\% del premio.\par
\noindent\textbf{DEC-160 --- Una única línea de tendencia modela el logaritmo del precio en función del rendimiento}\par
Después del cierre de las subastas de vehículos, el sistema debe ajustar una sola línea de tendencia global sobre todos los participantes computables. La variable explicativa es el rendimiento observado y la variable modelada es el logaritmo natural del precio oficial de subasta. El modelo log-lineal normativo es $\ln(C_i)=a+bS_i+\varepsilon_i$, con precio de tendencia $\widehat C(S_i)=\exp(a+bS_i)$, y se estima por mínimos cuadrados ordinarios cuando existe rango suficiente.\par\vspace{0.45em}
\noindent\textbf{DEC-161 --- La eficiencia mide subvaloración o sobrevaloración respecto de la tendencia}\par
La diferencia de eficiencia de un participante es el precio de tendencia asociado a su rendimiento menos el precio oficial de su vehículo: $D_i=\widehat C(S_i)-C_i$. Un valor positivo indica que el vehículo rindió como uno normalmente valorado por encima de su precio oficial y recibe una señal favorable de eficiencia; un valor negativo indica valoración superior a la tendencia para su rendimiento.\par\vspace{0.45em}
\noindent\textbf{DEC-162 --- No se publica información de pujas ni consecuencias derivadas antes del cierre}\par
Mientras la fase de subasta de vehículos del evento esté abierta, el sistema no debe publicar pujas de terceros, conteos, posiciones relativas, precio provisional, adjudicatario provisional, clasificación de eficiencia, puntos ni premio provisional. Cada postor puede consultar privadamente su propia puja vigente y la confirmación de sus propias operaciones. Los resultados de todos los vehículos se publican únicamente después del cierre de la fase.\par\vspace{0.45em}
\noindent\textbf{DEC-163 --- Una puja vigente por usuario y vehículo y sustitución solo al alza}\par
Cada usuario puede mantener como máximo una puja vigente por vehículo. Una puja vigente solo puede sustituirse por un importe estrictamente mayor; una puja igual o menor debe rechazarse sin modificar la vigente.\par\vspace{0.45em}
\noindent\textbf{DEC-164 --- Usuarios con verificación máxima pueden pujar y el dueño puede pujar por su propio vehículo}\par
La subasta de vehículos admite la participación de cualquier usuario con verificación reforzada activa, términos vigentes y sin exclusión, con independencia de sus demás papeles funcionales. El dueño puede presentar una puja sellada por su propio vehículo y esa puja participa normalmente en el ordenamiento de la subasta.\par\vspace{0.45em}
\noindent\textbf{DEC-165 --- La puja propia del dueño no retiene el principal}\par
La puja del dueño sobre su propio vehículo no exige retención del importe principal porque una eventual autoadjudicación no genera una obligación de transferirse a sí mismo ese principal. La excepción no se aplica a pujas del mismo usuario sobre vehículos ajenos y no exime de la comisión aplicable al comprador.\par\vspace{0.45em}
\noindent\textbf{DEC-166 --- La autoadjudicación conserva el vehículo y causa comisión al comprador}\par
Cuando el dueño previo presenta la mayor puja vigente, conserva el vehículo por autoadjudicación. El precio oficial sigue siendo el segundo precio válido. No se ejecuta cobro y pago circular del principal ni flujo de entrega o disputa consigo mismo. El dueño actúa como comprador adjudicatario y paga la comisión vigente calculada sobre el precio oficial.\par\vspace{0.45em}
\noindent\textbf{DEC-167 --- La puja del dueño y la autoadjudicación participan en la señal de eficiencia}\par
La puja válida del dueño puede determinar el segundo precio pagado por un tercero. Si el dueño se autoadjudica, el segundo precio oficial de ese resultado se usa como precio del vehículo en la muestra global precio--rendimiento. No existe tratamiento estadístico especial por la identidad del adjudicatario.\par\vspace{0.45em}
\noindent\textbf{DEC-168 --- El rendimiento observado se fija antes de la subasta}\par
El rendimiento usado por la eficiencia es la velocidad media de la vuelta válida más rápida de calificación, calculada como longitud oficial de vuelta dividida entre el menor tiempo válido del participante. El rendimiento queda fijado antes de abrir la subasta y no cambia por las pujas ni por el precio obtenido.\par\vspace{0.45em}
\noindent\textbf{DEC-169 --- El dueño inscrito es el participante y solo los vehículos participantes alcanzan la subasta}\par
El dueño del vehículo que registra la inscripción es el participante a efectos competitivos y económicos, aunque un piloto humano distinto pueda conducir. Solo los vehículos asociados a inscripciones computables que hayan participado en el evento pueden entrar en la subasta de vehículos, una vez cumplida la entrega al promotor.\par\vspace{0.45em}
\noindent\textbf{DEC-170 --- Durante la subasta solo se conoce públicamente el pozo inicial del premio de eficiencia}\par
Mientras la fase de subasta de vehículos permanezca abierta, el pozo económico oficial consolidado y los importes de aportaciones y contribuciones que lo integran permanecen reservados. Como referencia pública del premio de eficiencia se conserva únicamente el pozo inicial publicado; la publicación económica consolidada ocurre después del cierre conforme a las reglas vigentes.\par\vspace{0.45em}
\noindent\textbf{DEC-171 --- Pujar el mínimo real de venta es la estrategia robusta recomendada para el dueño}\par
Sea $v$ el mínimo real al que el dueño está dispuesto a vender, $H$ la mayor puja externa y $M$ la segunda mayor. Frente a pujar $v$, una puja $b>v$ solo eleva directamente el precio de una venta cuando $M<b<H$; intenta situarse dentro de la brecha oculta entre las dos mejores pujas externas. Como el sistema no revela $H$, $M$ ni señales provisionales, tampoco revela $H-M$ ni permite inferirla por una vía independiente. Pujar $v$ se considera la estrategia robusta recomendada bajo la información disponible en TOP Racing. Una puja superior es una estrategia especulativa que puede producir autoadjudicación, comisión y deterioro de eficiencia; no se afirma que $v$ sea estrategia dominante ni óptima bajo toda distribución externa imaginable.\par\vspace{0.45em}
\noindent\textbf{DEC-172 --- La adquisición de un vehículo puede tener finalidad de uso, estudio o aprendizaje técnico}\par
La transferencia de un vehículo por subasta puede ser valiosa para el adjudicatario por su uso competitivo, por su valor patrimonial o por el análisis técnico de su configuración física. TOP Racing no impone una confidencialidad tecnológica general por defecto sobre el vehículo transferido, sin perjuicio de límites legales, seguridad, derechos de terceros o términos específicos publicados.\par\vspace{0.45em}
\noindent\textbf{DEC-173 --- TOP Racing busca promover deporte de bajo costo y mérito tecnodeportivo}\par
El sistema debe entenderse como un mecanismo para promover deporte de pista accesible, desarrollo tecnológico costo--eficiente y destreza deportiva, no como una competencia de gasto absoluto. La eficiencia relativa premia la combinación de rendimiento medido y valoración económica inferior a la tendencia del evento.\par\vspace{0.45em}
\noindent\textbf{DEC-174 --- El vehículo de competición no se limita a vehículos motorizados caros}\par
La plataforma puede soportar categorías de pista con artefactos competitivos motorizados o no motorizados, caseros, abiertos, artesanales o de bajo costo, siempre que el evento publique reglas técnicas y de seguridad verificables y que exista rendimiento medible.\par\vspace{0.45em}
\noindent\textbf{DEC-175 --- La apertura tecnológica declarativa opcional queda superada por la obligación transversal posterior}\par
Esta decisión se conserva como procedencia histórica de la dimensión comunitaria de apertura tecnológica, pero queda reemplazada en su parte opcional por DEC-180. La apertura tecnológica no sustituye los resultados medidos, la seguridad ni el cálculo de eficiencia.\par\vspace{0.45em}
\noindent\textbf{DEC-176 --- El patrocinio público es rationale de política, no obligación funcional}\par
El patrocinio público o regional puede aumentar bolsas, infraestructura común o accesibilidad para amplificar incentivos a desarrollo tecnológico costo--eficiente y mérito deportivo. Esta decisión documenta la orientación de política pública y el programa de investigación; no obliga al software a garantizar resultados sociales ni reemplaza evaluación empírica.\par\vspace{0.45em}
\noindent\textbf{DEC-177 --- Los eventos comparables compiten por peso en los ámbitos territoriales y temporales que comparten}\par
La matriz de pesos debe interpretarse como una competencia por relevancia entre eventos que se solapan en los ámbitos comparativos definidos por las reglas territoriales y temporales del dominio. La entrada o mejora de P/C de un evento puede reducir el peso de otros eventos comparables en campeonatos superiores compartidos.\par\vspace{0.45em}
\noindent\textbf{DEC-178 --- La matriz de pesos genera una prima potencial de diferenciación territorial y temporal}\par
Un promotor situado en un rango inferior puede responder mejorando el P/C de su evento o eligiendo un territorio o periodo con menor solapamiento competitivo. La recuperación potencial hacia peso 100 constituye una prima de diferenciación antes de costos de entrada, demanda y economías de red. Esta decisión documenta el incentivo inducido; no obliga al promotor a trasladarse ni garantiza que la expansión sea rentable.\par\vspace{0.45em}
\noindent\textbf{DEC-179 --- La expansión descentralizada del ecosistema es rationale institucional y objeto de evaluación}\par
La presión de diferenciación producida por la matriz de pesos puede favorecer la creación de nuevos nodos territoriales y la extensión de la actividad a periodos menos congestionados, transportando hacia ellos los incentivos tecnodeportivos, comerciales y de difusión tecnológica del sistema. TOP Racing no presume cobertura óptima automática: la fragmentación, la demanda mínima, la seguridad, la calidad organizativa y las economías de red deben analizarse teórica y empíricamente.\par\vspace{0.45em}
\noindent\textbf{DEC-180 --- La tecnología abierta del vehículo es obligatoria y no implica acceso irrestricto al evento}\par
La tecnología abierta del vehículo debe ser una obligación transversal del dominio TOP Racing. La solución técnica del vehículo debe documentarse con información suficiente para ser consultada, estudiada, replicada o reutilizada conforme a los términos publicados aplicables. Esta obligación no impide que un promotor publique reglas de acceso, cupo, seguridad, sede, elegibilidad operativa o participación que limiten la admisión al evento cuando dichas reglas sean compatibles con las obligaciones normativas del dominio y estén dentro del alcance operativo de TOP Racing. La plataforma registra y aplica las reglas publicadas que pueda controlar, pero no garantiza por sí misma acceso físico o externo a un evento.\par\vspace{0.45em}
Estas decisiones se consideran vigentes mientras no sean reemplazadas por una nueva confirmación prioritaria del cliente debidamente registrada en esta SRS.\par \label{sec:01_05_vision_general}\par\vspace{0.45em}
\section{Resumen de cobertura de la refactorización}
\begin{longtable}{P{7.2cm}P{3.0cm}P{3.0cm}}
\toprule
\textbf{Elemento} & \textbf{Fuente heredada} & \textbf{Versión 2.5}\\
\midrule
Requisitos formales heredados & 295 IDs & Conservados mediante mapa de procedencia; los compuestos críticos se atomizan.\\
Requisitos canónicos v2 & No aplica & 384 obligaciones normalizadas con criterio de aceptación asociado.\\
Decisiones prioritarias & 150 entradas heredadas & DEC-001 a DEC-179; DEC-160 a DEC-179 consolidan subasta de vehículos, eficiencia, estrategia de puja del dueño, dimensión tecnodeportiva y expansión territorial--temporal inducida por la matriz de pesos.\\
Glosario & 157 términos heredados & 166 términos controlados: 157 conceptos heredados normalizados y 9 conceptos añadidos en v2.2--v2.5.\\
Calidad & 14 RNF + 12 RAC + RIE-015 & Consolidada en 15 QR sin duplicación normativa.\\
Cálculos & Dispersos en prosa y RF & Catálogo CALC con ecuaciones cerradas, NUM-001 y reconciliación monetaria CALC-MON-001.\\
Trazabilidad & Listas inversas extensas & Mapas tabulares heredado--canónico y canónico--procedencia.\\
Historial de iteraciones & En el cuerpo normativo & Referencia histórica externa al cuerpo normativo v2.5.\\
\bottomrule
\end{longtable}

\end{document}
